% comparison.tex -- master of the comparison note (Phases A and B of ugjs/PLAN.md).
% Orchestrator: Claude (Fable 5.1), 31 October 2026.  Parts in parts/, drafted by the council
% members named in each file.  Build: pdflatex x2 from this folder.
% preamble.tex -- shared preamble of the comparison note (orchestrator, 31 October 2026).
% Parts must not define macros; everything not listed here is written out explicitly.
\documentclass[11pt,a4paper]{article}
\usepackage[margin=2.9cm]{geometry}
\usepackage{amsmath,amssymb,amsthm,mathrsfs}
\usepackage[T1]{fontenc}
\usepackage{microtype}
\usepackage{booktabs}
\usepackage[colorlinks=true,linkcolor=blue,citecolor=blue]{hyperref}

% --- macros of part_i/main.tex, same names and meanings ---
\newcommand{\R}{\mathbb{R}}
\newcommand{\Z}{\mathbb{Z}}
\newcommand{\E}{\mathbb{E}}
\newcommand{\A}{\mathfrak{A}}
\newcommand{\G}{\mathcal{G}}
\newcommand{\CE}{\mathcal{C}_{\mathrm E}}
\newcommand{\CH}{\mathcal{C}_{\mathrm H}}
\newcommand{\HH}{\mathcal{H}}
\newcommand{\Om}{\Omega}
\newcommand{\ip}[2]{\langle #1,#2\rangle}
\DeclareMathOperator{\supp}{supp}
\DeclareMathOperator{\dist}{dist}
\DeclareMathOperator{\spec}{spec}
% --- macros of this note ---
\newcommand{\UGJS}{U_{\rm GJS}}
\newcommand{\GJS}[1]{\lVert #1\rVert_{\mathrm{GJS}}}   % the printed norm (1.1.7)-(1.1.9), eq. (1) of the record
\newcommand{\ev}[2]{\langle #1\rangle_{#2}}            % \ev{F}{h} = \int F dq_h
\DeclareMathOperator{\Cov}{Cov}                        % \Cov_h(F,G) = \ev{FG}{h}-\ev{F}{h}\ev{G}{h}
\newcommand{\Uc}{\mathcal U}                           % \Uc_h = \lambda\int h :P(\Phi): dz, the record's A3
\newcommand{\dmu}{d\mu_0}                              % free Euclidean measure
\newcommand{\dphi}{d\varphi_0}                         % free time-zero measure

% --- numbered statements ---
\theoremstyle{plain}
\newtheorem{dtheorem}{Theorem}
\renewcommand{\thedtheorem}{D.\arabic{dtheorem}}
\newtheorem{dlemma}[dtheorem]{Lemma}
\newtheorem{auxlemma}{Lemma}
\renewcommand{\theauxlemma}{N.\arabic{auxlemma}}
\newtheorem{auxprop}[auxlemma]{Proposition}
\newtheorem{auxcor}[auxlemma]{Corollary}
\theoremstyle{definition}
\newtheorem{auxdef}[auxlemma]{Definition}
\newtheorem{auxrem}[auxlemma]{Remark}
\newtheorem*{gapnote}{GAP}


\title{The uniform local moment from the printed clustering theorem\\[4pt]
\large Lemmas D.1--D.2 and Theorems D.3--D.4 of \texttt{ugjs/PLAN.md}: the comparison note
for Phases A and B (draft under audit; not integrated)}
\author{Alexander Stottmeister\\[3pt]
  \normalsize Institut f\"ur Theoretische Physik, Leibniz Universit\"at Hannover,\\
  \normalsize Appelstra\ss e 2, 30167 Hannover, Germany}
\date{31 October 2026}

\begin{document}
\maketitle

\begin{quote}\small
\textbf{Status (4 November 2026).} This note is the deliverable of Phases A and B of
\texttt{ugjs/PLAN.md}, drafted by the council members aqft, mp and cqft2 and assembled by the
orchestrator; it is AI-drafted and under audit. The status words inside
\S\S\ref{sec:notation}--\ref{sec:cons} (``proved'', ``proved conditional on \dots'') are the
drafters' claims under audit; the note as a whole has status \emph{conjectured} until the
audits V1 and V2 and the gates GA and GB of the plan have passed, and Appendix~\ref{app:gates}
records which have and what each claim is conditional on. Nothing in \texttt{part\_i/main.tex} has changed,
and integration as its Appendix~D is the author's decision (Checkpoint~2 of the plan). The
record's Theorem~12.2 is proved from the literal imports A1--A8 since 31 October 2026, with
A8(vi) read off GJS Theorem~1.1.11; this note does not change that status. Its content is the
mathematics of the removal: that the printed clustering theorem, with the record's Lemma~C.2,
Lemma~C.4 and Theorem~C.5, yields the uniform local moment, so that after integration GJS
Theorem~1.1.11 would no longer be used and one GJS theorem, Theorem~1.1.7, would remain imported.
\end{quote}

\section{Purpose, standard and status}\label{sec:purpose}

\emph{Purpose.} The record \cite{R} consumes, in its Theorem~5.1 and Lemmas~7.1--7.3 and 7.6,
a moment bound for sharp-time observables under the interacting Euclidean measures $dq_h$ at
slab cutoffs $h=h_{T,g}$, uniformly in $T>0$: Assumption A8(vi), equation $(\UGJS)$ of the
record, read since 31 October 2026 off GJS Theorem~1.1.11 at $B=I$, $n=0$. This note derives
that bound from the printed clustering theorem, GJS Theorem~1.1.7, (1.1.10)
\cite[p.~594]{GJS}, which the record imports as A8(iv), together with the record's own
Lemma~C.2 and Theorem~C.5 and the literal imports A3, A4 and A8(i)--(iii). Phase~A
(Theorem~\ref{thm:D3}) gives the bound at fixed spatial cutoff $g$, with constants depending on
$g$, which is all \S7 of the record needs; Phase~B (Theorem~\ref{thm:D4}) gives the bound of
Theorem~5.1 uniformly in the cutoff. Lemmas~\ref{lem:D1} and \ref{lem:D2} serve both.

\emph{Standard.} A printed statement is imported at the scope at which it is printed; this is
the standard under which the record imports A1--A8 (\texttt{ugjs/PLAN.md}, Part~7, ruling R9),
and this note uses no other. None of GJS's proofs is verified here. One case of
Theorem~1.1.7 is used that its printed proof does not reach, and it is disclosed wherever it
is used, in the plan's frozen sentence F-T, quoted verbatim (its bracketed clause is added only at
integration): GJS Theorem 1.1.7 is printed for localized monomials separated by a strip of any
orientation (p.~594). Its printed proof rotates the strip into the time direction (p.~597) and
so reaches sharp-time monomials only across strips parallel to the time-zero line. Lemma 7.4
uses only that case. In this note, Theorem~\ref{thm:D4} also uses the transverse case, at the
printed statement's scope, and every such use is flagged where it occurs;
Theorem~\ref{thm:D3} uses the parallel case only.

\emph{What the note does not do.} It does not prove the transverse case of Theorem~1.1.7; it
does not re-derive the cluster expansion; it renumbers no result of the record, and the
statements it changes are listed in \S\ref{sec:cons}: those whose constants contained $C_8$
(the $M$ of Theorem~5.1 and Corollary~5.2; Lemma~7.1, whose form changes; Lemma~7.6; the
radius $s_0$ in the proof of Theorem~7.7) and those that carry the count $N_I$, redefined on
GJS columns (Lemma~6.1 and the displays (7)--(8) of Proposition~6.2, with the notation table
and A8(ii)); it does not evaluate the constant of Theorem~C.5, which enters the constants below as $S_2(B)$; and it
proves nothing about the net of ground states beyond what the record already proves with
A8(vi). The open human readings it rests on, C5 items 1--4 and 10--13 of
\texttt{part\_i/C5\_reading\_record.md}, are named beside every status word at integration.

\emph{Inputs, by name.} From the record: A2 (the slab objects and the Feynman--Kac--Nelson
transfer, in \S\ref{sec:cons}), A3 (finite-volume stability), A5(i) (domain smoothing, for
$|V(h)|$ in the replaced Lemma~4.3(ii)), A4 (free hypercontractivity,
through Theorem~C.5 and the $3^{p'}$ comparison), A8(i)--(iv) (the cutoff class of p.~589, the
observables of p.~593, the norms of p.~594, Theorem~1.1.7), Definition~2.1 and Lemma~3.1 (the classes), Remark~2.5, \S2.6,
Proposition~3.2 and Remark~3.3 (the working coordinates), Lemmas~4.1--4.4, Lemma~6.1, Proposition~6.2 and Remark~6.3 (with the count $N_I$
redefined on GJS columns, \S\ref{sec:cons}), and from Appendix~C the normalization
Lemma~C.2, the ultraviolet Lemma~C.4 (through the identification of $\Uc_h$ with the
$L^2$-limit of $\Uc_{h,\kappa}$, which Lemma~\ref{lem:D1} and Lemma~\ref{lem:D2} use) and the
stability Theorem~C.5, each cited as proved conditional on the repairs of the V3 audit
(\texttt{ugjs/council/v3-appendix-c-report.md}: R8 for Lemma~C.2, R1--R10 for Lemma~C.4 and
Theorem~C.5); the Gaussian moments of the free field, written out in \S\ref{sec:d1}; and the
reading, named in \S\ref{sec:notation}, that identifies the note's localized pieces, the
one-square interaction terms and $V(h')(0)^2$ with GJS's monomials of p.~593. No new import is
made: every external statement used is one of A1--A8 as the record states it; every other
input is a result of the record or of this note, cited with its condition where it is used.

\section{Notation and standing conventions}\label{sec:notation}
% notation.tex -- drafted by cqft2 (council of ugjs/PLAN.md), 2 November 2026.
% Read for this part: part_i/main.tex (state of 2 November 2026; numbers and pages from
% labels-main.txt, re-checked against main.aux the same day); GJS pp. 589, 593, 594 on the
% page images gjs74/p-06, p-10, p-11 (p. 594 also at 300 dpi).  No other source is used.

This section fixes the notation of \texttt{ugjs/PLAN.md}, Part~1, and proves the elementary
facts that the later sections use: that every cutoff of the comparison lies in $\CE$
(Lemma~\ref{aux:notation-cutoffs}), the strip geometry (Lemma~\ref{aux:notation-strip}), the
ring count around the $3\times5$ box (Lemma~\ref{aux:notation-rings}), the finiteness and the
translation invariance of the stability constant $S_2$
(Proposition~\ref{aux:notation-S2}), the norm of the interaction in one square
(Lemma~\ref{aux:notation-cU}), the Wick isometry with the constant $V_*$ and the comparison
$\|Y\|_4\le3^{p'}\|Y\|_2$ (Lemmas~\ref{aux:notation-wick} and \ref{aux:notation-hyper}), and
the list of printed constants (\S\ref{sec:notation-printed}).  The record is \cite{R}; its
numbered statements are cited with number and page, and lines of
\texttt{part\_i/main.tex} refer to its state of 2 November 2026.  Every statement of this
section is \emph{conjectured} in the sense of the plan until audited.

{\sloppy\emph{Appendix~C of the record.}  The audit V3 of the plan (Part~5; report
\texttt{ugjs/council/}\allowbreak\texttt{v3-appendix-c-report.md}, 2 November 2026) found Lemmas~C.2--C.4 and
Theorem~C.5 of the record true, with every displayed computation re-derived, but not proved as
written: the ``classical Bessel facts'' of Lemma~C.4 are neither cited nor derived and $K_0$ is
undefined (\texttt{main.tex}, l.~2901, ll.~2989--2993); the constant $C_*$ of
$c_\kappa\le\frac1{2\pi}\log\kappa+C_*$ is defined nowhere (l.~2917); and the Wick-chaos
structure of $\Uc_h$ (``Wick isometry'', l.~2979; ``a sum \dots of Wick chaoses'', l.~3021;
``$\Gamma(e^{-t})Y_r=e^{-rt}Y_r$'', l.~3029) is used but neither stated nor imported.  (These
three descriptions were re-checked on \texttt{main.tex} for this section.)  The audit writes out
repairs R1--R10 that need no new mathematics.  Accordingly, every statement of this note that
uses one of these results is \emph{proved, conditional on repairs of the V3 audit}: on
R1--R10 where Lemma~C.4 or Theorem~C.5 is used, and on R8 alone where only Lemma~C.2 or the
comparison step in the proof of Theorem~C.5 is used (V3 report, \S2 and \S5, items
C5.3--C5.4).  In this section Lemma~C.4 is used in Lemma~\ref{aux:notation-cutoffs}(iv),
Lemma~\ref{aux:notation-transl}(iii)--(iv) and Lemma~\ref{aux:notation-cU}(iii)--(iv), and through
them in Proposition~\ref{aux:notation-S2}(iii) (R1--R10); Theorem~C.5 in
Proposition~\ref{aux:notation-S2}(i) (R1--R10); and the comparison step in
Lemma~\ref{aux:notation-hyper}(ii) (R8).  In Appendix~\ref{app:constants}, every constant
containing $S_2$, and the comparison constants $C_7c_U$ and $B'$ (Lemma~C.4 through
Lemma~\ref{aux:notation-cU}(iv), and the condition of Lemma~\ref{lem:D1}), rest on R1--R10, and
every constant containing only $Z$ or $9^{p'}$ on R8.\par}

\emph{Inputs of this section.}  From the record: \S2.1, Definition~2.1, Assumptions~A2, A3, A4
(its abstract form, in Lemma~\ref{aux:notation-hyper}) and A8(i)--(iv), Remark~2.5, \S2.6,
Remark~3.3, the proofs of Lemma~6.1 and Proposition~6.2 (run on GJS's columns in
Remark~\ref{aux:notation-counts}), the regularization of Appendix~C with Lemma~C.4 and
Theorem~C.5, and the comparison step in the proof of Theorem~C.5; from GJS, pp.~589, 593 and 594
as quoted.  No other source is used.  The note makes no new import; its conditional statuses rest
on the import paragraph R8 that the V3 audit proposes, and closing the GAP after
Lemma~\ref{aux:notation-wick} would need a further one.

\subsection{Coordinates, squares, rows and columns}\label{sec:notation-squares}

\emph{Working coordinates.}  We work in the working coordinates of the record, \S2.6
\cite[p.~14]{R} and Remark~3.3 \cite[p.~16]{R}: the bare mass is $m_0=M_*$ and
$\lambda<\lambda_*$, ``so that the \emph{unit-lattice} statements of Assumption~A8 apply
literally in the working coordinates'' (\S2.6).  A point of the Euclidean plane is
$z=(x^0,x^1)\in\R^2$, with $x^0$ the Euclidean time and $x^1$ the space coordinate; the
slab cutoffs are $h_{T,g}=\mathbf 1_{[-T,T]}\otimes g$, i.e.\
$h_{T,g}(z)=\mathbf 1_{[-T,T]}(x^0)\,g(x^1)$ (Assumption~A2 \cite[p.~8]{R}).  GJS also put
time first: ``$\Phi(\tau,f)\equiv\Phi(\delta_\tau\otimes f)=\int\Phi(\tau,\vec x)f(\vec
x)d\vec x$'' \cite[p.~589]{GJS}.

\emph{The squares of GJS.}  The observables of GJS are the monomials (1.1.5)--(1.1.6), and
p.~593 defines their localization as follows (read on the page image, quoted verbatim,
misprints kept): ``$Q$ is said to be localized if $w(x)$ is supported in a product
$\Delta_{j_1}\times\cdots\times\Delta_{j_n}\subset R^{2r}$ where $\Delta_{j_i}$ is the unit
lattice square in $R^2$ centered at the lattice side $j_i\in Z^2$. We assume a bound on the
$n_i$ in (1.1.5), $n_i<\bar n$. Any choice $\bar n\geqq\deg P$ is sufficient. Since there is
no restriction on $r$, $Q$ still may have an arbitrarily large degree.'' and ``For a
localized monomial $Q$ or $Q(x_0)$, let $N(\Delta)$ be the local degree, i.e.\ the sum of
the $n_i$ for which $\Delta_{j_i}=\Delta$.'' \cite[p.~593]{GJS}.  (``side'' stands for
``site''; the last index of the product is printed $n$ although the product lies in
$R^{2r}$.)  Page~593 prints no convention for points on the lines of the lattice; the
following definition supplies one, and it is used everywhere in this note.

\begin{auxdef}[squares, rows, columns; the convention on lattice lines]\label{aux:notation-squares}
For $(l,k)\in\Z^2$ let
\[
 \Delta_{l,k}:=[l-\tfrac12,\,l+\tfrac12)\times[k-\tfrac12,\,k+\tfrac12)\subset\R^2 ,
\]
the first factor being time.  \emph{Row} $l$ is $\bigcup_{k\in\Z}\Delta_{l,k}=[l-\frac12,
l+\frac12)\times\R$ and \emph{column} $k$ is $\bigcup_{l\in\Z}\Delta_{l,k}=\R\times
[k-\frac12,k+\frac12)$.  For a time $s\in\R$,
\[
 l_s:=\lfloor s+\tfrac12\rfloor
\]
is the row containing $s$, i.e.\ the unique integer with $l_s-\frac12\le s<l_s+\frac12$;
for $x^1\in\R$, $k(x^1):=\lfloor x^1+\frac12\rfloor$ is the column containing $x^1$.  In
particular a time on a lattice line, $s=l-\frac12$ with $l\in\Z$, lies in row $l_s=l$.  The
\emph{sup-index distance} of $\Delta_{l,k}$ and $\Delta_{l',k'}$ is
$\max\{|l-l'|,|k-k'|\}$.
\end{auxdef}

\begin{auxrem}[relation to the squares of p.~593 and to the record's cells]\label{aux:notation-gjslattice}
\emph{(i)} The squares $\Delta_{l,k}$ partition $\R^2$: the point $z$ lies in exactly one of
them, $\Delta_{l_{x^0},k(x^1)}$.  The closure
$\overline{\Delta}_{l,k}=[l-\frac12,l+\frac12]\times[k-\frac12,k+\frac12]$ is the unit
lattice square centred at the lattice site $(l,k)\in\Z^2$ of p.~593.  A kernel that vanishes
almost everywhere outside a product of squares $\Delta_{j_1}\times\cdots\times\Delta_{j_r}$
has its (closed) support in the product of their closures, so the monomial is localized in
the sense of p.~593; this is how ``localized'' is read in this note, for (1.1.5) and for
(1.1.6) alike (Assumption~A8(ii) \cite[p.~10]{R} defines ``localized'' for both forms).

\emph{(ii)} For a sharp-time factor at time $x_{i,0}$, localization in a closed square
requires $x_{i,0}\in[l-\frac12,l+\frac12]$; a time on a lattice line lies in two closed
squares of the same column, and Definition~\ref{aux:notation-squares} assigns the factor to
the square of row $l_{x_{i,0}}$.  Thus a sharp-time monomial at time $s$ whose spatial
kernel vanishes almost everywhere outside $\prod_i[k_i-\frac12,k_i+\frac12)$ is localized in
the squares $\Delta_{l_s,k_i}$, and these are \emph{its localization squares} in this note.
The choice of row changes neither the local degrees nor $\|w(x_0,\cdot)\|_2$, hence not the
norm (1.1.8) (eq.~(1) of the record, \cite[p.~10]{R}; cf.\ Remark~6.3 \cite[p.~20]{R}); the
strips of \S\ref{sec:notation-strips} are computed for the row $l_s$, and every strip
written there for that assignment separates the closed squares as well.

\emph{(iii)} The record counts in the cells $[j,j+1)$ (\S\S5--6 \cite[pp.~18--20]{R}),
``this lattice shifted by one half unit in space'' (Assumption~A8(ii) \cite[p.~10]{R}).
Following finding F3 of the plan, this note uses the squares of
Definition~\ref{aux:notation-squares} only.  The next remark records the one place where
the two lattices give different counts.

\emph{(iv) A reading of the source.}  GJS print (1.1.5)--(1.1.6) as formal integrals (p.~593),
define the action ``(1.8) $V(h)=\lambda\int_{R^2}:\mathscr P(\Phi(x)):h(x)dx$ as a function on
$\mathscr S'(R^2)$, where $:\ :$ denotes Wick order with respect to the measure $d\Phi$, as in
[9], [1]'' (p.~589), and give no construction.  That the random variables of this note --- the
one-square interactions $U_{\Delta,n}$ (Definition~\ref{aux:notation-UDelta}), the sharp-time
objects of Lemmas~\ref{aux:notation-wick}--\ref{aux:notation-sharpid}, the regularized limits of
\S\ref{sec:d1} and the pieces of $Q_V$ in \S\ref{sec:d4} --- are GJS's monomials of p.~593 under
GJS's $dq_h$, and that the record's $\Uc_h$ is GJS's $V(h)$, is a \emph{reading} of the source,
of the kind Assumption~A8, Scope~(d), of the record names \cite[p.~13]{R}.  It is named in the
status lines of Lemma~\ref{lem:D2} and Theorems~\ref{thm:D3}--\ref{thm:D4} and in
Appendix~\ref{app:gates}.  Item~10 of \texttt{part\_i/C5\_reading\_record.md}, still open, asks a
human reader to confirm the printed text of (1.1.5)--(1.1.6) and of the multi-time notation
(pp.~593, 595) on which the reading rests; it does not check the reading itself.  At sharp time,
Lemma~\ref{aux:notation-sharpid} shows that the two regularizations used in this note give the
same objects.
\end{auxrem}

\begin{auxrem}[counts on the squares of p.~593]\label{aux:notation-counts}
For a bounded interval $I\subset\R$ of length $L$ put
\[
 N^{\square}_I:=\#\{k\in\Z:\ |I\cap[k-\tfrac12,k+\tfrac12)|>0\},
\]
the number of columns meeting $I$ in positive measure, and recall the record's
$N_I:=\#\{j\in\Z:\ |I\cap[j,j+1)|>0\}$ (\S6 \cite[p.~19]{R}).

\emph{(i)} $N^{\square}_I\le\lfloor L\rfloor+2$.  The proof of Lemma~6.1 \cite[p.~19]{R}
uses only that the cells are consecutive unit intervals partitioning $\R$; run with the
columns: if $k_{\min}\le k_{\max}$ are the extreme columns meeting $I$ in positive measure,
then $\inf I<k_{\min}+\frac12$ and $\sup I>k_{\max}-\frac12$, so
$L>k_{\max}-k_{\min}-1$, whence $k_{\max}-k_{\min}-1\le\lceil L\rceil-1\le\lfloor
L\rfloor$ and $N^{\square}_I\le k_{\max}-k_{\min}+1\le\lfloor L\rfloor+2$.

\emph{(ii)} The two counts differ in general: $I=(-0.1,0.1)$ gives $N_I=2$ and
$N^{\square}_I=1$; $I=(0.4,0.6)$ gives $N_I=1$ and $N^{\square}_I=2$.

\emph{(iii)} Consequently the bounds that the record states with $N_I$ --- Proposition~6.2,
eqs.~(7)--(8) \cite[p.~20]{R}, and the constant $K'=K_2^2N_I^{1/2}\|f\|_2$ of
Lemma~7.6 \cite[p.~22]{R} --- are used in this note, for norms computed on the squares of
Definition~\ref{aux:notation-squares}, with $N_I$ replaced by $N^{\square}_I$.  The proof of
Proposition~6.2 uses the cells only through the decomposition
$f_i=\sum_jf_i\mathbf 1_{[j,j+1)}$ and the Cauchy--Schwarz step over the $N_I$ cells; with
the columns $[k-\frac12,k+\frac12)$ in place of the cells it runs verbatim and gives
\[
 \GJS{Q}\le K_1K_2^{2r}r!\,(N^{\square}_I)^{r/2}\prod_{i=1}^r\|f_i\|_2 ,\qquad
 \GJS{Q^2}\le K_1K_2^{4r}(2r)!\,(N^{\square}_I)^{r}\prod_{i=1}^r\|f_i\|_2^2
 \tag{N.a}
\]
for a cylinder monomial $Q=\prod_{i=1}^r\varphi_0(f_i)$ at one sharp time, $\supp
f_i\subseteq I$.  With $N_I$ itself the first bound can fail on these squares: for $r=1$ and
$f$ supported in $(0.4,0.6)$ and even about $\frac12$, the localized pieces of
$Q=\varphi_0(f)$ on the squares of Definition~\ref{aux:notation-squares} are
$\varphi_0(f\mathbf 1_{[-1/2,1/2)})$ and $\varphi_0(f\mathbf 1_{[1/2,3/2)})$, and (1.1.9)
gives $\GJS{Q}=K_1K_2^2(\|f\mathbf 1_{[-1/2,1/2)}\|_2+\|f\mathbf
1_{[1/2,3/2)}\|_2)=\sqrt2\,K_1K_2^2\|f\|_2$, while (7) with $N_I=1$ asserts
$\GJS{Q}\le K_1K_2^2\|f\|_2$.  Both counts obey the bound (i), so every constant written
with $\lfloor L\rfloor+2$ in place of the count is valid on both lattices.
\end{auxrem}

\subsection{Boxes, the restricted cutoff $h^B$ and the path $h_\alpha$}\label{sec:notation-boxes}

\begin{auxdef}[boxes; restriction; charged squares]\label{aux:notation-box}
A \emph{box} is a set
\[
 B=\bigcup_{l=l_1}^{l_2}\ \bigcup_{k=k_1}^{k_2}\Delta_{l,k}
 =[l_1-\tfrac12,\,l_2+\tfrac12)\times[k_1-\tfrac12,\,k_2+\tfrac12),
 \qquad l_1\le l_2,\ k_1\le k_2\ \text{integers};
\]
its \emph{shape} is the pair $(l_2-l_1+1,\;k_2-k_1+1)$ of the numbers of its rows and
columns, and $|B|:=(l_2-l_1+1)(k_2-k_1+1)$ is the number of its squares, which equals the
area of $B$ and of its closure $\overline B$.  Since the squares partition $\R^2$ and $B$ is
a union of squares, a square $\Delta$ satisfies $\Delta\not\subseteq B$ if and only if
$\Delta\cap B=\varnothing$.  For a function $h$ on $\R^2$ and a box $B$,
\[
 h^B:=h\,\mathbf 1_B,\qquad h_\alpha:=h^B+\alpha\,(h-h^B)\quad(\alpha\in[0,1]),
\]
and $\mathfrak S(h):=\{\Delta:\int_\Delta h\,dz>0\}$ is the set of squares \emph{charged}
by $h$.  For a measurable $h$, $\supp h$ denotes the closure of $\{z:h(z)\ne0\}$.
\end{auxdef}

The class $\CE$ is Definition~2.1 of the record \cite[p.~7]{R}:
\[
 \text{``}\CE:=\{h:\R^2\to[0,1]\ \text{measurable},\ \supp h\ \text{compact}\}
 \quad\text{(Euclidean class)''},
\]
``the printed class[\,] of \cite[p.~589]{GJS}''; on p.~589 GJS write ``Let $h$ be
a function with compact support in $R^2$, and let $0\leqq h\leqq1$.''  Assumption~A8(i)
\cite[p.~10]{R} reads: ``For $h\in\CE$, $dq_h=Z(h)^{-1}e^{-\mathcal U_h}d\mu_0$ as in
Assumption~A3''; and Assumption~A3 \cite[p.~8]{R} gives, for every $h\in\CE$,
$\Uc_h\in L^p(\dmu)$ for every $p<\infty$ and $0<Z(h)<\infty$.

\begin{auxlemma}[the cutoffs of the comparison lie in $\CE$]\label{aux:notation-cutoffs}
Let $h\in\CE$ and let $B$ be a box.
\begin{enumerate}
\item[\textup{(i)}] $h^B$, $h-h^B$ and $h\mathbf 1_\Delta$ (for every square $\Delta$) belong
to $\CE$, and $\supp h^B\subseteq\overline B\cap\supp h$.
\item[\textup{(ii)}] For every $\alpha\in[0,1]$, $h_\alpha=(1-\alpha)h^B+\alpha h=h\,(\mathbf
1_B+\alpha\mathbf 1_{\R^2\setminus B})$ belongs to $\CE$; moreover $h^B\le h_\alpha\le h$
pointwise, $h_0=h^B$, $h_1=h$, and $\supp h_\alpha\subseteq\supp h$.
\item[\textup{(iii)}] $\mathfrak S(h)$ is finite, every $\Delta\in\mathfrak S(h)$ meets
$\supp h$, and $h-h^B=\sum_{\Delta\in\mathfrak S(h),\ \Delta\cap B=\varnothing}h\mathbf
1_\Delta$ almost everywhere.
\item[\textup{(iv)}] If $h_1,\dots,h_m\in\CE$ and $c_1,\dots,c_m\in\R$ satisfy
$\sum_ic_ih_i=0$ almost everywhere, then $\sum_ic_i\Uc_{h_i}=0$ in $L^2(\dmu)$.  In
particular
\[
 \Uc_{h_\alpha}=\Uc_{h^B}+\alpha\,\Uc_{h-h^B}\quad(\alpha\in[0,1]),\qquad
 \Uc_{h-h^B}=\sum_{\Delta\in\mathfrak S(h),\ \Delta\cap B=\varnothing}\Uc_{h\mathbf 1_\Delta} .
\]
\end{enumerate}
\end{auxlemma}

\begin{proof}
(i) $B$ and $\Delta$ are Borel sets, so $h\mathbf 1_B$, $h\mathbf 1_{\R^2\setminus B}=h-h^B$
and $h\mathbf 1_\Delta$ are products of measurable functions, hence measurable; they take
values in $[0,1]$ because $h$ does and the indicators take values in $\{0,1\}$.  Each
vanishes wherever $h$ does, so its set of non-zeros lies in $\{h\ne0\}$ and its support is
a closed subset of the compact set $\supp h$, hence compact.  For $h^B$ the non-zeros lie in
$B\cap\{h\ne0\}$, whose closure lies in $\overline B\cap\supp h$.

(ii) Pointwise, $h^B+\alpha(h-h^B)=h\mathbf 1_B+\alpha h\mathbf 1_{\R^2\setminus B}$, which
equals $(1-\alpha)h^B+\alpha h$.  For $0\le\alpha\le1$ and $0\le h\le1$ this gives $0\le
h\mathbf 1_B\le h_\alpha\le h\le1$; $h_\alpha$ is measurable, and $\{h_\alpha\ne0\}\subseteq
\{h\ne0\}$ gives $\supp h_\alpha\subseteq\supp h$, compact.  At $\alpha=0,1$ the formula
gives $h^B$ and $h$.

(iii) $\supp h$ is compact, hence contained in $[-\rho,\rho]^2$ for some $\rho$, which meets
only finitely many squares.  If $\int_\Delta h>0$, the set $\Delta\cap\{h>0\}$ has positive
measure, so $\Delta$ meets $\{h\ne0\}\subseteq\supp h$; hence $\mathfrak S(h)$ is finite.
Since the squares partition $\R^2$ and $B$ is a union of squares, $\mathbf 1_{\R^2\setminus
B}=\sum_{\Delta\cap B=\varnothing}\mathbf 1_\Delta$ pointwise, so $h-h^B=\sum_{\Delta\cap
B=\varnothing}h\mathbf 1_\Delta$ pointwise, with only the finitely many squares meeting
$\{h\ne0\}$ contributing.  If $\Delta\notin\mathfrak S(h)$, then $h\ge0$ and $\int_\Delta
h=0$ give $h\mathbf 1_\Delta=0$ almost everywhere, and these terms may be dropped.

(iv) We use the record's regularization (Appendix~C, the notation paragraph before Lemma~C.2
\cite[pp.~39--40]{R}): $\Phi_\kappa=\delta_\kappa*\Phi$ with an even mollifier
$\delta_\kappa$, the variance $c_\kappa$ of $\Phi_\kappa(z)$, and
$\Uc_{h,\kappa}=\lambda\int h\,{:}\mathscr P(\Phi_\kappa){:}_{c_\kappa}dz$.  For every
$q\in\mathcal S'(\R^2)$ the function $z\mapsto\Phi_\kappa(z)(q)=(q*\delta_\kappa)(z)$ is
smooth, and ${:}\mathscr P(y){:}_{c}=\sum_na_nc^{n/2}\operatorname{He}_n(y/\sqrt c)$ is a
polynomial in $y$; so $\Uc_{h,\kappa}(q)$ is the integral of $h$ against a function bounded
on the compact $\supp h$, linear in $h$ and unchanged when $h$ is changed on a null set.
Hence $\sum_ic_i\Uc_{h_i,\kappa}=\lambda\int(\sum_ic_ih_i){:}\mathscr
P(\Phi_\kappa){:}_{c_\kappa}dz=0$ pointwise.  Lemma~C.4 \cite[p.~40]{R} states: ``There is
$c_2=c_2(\lambda,\mathscr P,\Lambda)$ such that for all $h\in\CE$ with
$\supp h\subseteq\Lambda$ and all $\kappa\ge2$, $\|\Uc_h-\Uc_{h,\kappa}\|_{L^2(d\mu_0)}^2\le
c_2\,\kappa^{-1}(\log\kappa)^{2p'}$.''  Its hypotheses hold for each $h_i$, with
$\Lambda:=\bigcup_i\supp h_i$ (compact) and every $\kappa\ge2$; it is used conditional on the
repairs R1--R10 of the V3 audit (introduction of this section), which also supply the meaning of
$\Uc_h$ as an element of $L^2(\dmu)$.  Therefore
\[
 \Bigl\|\sum_ic_i\Uc_{h_i}\Bigr\|_{L^2(\dmu)}
 \le\sum_i|c_i|\,\|\Uc_{h_i}-\Uc_{h_i,\kappa}\|_{L^2(\dmu)}
 \le\Bigl(\sum_i|c_i|\Bigr)\bigl(c_2\kappa^{-1}(\log\kappa)^{2p'}\bigr)^{1/2}
 \xrightarrow[\kappa\to\infty]{}0 .
\]
The two displayed identities are the cases $h_\alpha-h^B-\alpha(h-h^B)=0$ (all three in
$\CE$ by (i)--(ii)) and (iii).  The same identities follow without Lemma~C.4 from Step~1 of the
proof of Lemma~\ref{lem:D1} (\S\ref{sec:d1}), which proves from Lemma~\ref{aux:d1-moments}(a) that
the regularizations converge in $L^2(\dmu)$ for every bounded kernel of compact support and that
the limit is linear in the kernel, \emph{provided $\Uc_h$ is read as that limit}.  This is a
reading the record itself does not make: the record states no construction of $\Uc_h$ (V3 audit,
finding M-3); its notation paragraph asserts, without proof, that $\Uc_{h,\kappa}\to\Uc_h$
\cite[p.~39]{R} (V3 audit, item m-1); and Assumption~A3 imports $\Uc_h$ with Theorem~V.7 of
Simon's book as its source (\texttt{main.tex}, l.~535, \cite[p.~8]{R}).  Identifying the limit
with the $\Uc_h$ that A3 imports is the condition that \S\ref{sec:d1} states for Step~1 (the status
paragraph after Lemma~\ref{lem:D1}); on the reading, that identification moves into the import of
Assumption~A3.
\end{proof}

In Lemma~\ref{lem:D1} the path is used with $h_1=h^B$ and $h_2=h$; there
$h_2-h_1=h-h^B\ge0$, and it enters only as the observable $\Uc_{h-h^B}$, never as a cutoff.

\subsection{Strips and the strip width}\label{sec:notation-strips}

The exponent of (1.1.10) is a strip width, printed on p.~594 as follows (read on the page
image, quoted verbatim; the sentence ends without a period in the source): ``For two
localized monomials $Q_1$ and $Q_2$, with disjoint localization squares
$\{\Delta^{(1)}_{j_i}\}$ and $\{\Delta^{(2)}_{j_l}\}$ respectively, let $d$ be the width of
the widest strip bounded by two parallel lines, and separating the $\{\Delta^{(1)}_{j_i}\}$
from the $\{\Delta^{(2)}_{j_l}\}$'' \cite[p.~594]{GJS}.  No orientation is prescribed.

\begin{auxdef}[strips; separation; orientation; strip width]\label{aux:notation-stripdef}
Let $L_1,L_2$ be parallel lines in $\R^2$.  If $L_1\ne L_2$, the open region between them is
a \emph{strip} of width $\dist(L_1,L_2)$, and its complement is the union of two disjoint
closed half-planes $H_1,H_2$ with boundaries $L_1,L_2$.  If $L_1=L_2=:L$, the strip is the
line $L$ itself, of width $0$, and $H_1,H_2$ are the two closed half-planes bounded by $L$.
The strip \emph{separates} $X\subseteq\R^2$ from $Y\subseteq\R^2$ if $X\subseteq H_1$ and
$Y\subseteq H_2$, or the other way round; since $H_1,H_2$ are closed, it then also separates
$\overline X$ from $\overline Y$.  The strip is \emph{parallel} (to the time-zero line
$\{x^0=0\}$) if $L_1,L_2$ are lines $\{x^0=\text{const}\}$, and \emph{transverse}
otherwise.  For a square $\Delta$ and a non-empty finite set $R$ of squares with
$\Delta\notin R$,
\[
 d(\Delta,R):=\sup\Bigl\{\text{width of }S:\ S\text{ a strip separating }\Delta\text{ from }
 \textstyle\bigcup_{\Delta'\in R}\Delta'\Bigr\},
\]
defined whenever some strip separates them.
\end{auxdef}

If $Q_2$ is localized in $\Delta$ and $Q_1$ in squares of $R$, then every strip separating
$\Delta$ from $\bigcup R$ separates the localization squares of $Q_2$ from those of $Q_1$, so
its width $w$ is at most the $d$ of p.~594, and $e^{-m_1d}\le e^{-m_1w}$ because $m_1>0$.
Every application of (1.1.10) in this note is made in this weakened form, with the width of
one explicit separating strip.

\begin{auxlemma}[strip width]\label{aux:notation-strip}
Let $S:=\{\Delta_{l',k'}:l_1\le l'\le l_2,\ k_1\le k'\le k_2\}$ be a rectangle of squares, let
$R\subseteq S$ be non-empty, and let $\Delta=\Delta_{l,k}\notin S$.  Put
\[
 \rho_t:=\max\{l_1-l,\,0,\,l-l_2\},\qquad\rho_x:=\max\{k_1-k,\,0,\,k-k_2\},\qquad
 m:=\max\{\rho_t,\rho_x\},
\]
so that $m\ge1$ is the sup-index distance from $\Delta$ to the nearest square of $S$.
\begin{enumerate}
\item[\textup{(i)}] If $\rho_t\ge1$, the strip bounded by the lines $x^0=l_2+\frac12$ and
$x^0=l-\frac12$ (if $l>l_2$), resp.\ $x^0=l+\frac12$ and $x^0=l_1-\frac12$ (if $l<l_1$), is
\emph{parallel}, has width $\rho_t-1$, and separates $\Delta$ from $\bigcup S$.
\item[\textup{(ii)}] If $\rho_x\ge1$, the strip bounded by the lines $x^1=k_2+\frac12$ and
$x^1=k-\frac12$ (if $k>k_2$), resp.\ $x^1=k+\frac12$ and $x^1=k_1-\frac12$ (if $k<k_1$), is
\emph{transverse} (it is perpendicular to the time-zero line), has width $\rho_x-1$, and
separates $\Delta$ from $\bigcup S$.
\item[\textup{(iii)}] $d(\Delta,R)$ is defined and $d(\Delta,R)\ge m-1$.  If $m\ge2$, the
strip of \textup{(i)} or \textup{(ii)} realising $m-1$ has width $\ge1$.
\item[\textup{(iv)}] \textup{(Rows next to $R$.)}  If $R$ contains a square in a row $l'$
with $|l-l'|\le1$, then no parallel strip of positive width separates $\Delta$ from
$\bigcup R$: every strip of positive width that separates them is transverse.
\item[\textup{(v)}] \textup{(Monotonicity.)}  If $\varnothing\ne R_1\subseteq R$, then
$d(\Delta,R_1)\ge d(\Delta,R)$.
\end{enumerate}
\end{auxlemma}

\begin{proof}
Every square $\Delta_{l',k'}$ is contained in $\{l'-\frac12\le x^0<l'+\frac12\}$ and in
$\{k'-\frac12\le x^1<k'+\frac12\}$.

(i) Let $l>l_2$, so $\rho_t=l-l_2\ge1$.  Then $\bigcup S\subseteq\{x^0\le l_2+\frac12\}$ and
$\Delta\subseteq\{x^0\ge l-\frac12\}$.  If $\rho_t\ge2$, the lines $x^0=l_2+\frac12$ and
$x^0=l-\frac12$ are distinct, the open strip $\{l_2+\frac12<x^0<l-\frac12\}$ has width
$(l-\frac12)-(l_2+\frac12)=\rho_t-1$, and the two closed half-planes just named are the
complement of that strip; if $\rho_t=1$ the two lines coincide and the strip of width $0$
separates.  The case $l<l_1$ is the same with $\bigcup S\subseteq\{x^0\ge l_1-\frac12\}$,
$\Delta\subseteq\{x^0\le l+\frac12\}$ and width $(l_1-\frac12)-(l+\frac12)=l_1-l-1=\rho_t-1$.
The bounding lines are lines of constant $x^0$, so the strip is parallel.

(ii) The same computation in the coordinate $x^1$, with columns in place of rows; the
bounding lines $\{x^1=\text{const}\}$ are perpendicular to $\{x^0=0\}$, so the strip is
transverse.

(iii) $\Delta\notin S$ means $l\notin[l_1,l_2]$ or $k\notin[k_1,k_2]$, i.e.\ $m\ge1$.  If
$m=\rho_t$, (i) gives a strip of width $m-1$ separating $\Delta$ from $\bigcup S$;
otherwise $m=\rho_x$ and (ii) does.  As $\bigcup R\subseteq\bigcup S$, the same strip
separates $\Delta$ from $\bigcup R$.  So $d(\Delta,R)$ is defined and at least $m-1$, which
is $\ge1$ when $m\ge2$.

(iv) Let $\Delta'\in R$ lie in row $l'$ with $|l-l'|\le1$, and suppose that the parallel
strip $\{\alpha<x^0<\beta\}$, $\beta>\alpha$, separates $\Delta$ from $\bigcup R$.  If
$\Delta\subseteq\{x^0\le\alpha\}$ and $\Delta'\subseteq\{x^0\ge\beta\}$, then, because
$\Delta$ contains points with $x^0$ arbitrarily close to $l+\frac12$ and $\Delta'$ contains
points with $x^0=l'-\frac12$, we get $l+\frac12\le\alpha$ and $l'-\frac12\ge\beta$, whence
$l'-l\ge1+(\beta-\alpha)>1$, a contradiction.  The other case,
$\Delta\subseteq\{x^0\ge\beta\}$ and $\Delta'\subseteq\{x^0\le\alpha\}$, gives
$l-l'>1$ in the same way.

(v) A strip separating $\Delta$ from $\bigcup R$ separates $\Delta$ from the subset
$\bigcup R_1$, so the supremum defining $d(\Delta,R_1)$ is over a larger set.
\end{proof}

\begin{auxrem}[the bound fails for a general $R$]\label{aux:notation-counterex}
The rectangle $S\supseteq R$ in Lemma~\ref{aux:notation-strip} cannot be dropped.  Let
$R=\{\Delta_{0,0},\Delta_{0,4}\}$ and $\Delta=\Delta_{0,2}$ (verifier, item m2 of the R3
report).  Each square of $R$ is at sup-index distance $2$ from $\Delta$, but no strip, of
positive or of zero width, separates $\Delta$ from $\Delta_{0,0}\cup\Delta_{0,4}$: a closed
half-plane containing $\Delta_{0,0}\cup\Delta_{0,4}$ contains the closures and their convex
hull $[-\frac12,\frac12]\times[-\frac12,\frac92]$, which contains
$\overline\Delta_{0,2}=[-\frac12,\frac12]\times[\frac32,\frac52]$; so a separating strip would
place $\Delta$ in $H_1\cap H_2$, which is empty for positive width and a line for width
$0$, and a square lies in neither.  Thus $d(\Delta,R)$ is undefined, (1.1.10) does not apply
to a pair localized in $R$ and in $\Delta$, and the inequality
``$d(\Delta,R)\ge(\text{sup-index distance to the nearest square of }R)-1$'' is false.
Lemma~\ref{aux:notation-strip} measures the distance to a rectangle containing $R$; here
every such rectangle contains $\Delta$.
\end{auxrem}

\begin{auxrem}[tilted strips]\label{aux:notation-tilted}
Lemma~\ref{aux:notation-strip} uses axis-parallel strips and gives lower bounds only; the
widest separating strip may be tilted and wider.  For $R=\{\Delta_{0,0}\}$ and
$\Delta=\Delta_{3,3}$ the axis-parallel strips have width $2$, while the strip between the
lines $x^0+x^1=1$ and $x^0+x^1=5$ separates $\overline\Delta_{0,0}$ (on which $x^0+x^1\le1$)
from $\overline\Delta_{3,3}$ (on which $x^0+x^1\ge5$) and has width $4/\sqrt2=2\sqrt2$.  A
tilted strip is transverse in the sense of Definition~\ref{aux:notation-stripdef}; the note
never needs one.
\end{auxrem}

\subsection{The $3\times5$ box and its rings}\label{sec:notation-rings}

\begin{auxdef}[the box $B_{3\times5}$ and its rings]\label{aux:notation-ringdef}
For $a\in\R$ let $k_a:=\lfloor a+\frac12\rfloor$ and
\[
 S_a:=\{\Delta_{0,k}:\ k_a\le k\le k_a+2\},
\]
\[
 B^{(j)}:=\bigcup\{\Delta_{l,k}:\ |l|\le1+j,\ k_a-1-j\le k\le k_a+3+j\}\qquad(j\ge0),
\]
and $B_{3\times5}=B_{3\times5}(a):=B^{(0)}$, the rows $-1,0,1$ times the columns
$k_a-1,\dots,k_a+3$.  For $j\ge1$ the \emph{ring} $\mathcal L_j$ is the set of squares
contained in $B^{(j)}$ and not in $B^{(j-1)}$.  For a square $\Delta_{l,k}$ put
$\delta(\Delta_{l,k}):=\max\{|l|-1,\ k_a-1-k,\ k-k_a-3\}$.
\end{auxdef}

\begin{auxlemma}[ring count and strip widths around $B_{3\times5}$]\label{aux:notation-rings}
Let $a\in\R$ and let $R\subseteq S_a$ be non-empty.
\begin{enumerate}
\item[\textup{(i)}] The columns $k$ with $[k-\frac12,k+\frac12)\cap[a,a+2]\ne\varnothing$ are
exactly $k_a,k_a+1,k_a+2$; the columns meeting $[a,a+2]$ in positive measure are among
them.  Hence every square of row $0$ whose column meets $[a,a+2]$ lies in $S_a$, and
$|B_{3\times5}|=15$.
\item[\textup{(ii)}] $\Delta\subseteq B^{(j)}$ if and only if $\delta(\Delta)\le j$.  The rings
$\mathcal L_j=\{\Delta:\delta(\Delta)=j\}$, $j\ge1$, are pairwise disjoint, their union is
the set of all squares $\Delta\not\subseteq B_{3\times5}$, and $|\mathcal L_j|=8j+12$.
\item[\textup{(iii)}] The $20$ squares of $\mathcal L_1$ are $\Delta_{2,k}$ and $\Delta_{-2,k}$
for $k_a-2\le k\le k_a+4$ (fourteen squares) and $\Delta_{l,k_a-2}$, $\Delta_{l,k_a+4}$ for
$|l|\le1$ (six squares).
\item[\textup{(iv)}] Every $\Delta=\Delta_{l,k}\in\mathcal L_j$ satisfies $d(\Delta,R)\ge j$.  If
$|l|=1+j$, a parallel strip of width $j$ separates $\Delta$ from $\bigcup R$.  If $|l|\le j$,
then $\Delta$ lies in column $k_a-1-j$ or $k_a+3+j$, a transverse strip of width $j$
separates it, and every parallel strip separating it has width at most $|l|-1\le j-1$ (none
of positive width if $|l|\le1$).  The four corner squares ($|l|=1+j$ in the two side
columns) admit both strips of width $j$.
\item[\textup{(v)}] For the squares of rows $-1,0,1$ outside $B_{3\times5}$, i.e.\
$\Delta_{l,k}$ with $|l|\le1$ and $k\le k_a-2$ or $k\ge k_a+4$, no parallel strip of positive
width separates $\Delta_{l,k}$ from $\bigcup R$: every strip of positive width separating
them is transverse.
\item[\textup{(vi)}] For every $m_1>0$, with $x:=e^{-m_1}\in(0,1)$, put
\[
 \overline\Sigma_{3\times5}:=\sum_{j\ge1}(8j+12)\,e^{-m_1j}=\frac{4x\,(5-3x)}{(1-x)^2}
 =\frac{4e^{-m_1}\,(5-3e^{-m_1})}{(1-e^{-m_1})^2} .
\]
With the strips of \textup{(iv)}, of width $w_\Delta:=j$ for $\Delta\in\mathcal L_j$,
$\sum_{\Delta\not\subseteq B_{3\times5}}e^{-m_1w_\Delta}=\overline\Sigma_{3\times5}$; and the sum
over the widest strips, which depends on $R$, satisfies
\[
 \Sigma_{3\times5}:=\sum_{\Delta\not\subseteq B_{3\times5}}e^{-m_1d(\Delta,R)}
 \le\overline\Sigma_{3\times5} .
\]
\end{enumerate}
\end{auxlemma}

\begin{proof}
(i) $[k-\frac12,k+\frac12)\cap[a,a+2]\ne\varnothing$ if and only if $k-\frac12\le a+2$ and
$k+\frac12>a$, i.e.\ $a-\frac12<k\le a+\frac52$.  The interval $(a-\frac12,a+\frac52]$ is
half-open of length $3$, so it contains exactly three integers; the smallest integer
$>a-\frac12$ is $\lfloor a-\frac12\rfloor+1=\lfloor a+\frac12\rfloor=k_a$.  A column meeting
$[a,a+2]$ in positive measure meets it.  The rows of $S_a$ and the columns of $B^{(0)}$ give
$|B_{3\times5}|=3\cdot5=15$.

(ii) $\Delta_{l,k}\subseteq B^{(j)}$ iff $|l|\le1+j$, $k\ge k_a-1-j$ and $k\le k_a+3+j$, i.e.\
iff $|l|-1\le j$, $k_a-1-k\le j$ and $k-k_a-3\le j$, i.e.\ iff $\delta(\Delta_{l,k})\le j$.
Since $\delta$ takes integer values, $\mathcal L_j=\{\delta=j\}$; these sets are disjoint
for different $j$, and $\Delta\not\subseteq B_{3\times5}=B^{(0)}$ iff $\delta(\Delta)\ge1$.
$B^{(j)}$ has $2(1+j)+1=3+2j$ rows and $(k_a+3+j)-(k_a-1-j)+1=5+2j$ columns, and
$B^{(j-1)}\subseteq B^{(j)}$, so
\[
 |\mathcal L_j|=(3+2j)(5+2j)-(1+2j)(3+2j)=(3+2j)\bigl[(5+2j)-(1+2j)\bigr]=4(3+2j)=8j+12 .
\]

(iii) $\delta(\Delta_{l,k})=1$ means: either $|l|=2$ and $k_a-2\le k\le k_a+4$ ($2\times7$
squares), or $|l|\le1$ and $k\in\{k_a-2,k_a+4\}$ ($3\times2$ squares); $14+6=20=8\cdot1+12$.

(iv) $S_a$ is the rectangle of Lemma~\ref{aux:notation-strip} with $l_1=l_2=0$, $k_1=k_a$,
$k_2=k_a+2$, and $\Delta\notin S_a$ since $\delta(\Delta)\ge1$.  Let $\Delta=\Delta_{l,k}\in
\mathcal L_j$.  If $|l|-1=j$, then $\rho_t=|l|=j+1\ge2$ and Lemma~\ref{aux:notation-strip}(i)
gives a parallel strip of width $j$.  If $k_a-1-k=j$, then $\rho_x=k_a-k=j+1$ and
Lemma~\ref{aux:notation-strip}(ii) gives a transverse strip of width $j$; if $k-k_a-3=j$, then
$\rho_x=k-(k_a+2)=j+1$, and again a transverse strip of width $j$.  At least one of the three
cases holds because $\delta(\Delta)=j$; if $|l|\le j$ the first fails, so $\Delta$ lies in a
side column.  Each strip separates $\Delta$ from $\bigcup S_a\supseteq\bigcup R$, so
$d(\Delta,R)\ge j$.  For a parallel strip $\{\alpha<x^0<\beta\}$ separating $\Delta$ from a
square $\Delta'\in R$ of row $0$, the computation in the proof of
Lemma~\ref{aux:notation-strip}(iv) gives $\beta-\alpha\le|l|-1$; for $|l|\le1$ this excludes
positive width.  The corners have $|l|=1+j$ and $k\in\{k_a-1-j,k_a+3+j\}$.

(v) $R$ is a non-empty subset of row $0$, and $|l-0|\le1$; apply
Lemma~\ref{aux:notation-strip}(iv).

(vi) By (ii) the sum over $\Delta\not\subseteq B_{3\times5}$ is the sum over $j\ge1$ of the
sums over $\mathcal L_j$, which have $8j+12$ terms each, all equal to $e^{-m_1w_\Delta}=x^j$;
this gives $\sum_\Delta e^{-m_1w_\Delta}=\sum_{j\ge1}(8j+12)x^j$.  By (iv), $d(\Delta,R)\ge
w_\Delta=j$ for $\Delta\in\mathcal L_j$, and $m_1>0$, so $e^{-m_1d(\Delta,R)}\le x^j$ and
$\Sigma_{3\times5}\le\overline\Sigma_{3\times5}$.  The closed form: for $0\le x<1$ and $N\ge1$,
\[
 (1-x)\sum_{j=1}^Nx^j=x-x^{N+1},\qquad
 (1-x)\sum_{j=1}^Njx^j=\sum_{j=1}^Nx^j-Nx^{N+1},
\]
and $x^{N+1}\to0$, $Nx^{N+1}\to0$ as $N\to\infty$; hence $\sum_{j\ge1}x^j=x/(1-x)$ and
$\sum_{j\ge1}jx^j=x/(1-x)^2$, and
\[
 \sum_{j\ge1}(8j+12)x^j=\frac{8x}{(1-x)^2}+\frac{12x}{1-x}
 =\frac{8x+12x-12x^2}{(1-x)^2}=\frac{4x(5-3x)}{(1-x)^2} .
\]
Test by hand: $4x(5-3x)(1-x)^{-2}=4x(5-3x)(1+2x+3x^2+\cdots)=20x+28x^2+\cdots$, and
$8\cdot1+12=20$, $8\cdot2+12=28$.
\end{proof}

\begin{auxrem}[comparison and check]\label{aux:notation-ringrem}
(a) The bound $(2j+3)(2j+5)$ of the plan (\texttt{aqft}) is $|B^{(j)}|$, the number of
squares of the box that contains $\mathcal L_j$; it is valid and cruder.
(b) \emph{Where the transverse case enters.}  The squares of rows $-1,0,1$ outside
$B_{3\times5}$ admit no parallel strip of positive width
(Lemma~\ref{aux:notation-rings}(v)).  The other side squares of $\mathcal L_j$, in rows
$2\le|l|\le j$, admit parallel strips, but only of width $\le|l|-1<j$; the bound
$d(\Delta,R)\ge j$ uses the transverse strip for them too.  It cannot be replaced by the
parallel width without losing the uniformity: for fixed $|l|\ge2$ the parallel width is the
same in every column, and the number of squares of row $l$ charged by $h$ is not bounded
uniformly in $h\in\CE$.  So the uniform bound $\overline\Sigma_{3\times5}$ of (vi), with the
strips of (iv), uses the transverse case of Theorem~1.1.7
for every square of $\mathcal L_j$ with $|l|\le j$, not only for rows $-1,0,1$; each such
use falls under sentence F-T of \S\ref{sec:purpose}.
(c) \emph{Numerical check} (\texttt{parts/cqft2\_ringcheck.py}; nothing is deduced from
it): for $j=1,\dots,12$ and several $a$, the script lists the squares with $\delta=j$,
confirms $8j+12$ of them, computes for each the axis-parallel width of
Lemma~\ref{aux:notation-strip} and confirms it is $\ge j$, and compares partial sums of
$\sum(8j+12)x^j$ with the closed form in exact rational arithmetic.  Run record: 3 November 2026,
\texttt{/usr/local/bin/python3} (Python 3.14.7, numpy 2.5.2, scipy 1.18.1), output
\texttt{parts/cqft2\_ringcheck.output.txt}, ``ALL CHECKS PASSED''; it also checks the kernel
bound of Lemma~\ref{aux:notation-kernel} against quadrature (relative tolerance $10^{-11}$) and
the inequality $(4n-1)!!\le4^n(2n)!$ of Appendix~\ref{app:constants} for $n\le40$.
\end{auxrem}

\subsection{Translation covariance and the stability constant $S_2$}\label{sec:notation-S2}

The free Euclidean measure $\dmu$ is the centred Gaussian measure on $\mathcal
S'(\R^2)$, with the $\sigma$-algebra generated by the coordinate functions
$\Phi(f)(q)=\langle q,f\rangle$, and covariance $G:=(-\Delta_{\R^2}+m_0^2)^{-1}$
(Assumption~A2 \cite[p.~8]{R}; \cite[p.~589]{GJS}: ``$d\Phi$ is the Gaussian measure with mean
zero and covariance $(-\Delta+m_0^2)^{-1}$'').

\begin{auxlemma}[translation covariance]\label{aux:notation-transl}
For $a\in\R^2$ let $(\tau_af)(z):=f(z-a)$ for functions on $\R^2$, and $(T_aq)(f):=q(\tau_af)$
for $q\in\mathcal S'(\R^2)$, $f\in\mathcal S(\R^2)$.
\begin{enumerate}
\item[\textup{(i)}] $T_a$ is a measurable bijection of $\mathcal S'(\R^2)$ with inverse
$T_{-a}$, and $(T_a)_*\mu_0=\mu_0$.  Hence $\int F\circ T_a\,\dmu=\int F\,\dmu$ for every
measurable $F\ge0$, and $F\mapsto F\circ T_a$ is an isometry of every $L^p(\dmu)$.
\item[\textup{(ii)}] $\tau_a$ maps $\CE$ bijectively onto $\CE$, with $\supp\tau_ah=\supp
h+a$.
\item[\textup{(iii)}] For every $h\in\CE$, $\Uc_{\tau_ah}=\Uc_h\circ T_a$ almost everywhere.
\item[\textup{(iv)}] For every $h\in\CE$, $Z(\tau_ah)=Z(h)$, and $\int F\circ T_a\,dq_{\tau_ah}=
\int F\,dq_h$ for every measurable $F\ge0$ and every $F\in L^1(dq_h)$.
\end{enumerate}
\end{auxlemma}

\begin{proof}
(i) $\Phi(f)\circ T_a=\Phi(\tau_af)$, since $\Phi(f)(T_aq)=(T_aq)(f)=q(\tau_af)$; so the
preimage of a cylinder set is a cylinder set, and $T_a$ is measurable, with inverse $T_{-a}$.
For $f_1,\dots,f_m\in\mathcal S(\R^2)$, the vector $(\Phi(f_i)\circ T_a)_i=(\Phi(\tau_af_i))_i$
is centred Gaussian under $\mu_0$ with covariance matrix $\ip{\tau_af_i}{G\,\tau_af_j}$.  In
Fourier space $\widehat{\tau_af}(p)=e^{-ip\cdot a}\hat f(p)$ and $G$ is multiplication by
$(p^2+m_0^2)^{-1}$, so the phases cancel and $\ip{\tau_af_i}{G\,\tau_af_j}=\ip{f_i}{Gf_j}$.
Thus $(T_a)_*\mu_0$ and $\mu_0$ have the same finite-dimensional distributions; they agree on
the cylinder sets, a $\pi$-system generating the $\sigma$-algebra, hence everywhere.  The
integral identity follows for indicators, then for simple functions, then for $F\ge0$ by
monotone convergence; applied to $|F|^p$ it gives the isometry.

(ii) $\tau_ah$ is measurable with values in $[0,1]$, and $\{\tau_ah\ne0\}=\{h\ne0\}+a$, so
$\supp\tau_ah=\supp h+a$ is compact; $\tau_{-a}$ is the inverse.

(iii) In the record's regularization (Appendix~C \cite[pp.~39--40]{R}),
$\Phi_\kappa(z)=\Phi(\delta_\kappa(z-\cdot))$ with $\delta_\kappa$ even, so
$\delta_\kappa(z-\cdot)=\tau_z\delta_\kappa$ and, by (i),
$\Phi_\kappa(z)\circ T_a=\Phi(\tau_a\tau_z\delta_\kappa)=\Phi_\kappa(z+a)$.  The Wick constant
$c_\kappa=\ip{\tau_z\delta_\kappa}{G\,\tau_z\delta_\kappa}=\ip{\delta_\kappa}{G\delta_\kappa}$
does not depend on $z$ (by the computation in (i); the record notes ``translation
invariance'').  Hence, for every $q$,
\begin{align*}
 \Uc_{h,\kappa}(T_aq)&=\lambda\int h(z)\,{:}\mathscr P(\Phi_\kappa(z+a)(q)){:}_{c_\kappa}\,dz\\
 &=\lambda\int h(w-a)\,{:}\mathscr P(\Phi_\kappa(w)(q)){:}_{c_\kappa}\,dw
 =\Uc_{\tau_ah,\kappa}(q) .
\end{align*}
Lemma~C.4 \cite[p.~40]{R} (quoted in the proof of Lemma~\ref{aux:notation-cutoffs}; used
conditional on the repairs R1--R10 of the V3 audit), applied to $h$ with $\Lambda=\supp h$ and
to $\tau_ah\in\CE$ with $\Lambda=\supp h+a$, gives
$\Uc_{h,\kappa}\to\Uc_h$ and $\Uc_{\tau_ah,\kappa}\to\Uc_{\tau_ah}$ in $L^2(\dmu)$ as
$\kappa\to\infty$; by (i), $\Uc_{h,\kappa}\circ T_a\to\Uc_h\circ T_a$ in $L^2(\dmu)$.  The
sequence $\Uc_{h,\kappa}\circ T_a=\Uc_{\tau_ah,\kappa}$ has both limits, so they agree almost
everywhere.

(iv) By (iii), $e^{-\Uc_{\tau_ah}}=e^{-\Uc_h}\circ T_a$ almost everywhere, and (i) gives
$Z(\tau_ah)=Z(h)$ and, for $F\ge0$, $\int(F\circ T_a)e^{-\Uc_{\tau_ah}}\dmu=\int(F
e^{-\Uc_h})\circ T_a\,\dmu=\int Fe^{-\Uc_h}\dmu$.  Divide by $Z(h)$; for $F\in L^1(dq_h)$
apply this to $F^\pm$.
\end{proof}

\begin{auxdef}[the stability constant]\label{aux:notation-S2def}
For a compact set $\Lambda\subset\R^2$,
\[
 S_2(\Lambda):=\sup\bigl\{\|e^{-\Uc_{h'}}\|_{L^2(\dmu)}:\ h'\in\CE,\ \supp h'\subseteq\Lambda\bigr\},
\]
and for a box $B$, $S_2(B):=S_2(\overline B)$.
\end{auxdef}

\begin{auxprop}[$S_2$ is finite and depends on the shape only]\label{aux:notation-S2}
\begin{enumerate}
\item[\textup{(i)}] $1\le S_2(\Lambda)<\infty$ for every compact $\Lambda$.
\item[\textup{(ii)}] $\Lambda\subseteq\Lambda'$ implies $S_2(\Lambda)\le S_2(\Lambda')$.
\item[\textup{(iii)}] $S_2(\Lambda+a)=S_2(\Lambda)$ for every $a\in\R^2$.  In particular, two
boxes of the same shape have the same $S_2$: $S_2(B)$ depends on the box $B$ only through its
shape.
\item[\textup{(iv)}] For every box $B$ and every $h\in\CE$,
$\|e^{-\Uc_{h^B}}\|_{L^2(\dmu)}\le S_2(B)$.
\end{enumerate}
\end{auxprop}

\begin{proof}
(i) Theorem~C.5 of the record \cite[p.~41]{R} reads: ``For every $p<\infty$ and compact
$\Lambda$, $\sup\{\|e^{-\mathcal U_h}\|_{L^p(d\mu_0)}:h\in\CE,\ \supp h\subseteq\Lambda\}<\infty$.''
Its hypotheses, checked: $p=2<\infty$; $\Lambda$ compact.  Its standing data are those of
the record's \S2.1 \cite[p.~6]{R} and Appendix~C, which hold here: $m_0>0$; $\mathscr
P(\xi)=\sum_{n=0}^{2p'}a_n\xi^n$ with $a_{2p'}>0$ and $p'\ge1$ (used in Lemma~C.3); $\lambda\ge0$
(its proof treats $\lambda=0$ separately).  The supremum in
Definition~\ref{aux:notation-S2def} is the supremum of Theorem~C.5 at $p=2$, so
$S_2(\Lambda)<\infty$; this is proved conditional on the repairs R1--R10 of the V3 audit.  (The
V3 report, \S5, item C5.1, also observes that the bound produced by the proof of Theorem~C.5
depends on $\Lambda$ only through $|\Lambda|$ and $\operatorname{diam}\Lambda$, which the
statement does not record; (iii) below proves the stronger fact that the supremum itself is
translation invariant.)  The function $h'\equiv0$ is measurable, $[0,1]$-valued and has empty,
hence compact, support contained in $\Lambda$; $\Uc_0=0$ by
Lemma~\ref{aux:notation-cutoffs}(iv) ($m=1$, $c_1=1$, $h_1=0$), so $\|e^{-\Uc_0}\|_2=1$ and
$S_2(\Lambda)\ge1$.

(ii) The supremum is over a larger set.

(iii) By Lemma~\ref{aux:notation-transl}(ii), $h'\mapsto\tau_ah'$ is a bijection from
$\{h'\in\CE:\supp h'\subseteq\Lambda\}$ onto $\{h''\in\CE:\supp h''\subseteq\Lambda+a\}$, and
by (iii) and (i) of that lemma
$\|e^{-\Uc_{\tau_ah'}}\|_{L^2(\dmu)}=\|e^{-\Uc_{h'}}\circ T_a\|_{L^2(\dmu)}=
\|e^{-\Uc_{h'}}\|_{L^2(\dmu)}$.  Two boxes of the same shape differ by a translation
$a\in\Z^2$, and so do their closures.

(iv) $h^B\in\CE$ and $\supp h^B\subseteq\overline B$ by
Lemma~\ref{aux:notation-cutoffs}(i).
\end{proof}

\begin{auxrem}[what $S_2$ is not]\label{aux:notation-S2rem}
$S_2(B)$ is not evaluated.  The proof of Theorem~C.5 ends ``with a bound depending only on
$(p,\lambda,\mathscr P,\Lambda)$'' \cite[p.~41]{R} and produces no number, so every constant
below that contains $S_2$ is finite and unevaluated, and the constant $M_B$ of
Theorem~\ref{thm:D4} is not explicit (plan, Part~2).  Theorem~C.5 is a bound under $\dmu$ for
a fixed compact region; it is not the stability estimate (2.4.34) of GJS under interpolated
measures and has no growth linear in the volume (plan, trap~14).  By
Proposition~\ref{aux:notation-S2}(iii), $S_2(B_s(g,I))$ in Theorem~\ref{thm:D3} depends only
on $N_B$, and $S_2(B_{3\times5}(a))$ in Theorem~\ref{thm:D4} is one number for all $a$.
\end{auxrem}

\subsection{The interaction in one square and the constant $c_U$}\label{sec:notation-cU}

The printed norms are used exactly as the record transcribes them in its eq.~(1),
Assumption~A8(iii) \cite[p.~10]{R}:
\[
\begin{gathered}
 \text{(1.1.7)}\ \ \GJS{Q}=K_1\Bigl[\prod_\Delta K_2^{2N(\Delta)}N(\Delta)!\Bigr]\|w\|_2,\\
 \text{(1.1.8)}\ \ \GJS{Q(x_0)}=K_1\Bigl[\prod_\Delta K_2^{2N(\Delta)}N(\Delta)!\Bigr]\|w(x_0,\cdot)\|_2,
\end{gathered}
\]
``where $N(\Delta)$ is the total degree localized in the unit square $\Delta$, and (1.1.9): a
nonlocalized polynomial is normed by summing its localized pieces'', with $K_2\ge1$ fixed
there; $K_1$ enters to the first power and the local weight is $K_2^{2N(\Delta)}N(\Delta)!$.
On p.~594 GJS print (1.1.9) as: ``If $Q$ is not localized, or is a general polynomial, it is
written as a sum of localized monomials $Q_i$, and then we set (1.1.9) $\|Q\|=\sum_i\|Q_i\|$.''
\cite[p.~594]{GJS}.  The interaction is $\Uc_h=\lambda\int h(z){:}\mathscr P(\Phi(z)){:}\,dz$
with $\mathscr P(\xi)=\sum_{n=0}^{2p'}a_n\xi^n$ (Assumption~A2 \cite[p.~8]{R}, \S2.1
\cite[p.~6]{R}).

\begin{auxdef}[the interaction in one square]\label{aux:notation-UDelta}
For $h\in\CE$, a square $\Delta$ and $0\le n\le2p'$, let
\[
 U_{\Delta,n}:=\lambda a_n\int_\Delta h(z)\,{:}\Phi(z)^n{:}\,dz ,\qquad
 U_\Delta:=\sum_{n=1}^{2p'}U_{\Delta,n} ,
\]
so that $U_{\Delta,n}$ ($n\ge1$) is the monomial of type (1.1.6) with $r=1$, one Wick factor
of degree $n$, and kernel $w_{\Delta,n}:=\lambda a_n\,h\mathbf 1_\Delta\in L^2(\R^2)$, while
$U_{\Delta,0}=\lambda a_0\int_\Delta h\,dz$ is a real number.  Expanding ${:}\mathscr
P(\Phi){:}=\sum_na_n{:}\Phi^n{:}$ in the definition of $\Uc$,
$\Uc_{h\mathbf 1_\Delta}=U_{\Delta,0}+U_\Delta$.
\end{auxdef}

\begin{auxlemma}[the constant $c_U$]\label{aux:notation-cU}
Let $h\in\CE$ and let $\Delta$ be a square.
\begin{enumerate}
\item[\textup{(i)}] For $1\le n\le2p'$, $U_{\Delta,n}$ is a localized monomial of type (1.1.6)
with localization square $\Delta$, local degree $N(\Delta)=n<\bar n=2p'+1$, and
\[
 \GJS{U_{\Delta,n}}=K_1K_2^{2n}n!\;\lambda|a_n|\;\|h\mathbf 1_\Delta\|_{L^2(\R^2)}
 \le\lambda|a_n|K_1K_2^{2n}n! .
\]
\item[\textup{(ii)}] With (1.1.9) applied to $U_\Delta=\sum_{n=1}^{2p'}U_{\Delta,n}$,
\[
 \GJS{U_\Delta}=\sum_{n=1}^{2p'}\GJS{U_{\Delta,n}}\le c_U:=\lambda\sum_{n=1}^{2p'}|a_n|\,K_1K_2^{2n}\,n! .
\]
The constant $c_U$ depends on $\lambda$, $a_1,\dots,a_{2p'}$, $K_1$, $K_2$ only; it is
uniform in $\Delta$ and in $h\in\CE$.
\item[\textup{(iii)}] \textup{(The term $n=0$.)}  For every $h'\in\CE$ and every $F\in
L^1(dq_{h'})$, $\Cov_{h'}(F,U_{\Delta,0})=0$.  The measures $dq_h$, $h\in\CE$, do not depend
on $a_0$.  As a coefficient of the interaction, $a_0$ enters the estimates of this note only
through the lower bound of Lemma~C.2 \cite[p.~40]{R}, $Z(h')\ge e^{-\lambda|a_0||\Lambda|}$, and
through $\Uc_{h'}$ inside $S_2$, i.e.\ through the base terms of Theorems~\ref{thm:D3}
and~\ref{thm:D4}.  As a coefficient of the observable $V(h')$ of Theorem~\ref{thm:D4} it also
enters $V_*$ (the term $n=0$ of Lemma~\ref{aux:notation-wick}(iv)) and the bound $\mathfrak Q$ on
$\GJS{Q_V}$ (the sum over $n\ge0$, Appendix~\ref{app:constants}).  It does not enter $c_U$.
\item[\textup{(iv)}] For every box $B$,
$\Uc_{h-h^B}=\sum_{\Delta\in\mathfrak S(h),\,\Delta\cap B=\varnothing}(U_{\Delta,0}+U_\Delta)$
in $L^2(\dmu)$; and for every $h'\in\CE$ and every $F$ such that $F$ and the finitely many
$U_{\Delta,n}$ involved lie in $L^2(dq_{h'})$,
\[
 \Cov_{h'}\bigl(F,\Uc_{h-h^B}\bigr)=\sum_{\Delta\in\mathfrak S(h),\ \Delta\cap B=\varnothing}\
 \sum_{n=1}^{2p'}\Cov_{h'}\bigl(F,U_{\Delta,n}\bigr) .
\]
\end{enumerate}
\end{auxlemma}

\begin{proof}
(i) In the notation of p.~593, $U_{\Delta,n}=\int dx_0\,Q(x_0)$ with
$Q(x_0)=\int d\vec x\,w_{\Delta,n}(x)\,{:}\Phi(x)^n{:}$, $r=1$, $n_1=n$, and the kernel
$w_{\Delta,n}=\lambda a_nh\mathbf 1_\Delta$ vanishes outside $\Delta$; so the monomial is
localized in $\Delta$ (Remark~\ref{aux:notation-gjslattice}(i)), with $N(\Delta)=n$ and
$N(\Delta')=0$ for every other square $\Delta'$, whose factors in (1.1.7) are
$K_2^0\,0!=1$.  The degree bound: $n\le2p'<2p'+1=\bar n$ (Assumption~A8(ii)
\cite[p.~10]{R}).  By (1.1.7),
$\GJS{U_{\Delta,n}}=K_1K_2^{2n}n!\,\|w_{\Delta,n}\|_2$, and
$\|w_{\Delta,n}\|_2=\lambda|a_n|\,(\int_\Delta h^2\,dz)^{1/2}\le\lambda|a_n|\,|\Delta|^{1/2}=
\lambda|a_n|$ because $0\le h\le1$ and $|\Delta|=1$.  (Test: for $h=\mathbf 1_\Delta$ equality
holds, $\GJS{U_{\Delta,n}}=\lambda|a_n|K_1K_2^{2n}n!$.)

(ii) The decomposition $U_\Delta=\sum_nU_{\Delta,n}$ writes the polynomial $U_\Delta$ as a sum
of the localized monomials of (i), and (1.1.9) sets its norm to the sum of theirs; sum the
bounds of (i).  The bound contains no $\Delta$ and no $h$.

(iii) $dq_{h'}$ is a probability measure (Assumption~A3 \cite[p.~8]{R}), so for a constant $c$,
$\int Fc\,dq_{h'}-\int F\,dq_{h'}\int c\,dq_{h'}=c\int F\,dq_{h'}-c\int F\,dq_{h'}=0$.  At the
regularized level ${:}\mathscr P(y){:}_c=a_0+{:}(\mathscr P-a_0)(y){:}_c$, since the Wick
polynomial of degree $0$ is the constant; hence $\Uc_{h,\kappa}=\lambda a_0\int h\,dz+
\Uc^{(0)}_{h,\kappa}$, where $\Uc^{(0)}$ is the interaction of $\mathscr P-a_0$.  Lemma~C.4,
which applies to $\mathscr P-a_0$ as well (same degree $2p'$, same leading coefficient),
passes this to the limit: $\Uc_h=\lambda a_0\int h\,dz+\Uc^{(0)}_h$ in $L^2(\dmu)$.  The
factor $e^{-\lambda a_0\int h}$ of $e^{-\Uc_h}$ then cancels against the same factor of
$Z(h)$.  The two places where $a_0$ survives as a coefficient of the interaction are the separate
bounds of the numerator $\|e^{-\Uc_{h^B}}\|_2\le S_2(B)$ and of the denominator
$Z(h^B)\ge e^{-\lambda|a_0||B|}$ in the base term; the ratio itself does not depend on $a_0$.  In
$V_*$ and $\mathfrak Q$ it is a coefficient of the observable $V(h')=\lambda\sum_{n\ge0}a_nW_n$,
not of the measure; $c_U$ sums over $n\ge1$ only.

(iv) Lemma~\ref{aux:notation-cutoffs}(iv) and Definition~\ref{aux:notation-UDelta} give the
first identity.  The covariance is bilinear on $L^2(dq_{h'})$, and the constants
$U_{\Delta,0}$ drop out by (iii).  The integrability of $F$ and of the $U_{\Delta,n}$ under
$dq_{h'}$ is part of the chain of Lemma~\ref{lem:D1} (\S\ref{sec:d1}).
\end{proof}

\subsection{The kernel $C$, the Wick isometry and $V_*$}\label{sec:notation-wick}

The time-zero measure $\dphi$ is ``Gaussian with covariance $C=(2\mu)^{-1}$,
$\mu=(-\partial_x^2+m_0^2)^{1/2}$'' (\S2.1 \cite[p.~6]{R}).  Fourier convention:
$\hat f(k)=\int e^{-ikx}f(x)\,dx$, so that $\ip fg=(2\pi)^{-1}\int\overline{\hat f}\hat g\,dk$
and $\mu$ acts as multiplication by $(k^2+m_0^2)^{1/2}$.

\begin{auxlemma}[the kernel of $(2\mu)^{-1}$]\label{aux:notation-kernel}
For $u\in\R\setminus\{0\}$ put
\[
 C(u):=\frac1{4\pi}\int_0^\infty t^{-1}\,e^{-tm_0^2-u^2/(4t)}\,dt ,\qquad
 b_0:=\log4+e^{-1}+m_0^{-2} .
\]
\begin{enumerate}
\item[\textup{(i)}] $0<C(u)<\infty$; $C$ is even, nonincreasing in $|u|$, and continuous on
$\R\setminus\{0\}$.
\item[\textup{(ii)}] For real $f,g\in L^1(\R)\cap L^2(\R)$, $\ip f{(2\mu)^{-1}g}=\iint
f(x)\,C(x-y)\,g(y)\,dx\,dy$, the double integral converging absolutely.  Thus $C$ is the
kernel of $(2\mu)^{-1}$, and $\E_{\phi_0}[\varphi_0(f)\varphi_0(g)]=\iint f(x)C(x-y)g(y)\,dx\,dy$.
\item[\textup{(iii)}] $C(u)\le(4\pi)^{-1}\bigl(b_0+2|\log|u||\bigr)$ for every $u\ne0$.
\item[\textup{(iv)}] For $n\ge0$ let $I_n:=\iint_{[-1,1]^2}C(x-y)^n\,dx\,dy=\int_{-2}^2(2-|u|)\,
C(u)^n\,du$.  Then $I_0=4$, and for $n\ge1$
\[
 I_n\le(4\pi)^{-n}\,2^{n+1}\bigl(2b_0^n+2^n(n!+(\log2)^n)\bigr)<\infty .
\]
\item[\textup{(v)}] \textup{(The Euclidean kernel.)}  For real $f,g\in L^1(\R^2)\cap L^2(\R^2)$,
$\ip f{(-\Delta_{\R^2}+m_0^2)^{-1}g}=\iint f(z)\,C(|z-z'|)\,g(z')\,dz\,dz'$, the double integral
converging absolutely.  So the covariance $G$ of $\dmu$ has the kernel $G(\zeta)=C(|\zeta|)$,
$\zeta\ne0$, which is positive, radial, nonincreasing in $|\zeta|$, continuous on
$\R^2\setminus\{0\}$, bounded by $(4\pi)^{-1}(b_0+2|\log|\zeta||)$, and satisfies
$\int_{\R^2}G=m_0^{-2}$.
\end{enumerate}
\end{auxlemma}

\begin{proof}
(ii) For $A>0$, the substitution $s=tA$ gives $\int_0^\infty t^{-1/2}e^{-tA}dt=\Gamma(\frac12)
A^{-1/2}=\sqrt\pi A^{-1/2}$; with $A=k^2+m_0^2$,
\[
 \bigl(2(k^2+m_0^2)^{1/2}\bigr)^{-1}=(2\sqrt\pi)^{-1}\int_0^\infty t^{-1/2}e^{-tm_0^2}e^{-tk^2}dt .
\]
The heat kernel $p_t(u):=(4\pi t)^{-1/2}e^{-u^2/(4t)}$ has $\hat p_t(k)=e^{-tk^2}$ (complete the
square), so for $f\in L^1\cap L^2$, by Plancherel and the convolution theorem,
\[
 (2\pi)^{-1}\int|\hat f(k)|^2e^{-tk^2}dk=\iint f(x)\,p_t(x-y)\,f(y)\,dx\,dy .
\]
Insert the $t$-integral into
$\ip f{(2\mu)^{-1}f}=(2\pi)^{-1}\int|\hat f(k)|^2\bigl(2(k^2+m_0^2)^{1/2}\bigr)^{-1}dk$ and
exchange $k$ and $t$ (Tonelli; the integrand is $\ge0$):
\begin{align*}
 \ip f{(2\mu)^{-1}f}
 &=(2\sqrt\pi)^{-1}\int_0^\infty t^{-1/2}e^{-tm_0^2}\,
   (2\pi)^{-1}\!\int|\hat f(k)|^2e^{-tk^2}dk\,dt\\
 &=(2\sqrt\pi)^{-1}\int_0^\infty t^{-1/2}e^{-tm_0^2}\iint f(x)\,p_t(x-y)\,f(y)\,dx\,dy\,dt .
\end{align*}
For $f\ge0$ the integrand in $(t,x,y)$ is nonnegative and Tonelli exchanges again; as
\[
 (2\sqrt\pi)^{-1}\,t^{-1/2}\,(4\pi t)^{-1/2}=(4\pi t)^{-1},
\]
the result is
\[
 \ip f{(2\mu)^{-1}f}=\iint f(x)\,C(x-y)\,f(y)\,dx\,dy ,
\]
which is therefore at most $\|f\|_2^2/(2m_0)<\infty$, since $\|(2\mu)^{-1}\|\le(2m_0)^{-1}$.  For real $f,g$ apply this to
$F:=|f|+|g|\ge0$:
\[
 \iint F(x)\,C(x-y)\,F(y)\,dx\,dy<\infty
\]
dominates every double integral of $|f|,|g|$ against $C$, and hence (by the same Tonelli
step) against $t^{-1/2}e^{-tm_0^2}p_t$ in $(t,x,y)$; so Fubini applies to the display for
signed $f$, giving $\ip f{(2\mu)^{-1}f}=\iint fCf$; polarization (both sides are
symmetric bilinear forms, $C$ being even) gives the statement for $f,g$.

(i) The integrand is positive, so $C(u)>0$; finiteness is (iii).  $C$ depends on $u^2$
only, and $e^{-u^2/(4t)}$ is nonincreasing in $|u|$.  For $|u|\ge u_0>0$ the integrand is
dominated by $t^{-1}e^{-tm_0^2-u_0^2/(4t)}$, whose integral is $4\pi C(u_0)<\infty$; dominated
convergence gives continuity on $\{|u|>u_0\}$ for every $u_0>0$.

(iii) Drop $e^{-tm_0^2}\le1$ where convenient.  The substitution $s=u^2/(4t)$ gives,
for $\tau>0$,
\[
\begin{gathered}
 \int_0^{\tau}t^{-1}e^{-u^2/(4t)}dt=\int_{u^2/(4\tau)}^\infty s^{-1}e^{-s}ds,\\
 \int_{1/4}^\infty s^{-1}e^{-s}ds\le\int_{1/4}^1\frac{ds}s+\int_1^\infty e^{-s}ds=\log4+e^{-1};
\end{gathered}
\]
also $\int_1^\infty t^{-1}e^{-tm_0^2}dt\le\int_1^\infty e^{-tm_0^2}dt\le m_0^{-2}$.  If
$0<|u|\le1$, split $\int_0^\infty=\int_0^{u^2}+\int_{u^2}^1+\int_1^\infty$: the three pieces
are at most $\log4+e^{-1}$ ($\tau=u^2$), $\int_{u^2}^1t^{-1}dt=2|\log|u||$, and $m_0^{-2}$.  If
$|u|>1$, split $\int_0^1+\int_1^\infty$: at most $\int_{u^2/4}^\infty s^{-1}e^{-s}ds\le
\log4+e^{-1}$ (as $u^2/4>\frac14$) and $m_0^{-2}$.  In both cases $4\pi C(u)\le
b_0+2|\log|u||$.

(iv) For $\varphi\ge0$ measurable, $\iint_{[-1,1]^2}\varphi(x-y)\,dx\,dy=\int_{-2}^2(2-|u|)
\varphi(u)\,du$, since $\{x\in[-1,1]:x-u\in[-1,1]\}$ has length $2-|u|$ for $|u|\le2$.  For
$n=0$ this is $4$.  For $n\ge1$, by (iii) and $(\alpha+\beta)^n\le2^{n-1}(\alpha^n+\beta^n)$
($\alpha,\beta\ge0$, convexity), and $2-|u|\le2$,
\[
 I_n\le4\int_0^2C(u)^n\,du\le4(4\pi)^{-n}2^{n-1}\int_0^2\bigl(b_0^n+2^n|\log u|^n\bigr)du ,
\]
and $\int_0^2|\log u|^ndu\le\int_0^1|\log u|^ndu+(\log2)^n=n!+(\log2)^n$ (substitute $u=e^{-s}$
in the first integral).  Collect: $4\cdot2^{n-1}(2b_0^n+2^n(n!+(\log2)^n))(4\pi)^{-n}$.

(v) The computation of (ii) in two dimensions.  With $\hat f(p):=\int e^{-ip\cdot z}f(z)\,dz$,
$\ip f{(-\Delta+m_0^2)^{-1}f}=(2\pi)^{-2}\int|\hat f(p)|^2(|p|^2+m_0^2)^{-1}dp$, and
$(|p|^2+m_0^2)^{-1}=\int_0^\infty e^{-tm_0^2}e^{-t|p|^2}dt$.  The heat kernel
$P_t(\zeta):=(4\pi t)^{-1}e^{-|\zeta|^2/(4t)}$ has $\hat P_t(p)=e^{-t|p|^2}$, so
$(2\pi)^{-2}\int|\hat f|^2e^{-t|p|^2}dp=\iint f(z)P_t(z-z')f(z')\,dz\,dz'$.  Tonelli for $f\ge0$,
then $|f|+|g|$ for absolute convergence (with $\|(-\Delta+m_0^2)^{-1}\|\le m_0^{-2}$), Fubini and
polarization, exactly as in (ii), give the kernel
$\int_0^\infty e^{-tm_0^2}(4\pi t)^{-1}e^{-|\zeta|^2/(4t)}dt=C(|\zeta|)$.  The properties follow
from (i) and (iii), and by Tonelli
$\int_{\R^2}G=(4\pi)^{-1}\int_0^\infty t^{-1}e^{-tm_0^2}(4\pi t)\,dt=m_0^{-2}$.
\end{proof}

\emph{Chaoses.}  Let $\mathcal H_{\le n}$ be the closure in $L^2(\dphi)$ of the polynomials
of degree $\le n$ in the variables $\varphi_0(f)$, and $\mathcal H_n:=\mathcal
H_{\le n}\ominus\mathcal H_{\le n-1}$ (the $n$-th Wick chaos); the same definitions are made in
$L^2(\dmu)$ with the variables $\Phi(f)$.  These are definitions and need no import.  The fact
that the second quantization $\Gamma(e^{-t}\mathbf 1)$ acts on the $n$-th chaos as
multiplication by $e^{-nt}$, which the record's proof of Theorem~C.5 uses
(``$\Gamma(e^{-t})Y_r=e^{-rt}Y_r$'', \cite[p.~41]{R}) and does not import (V3 audit, M-3;
repair R8), is used here only in Lemma~\ref{aux:notation-hyper}.  R8 imports that action on
Simon's $\Gamma_n$, the closed span of the Wick monomials ${:}\Phi(f_1)\cdots\Phi(f_n){:}$; every
object of degree $n$ to which this note applies it is an $L^2$-limit of integrals of Wick
monomials ${:}\Phi(g)^n{:}$ with all $n$ test functions equal (Lemma~\ref{aux:notation-hyper}(i)),
so it lies in $\Gamma_n$ directly.  For a centred
Gaussian variable $X$, ${:}X^n{:}:=c^{n/2}\operatorname{He}_n(X/\sqrt c)$, $c=\E[X^2]$, is
the coefficient of $s^n/n!$ in ${:}e^{sX}{:}:=e^{sX-s^2c/2}$ (as in the proof of
Proposition~3.2 \cite[p.~16]{R}).  Fix $\rho\in C_c^\infty(\R)$, $\rho\ge0$, even, $\supp
\rho\subseteq[-1,1]$, $\int\rho=1$, put $\rho_\kappa(x):=\kappa\rho(\kappa x)$ and
$\varphi_\kappa(x):=\varphi_0(\rho_\kappa(x-\cdot))$ for $\kappa\ge1$.

\begin{auxlemma}[Wick isometry; the constant $V_*$]\label{aux:notation-wick}
Let $h'$ be real and measurable with $\|h'\|_\infty\le1$ and $\supp h'\subseteq[a,a+2]$,
$a\in\R$, and let $0\le n,m\le2p'$.
\begin{enumerate}
\item[\textup{(i)}] For $\kappa,\kappa'\ge1$:
$\E_{\phi_0}\bigl[{:}\varphi_\kappa(x)^n{:}\,{:}\varphi_{\kappa'}(y)^m{:}\bigr]=\delta_{nm}\,n!\,
C_{\kappa,\kappa'}(x-y)^n$, where $C_{\kappa,\kappa'}(u):=\int\nu(v)C(u-v)\,dv>0$ and
$\nu:=\rho_\kappa*\rho_{\kappa'}$.
\item[\textup{(ii)}] If $\min\{\kappa,\kappa'\}\ge6$, then $C_{\kappa,\kappa'}(u)\le
c_\rho\,(b_0+2+2|\log|u||)$ for all $u\ne0$, with $c_\rho:=3\|\rho\|_\infty/\pi$; and
$C_{\kappa,\kappa'}(u)\to C(u)$ for every $u\ne0$ as $\min\{\kappa,\kappa'\}\to\infty$.
\item[\textup{(iii)}] $W_{n,\kappa}:=\int h'(x)\,{:}\varphi_\kappa(x)^n{:}\,dx$ lies in
$\mathcal H_n$, and converges in $L^2(\dphi)$, as $\kappa\to\infty$, to some $W_n\in\mathcal
H_n$ with $\ip{W_n}{W_m}=\delta_{nm}\,n!\iint h'(x)h'(y)C(x-y)^n\,dx\,dy$.
\item[\textup{(iv)}] $V(h'):=\lambda\sum_{n=0}^{2p'}a_nW_n$ satisfies
\[
 \|V(h')\|^2_{L^2(\dphi)}=\lambda^2\sum_{n=0}^{2p'}a_n^2\,n!\iint h'(x)h'(y)\,C(x-y)^n\,dx\,dy
 \le\overline V_*:=\lambda^2\sum_{n=0}^{2p'}a_n^2\,n!\,I_n .
\]
Hence $V_*:=\sup_{h'}\|V(h')\|^2_{L^2(\dphi)}\le\overline V_*<\infty$, the supremum over all
$h'$ as above and all $a\in\R$.
\end{enumerate}
\end{auxlemma}

\begin{proof}
(i) $X:=\varphi_\kappa(x)$ and $Y:=\varphi_{\kappa'}(y)$ are jointly centred Gaussian, with
$\E[XY]=\ip{\rho_\kappa(x-\cdot)}{(2\mu)^{-1}\rho_{\kappa'}(y-\cdot)}$; by
Lemma~\ref{aux:notation-kernel}(ii) and the substitutions $x'=x-\alpha$, $y'=y-\beta$ this is
$\iint\rho_\kappa(\alpha)\rho_{\kappa'}(\beta)C(x-y-\alpha+\beta)\,d\alpha\,d\beta=
C_{\kappa,\kappa'}(x-y)$, with $\nu$ the density of $\alpha-\beta$, which is
$\rho_\kappa*\rho_{\kappa'}$ because $\rho_{\kappa'}$ is even; $C>0$ gives
$C_{\kappa,\kappa'}>0$.  The Gaussian moment generating function gives
$\E[{:}e^{sX}{:}\,{:}e^{tY}{:}]=e^{-(s^2\E X^2+t^2\E Y^2)/2}\E[e^{sX+tY}]=e^{st\,\E[XY]}$.  Both
sides are entire in $(s,t)$; the left side may be differentiated under $\E$ (for $|s|,|t|\le1$
the derivatives are dominated by polynomials in $|X|,|Y|$ times $e^{|X|+|Y|}$, which are
integrable), and its coefficient of $s^nt^m/(n!\,m!)$ is $\E[{:}X^n{:}{:}Y^m{:}]$; on the right
it is $\delta_{nm}n!\,\E[XY]^n$.

(ii) $\nu\ge0$, $\int\nu=1$, $\supp\nu\subseteq[-\varepsilon,\varepsilon]$ with
$\varepsilon:=\kappa^{-1}+\kappa'^{-1}\le2/K$, $K:=\min\{\kappa,\kappa'\}$, and
$\|\nu\|_\infty\le\min\{\|\rho_\kappa\|_\infty,\|\rho_{\kappa'}\|_\infty\}=K\|\rho\|_\infty$.
\emph{Case $|u|\ge2\varepsilon$.}  For $|v|\le\varepsilon$, $|u-v|\ge|u|/2$, so by
Lemma~\ref{aux:notation-kernel}(i),(iii) $C(u-v)\le C(u/2)\le(4\pi)^{-1}(b_0+2\log2+
2|\log|u||)$; average over $\nu$.  \emph{Case $0<|u|<2\varepsilon$.}  Then
$C_{\kappa,\kappa'}(u)\le K\|\rho\|_\infty\int_{|w|\le3\varepsilon}C(w)\,dw$, and
$3\varepsilon\le6/K\le1$.  For $0<r\le1$, Lemma~\ref{aux:notation-kernel}(iii) and
$\int_0^r|\log w|\,dw=r(1+|\log r|)$ give $\int_{|w|\le r}C\le(4\pi)^{-1}2r(b_0+2+2|\log
r|)$, which is increasing in $r\in(0,1]$ (derivative $(4\pi)^{-1}2(b_0-2\log r)>0$); at
$r=6/K$ it is $(4\pi)^{-1}(12/K)(b_0+2+2\log(K/6))$.  Since $|u|<4/K$,
$\log(K/6)<\log(2/3)-\log|u|\le|\log|u||$.  So $C_{\kappa,\kappa'}(u)\le(3\|\rho\|_\infty/\pi)
(b_0+2+2|\log|u||)$.  As $\int\rho=1$ on $[-1,1]$, $\|\rho\|_\infty\ge\frac12$, so
$c_\rho\ge(4\pi)^{-1}$ and $2\log2<2$ make $c_\rho(b_0+2+2|\log|u||)$ a bound in both cases.
Convergence: $C_{\kappa,\kappa'}(u)-C(u)=\int\nu(v)(C(u-v)-C(u))dv\to0$ as $\varepsilon\to0$, by
the continuity of $C$ at $u\ne0$.

(iii) $(x,q)\mapsto{:}\varphi_\kappa(x)^n{:}(q)$ is continuous in $x$ and measurable in $q$,
and its moments are bounded uniformly in $x$; Fubini therefore applies below.  For a
polynomial $P$ of degree $<n$ in variables $\varphi_0(g_i)$, the coefficient of $s^n$ in
$\E[{:}e^{sX}{:}\,e^{\sum_it_i\varphi_0(g_i)}]=e^{s\sum_it_i\E[X\varphi_0(g_i)]}\,
\E[e^{\sum_it_i\varphi_0(g_i)}]$ is a power series in the $t_i$ all of whose terms have
degree $\ge n$, so $\E[{:}X^n{:}P]=0$; with Fubini, $W_{n,\kappa}\perp\mathcal H_{\le n-1}$.
For $Q\perp\mathcal H_{\le n}$, $\E[W_{n,\kappa}Q]=\int h'(x)\E[{:}\varphi_\kappa(x)^n{:}Q]dx=0$,
so $W_{n,\kappa}\in\mathcal H_{\le n}$; hence $W_{n,\kappa}\in\mathcal H_n$.  By (i) and
Fubini,
\[
 \|W_{n,\kappa}-W_{n,\kappa'}\|_2^2=n!\iint h'(x)h'(y)\bigl[C_{\kappa,\kappa}^n-2C_{\kappa,\kappa'}^n
 +C_{\kappa',\kappa'}^n\bigr](x-y)\,dx\,dy .
\]
By (ii), for $\min\{\kappa,\kappa'\}\ge6$ each of the three integrands is dominated by
\[
 \mathbf 1_{[a,a+2]}(x)\,\mathbf 1_{[a,a+2]}(y)\;c_\rho^n\bigl(b_0+2+2|\log|x-y||\bigr)^n ,
\]
which is integrable (as in Lemma~\ref{aux:notation-kernel}(iv)), and each converges pointwise off the
diagonal to $h'(x)h'(y)C(x-y)^n$.  By dominated convergence the three integrals have the same
limit, so $(W_{n,\kappa})$ is Cauchy; its limit $W_n$ lies in the closed subspace $\mathcal
H_n$, and $\ip{W_n}{W_m}=\lim_\kappa\ip{W_{n,\kappa}}{W_{m,\kappa}}$ is given by (i) and the
same dominated convergence.

(iv) Expand the square with (iii).  Since $|h'|\le\mathbf 1_{[a,a+2]}$ and $C^n\ge0$
(Lemma~\ref{aux:notation-kernel}(i): here the positivity of $C$ is used, for odd $n$),
$|\iint h'(x)h'(y)C(x-y)^n|\le\iint_{[a,a+2]^2}C(x-y)^n=I_n$, the last step by the
translation $x\mapsto x-a-1$, $y\mapsto y-a-1$.  The bound does not depend on $h'$ or $a$.
\end{proof}

\begin{gapnote}
The record's $V(h)$ (\S2.1 \cite[p.~6]{R}: ``$V(h):=\lambda\int
h(x)\,{:}\mathscr P(\varphi_0(x)){:}_C\,dx$, a self-adjoint multiplication operator in
$Q$-space'', imported with Assumption~A1) is identified here with the limit
$\lambda\sum_na_nW_n$ of Lemma~\ref{aux:notation-wick}.  Missing: the record states no
construction of the Wick powers in $V(h)$.  Its Appendix~C makes the corresponding
identification for $\Uc_h$ (``so that $\mathcal U_{h,\kappa}\to\mathcal U_h$'',
\cite[pp.~39--40]{R}), and the V3 audit's repair R8 would import the construction for $\Uc_h$
from Simon's book (chapters~I and~V, by the errata header of the V3 report); the same import for
the time-zero field (the V3 report names Simon's Theorem~V.2(b); not read for this note) would
close this gap.  \emph{Where it is used:} only where the record's $Q$-space operator $V(h)$ is
replaced by the constructed limit, i.e.\ in the record's Theorem~5.1, Step~1 (``$Q_V$ is exactly
the multiplication function $V(h)^2$''), and in the transfer of Theorem~\ref{thm:D4} to the
record's Theorem~5.1 (\S\ref{sec:cons}).  Lemmas~\ref{aux:notation-wick}--\ref{aux:notation-sharpid},
the constants $V_*$, $9^{p'}$ and $M_B$ of Appendix~\ref{app:constants}, and the Euclidean bound
of Theorem~\ref{thm:D4} concern the constructed objects (Remark~\ref{aux:notation-sharpidrem}(b))
and do not need it.
\end{gapnote}

\begin{auxlemma}[the comparison $\|Y\|_4\le3^{p'}\|Y\|_2$]\label{aux:notation-hyper}
Let $h'$ be as in Lemma~\ref{aux:notation-wick}, and let
\[
 Y:=V(h')(0)=\lambda\sum_{n=0}^{2p'}a_n\int h'(x)\,{:}\Phi(0,x)^n{:}\,dx
\]
be the sharp-time observable at time $0$, as an element of $L^2(\dmu)$.
\begin{enumerate}
\item[\textup{(i)}] $Y=\sum_{r=0}^{2p'}Y_r$, where each $Y_r$ lies in the $r$-th Wick chaos
of $L^2(\dmu)$; moreover $\|Y\|_{L^2(\dmu)}=\|V(h')\|_{L^2(\dphi)}$.
\item[\textup{(ii)}] $\|Y\|_{L^4(\dmu)}\le3^{p'}\|Y\|_{L^2(\dmu)}$.  Consequently
\[
 \|Y^2\|_{L^2(\dmu)}=\|Y\|_{L^4(\dmu)}^2\le9^{p'}\|V(h')\|_{L^2(\dphi)}^2\le9^{p'}V_* ,
\]
and the same bound holds for $\|V(h')^2\|_{L^2(\dphi)}$, which equals $\|Y^2\|_{L^2(\dmu)}$.
\end{enumerate}
\end{auxlemma}

\begin{proof}
(i) Remark~2.5 of the record \cite[p.~14]{R}: ``The sharp-time restriction of the Euclidean
covariance equals $C$'' and ``sharp-time $d\mu_0$-Wick ordering coincides, constant by constant,
with the $C$-Wick ordering used in $V(h)$''.  GJS write the sharp-time field as
``$\Phi(\tau,f)\equiv\Phi(\delta_\tau\otimes f)$'' \cite[p.~589]{GJS} and print no construction;
we read $\Phi(0,f)$ as the extension of the linear map $f\mapsto\Phi(f)$ by continuity in the
covariance norm, which is finite at $\delta_0\otimes f$ because
$\ip{\delta_0\otimes f}{G\,\delta_0\otimes f}=\ip f{Cf}$ (Remark~2.5).  So $\Phi(0,f)$ is the
$L^2(\dmu)$-limit of $\Phi(f_\varepsilon)$ along test functions $f_\varepsilon\to\delta_0\otimes
f$ in that norm, a centred Gaussian variable in the closed first chaos of $L^2(\dmu)$, and the
family $(\Phi(0,f))_f$ has the covariance $C$ of $(\varphi_0(f))_f$ under $\dphi$.  So $J_0:F(\varphi_0)\mapsto F(\Phi(0,\cdot))$
preserves the joint laws of finitely many field variables and is an isometry
$L^2(\dphi)\to L^2(\dmu)$; it maps ${:}\varphi_\kappa(x)^n{:}$ to the Wick power, at its own
variance (the same $c$), of the Gaussian variable $\Phi(0,\rho_\kappa(x-\cdot))$.  The
generating-function argument of Lemma~\ref{aux:notation-wick}(iii) holds in any Gaussian space:
this Wick power is orthogonal to every polynomial of degree $<n$ in the variables $\Phi(f)$, and
it is a polynomial of degree $n$ in a first-chaos variable, hence (Gaussian $L^2$-limits converge
in every $L^q$) in the closure of the polynomials of degree $\le n$.  So $J_0W_{n,\kappa}$ lies in
the $n$-th Wick chaos of $L^2(\dmu)$, and so does its limit $J_0W_n$.  Moreover $J_0W_{n,\kappa}$
is the integral over $x$ of the Wick monomials ${:}\Phi(0,\rho_\kappa(x-\cdot))^n{:}$, each the
$L^2$-limit of ${:}\Phi(f_\varepsilon)^n{:}$ for test functions $f_\varepsilon\to\delta_0\otimes
\rho_\kappa(x-\cdot)$ in the covariance norm; so $J_0W_n$ also lies in Simon's $\Gamma_n$, the
closed span of the Wick monomials, on which repair R8 imports the action of $\Gamma$.  Put
$Y_r:=\lambda a_rJ_0W_r$; then $Y=J_0V(h')=\sum_rY_r$, and $\|Y\|_{L^2(\dmu)}=\|V(h')\|_{L^2(\dphi)}$.

(ii) The record proves the comparison inside the proof of Theorem~C.5 \cite[p.~41]{R}
(\texttt{part\_i/main.tex}, ll.~3021--3031), under $\dmu$: ``$Y:=\mathcal U_h-\mathcal
U_{h,\kappa}$ is a sum $Y=\sum_{r=0}^{2p'}Y_r$ of Wick chaoses, so Nelson's hypercontractive
moment comparison applies in the form $\|Y\|_{L^q}\le(q-1)^{p'}\|Y\|_{L^2}$ $(q\ge2)$.  (Proof
from the abstract form of Assumption~A4 [Thm.~I.17]: the second-quantized contraction
$\Gamma(e^{-t})$ on $L^2(d\mu_0)$ satisfies $\|\Gamma(e^{-t})X\|_q\le\|X\|_2$ for $q=1+e^{2t}$,
and $\Gamma(e^{-t})Y_r=e^{-rt}Y_r$.  Write $Y=\Gamma(e^{-t})Z$ with $Z=\sum_re^{rt}Y_r$; by
orthogonality of the chaoses, $\|Z\|_2\le e^{2p't}\|Y\|_2$, whence $\|Y\|_q\le\|Z\|_2\le
e^{2p't}\|Y\|_2=(q-1)^{p'}\|Y\|_2$.)''  The argument uses of its $Y$ only that it is a sum of
Wick chaoses of degree $r\le2p'$ in $L^2(\dmu)$, which our $Y$ is by (i); so it applies verbatim,
conditional on the repair R8 of the V3 audit, which imports the action
$\Gamma(e^{-t})Y_r=e^{-rt}Y_r$ (V3 report, \S5, items C5.3--C5.4; part (i) above does not need
it).  With $q=4$: $t=\frac12\log3$, the
contraction $e^{-t}\mathbf 1$ has $\|e^{-t}\mathbf 1\|^2=\frac13=(2-1)/(4-1)$ as the abstract
form of Assumption~A4 \cite[p.~9]{R} requires, and $e^{2p't}=3^{p'}$:
\[
 \|Y\|_4=\|\Gamma(e^{-t})Z\|_4\le\|Z\|_2=\Bigl(\sum_re^{2rt}\|Y_r\|_2^2\Bigr)^{1/2}\le
 3^{p'}\|Y\|_2 .
\]
(Test: $p'=1$, $Y$ centred Gaussian with variance $\sigma^2$: $\|Y\|_4=3^{1/4}\sigma\le3\sigma$.)
Finally $\|Y^2\|_2=(\E Y^4)^{1/2}=\|Y\|_4^2$, Lemma~\ref{aux:notation-wick}(iv) bounds
$\|V(h')\|_{L^2(\dphi)}^2$ by $V_*$, and $J_0$ maps $V(h')^2$ to $Y^2$ isometrically.
\end{proof}

By Lemma~\ref{aux:notation-sharpid} below and Remark~\ref{aux:notation-sharpidrem}(b), $Y$ is also
the sharp-time observable built with the record's regularization, so the bound applies to $Q_V$ in
the form in which Lemma~\ref{lem:D2} uses it.

\begin{auxlemma}[the two regularizations of a sharp-time Wick power agree]\label{aux:notation-sharpid}
Let $t\in\R$, let $f$ be real, bounded and measurable with compact support in $\R$, and let
$0\le n\le2p'$.  Let $\delta_\kappa$ be the record's mollifier (Appendix~C, notation paragraph
\cite[pp.~39--40]{R}: $\delta\in C_c^\infty(\R^2)$, $\delta\ge0$, $\delta(z)=\delta(-z)$,
$\supp\delta\subseteq\{|z|\le1\}$, $\int\delta=1$, $\delta_\kappa(z)=\kappa^2\delta(\kappa z)$),
let $\Phi_\kappa(\zeta)=\Phi(\delta_\kappa(\zeta-\cdot))$, with variance $c_\kappa$, and
\[
 X_\kappa:=\int f(x)\,{:}\Phi_\kappa(t,x)^n{:}_{c_\kappa}\,dx\qquad(\kappa\ge1),
\]
the integral taken pointwise in $q\in\mathcal S'(\R^2)$.  Let $W_n(f)$ be the limit of
Lemma~\ref{aux:notation-wick}(iii) with $f$ in place of $h'$ (its proof uses only that the kernel
is bounded with compact support), and let $J_t$ be the sharp-time map at time $t$,
$F(\varphi_0)\mapsto F(\Phi(t,\cdot))$; $J_0$ is the map of Lemma~\ref{aux:notation-hyper}.
Then $X_\kappa\to J_tW_n(f)$ in $L^2(\dmu)$ as $\kappa\to\infty$.
\end{auxlemma}

\begin{proof}
\emph{The map $J_t$.}  $J_t$ is defined on cylinder functions,
\[
 J_t:\ F(\varphi_0(g_1),\dots,\varphi_0(g_k))\longmapsto F(\Phi(t,g_1),\dots,\Phi(t,g_k)),
\]
and preserves their laws: by Lemma~\ref{aux:notation-transl}(i) with $a=(t,0)$, $\Phi(g)\circ T_a=\Phi(\tau_ag)$
and $(T_a)_*\mu_0=\mu_0$; the extension by continuity in the covariance norm, which $\tau_a$
preserves, gives $\Phi(0,g)\circ T_a=\Phi(t,g)$, so $(\Phi(t,g))_g$ has the law of $(\Phi(0,g))_g$,
which is that of $(\varphi_0(g))_g$ (proof of Lemma~\ref{aux:notation-hyper}(i)).  Hence $J_t$
extends to an isometry of every $L^p(\dphi)$ into $L^p(\dmu)$, is multiplicative, and maps a Wick
power at its own variance to the Wick power at its own variance of the image variable.

\emph{Covariances.}  By Lemma~\ref{aux:notation-kernel}(v) and the evenness of $\delta$,
$\E[\Phi_\kappa(t,x)\Phi_\kappa(t,y)]=(\delta_\kappa*\delta_\kappa*G)((0,x-y))$.  The variable
$\Phi(t,g)$, $g:=\rho_{\kappa'}(y-\cdot)$, is the $L^2(\dmu)$-limit of $\Phi(\eta_\varepsilon\otimes
g)$ for $\eta_\varepsilon\in C_c^\infty(\R)$, $\eta_\varepsilon\ge0$, $\int\eta_\varepsilon=1$,
$\supp\eta_\varepsilon\subseteq[t-\varepsilon,t+\varepsilon]$ (the reading of
Lemma~\ref{aux:notation-hyper}(i)); since $\delta_\kappa*G$ is continuous,
$\E[\Phi_\kappa(t,x)\Phi(\eta_\varepsilon\otimes g)]=\int(\delta_\kappa*G)((t,x)-z')\,
\eta_\varepsilon(z'^0)\,g(z'^1)\,dz'\to\int(\delta_\kappa*G)((0,x-y'))\,g(y')\,dy'$, so
\[
 \E\bigl[\Phi_\kappa(t,x)\,\Phi(t,\rho_{\kappa'}(y-\cdot))\bigr]
 =\int(\delta_\kappa*G)((0,x-y'))\,\rho_{\kappa'}(y-y')\,dy'=:\Gamma_{\kappa,\kappa'}(x,y) .
\]

\emph{Inner products.}  Put $Y_{\kappa'}:=J_tW_{n,\kappa'}(f)=\int f(y)\,{:}\Phi(t,\rho_{\kappa'}
(y-\cdot))^n{:}\,dy$ and $Y:=J_tW_n(f)$, so that $Y_{\kappa'}\to Y$ in $L^2(\dmu)$
(Lemma~\ref{aux:notation-wick}(iii) and the isometry).  The generating-function identity in the
proof of Lemma~\ref{aux:notation-wick}(i) holds for any two jointly centred Gaussian variables, and
Fubini applies as there: the moments are bounded uniformly in $x,y$, and $|f|\le\|f\|_\infty
\mathbf 1_K$ with $K$ bounded.  Hence
\begin{gather*}
 \E[X_\kappa^2]=n!\iint f(x)f(y)\,(\delta_\kappa*\delta_\kappa*G)((0,x-y))^n\,dx\,dy,\\
 \E[X_\kappa Y_{\kappa'}]=n!\iint f(x)f(y)\,\Gamma_{\kappa,\kappa'}(x,y)^n\,dx\,dy ,
\end{gather*}
and $\E[Y^2]=I:=n!\iint f(x)f(y)\,C(x-y)^n\,dx\,dy$ by Lemma~\ref{aux:notation-wick}(iii).

\emph{A majorant.}  Let $\kappa\ge6$, and let $\nu$ be a probability density on $\R^2$ with
$\supp\nu\subseteq\{|v|\le\varepsilon\}$, $\varepsilon\le2/\kappa$, and
$\|\nu\|_\infty\le\kappa^2\|\delta\|_\infty$.  Then, for $u\ne0$,
\[
 (\nu*G)((0,u))\le9\|\delta\|_\infty\bigl(b_0+2+2|\log|u||\bigr),
\]
and $(\nu*G)((0,u))\to G((0,u))=C(u)$ as $\varepsilon\to0$.  Indeed, if $|u|\ge2\varepsilon$, then
$|(0,u)-v|\ge|u|/2$ on $\supp\nu$, and by Lemma~\ref{aux:notation-kernel}(v), (i) and (iii),
$G((0,u)-v)\le C(|u|/2)\le(4\pi)^{-1}(b_0+2\log2+2|\log|u||)$.  If $0<|u|<2\varepsilon$, then
$|(0,u)-v|<3\varepsilon\le6/\kappa\le1$ on $\supp\nu$, so, in polar coordinates and by (iii),
\begin{align*}
 (\nu*G)((0,u))&\le\kappa^2\|\delta\|_\infty\,2\pi\int_0^{6/\kappa}rC(r)\,dr
 \le\kappa^2\|\delta\|_\infty\,\frac{2\pi}{4\pi}\,\frac{(6/\kappa)^2}{2}
 \Bigl(b_0+1+2\log\frac\kappa6\Bigr)\\
 &=9\|\delta\|_\infty\Bigl(b_0+1+2\log\frac\kappa6\Bigr),
\end{align*}
using $\int_0^{r_0}r(b_0+2|\log r|)\,dr=\frac{r_0^2}2(b_0+1+2|\log r_0|)$, which increases in
$r_0\in(0,1]$ (derivative $r_0(b_0-2\log r_0)>0$); and $\log(\kappa/6)<|\log|u||$ because
$|u|<4/\kappa$.  Both cases are at most $9\|\delta\|_\infty(b_0+2+2|\log|u||)$, since
$\|\delta\|_\infty\ge1/\pi$ ($\delta$ is a probability density on the unit disc) and $2\log2<2$.
The convergence holds because $G$ is continuous at $(0,u)\ne0$.  The majorant applies to
$\nu=\delta_\kappa$ ($\varepsilon=1/\kappa$) and to $\nu=\delta_\kappa*\delta_\kappa$
($\varepsilon=2/\kappa$, $\|\delta_\kappa*\delta_\kappa\|_\infty\le\|\delta_\kappa\|_\infty=
\kappa^2\|\delta\|_\infty$).

\emph{Limits.}  (1) For fixed $\kappa$, $\delta_\kappa*G$ is continuous and bounded by
$\|\delta_\kappa\|_\infty\int G=\kappa^2\|\delta\|_\infty m_0^{-2}$
(Lemma~\ref{aux:notation-kernel}(v)), so $\Gamma_{\kappa,\kappa'}(x,y)\to(\delta_\kappa*G)((0,x-y))$
as $\kappa'\to\infty$, boundedly; by dominated convergence and $Y_{\kappa'}\to Y$,
\[
 \E[X_\kappa Y]=n!\iint f(x)f(y)\,(\delta_\kappa*G)((0,x-y))^n\,dx\,dy .
\]
(2) As $\kappa\to\infty$, the integrands of $\E[X_\kappa^2]$ and of $\E[X_\kappa Y]$ converge to
$f(x)f(y)C(x-y)^n$ off the diagonal and are dominated by
$\|f\|_\infty^2\mathbf 1_K(x)\mathbf 1_K(y)(9\|\delta\|_\infty)^n(b_0+2+2|\log|x-y||)^n$, which is
integrable (as in Lemma~\ref{aux:notation-kernel}(iv)).  So $\E[X_\kappa^2]\to I$ and
$\E[X_\kappa Y]\to I$, and
\[
 \|X_\kappa-Y\|_{L^2(\dmu)}^2=\E[X_\kappa^2]-2\,\E[X_\kappa Y]+\E[Y^2]\to I-2I+I=0 .\qedhere
\]
\end{proof}

\begin{auxrem}[what Lemma~\ref{aux:notation-sharpid} identifies]\label{aux:notation-sharpidrem}
(a) The localized monomial of type (1.1.5) with one factor of degree $n\le2p'$ and a bounded
kernel $f$ at time $t$, which Definition~\ref{aux:d1-obs}(b) of \S\ref{sec:d1} constructs with the
record's mollifier, is the limit of $X_\kappa$ (Lemma~\ref{aux:d1-moments}(a)), hence equals
$J_tW_n(f)$.  For a product kernel $f_1\otimes\cdots\otimes f_r$ at one time $t$, with every $f_i$
bounded and vanishing outside one column, Lemma~\ref{aux:d1-moments}(c) makes the monomial the
product of its one-factor monomials, hence equal to
$\prod_iJ_tW_{n_i}(f_i)=J_t(\prod_iW_{n_i}(f_i))$ almost everywhere.
(b) In particular the pieces $Q_{n,n',k,k'}$ of $Q_V$ (\S\ref{sec:d4}), to which
Theorem~\ref{thm:D4} applies Lemma~\ref{lem:D2}, are
$\lambda^2a_na_{n'}\,J_0W_n(h'_k)\,J_0W_{n'}(h'_{k'})$, and their sum is $Q_V=(J_0V(h'))^2=Y^2$ with
$Y$ of Lemma~\ref{aux:notation-hyper}, because $W_n(h')=\sum_kW_n(h'_k)$ (the regularized integrals,
hence their limits, are linear in the kernel).  So the base term of Theorem~\ref{thm:D4}, bounded
through Lemma~\ref{aux:notation-hyper}, and its comparison term, bounded through
Lemma~\ref{lem:D2}, concern one and the same element of $L^2(\dmu)$.
(c) At degree $1$ and $f\in C_c^\infty(\R)$, $X_\kappa=\Phi(\delta_\kappa*(\delta_t\otimes f))\to
\Phi(t,f)=J_tW_1(f)$; this is also immediate from $\delta_\kappa*(\delta_t\otimes f)\to
\delta_t\otimes f$ in the covariance norm.  It covers the unsplit factors $\Phi(s,f_i)$,
$f_i\in C_c^\infty(\R)$, of the cylinder monomials of Theorem~\ref{thm:D3}, but not the pieces to
which Lemma~\ref{lem:D2} is applied there (Lemma~\ref{aux:d3-cyl}(vi)): their kernels
$f_i\mathbf 1_{[k-\frac12,k+\frac12)}$ are bounded, measurable and of compact support, but need
not be smooth.  Those pieces are covered by Lemma~\ref{aux:notation-sharpid} with $n=1$, through
(a), which identifies a one-factor piece of degree~$1$ with kernel $g$ at time $t$ with
$J_tW_1(g)$, and by Lemma~\ref{aux:d1-moments}(f)(1), which identifies it with $\Phi(t,g)$;
Lemma~\ref{aux:d3-cyl}(vi) uses the latter, with Lemma~\ref{aux:d1-moments}(c) for the
products.
(d) Lemma~\ref{aux:notation-sharpid} uses of Appendix~C only its regularization; it needs no repair
of the V3 audit.
\end{auxrem}

\subsection{The printed constants}\label{sec:notation-printed}

The constants $C_7$, $m_1$, $K_1$, $K_2$ and $\bar n$ are fixed once, with the working
constants of \S2.6 of the record \cite[p.~14]{R}: ``$m_1>0$ and $C_7,C_8,K_1,K_2$ denote the
printed constants''.  The source prints $C_7$ and $m_1$ as ``independent of $h$, $Q_1$ and
$Q_2$''; how they depend on $(m_0,\lambda,\mathscr P,K_1,K_2)$ is not printed, and this note
uses no such dependence.  The constant $C_8$ of the record (Assumption~A8(v)--(vi)) is not
used in this note.

\begin{center}
\small
\begin{tabular}{@{}p{0.11\textwidth}p{0.42\textwidth}p{0.39\textwidth}@{}}
\toprule
constant & printed in \cite{GJS} (read on the page images) & in the record \cite{R}\\
\midrule
$\bar n$ & p.~593: ``We assume a bound on the $n_i$ in (1.1.5), $n_i<\bar n$. Any choice
$\bar n\geqq\deg P$ is sufficient.'' & A8(ii), p.~10: the fixed choice $\bar n=2p'+1$; every
Wick factor used here has degree $\le2p'<\bar n$\\
$K_1$, $K_2$ & (1.1.7)--(1.1.8), p.~594; Theorem~1.1.7, p.~594: ``Let $m_0$, $K_2$ and $K_1$ be
sufficiently large and let $\lambda$ be sufficiently small.'' & A8(iii), eq.~(1), p.~10;
``we fix $K_2\ge1$''\\
$C_7$, $m_1$ & Theorem~1.1.7, p.~594: ``There are constants $0<m_1$ and $O(1)$ independent of
$h$, $Q_1$ and $Q_2$ such that'' (1.1.10); $C_7$ is the record's name for this $O(1)$ &
Source Fact~2.2 (A8(iv)), p.~10: ``with $m_1>0$ and $C_7$ independent of $h$, $Q_1$, $Q_2$
(so stated in the source)''\\
$m_0$, $\lambda$ & hypotheses of Theorem~1.1.7, p.~594 (quoted above) & \S2.6, p.~14: $m_0=M_*$,
$0\le\lambda<\lambda_*$ in working coordinates\\
$\mathscr P$, $p'$, $a_n$ & (1.8), p.~589: ``$V(h)=\lambda\int_{R^2}:\mathscr P(\Phi(x)):h(x)dx$'' &
\S2.1, p.~6: $\mathscr P(\xi)=\sum_{n=0}^{2p'}a_n\xi^n$, $a_{2p'}>0$, $p'\ge1$\\
\bottomrule
\end{tabular}
\end{center}

The derived constants of this section are $c_U$ (Lemma~\ref{aux:notation-cU}), $S_2$
(Definition~\ref{aux:notation-S2def}; finite and unevaluated), $b_0$ and $I_n$
(Lemma~\ref{aux:notation-kernel}), $\overline V_*\ge V_*$ (Lemma~\ref{aux:notation-wick}),
and the ring constant $\overline\Sigma_{3\times5}\ge\Sigma_{3\times5}$
(Lemma~\ref{aux:notation-rings}(vi)); Appendix~\ref{app:constants}
tabulates them with the constants of Theorems~\ref{thm:D3}--\ref{thm:D4} and of
\S\ref{sec:cons}.


\section{The affine-path identity}\label{sec:d1}
% d1.tex -- drafted by mp (council member of ugjs/PLAN.md), 31 October and 2 November 2026.
% Section 3 of the comparison note: Lemma D.1 with the integrability chain of gate GA2.

This section proves Lemma~\ref{lem:D1}. It is the identity that GJS print as the first two
equalities of (1.1.12) \cite[p.~595]{GJS},
\[
 \text{``}\langle Q\rangle_{h_2}-\langle Q\rangle_{h_1}
 =\int_0^1d\alpha\,\frac{d}{d\alpha}\langle Q\rangle_{h_1+\alpha g}
 =-\int_0^1d\alpha\bigl[\langle QV(g)\rangle_{h_1+\alpha g}
 -\langle Q\rangle_{h_1+\alpha g}\langle V(g)\rangle_{h_1+\alpha g}\bigr]\text{''},
\]
where ``To prove (1.1.11) let $g=h_2-h_1$'' (p.~594) and $\langle F\rangle_h=\int F\,dq_h$
(ibid.); the pages give no integrability argument, and this section supplies one. Its inputs
are: Assumption~A3 of the record, as stated there; the linearity of $h\mapsto\Uc_h$ (Step~1
of the proof below), which rests on repair R8 of the V3 audit or, on the route written in
Step~1, on the record's identification of $\Uc_h$ with the limit of its regularizations, that is
on Lemma~C.4 of the record with the repairs R1--R10 (the status paragraph after
Lemma~\ref{lem:D1}); the Gaussian moments of the observables, computed explicitly in
Lemmas~\ref{aux:d1-wick}, \ref{aux:d1-cov} and~\ref{aux:d1-moments} from Assumption~A2 of the
record alone (neither the decomposition of $L^2(\dmu)$ into Wick chaoses nor
hypercontractivity is used); H\"older's inequality, with the exponents written at each use;
dominated convergence; Fatou's lemma; and the fundamental theorem of calculus. No smallness
hypothesis is used: everything in this section holds for every $\lambda\ge0$ and every
$m_0>0$.

\subsection*{The observables and their Gaussian moments}

\emph{Gaussian variables.} As in \S\ref{sec:notation}, $\dmu$ is the centred Gaussian measure
on $\mathcal S'(\R^2)$ with covariance $(-\Delta_{\R^2}+m_0^2)^{-1}$ (Assumption~A2 of the
record \cite[p.~8]{R}), and $\Phi(f)(q)=\langle q,f\rangle$. Let $\mathscr H_1\subset
L^2(\dmu)$ be the closed linear span of the $\Phi(f)$, $f\in C_c^\infty(\R^2)$ real. Its
elements are jointly Gaussian and centred, and
$\E[\Phi(f)\Phi(g)]=\ip{f}{(-\Delta_{\R^2}+m_0^2)^{-1}g}_{L^2(\R^2)}$.

\emph{Regularization and Wick powers.} Appendix~C of the record \cite[p.~39]{R} fixes
$\delta\in C_c^\infty(\R^2)$ with $\delta\ge0$, $\delta(z)=\delta(-z)$,
$\supp\delta\subseteq\{z:|z|\le1\}$ and $\int\delta=1$, and puts
$\delta_\kappa(z):=\kappa^2\delta(\kappa z)$. For $\kappa\ge2$ and $\zeta\in\R^2$ let
$\Phi_\kappa(\zeta):=\Phi(\delta_\kappa(\zeta-\cdot))\in\mathscr H_1$, the record's
$\delta_\kappa*\Phi$ at $\zeta$; for every $q\in\mathcal S'(\R^2)$ the function
$\zeta\mapsto\Phi_\kappa(\zeta)(q)=\langle q,\delta_\kappa(\zeta-\cdot)\rangle$ is smooth, being
the convolution of a tempered distribution with a test function. Let
$c_\kappa:=\E[\Phi_\kappa(\zeta)^2]$, which does not depend on $\zeta$
(Lemma~\ref{aux:d1-cov}(iii)). For a centred Gaussian variable $X$ and $c\ge0$ the Wick powers
${:}X^n{:}_c$ are the coefficients in
$\sum_{n\ge0}\frac{s^n}{n!}{:}X^n{:}_c=\exp(sX-\frac12s^2c)$, the convention of the record
(proof of Proposition~3.2, \cite[p.~16]{R}); expanding the two exponentials,
\begin{equation}\label{eq:d1-hermite}
 {:}X^n{:}_c=\sum_{0\le j\le n/2}\frac{n!}{j!\,(n-2j)!}\Bigl(-\frac c2\Bigr)^jX^{n-2j} .
\end{equation}

\begin{auxlemma}[moments of products of Wick powers]\label{aux:d1-wick}
Let $X_1,\dots,X_\ell\in\mathscr H_1$, $c_a:=\E[X_a^2]$, $\gamma_{ab}:=\E[X_aX_b]$, and let
$m_1,\dots,m_\ell\ge0$ be integers. Then
\begin{equation}\label{eq:d1-wick}
 \E\Bigl[\prod_{a=1}^\ell{:}X_a^{m_a}{:}_{c_a}\Bigr]=\mathcal P_m(\gamma)
 :=\sum_{(k_{ab})}\ \prod_{a=1}^\ell m_a!\prod_{1\le a<b\le\ell}
 \frac{\gamma_{ab}^{k_{ab}}}{k_{ab}!} ,
\end{equation}
the sum running over all families $(k_{ab})_{1\le a<b\le\ell}$ of nonnegative integers with
$\sum_{b\ne a}k_{ab}=m_a$ for every $a$, where $k_{ba}:=k_{ab}$ (an empty sum is $0$). Thus
$\mathcal P_m$ is a polynomial with nonnegative coefficients in the off-diagonal covariances
$\gamma_{ab}$, $a<b$, alone, and each of its monomials has degree $\frac12\sum_am_a$. For
$\ell=2$: $\E[{:}X_1^n{:}_{c_1}{:}X_2^{n'}{:}_{c_2}]=\delta_{nn'}\,n!\,\gamma_{12}^n$ (Wick
orthogonality and the Wick isometry).
\end{auxlemma}

\begin{proof}
For real $s_1,\dots,s_\ell$ the variable $\sum_as_aX_a$ is centred Gaussian with variance
$\sum_as_a^2c_a+2\sum_{a<b}s_as_b\gamma_{ab}$, so
\[
 \E\Bigl[\prod_a\exp\bigl(s_aX_a-\tfrac12s_a^2c_a\bigr)\Bigr]
 =\exp\Bigl(\sum_{a<b}s_as_b\gamma_{ab}\Bigr)
 =\prod_{a<b}\ \sum_{k\ge0}\frac{(s_as_b\gamma_{ab})^k}{k!} .
\]
On the left expand each factor by the generating function. The expansion may be integrated
term by term: by \eqref{eq:d1-hermite},
$\sum_n\frac{|s|^n}{n!}|{:}X^n{:}_c|
\le\sum_n\sum_{j\le n/2}\frac{|s|^n(c/2)^j|X|^{n-2j}}{j!\,(n-2j)!}=e^{|s||X|+s^2c/2}$, and
$\E\prod_ae^{|s_a||X_a|}\le\prod_a(\E e^{\ell|s_a||X_a|})^{1/\ell}<\infty$ by H\"older's
inequality with $\ell$ factors, each with exponent $\ell$, and the Gaussian exponential
moments. Fubini's theorem then turns the left side into
$\sum_m\prod_a\frac{s_a^{m_a}}{m_a!}\,\E[\prod_a{:}X_a^{m_a}{:}_{c_a}]$, a power series
converging for all $s$, and comparing the coefficients of $\prod_as_a^{m_a}$ gives
\eqref{eq:d1-wick}. No $\gamma_{aa}$ occurs, and every term has degree
$\sum_{a<b}k_{ab}=\frac12\sum_am_a$. For $\ell=2$ the only family is $k_{12}=n=n'$.
\end{proof}

\begin{auxlemma}[the free covariance]\label{aux:d1-cov}
Let $G$ be the kernel of $(-\Delta_{\R^2}+m_0^2)^{-1}$. Put $g_0:=(4+m_0^{-2})/(4\pi)$,
$\log_+t:=\max\{\log t,0\}$ and $M_0(\varrho):=g_0+\frac1{2\pi}\log_+\frac1\varrho$ for
$\varrho>0$.
\begin{itemize}
\item[(i)] For $\zeta\ne0$,
$G(\zeta)=\int_0^\infty(4\pi s)^{-1}\exp\bigl(-\frac{|\zeta|^2}{4s}-m_0^2s\bigr)\,ds$, and
$\E[\Phi(f)\Phi(g)]=\iint f(\zeta)\,G(\zeta-\zeta')\,g(\zeta')\,d\zeta\,d\zeta'$ for real
$f,g\in C_c^\infty(\R^2)$. Likewise the time-zero covariance $C=(2\mu)^{-1}$ of the record
has the kernel $C(x)=G(0,x)$, $x\in\R\setminus\{0\}$:
$\ip{f}{Cg}_{L^2(\R)}=\iint f(x)\,G(0,x-y)\,g(y)\,dx\,dy$ for real $f,g\in C_c^\infty(\R)$.
\item[(ii)] $G$ is positive, continuous on $\R^2\setminus\{0\}$, a nonincreasing function of
$|\zeta|$, and $G(\zeta)\le M_0(|\zeta|)$. In particular the kernel $C(x)=G(0,x)$ of
$(2\mu)^{-1}$ is positive and $C(x)\le M_0(|x|)$; this is the kernel form of the record's
conventions lock, Remark~2.5 \cite[p.~14]{R}.
\item[(iii)] For $\kappa,\kappa'\ge2$ and $\zeta,\zeta'\in\R^2$:
$\E[\Phi_\kappa(\zeta)\Phi_{\kappa'}(\zeta')]=G_{\kappa\kappa'}(\zeta-\zeta')$ with
$G_{\kappa\kappa'}:=\delta_\kappa*\delta_{\kappa'}*G$, a continuous function on $\R^2$;
$0<G_{\kappa\kappa'}(\zeta)\le M_1(|\zeta|):=g_1+b_1\log_+\frac1{|\zeta|}$ for $\zeta\ne0$,
where
$g_1:=\max\bigl\{g_0+\frac{\log2}{2\pi},\ \|\delta\|_\infty(4\pi g_0+1+2\log4)\bigr\}$ and
$b_1:=\max\{\frac1{2\pi},\,2\|\delta\|_\infty\}$; and
$G_{\kappa\kappa'}(\zeta)\to G(\zeta)$ as $\min\{\kappa,\kappa'\}\to\infty$, for every
$\zeta\ne0$. In particular $c_\kappa=G_{\kappa\kappa}(0)$.
\item[(iv)] For $R>0$, $q\in[1,\infty)$ and $d\in\{1,2\}$:
$\int_{\{t\in\R^d:\,|t|\le R\}}M_1(|t|)^q\,dt<\infty$.
\end{itemize}
\end{auxlemma}

\begin{proof}
(ii) The integrand of the formula in (i), taken as the definition of $G$ on
$\R^2\setminus\{0\}$, is positive and nonincreasing in $|\zeta|$; on $\{|\zeta|\ge\varepsilon\}$
it is bounded by $(4\pi s)^{-1}e^{-\varepsilon^2/(4s)-m_0^2s}$, which is integrable over
$(0,\infty)$, so $G$ is continuous there by dominated convergence. Let $0<|\zeta|\le1$ and split
the $s$-integral at $s=|\zeta|^2$ and $s=1$. With $s=|\zeta|^2v$ and $e^{-y}\le y^{-1}$
($y>0$),
$\int_0^{|\zeta|^2}\frac{e^{-|\zeta|^2/(4s)}}{4\pi s}ds=\frac1{4\pi}\int_0^1\frac{e^{-1/(4v)}}{v}dv
\le\frac1{4\pi}\int_0^14\,dv=\frac1\pi$; next
$\int_{|\zeta|^2}^1\frac{ds}{4\pi s}=\frac1{2\pi}\log\frac1{|\zeta|}$; and
$\int_1^\infty\frac{e^{-m_0^2s}}{4\pi s}ds\le\frac1{4\pi m_0^2}$. Dropping the factors
$e^{-m_0^2s}\le1$ and $e^{-|\zeta|^2/(4s)}\le1$ where they were not used,
$G(\zeta)\le\frac1\pi+\frac1{4\pi m_0^2}+\frac1{2\pi}\log\frac1{|\zeta|}=M_0(|\zeta|)$. For
$|\zeta|>1$, monotonicity gives $G(\zeta)\le\sup_{|\eta|=1}G(\eta)\le g_0=M_0(|\zeta|)$. The
statement about $C$ follows, since $C(x)=G(0,x)$ by (i).

(iv) Since $\log_+(1/|t|)=0$ for $|t|\ge1$ and $(a+b)^q\le2^{q-1}(a^q+b^q)$ for $a,b\ge0$, it
suffices that $\int_0^1(\log\frac1t)^qdt=\Gamma(q+1)$ and
$\int_0^1t(\log\frac1t)^qdt=2^{-q-1}\Gamma(q+1)$ are finite (substitute $t=e^{-u}$).

(i) By Assumption~A2 and Parseval's identity, with $\hat f(p):=\int e^{-ip\cdot\zeta}f(\zeta)\,d\zeta$,
$\E[\Phi(f)\Phi(g)]=(2\pi)^{-2}\int\overline{\hat f(p)}\,\hat g(p)\,(|p|^2+m_0^2)^{-1}dp$.
Insert $(|p|^2+m_0^2)^{-1}=\int_0^\infty e^{-s(|p|^2+m_0^2)}ds$; Fubini's theorem applies,
because the double integral of $|\hat f\hat g|e^{-s(|p|^2+m_0^2)}$ is
$\int|\hat f\hat g|(|p|^2+m_0^2)^{-1}dp<\infty$. For fixed $s>0$, expanding $\hat f,\hat g$
(all integrals absolutely convergent) and using
$(2\pi)^{-2}\int e^{ip\cdot\zeta-s|p|^2}dp=(4\pi s)^{-1}e^{-|\zeta|^2/(4s)}$ (two
one-dimensional Gaussian integrals) gives
$(2\pi)^{-2}\int\overline{\hat f}\hat g\,e^{-s|p|^2}dp
=\iint f(\zeta)(4\pi s)^{-1}e^{-|\zeta-\zeta'|^2/(4s)}g(\zeta')\,d\zeta\,d\zeta'$. A second use
of Fubini's theorem, legitimate since by (ii) and (iv) the absolute integrand integrates to
$\iint|f(\zeta)|\,G(\zeta-\zeta')\,|g(\zeta')|<\infty$ ($f,g$ bounded with bounded support),
gives the kernel formula with $G$ as displayed in (i). The time-zero statement is the same
computation in one dimension: $(2\mu(k))^{-1}=\frac12\pi^{-1/2}\int_0^\infty
s^{-1/2}e^{-s(k^2+m_0^2)}ds$ (since $\int_0^\infty s^{-1/2}e^{-sa}ds=\sqrt\pi\,a^{-1/2}$),
$(2\pi)^{-1}\int e^{ikx-sk^2}dk=(4\pi s)^{-1/2}e^{-x^2/(4s)}$, and
$\frac12\pi^{-1/2}s^{-1/2}(4\pi s)^{-1/2}=(4\pi s)^{-1}$; the two uses of Fubini's theorem are
justified as before, by $\int|\hat f\hat g|(2\mu)^{-1}dk<\infty$ and by (ii) and (iv) with
$d=1$.

(iii) By (i) with $f=\delta_\kappa(\zeta-\cdot)$, $g=\delta_{\kappa'}(\zeta'-\cdot)$, the
substitutions $\eta=\zeta-a$, $\eta'=\zeta'+b$ and $\delta_{\kappa'}(-b)=\delta_{\kappa'}(b)$,
$\E[\Phi_\kappa(\zeta)\Phi_{\kappa'}(\zeta')]=\iint\delta_\kappa(a)\delta_{\kappa'}(b)\,
G(\zeta-\zeta'-a-b)\,da\,db=G_{\kappa\kappa'}(\zeta-\zeta')$; it is continuous as the
convolution of the locally integrable $G$ (by (ii) and (iv)) with
$\tilde\delta:=\delta_\kappa*\delta_{\kappa'}\in C_c^\infty(\R^2)$. Now $\tilde\delta\ge0$,
$\int\tilde\delta=1$, $\supp\tilde\delta\subseteq\{|c|\le\varrho\}$ with
$\varrho:=\kappa^{-1}+\kappa'^{-1}\le1$, and
$\|\tilde\delta\|_\infty\le\min\{\|\delta_\kappa\|_\infty\|\delta_{\kappa'}\|_1,
\|\delta_{\kappa'}\|_\infty\|\delta_\kappa\|_1\}=\kappa_m^2\|\delta\|_\infty$ with
$\kappa_m:=\min\{\kappa,\kappa'\}$, and $\kappa_m^{-1}\le\varrho\le2\kappa_m^{-1}$.
Positivity follows from $G>0$. If $|\zeta|\ge2\varrho$, every $c$ with $|c|\le\varrho$ has
$|\zeta-c|\ge|\zeta|/2$, so by (ii)
$G_{\kappa\kappa'}(\zeta)\le M_0(|\zeta|/2)\le g_0+\frac{\log2}{2\pi}
+\frac1{2\pi}\log_+\frac1{|\zeta|}\le M_1(|\zeta|)$. If $0<|\zeta|<2\varrho$, then
$G_{\kappa\kappa'}(\zeta)\le\kappa_m^2\|\delta\|_\infty\int_{|\theta-\zeta|\le\varrho}
M_0(|\theta|)\,d\theta\le\kappa_m^2\|\delta\|_\infty\int_{|\theta|\le\varrho}M_0(|\theta|)\,d\theta$.
The second inequality holds for every nonincreasing $M_0$: the sets
$A:=\{|\theta-\zeta|\le\varrho<|\theta|\}$ and $A':=\{|\theta|\le\varrho<|\theta-\zeta|\}$ have
equal measure, $M_0(|\theta|)\le M_0(\varrho)$ on $A$ and $M_0(|\theta|)\ge M_0(\varrho)$ on
$A'$. Since $\varrho\le1$,
$\int_{|\theta|\le\varrho}M_0(|\theta|)\,d\theta=\pi\varrho^2g_0+\int_0^\varrho
t\log\frac1t\,dt=\pi\varrho^2g_0+\frac{\varrho^2}2\log\frac1\varrho+\frac{\varrho^2}4$, and with
$\kappa_m^2\varrho^2\le4$ and $\log\frac1\varrho\le\log\kappa_m$ this gives
$G_{\kappa\kappa'}(\zeta)\le\|\delta\|_\infty(4\pi g_0+1+2\log\kappa_m)$. Finally
$|\zeta|<2\varrho\le4/\kappa_m$ gives $\log\kappa_m<\log4+\log_+\frac1{|\zeta|}$, so
$G_{\kappa\kappa'}(\zeta)\le M_1(|\zeta|)$ in this case too. For $\zeta\ne0$ and
$\varrho\le|\zeta|/2$,
$|G_{\kappa\kappa'}(\zeta)-G(\zeta)|\le\sup_{|c|\le\varrho}|G(\zeta-c)-G(\zeta)|$, which tends
to $0$ as $\varrho\to0$, i.e.\ as $\kappa_m\to\infty$, by the continuity of $G$ at $\zeta$.
Taking $\zeta=\zeta'$ gives $c_\kappa=G_{\kappa\kappa}(0)$.
\end{proof}

\begin{auxdef}[localized monomials, made precise]\label{aux:d1-obs}
Let $r\ge0$ and let $n_1,\dots,n_r$ be integers with $0\le n_i<\bar n$; $N:=n_1+\dots+n_r$ is
the total degree, and $\Delta_{j_1},\dots,\Delta_{j_r}$ are GJS squares. For $\kappa\ge2$:
\begin{itemize}
\item[(a)] \emph{type (1.1.6):} for real $w\in L^2(\R^{2r})$ vanishing outside
$\Delta_{j_1}\times\dots\times\Delta_{j_r}$, and $x=(x_1,\dots,x_r)$ with $x_i\in\R^2$,
\[
 Q^{(\kappa)}:=\int_{\R^{2r}}w(x)\,A_\kappa(x)\,dx,\qquad
 A_\kappa(x):=\prod_{i=1}^r{:}\Phi_\kappa(x_i)^{n_i}{:}_{c_\kappa} ;
\]
\item[(b)] \emph{type (1.1.5):} for times $x_0=(x_{1,0},\dots,x_{r,0})\in\R^r$, equal or not,
and real $w(x_0,\cdot)\in L^2(\R^r)$ vanishing at $\vec x=(x_1,\dots,x_r)$ unless
$(x_{i,0},x_i)\in\Delta_{j_i}$ for every $i$,
\[
 Q^{(\kappa)}(x_0):=\int_{\R^r}w(x_0,\vec x)\,A_\kappa(\vec x)\,d\vec x,\qquad
 A_\kappa(\vec x):=\prod_{i=1}^r{:}\Phi_\kappa(x_{i,0},x_i)^{n_i}{:}_{c_\kappa} .
\]
\end{itemize}
The integrals are taken pointwise in $q\in\mathcal S'(\R^2)$: for each $q$ the integrand
$A_\kappa(\cdot)(q)$ is continuous, by the smoothness of $\zeta\mapsto\Phi_\kappa(\zeta)(q)$ and
\eqref{eq:d1-hermite}, and $w$ is integrable, its support being bounded; the result is
measurable in $q$ because $(q,x)\mapsto A_\kappa(x)(q)$ is measurable in $q$ and continuous in
$x$. The localized monomial $Q$, resp.\ $Q(x_0)$, is the limit of $Q^{(\kappa)}$, resp.\
$Q^{(\kappa)}(x_0)$, in $L^2(\dmu)$ as $\kappa\to\infty$, which exists by
Lemma~\ref{aux:d1-moments}(a). The class $\mathfrak F$ consists of the finite real linear
combinations of localized monomials of either type.
\end{auxdef}

\begin{auxrem}[provenance of the definition]\label{aux:d1-def}
GJS print (1.1.5)--(1.1.6) as formal integrals, ``$Q(x_0)=\int d\vec x\,w(x)\prod_{i=1}^r
{:}\Phi(x_i)^{n_i}{:}$'' and ``$Q=\int dx_0Q(x_0)$'', with ``localized'' meaning that ``$w(x)$
is supported in a product $\Delta_{j_1}\times\cdots\times\Delta_{j_n}\subset R^{2r}$ where
$\Delta_{j_i}$ is the unit lattice square in $R^2$ centered at the lattice side [sic]
$j_i\in Z^2$'' (p.~593; the index $n$ of the product is printed against $r$, also [sic]); the
record imports the class as printed (A8(ii), \cite[p.~10]{R}) and gives no construction.
Definition~\ref{aux:d1-obs} makes the integrals precise in the way the record treats the
interaction: for $h\in\CE$, the degree-$n$ part of the record's $\Uc_{h,\kappa}$ is the sum,
over the finitely many squares $\Delta$ meeting $\supp h$, of the $Q^{(\kappa)}$ with $r=1$ and
kernel $\lambda a_nh\mathbf 1_\Delta$, and the record identifies $\Uc_h$ with the limit of
$\Uc_{h,\kappa}$ (Appendix~C, notation paragraph, \cite[p.~39]{R}; quantified by
Lemma~C.4, \cite[p.~40]{R}). For a product kernel, $Q$ is the product of the individually
Wick-ordered smeared powers (Lemma~\ref{aux:d1-moments}(c)), which is how A8(ii) of the record
reads p.~593 (``products of individually Wick-ordered powers are admissible without
reordering''). Case (b) admits distinct times $x_{i,0}$, the reading of p.~595 (``In the
monomial $Q(x_0)$ of (1.1.5), suppose $x_0=(x_{1,0},\cdots,x_{r,0})$''); nothing below
depends on whether the times are equal. The bound $n_i<\bar n$ is carried because p.~593
prints it (``We assume a bound on the $n_i$ in (1.1.5), $n_i<\bar n$''); nothing in this
section uses it.  That these limits are the monomials of p.~593, which GJS print as formal
integrals with ``${:}\ {:}$'' denoting ``Wick order with respect to the measure $d\Phi$, as in
[9], [1]'' (p.~589), is a reading of the source of the kind the record names in
Assumption~A8, Scope (d); so is the identification of the record's $\Uc_h$ with GJS's
$V(h)$ of (1.8), p.~589, which the V3 audit (\S12) leaves unchecked.  The limits do not depend on
the mollifier $\delta$ (Lemma~\ref{aux:d1-moments}(e)), and at sharp times they agree with the
constructions of \S\ref{sec:notation} (Lemma~\ref{aux:d1-moments}(f) and
Lemma~\ref{aux:notation-sharpid}).
\end{auxrem}

\begin{auxlemma}[Gaussian moments of the observables]\label{aux:d1-moments}
In the setting of Definition~\ref{aux:d1-obs} write $x$ for the integration variable ($x\in
\R^{2r}$ in case (a), $x=\vec x\in\R^r$ in case (b)), $w(x)$ for the kernel and $K$ for a
bounded set outside which $w$ vanishes.
\begin{itemize}
\item[(a)] $Q^{(\kappa)}\in L^2(\dmu)$ for $\kappa\ge2$, and $Q^{(\kappa)}$ converges in
$L^2(\dmu)$ as $\kappa\to\infty$; the limit is $Q$.
\item[(b)] For every integer $j\ge1$,
$\E[Q^{2j}]\le\|\Pi_{Q,j}\|_{L^2(K^{2j})}\,\|w\|_{L^2}^{2j}<\infty$, with $\Pi_{Q,j}$ as in
\eqref{eq:d1-major}; $\|\Pi_{Q,j}\|_{L^2(K^{2j})}$ depends only on $j$, $n_1,\dots,n_r$, $K$,
$m_0$ and $\delta$. Hence $Q\in L^q(\dmu)$ for every $q<\infty$, and so is every
$F\in\mathfrak F$, by Minkowski's inequality.
\item[(c)] If $w(x)=\prod_{i=1}^rf_i(x_i)$, with each $f_i$ vanishing outside (the spatial
projection of) $\Delta_{j_i}$, then $Q=\prod_{i=1}^rX_i$ almost everywhere, where $X_i$ is the
monomial of Definition~\ref{aux:d1-obs} with one factor, degree $n_i$ and kernel $f_i$ (in case
(b) at the time $x_{i,0}$), and
$\|X_i\|_{L^2(\dmu)}^2=n_i!\iint f_i(u)f_i(v)\,\Gamma(u-v)^{n_i}\,du\,dv$ with $\Gamma=G$ in
case (a) and $\Gamma=C$ in case (b).
\item[(d)] \emph{(Wick orthogonality and isometry for the limits.)} If $X$ and $X'$ are monomials
with one factor, of degrees $n,n'$ and kernels $f,f'$, both of type (a), or both of type (b) at
one time $t$, then $\E[XX']=\delta_{nn'}\,n!\iint f(u)f'(v)\,\Gamma(u-v)^n\,du\,dv$, with
$\Gamma$ as in (c).
\item[(e)] \emph{(Independence of the mollifier.)} If $\delta'$ has the properties of $\delta$
($\delta'\in C_c^\infty(\R^2)$, $\delta'\ge0$, $\delta'(z)=\delta'(-z)$,
$\supp\delta'\subseteq\{|z|\le1\}$, $\int\delta'=1$), the limits of
Definition~\ref{aux:d1-obs} formed with $\delta$ and with $\delta'$ coincide almost everywhere.
\item[(f)] \emph{(Agreement with the sharp-time constructions of \S\ref{sec:notation}.)} Let
$\Phi(t,f)$ be the sharp-time field defined in the proof of Lemma~\ref{aux:notation-hyper}(i)
(the extension of $f\mapsto\Phi(f)$ by continuity in the covariance norm, there at $t=0$).
\begin{itemize}
\item[(1)] For real $f\in L^2(\R)$ with bounded support and every $t\in\R$, the monomial of
Definition~\ref{aux:d1-obs}(b) with one factor of degree $1$ and kernel $f$ at the time $t$ is
$\Phi(t,f)$.
\item[(2)] For real bounded measurable $f$ vanishing outside one column, $0\le n\le2p'$ and every
$t\in\R$, the monomial with one factor of degree $n$ and kernel $f$ at the time $t$ equals
$J_tW_n(f)$ almost everywhere, and a product-kernel monomial at one time $t$ with factors of
this kind equals $J_t(\prod_iW_{n_i}(f_i))$ almost everywhere: this is
Lemma~\ref{aux:notation-sharpid} with Remark~\ref{aux:notation-sharpidrem}(a) of
\S\ref{sec:notation} ($W_n(f)$ the limit of Lemma~\ref{aux:notation-wick}(iii) for $h'=f$, $J_t$
the sharp-time map there), the monomial being, by (a), the $L^2(\dmu)$-limit of that lemma's
$X_\kappa$.
\end{itemize}
Only (1) is proved here, since Lemma~\ref{aux:notation-sharpid} does not cover it: that lemma
needs a bounded kernel and names the degree-$1$ limit $J_tW_1(f)$, while (1) holds for every
kernel in $L^2$ with bounded support and names it $\Phi(t,f)$ (for $n\ge2$ and an unbounded
kernel, \S\ref{sec:notation} has no $W_n(f)$ to compare with).
\end{itemize}
\end{auxlemma}

\begin{proof}
\emph{The moment formula.} For $x\in K$ let $\zeta_i(x):=x_i$ in case (a) and
$\zeta_i(x):=(x_{i,0},x_i)$ in case (b). Let $j\ge1$, let $\kappa_1,\dots,\kappa_{2j}\ge2$ and
let $x^1,\dots,x^{2j}$ be configurations. Lemma~\ref{aux:d1-wick} with the $\ell=2jr$ variables
$\Phi_{\kappa_c}(\zeta_i(x^c))$, $c\le2j$, $i\le r$, the degrees $m=(n,\dots,n)$ and the
variances $c_{\kappa_c}$, together with Lemma~\ref{aux:d1-cov}(iii), gives
\begin{equation}\label{eq:d1-two}
 \E\Bigl[\prod_{c=1}^{2j}A_{\kappa_c}(x^c)\Bigr]
 =\mathcal P_{(n,\dots,n)}\bigl(\gamma(x^1,\dots,x^{2j})\bigr),
\end{equation}
where the entries of $\gamma$ are $G_{\kappa_c\kappa_{c'}}(\zeta_i(x^c)-\zeta_{i'}(x^{c'}))$,
for $c<c'$ and all $i,i'$, and for $c=c'$ and $i<i'$. Each entry is a continuous function of
the configurations (Lemma~\ref{aux:d1-cov}(iii)), so
$x\mapsto\E[A_\kappa(x)^{2j}]$ is continuous and bounded on the compact set $\bar K$.

\emph{The majorant.} By Lemma~\ref{aux:d1-cov}(iii) every entry of $\gamma$ is positive and
at most $M_1$ of the distance of its two points; in case (b) that distance is at least the
distance $|u-v|$ of the two spatial integration variables $u,v$ attached to the entry, and
$M_1$ is nonincreasing. As $\mathcal P$ has nonnegative coefficients, for all
$\kappa_1,\dots,\kappa_{2j}\ge2$
\begin{equation}\label{eq:d1-major}
 0\le\mathcal P_{(n,\dots,n)}\bigl(\gamma(x^1,\dots,x^{2j})\bigr)\le
 \Pi_{Q,j}(x^1,\dots,x^{2j}):=\mathcal P_{(n,\dots,n)}\bigl(\mu(x^1,\dots,x^{2j})\bigr),
\end{equation}
where $\mu$ has the entry $M_1(|u-v|)$ in place of each entry of $\gamma$, $(u,v)$ the pair of
integration variables (in $\R^2$ in case (a), in $\R$ in case (b)) attached to it. Each
monomial of $\Pi_{Q,j}$ is a product of $jN$ factors $M_1(|u_e-v_e|)$, $e=1,\dots,jN$, in which
$u_e$ and $v_e$ are two distinct integration variables, because $\mathcal P$ contains no
diagonal covariance. By H\"older's inequality with the $jN$ exponents $2jN$,
$\|\prod_eM_1(|u_e-v_e|)\|_{L^2(K^{2j})}\le\prod_e\|M_1(|u_e-v_e|)\|_{L^{2jN}(K^{2j})}$, and each
factor is finite: integrate first over $u_e$, which ranges over a bounded set, so that
$|u_e-v_e|\le R$ for some $R<\infty$, and use Lemma~\ref{aux:d1-cov}(iv); the other variables
range over bounded sets. Hence $\Pi_{Q,j}\in L^2(K^{2j})$ (if $N=0$, $\Pi_{Q,j}=1$).

\emph{Moments of the regularized observables.} For each $q$, $Q^{(\kappa)}(q)^{2j}$ is the
product of $2j$ absolutely convergent integrals, so
\[
 Q^{(\kappa)}(q)^{2j}=\int_{K^{2j}}\prod_{c=1}^{2j}w(x^c)\prod_{c=1}^{2j}A_\kappa(x^c)(q)\;
 dx^1\cdots dx^{2j} .
\]
Fubini's theorem on $\mathcal S'(\R^2)\times K^{2j}$ applies: by H\"older's inequality with the
$2j$ exponents $2j$, $\E\prod_c|A_\kappa(x^c)|\le\sup_{x\in\bar K}\E[A_\kappa(x)^{2j}]<\infty$,
and $\int\prod_c|w(x^c)|\le(|K|^{1/2}\|w\|_2)^{2j}$. With \eqref{eq:d1-two},
\eqref{eq:d1-major} and the Cauchy--Schwarz inequality on $K^{2j}$,
\begin{align}
 \E\bigl[(Q^{(\kappa)})^{2j}\bigr]
 &=\int_{K^{2j}}\prod_cw(x^c)\,\mathcal P\bigl(\gamma(x^1,\dots,x^{2j})\bigr)\,
   dx^1\cdots dx^{2j}\notag\\
 &\le\|w^{\otimes2j}\|_{L^2}\,\|\Pi_{Q,j}\|_{L^2(K^{2j})}
 =\|\Pi_{Q,j}\|_{L^2(K^{2j})}\,\|w\|_2^{2j} .\label{eq:d1-moment}
\end{align}

(a) The case $j=1$ of \eqref{eq:d1-moment} gives $Q^{(\kappa)}\in L^2(\dmu)$, and the same
computation with two regularizations $\kappa,\kappa'$ gives
$\E[Q^{(\kappa)}Q^{(\kappa')}]=\iint w(x)w(y)\,\mathcal P_{(n,n)}(\gamma^{\kappa\kappa'}(x,y))\,dx\,dy$,
the entries of $\gamma^{\kappa\kappa'}$ being as in \eqref{eq:d1-two} with
$(\kappa_1,\kappa_2)=(\kappa,\kappa')$, dominated by $|w(x)||w(y)|\Pi_{Q,1}(x,y)$, an
integrable function. As $\kappa,\kappa'\to\infty$, every entry converges to the corresponding
value of $G$, for every $(x,y)$ at which the relevant differences are nonzero, that is off a
null set (Lemma~\ref{aux:d1-cov}(iii)); by dominated convergence
$\E[Q^{(\kappa)}Q^{(\kappa')}]$ converges to a limit $I_Q$ independent of how
$\kappa,\kappa'\to\infty$. Therefore
$\|Q^{(\kappa)}-Q^{(\kappa')}\|_2^2=\E[Q^{(\kappa)}Q^{(\kappa)}]-2\E[Q^{(\kappa)}Q^{(\kappa')}]
+\E[Q^{(\kappa')}Q^{(\kappa')}]\to I_Q-2I_Q+I_Q=0$: the family $Q^{(\kappa)}$ is Cauchy in
$L^2(\dmu)$ and its limit $Q$ exists, with $\|Q\|_2^2=I_Q$.

(b) Choose $\kappa_k\to\infty$ with $Q^{(\kappa_k)}\to Q$ almost everywhere (possible by (a)).
Fatou's lemma and \eqref{eq:d1-moment} give
$\E[Q^{2j}]\le\liminf_k\E[(Q^{(\kappa_k)})^{2j}]\le\|\Pi_{Q,j}\|_{L^2(K^{2j})}\|w\|_2^{2j}$.
For $q<\infty$ take $2j\ge q$ and use $\|Q\|_q\le\|Q\|_{2j}$, $\dmu$ being a probability
measure.

(c) For each $q$ the integral defining $Q^{(\kappa)}(q)$ factorizes by Fubini's theorem,
$Q^{(\kappa)}(q)=\prod_iX_i^{(\kappa)}(q)$ with
$X_i^{(\kappa)}(q):=\int f_i(x_i)\,{:}\Phi_\kappa(\zeta_i)^{n_i}{:}_{c_\kappa}(q)\,dx_i$. Take any
$\kappa_k\to\infty$. By (a), $X_i^{(\kappa_k)}\to X_i$ and $Q^{(\kappa_k)}\to Q$ in $L^2(\dmu)$,
hence in probability; a finite product of sequences converging in probability converges in
probability to the product of the limits, and limits in probability are unique almost
everywhere; so $Q=\prod_iX_i$ almost everywhere. For one factor,
$\mathcal P_{(n_i,n_i)}(\gamma)=n_i!\,\gamma^{n_i}$ (Lemma~\ref{aux:d1-wick}, $\ell=2$), and the
limiting covariance is $G(u-v)$ in case (a) and $G((x_{i,0},u)-(x_{i,0},v))=G(0,u-v)=C(u-v)$ in
case (b) (Lemma~\ref{aux:d1-cov}(i)); so $\|X_i\|_2^2=I_{X_i}$ is the displayed integral.

(d) As in (a), Lemma~\ref{aux:d1-wick} with $\ell=2$ gives
\[
 \E\bigl[X^{(\kappa)}X'^{(\kappa')}\bigr]=\delta_{nn'}\,n!\iint f(u)f'(v)\,
 G_{\kappa\kappa'}(\zeta(u)-\zeta(v))^n\,du\,dv ,
\]
with $\zeta(u)=u$ in case (a) and $\zeta(u)=(t,u)$ in case (b); the integrand is dominated by $n!\,|f(u)||f'(v)|M_1(|u-v|)^n$, which is integrable by the
Cauchy--Schwarz inequality and Lemma~\ref{aux:d1-cov}(iv). Let $\kappa,\kappa'\to\infty$: the
left side tends to $\E[XX']$ because $X^{(\kappa)}\to X$ and $X'^{(\kappa')}\to X'$ in
$L^2(\dmu)$, and the right side to the stated integral by dominated convergence and
Lemma~\ref{aux:d1-cov}(i),(iii).

(e) Write $\Phi'_{\kappa'}(\zeta):=\Phi(\delta'_{\kappa'}(\zeta-\cdot))$.  As in the proof of
Lemma~\ref{aux:d1-cov}(iii), $\E[\Phi_\kappa(\zeta)\Phi'_{\kappa'}(\zeta')]=(\tilde\delta*G)
(\zeta-\zeta')$ with $\tilde\delta:=\delta_\kappa*\delta'_{\kappa'}$, which is nonnegative, has
integral $1$ and support in $\{|c|\le\kappa^{-1}+\kappa'^{-1}\}$, and satisfies
$\|\tilde\delta\|_\infty\le\min\{\kappa^2\|\delta\|_\infty,\kappa'^2\|\delta'\|_\infty\}\le
\min\{\kappa,\kappa'\}^2\max\{\|\delta\|_\infty,\|\delta'\|_\infty\}$.  Those are the only
properties of $\tilde\delta$ that proof uses, so $\tilde\delta*G$ is positive, continuous, bounded
by $M_1$ with $\|\delta\|_\infty$ replaced by $\max\{\|\delta\|_\infty,\|\delta'\|_\infty\}$, and
tends to $G$ off $0$.  The computation of (a), with one regularization formed with $\delta$ and
the other with $\delta'$, each Wick ordered at its own variance, then gives
$\E[Q^{(\kappa)}Q'^{(\kappa')}]\to I_Q$, the common limit of $\E[(Q^{(\kappa)})^2]$ and
$\E[(Q'^{(\kappa')})^2]$, since $I_Q$ involves $G$ only.  Hence
$\|Q-Q'\|_2^2=I_Q-2I_Q+I_Q=0$.

(f)(1) For each $q$, $\int f(x)\,\Phi_\kappa(t,x)(q)\,dx=\langle q,\delta_\kappa*(\delta_t\otimes
f)\rangle$, the integral of the test functions $f(x)\delta_\kappa((t,x)-\cdot)$ being one in
$\mathcal S(\R^2)$; so the regularized monomial is $\Phi(\delta_\kappa*(\delta_t\otimes f))$.  In
the covariance norm $\delta_\kappa*(\delta_t\otimes f)\to\delta_t\otimes f$: the squared distance
is $(2\pi)^{-2}\int|\hat\delta(p/\kappa)-1|^2|\hat f(p_1)|^2(|p|^2+m_0^2)^{-1}dp$, with
$|\hat\delta|\le1$ and the integrable dominator $4|\hat f(p_1)|^2(p_0^2+p_1^2+m_0^2)^{-1}$
(its $p_0$-integral is $4\pi|\hat f(p_1)|^2(p_1^2+m_0^2)^{-1/2}$), so it tends to $0$.  Hence the
limit is $\Phi(t,f)$.

(f)(2) is Lemma~\ref{aux:notation-sharpid} with Remark~\ref{aux:notation-sharpidrem}(a), as
stated; it is not proved again here.
\end{proof}

\subsection*{The identity}

Throughout this subsection $h_1,h_2\in\CE$ are fixed, and for $\alpha\in[0,1]$
\[
 h_\alpha:=h_1+\alpha(h_2-h_1)=(1-\alpha)h_1+\alpha h_2,\qquad
 \Uc_{h_2-h_1}:=\Uc_{h_2}-\Uc_{h_1} .
\]
The symbol $h_\alpha$ is used only with the letter $\alpha$: the numerals in $h_1,h_2$ always
denote the two given cutoffs, and the path equals $h_1$ at $\alpha=0$ and $h_2$ at $\alpha=1$.
At $h_1=h^B$, $h_2=h$ it is the path $h_\alpha=h^B+\alpha(h-h^B)$ of \S\ref{sec:notation}.
$\Uc_{h_2-h_1}$ lies in $L^q(\dmu)$ for every $q<\infty$ (Step~2 below); it is an
\emph{observable}, the covariance partner of $F$, and the function $h_2-h_1$, which may take
both signs, is nowhere used as a cutoff: every measure in this section is $dq_{h_\alpha}$
with $h_\alpha\in\CE$.

\begin{dlemma}[affine-path identity]\label{lem:D1}
Let $h_1,h_2\in\CE$, and let $F$ be a finite sum of localized monomials of type $(1.1.5)$ or
$(1.1.6)$, with kernels in $L^2$ and $n_i<\bar n$, that is, $F\in\mathfrak F$
(Definition~\ref{aux:d1-obs}). Then $h_\alpha\in\CE$ for every $\alpha\in[0,1]$; the functions
$F$, $\Uc_{h_2-h_1}$ and $F\,\Uc_{h_2-h_1}$ are $dq_{h_\alpha}$-integrable for every
$\alpha\in[0,1]$; the map $\alpha\mapsto\ev{F}{h_1+\alpha(h_2-h_1)}$ is $C^1$ on $[0,1]$
(one-sided derivatives at the endpoints), with
\begin{equation}\label{eq:d1-deriv}
 \frac{d}{d\alpha}\ev{F}{h_\alpha}=-\Cov_{h_\alpha}\bigl(F,\Uc_{h_2-h_1}\bigr);
\end{equation}
and
\begin{equation}\label{eq:d1-identity}
 \ev{F}{h_2}-\ev{F}{h_1}=-\int_0^1\Cov_{h_\alpha}\bigl(F,\Uc_{h_2-h_1}\bigr)\,d\alpha .
\end{equation}
\end{dlemma}

Here $\ev{F}{h}=\int F\,dq_h$ and $\Cov_h(F,G)=\ev{FG}{h}-\ev{F}{h}\ev{G}{h}$; the measures
$dq_h=Z(h)^{-1}e^{-\Uc_h}\dmu$ are those of Assumption~A3 of the record.

\emph{Status of Lemma~\ref{lem:D1}.} \emph{Proved, conditional on repair R8 of the V3 audit.}
Step~1 of the proof needs only that $h\mapsto\Uc_h$ is linear, for the record's $\Uc_h$, the
object of Assumption~A3.  The record states no construction of $\Uc_h$ (V3 audit, finding M-3).
Repair R8, as the V3 report writes it (its \S9; Simon's book was not read by this drafter),
defines $\Uc_h:=\lambda\sum_na_n{:}\Phi^n(h){:}$ with ${:}\Phi^n(g){:}$ fixed by pairings that
are linear in $g$, so the linearity holds under R8.  The route written in Step~1 obtains it
instead from the record's identification of $\Uc_h$ with the limit in $L^2(\dmu)$ of its
regularizations $\Uc_{h,\kappa}$, which the record asserts (Appendix~C, notation paragraph,
\cite[p.~39]{R}) without proof (V3 audit, item m-1) and which Lemma~C.4 of the record
\cite[p.~40]{R} gives under the repairs R1--R10; on that route Lemma~\ref{lem:D1} is proved
conditional on Lemma~C.4.  Step~1 proves from Lemma~\ref{aux:d1-moments}(a) alone that the
regularizations converge in $L^2(\dmu)$ and that the limit is linear in $h$; the condition is
what identifies that limit with the record's $\Uc_h$, the object that Assumption~A3 imports with
Theorem~V.7 of Simon's book as its source \cite[p.~8]{R}, and that the notation paragraph's
``so that'' asserts, without defining it, to be the limit.  If $\Uc_h$ is instead defined as
that limit (the reading of the last paragraph of the proof of
Lemma~\ref{aux:notation-cutoffs}), Step~1 needs no condition, and the identification moves into
the import of Assumption~A3.  Every other step is proved here from Assumptions~A2 and~A3.

\begin{proof}
\emph{Step 0: the path stays in $\CE$.} For $\alpha\in[0,1]$, $h_\alpha$ is measurable;
$0\le h_\alpha\le1$ pointwise, as a convex combination of two numbers in $[0,1]$; and, since
$h_1,h_2\ge0$, $\{h_\alpha\ne0\}\subseteq\{h_1\ne0\}\cup\{h_2\ne0\}$, so
$\supp h_\alpha\subseteq\Lambda:=\supp h_1\cup\supp h_2$, which is compact ($\supp$ the
closure of the set of non-zeros, as in \S\ref{sec:notation}; the inclusion holds equally for
essential supports). By Definition~2.1 of the record \cite[p.~7]{R} ($\CE$: measurable
$h:\R^2\to[0,1]$ with compact support), $h_\alpha\in\CE$. For $h_1=h^B$, $h_2=h$ this is
Lemma~\ref{aux:notation-cutoffs}(ii) of \S\ref{sec:notation}.

\emph{Step 1: linearity, and representatives.} The record identifies $\Uc_h$, $h\in\CE$,
with the limit in $L^2(\dmu)$ of $\Uc_{h,\kappa}=\lambda\int h\,{:}\mathscr
P(\Phi_\kappa){:}_{c_\kappa}dz$ (Appendix~C, notation paragraph, \cite[p.~39]{R}: ``so that
$\Uc_{h,\kappa}\to\Uc_h$ in every $L^p(d\mu_0)$''; Lemma~C.4, \cite[p.~40]{R}, quantifies the
convergence in $L^2$) and gives no other construction of $\Uc_h$; the V3 audit of Appendix~C
records this as its finding M-3. For every bounded measurable $w$ of compact support, of
either sign, $\lambda\int w\,{:}\mathscr P(\Phi_\kappa){:}_{c_\kappa}dz$ converges in
$L^2(\dmu)$ as $\kappa\to\infty$: write $w=\sum_\Delta w\mathbf 1_\Delta$ over the finitely many
squares meeting its support, and apply Lemma~\ref{aux:d1-moments}(a) with $r=1$ to each kernel
$\lambda a_nw\mathbf 1_\Delta$, $0\le n\le2p'$. The regularized integral is linear in $w$
(pointwise in $q$), hence so is the limit. With $w=h_\alpha,h_1,h_2$ and
$h_\alpha-(1-\alpha)h_1-\alpha h_2=0$,
\begin{equation}\label{eq:d1-affine}
 \Uc_{h_\alpha}=(1-\alpha)\Uc_{h_1}+\alpha\Uc_{h_2}=\Uc_{h_1}+\alpha\,\Uc_{h_2-h_1}
 \qquad\dmu\text{-a.e., for each }\alpha\in[0,1],
\end{equation}
where $\Uc_{h_2-h_1}$, the limit for $w=h_2-h_1$, equals $\Uc_{h_2}-\Uc_{h_1}$. This is the
linearity of $h\mapsto\Uc_h$ recorded in the proof of Lemma~3.1 of the record
\cite[p.~15]{R}; for $h_1=h^B$, $h_2=h$ it is Lemma~\ref{aux:notation-cutoffs}(iv) of
\S\ref{sec:notation}, proved there from Lemma~C.4.  Under repair R8 of the V3 audit, which
defines $\Uc_h$ linearly in $h$, \eqref{eq:d1-affine} holds directly (the status paragraph after
Lemma~\ref{lem:D1}).
The null set in \eqref{eq:d1-affine} may depend on $\alpha$. Fix therefore
finite real-valued representatives $v$ of $\Uc_{h_1}$ and $d$ of $\Uc_{h_2-h_1}$, and put
$u_\alpha:=v+\alpha d$ pointwise, for all $\alpha\in[0,1]$ at once. By
\eqref{eq:d1-affine}, $u_\alpha$ represents $\Uc_{h_\alpha}$ for every $\alpha$; $u_0=v$
represents $\Uc_{h_1}$ and $u_1=v+d$ represents $\Uc_{h_2}$. Every pointwise statement below
concerns these representatives.

\emph{Step 2: integrability, with the exponents.} Assumption~A3 of the record
\cite[p.~8]{R} reads: ``For every $h\in\CE$: $\Uc_h\in L^p(d\mu_0)$ for every $p<\infty$, and
$e^{-\varrho\,\Uc_h}\in L^1(d\mu_0)$ for every finite $\varrho>0$. In particular
$0<Z(h):=\int e^{-\Uc_h}d\mu_0<\infty$''. We use it at $h_1$, $h_2$ and $h_\alpha$, all in
$\CE$ by Step~0.
\begin{itemize}
\item[(I1)] $F\in L^2(\dmu)$, by Lemma~\ref{aux:d1-moments}(b).
\item[(I2)] $d\in L^q(\dmu)$ for every $q<\infty$: A3 gives $\Uc_{h_1},\Uc_{h_2}\in L^q$, and
$\|d\|_q\le\|\Uc_{h_1}\|_q+\|\Uc_{h_2}\|_q$.
\item[(I3)] $W:=e^{-u_0}+e^{-u_1}\in L^q(\dmu)$ for every $q<\infty$: A3 with $\varrho:=q$ gives
$\|e^{-\Uc_{h_i}}\|_q^q=\int e^{-q\,\Uc_{h_i}}d\mu_0<\infty$ for $i=1,2$, and Minkowski's
inequality.
\item[(I4)] For every point and every $\beta\in[0,1]$, by the convexity of the exponential
function, $e^{-u_\beta}=e^{-(1-\beta)u_0-\beta u_1}\le(1-\beta)e^{-u_0}+\beta e^{-u_1}\le W$.
\item[(I5)] H\"older's inequality with the exponents $(2,4,4)$, where
$\frac12+\frac14+\frac14=1$, and with the exponents $(2,2)$:
\[
 \|F\,d\,W\|_1\le\|F\|_2\|d\|_4\|W\|_4,\quad
 \|F\,W\|_1\le\|F\|_2\|W\|_2,\quad \|d\,W\|_1\le\|d\|_2\|W\|_2,
\]
and $\|W\|_1\le\|W\|_2$.
\end{itemize}
Put $Z(\alpha):=\int e^{-u_\alpha}d\mu_0=Z(h_\alpha)$ and $N(\alpha):=\int F\,e^{-u_\alpha}d\mu_0$.
By (I4), $|F|e^{-u_\alpha}\le|F|W$, $|d|e^{-u_\alpha}\le|d|W$ and
$|F\,d|e^{-u_\alpha}\le|F||d|W$, all in $L^1(\dmu)$ by (I5); dividing by $Z(h_\alpha)>0$ (A3)
shows that $F$, $d=\Uc_{h_2-h_1}$ and $F\,d$ are $dq_{h_\alpha}$-integrable for every $\alpha$,
and $\ev{F}{h_\alpha}=N(\alpha)/Z(\alpha)$.

\emph{Step 3: differentiation under the integral.} For $\alpha,\alpha+\varepsilon\in[0,1]$,
$\varepsilon\ne0$, pointwise
\begin{equation}\label{eq:d1-quot}
 \frac{e^{-u_{\alpha+\varepsilon}}-e^{-u_\alpha}}{\varepsilon}
 =-\,d\int_0^1e^{-u_{\alpha+\theta\varepsilon}}\,d\theta ,
\end{equation}
since $u_{\alpha+\theta\varepsilon}=u_\alpha+\theta\varepsilon d$ and
$\frac{d}{d\theta}e^{-u_\alpha-\theta\varepsilon d}=-\varepsilon\,d\,
e^{-u_{\alpha+\theta\varepsilon}}$. By (I4) the modulus of \eqref{eq:d1-quot} is at most
$|d|\,W$, and \eqref{eq:d1-quot} tends to $-d\,e^{-u_\alpha}$ pointwise as $\varepsilon\to0$.
Multiply by $F$, resp.\ by $1$, and apply dominated convergence with the dominator
$|F||d|W$, resp.\ $|d|W$, integrable by (I5): $N$ and $Z$ are differentiable on $[0,1]$
(one-sided at the endpoints), with
\[
 N'(\alpha)=-\int F\,d\,e^{-u_\alpha}\,d\mu_0,\qquad Z'(\alpha)=-\int d\,e^{-u_\alpha}\,d\mu_0 .
\]
If $\alpha_k\to\alpha$ in $[0,1]$, then $F\,d\,e^{-u_{\alpha_k}}\to F\,d\,e^{-u_\alpha}$ and
$d\,e^{-u_{\alpha_k}}\to d\,e^{-u_\alpha}$ pointwise, with the same dominators; so $N'$ and $Z'$
are continuous, and $N$, $Z$ are $C^1$ on $[0,1]$.

\emph{Step 4: normalization.} $Z(\alpha)=Z(h_\alpha)>0$ for every $\alpha\in[0,1]$, by A3 at
$h_\alpha\in\CE$; $Z$ is continuous on the compact interval $[0,1]$ (Step~3), so
$\min_{[0,1]}Z>0$. (Lemma~C.2 of the record \cite[p.~40]{R} gives the explicit bound
$Z(\alpha)\ge e^{-\lambda|a_0||\Lambda|}$, since $\supp h_\alpha\subseteq\Lambda$ by Step~0; it is
proved conditional on repair R8 of the V3 audit, and this proof does not use it.)

\emph{Step 5: the quotient rule.} By Steps~3 and~4, $\alpha\mapsto\ev{F}{h_\alpha}=N/Z$ is
$C^1$ on $[0,1]$, and, since $Z(\alpha)^{-1}e^{-u_\alpha}d\mu_0=dq_{h_\alpha}$ ($u_\alpha$
represents $\Uc_{h_\alpha}$),
\[
 \frac{d}{d\alpha}\ev{F}{h_\alpha}=\frac{N'}{Z}-\frac{N\,Z'}{Z^2}
 =-\int F\,d\;dq_{h_\alpha}+\int F\,dq_{h_\alpha}\int d\;dq_{h_\alpha}
 =-\Cov_{h_\alpha}\bigl(F,\Uc_{h_2-h_1}\bigr),
\]
which is \eqref{eq:d1-deriv}.

\emph{Step 6: integration.} The fundamental theorem of calculus for the $C^1$ function
$\alpha\mapsto\ev{F}{h_\alpha}$ on $[0,1]$, whose values at $\alpha=0$ and $\alpha=1$ are
$\ev{F}{h_1}$ and $\ev{F}{h_2}$ ($u_0$ and $u_1$ represent $\Uc_{h_1}$ and $\Uc_{h_2}$), gives
\eqref{eq:d1-identity}.
\end{proof}

\begin{auxrem}[what the proof uses, and what it does not]\label{aux:d1-uses}
(a) \emph{The class of $F$.} $F$ enters only through $F\in L^2(\dmu)$. For any
$F\in L^{p_0}(\dmu)$ with $1<p_0<\infty$ the proof runs verbatim with the exponents
$(p_0,2p_0',2p_0')$, $p_0':=p_0/(p_0-1)$, in place of $(2,4,4)$ in (I5)
($\frac1{p_0}+\frac1{2p_0'}+\frac1{2p_0'}=1$), and $(p_0,p_0')$ in place of $(2,2)$; so
Lemma~\ref{lem:D1} holds for every $F\in\bigcup_{p_0>1}L^{p_0}(\dmu)$. Neither $n_i<\bar n$,
nor the localization squares, nor the GJS norm enters.

(b) \emph{No constant, no uniformity in $h$.} Neither Theorem~C.5 nor Lemma~C.2 of the
record nor Assumption~A4 is used: the dominator involves only the endpoints $h_1,h_2$, at which
A3 applies, and positivity of $Z$ on the path comes from A3 and continuity (Step~4).
Lemma~\ref{lem:D1} is qualitative and carries no constant. The Gaussian input is computed
explicitly (Lemmas~\ref{aux:d1-wick}, \ref{aux:d1-cov} and~\ref{aux:d1-moments}); it does
not use the chaos structure of $\Uc_h$ or the Wick isometry for $\Uc_h$ that the V3 audit found
unimported in Appendix~C of the record (its finding M-3).  The proof is not independent of
Appendix~C, though: the linearity used in Step~1 is the path-space face of M-3, and it rests on
repair R8 or, on the route written in Step~1, on Lemma~C.4 (the status paragraph after
Lemma~\ref{lem:D1}).

(c) \emph{The partner is an observable.} $\Uc_{h_2-h_1}$ is an element of
$\bigcap_{q<\infty}L^q(\dmu)$, built from $\Uc_{h_1}$ and $\Uc_{h_2}$. A3 is applied only at
$h_1$, $h_2$ and $h_\alpha$, all in $\CE$; Step~1 takes the regularized limit for the kernel
$h_2-h_1$ as for any observable; and no statement about measures is applied at $h_2-h_1$,
for which, as Remark~C.6 of the record states \cite[p.~41]{R}, the Wick lower bound of
Lemma~C.3 would fail when $h_2-h_1$ takes both signs. In the application in
Lemma~\ref{lem:D2}, $h_1=h^B$ and $h_2=h$, so that $h_2-h_1=h\,\mathbf 1_{\R^2\setminus B}\ge0$.

(d) \emph{Relation to the record and to GJS.} Proposition~8.1 of the record, Steps~2--3
\cite[pp.~24--25]{R}, runs the same argument for slab cutoffs, with the log-convexity bound
(13) of the record in place of the convexity bound (I4) and with the difference quotient (14)
of the record, which is \eqref{eq:d1-quot}. GJS's (1.1.12) states the identity without the
integrability argument; its third equality, the decomposition of $V(g)$ into the pieces
$V(g_i)$, is the business of Lemma~\ref{lem:D2}.
\end{auxrem}

\begin{auxrem}[a check of the bookkeeping]\label{aux:d1-check}
In the one-variable model in which $\dmu$ is replaced by the standard Gaussian law of a real
variable $X$, $\Uc_h$ by $hX^2/2$ for $h\in[0,1]$, $F=X^2$, $h_1=0$ and $h_2=1$, one has
$dq_h=N(0,(1+h)^{-1})$, $\ev{F}{h}=(1+h)^{-1}$ and
$\Cov_h(X^2,X^2/2)=\frac12\operatorname{Var}_h(X^2)=(1+h)^{-2}$, so that
$-\int_0^1(1+\alpha)^{-2}d\alpha=-\frac12=\ev{F}{h_2}-\ev{F}{h_1}$, as \eqref{eq:d1-identity}
asserts. This tests the signs and the normalization only; it is not part of the proof.

\emph{Numerical record} (\emph{numerical}; nothing in this section is deduced from it).  The script
\texttt{parts/check\_d1.py}, run on 3 November 2026 with the system \texttt{python3} 3.14.7, numpy
2.5.2 and scipy 1.18.1 (output and run header in \texttt{parts/}\allowbreak\texttt{check\_d1.output.txt}), checks: the moment
formula of Lemma~\ref{aux:d1-wick} against Isserlis's theorem in exact rational arithmetic, for
twelve degree patterns with random rational covariance matrices (seed 20261102), all equal; the
bound $G\le M_0$ of Lemma~\ref{aux:d1-cov}(ii) on 4001 log-spaced points
$|\zeta|\in[10^{-10},10^3]$ for six values of $m_0\in[0.05,10]$, with margin at least $0.319$; the
formula of Lemma~\ref{aux:d1-cov}(i) against the value $K_0(m_0|\zeta|)/(2\pi)$ at three points, to
twelve digits (default tolerances of \texttt{scipy.integrate.quad}); and the model of this remark by
quadrature.
\end{auxrem}


\section{Localized comparison}\label{sec:d2}
% d2.tex -- drafted by aqft (council member of ugjs/PLAN.md), 2 November 2026 (begun 31 October 2026);
% revised 3 November 2026 after the audits V1 (cqft, mp) and V2, and 4 November after V1b/V2b (Step 3 and the
% status line now rest on Lemma C.4, R1--R10): status line; Step 3 through Step 1
% of the proof of Lemma D.1; the reading of the source named in Step 4; form (b)
% described as a form of the plan's statement, and F-T attached to its uncontrolled orientations.
% Section 4: Lemma D.2 (localized comparison).  Merged with \S2 (notation.tex, cqft2) on
% 2 November 2026: squares, boxes, the path h_alpha, the charged squares, strips, U_{Delta,n} and
% c_U are those of \S2 and are not restated here.  Auxiliary statements of this part:
% aux:d2-widest (attainment of the strip width), aux:d2-int (square integrability), aux:d2-uses.

This section proves Lemma~\ref{lem:D2}: the expectation of a localized observable under $dq_h$
differs from its expectation under the cutoff restricted to a box of squares containing the
observable's squares by a sum of clustering terms, one for each square outside the box, and
each term is bounded by GJS Theorem~1.1.7.  The device is the source's own: in the proof of its
(1.1.11), GJS write the difference of two expectations as an integral along an affine path of
cutoffs of a sum over unit squares and state ``We apply Theorem 1.1.7 to each term in the sum to
obtain'' \cite[p.~595]{GJS}.

\subsection{Standing assumption and notation}\label{sub:d2-conv}

\emph{Standing assumption.}  In \S\S\ref{sec:d2}--\ref{sec:d4} and \S\ref{sec:cons},
$0\le\lambda/m_0^2<\varepsilon_{\mathscr P}$, and we work in the working coordinates of \S2.6 of
the record \cite[p.~14]{R} (Remark~3.3 of the record \cite[p.~16]{R} transports the conclusions):
the pair $(m_0,\lambda)$ satisfies the hypotheses printed for GJS Theorem~1.1.7, ``Let $m_0$,
$K_2$ and $K_1$ be sufficiently large and let $\lambda$ be sufficiently small''
\cite[p.~594]{GJS}, and $m_1>0$, $C_7$, $K_1$ and $K_2\ge1$ are the printed constants of
Assumption~A8(iii)--(iv) of the record \cite[p.~10]{R} (\S\ref{sec:notation-printed}).  The
interaction polynomial is $\mathscr P(\xi)=\sum_{n=0}^{2p'}a_n\xi^n$, $a_{2p'}>0$, $p'\ge1$, and
$\bar n:=2p'+1$ (Assumption~A8(ii) of the record).  For $h\in\CE$,
$\Uc_h=\lambda\int h\,{:}\mathscr P(\Phi){:}\,dz$, $Z(h)=\int e^{-\Uc_h}\dmu$,
$dq_h=Z(h)^{-1}e^{-\Uc_h}\dmu$, $\ev{F}{h}=\int F\,dq_h$ (Assumption~A3 of the record
\cite[p.~8]{R}) and $\Cov_h(F,G)=\ev{FG}{h}-\ev{F}{h}\ev{G}{h}$.

\emph{Notation.}  We use the notation of \S\ref{sec:notation}: the squares $\Delta_{l,k}$, rows
and columns (Definition~\ref{aux:notation-squares}); boxes $B$, $|B|$, the restriction $h^B$,
the path $h_\alpha$ and the set $\mathfrak S(h)$ of squares charged by $h$
(Definition~\ref{aux:notation-box}); strips, separation, the parallel and the transverse
orientation, and $d(\Delta,R)$ (Definition~\ref{aux:notation-stripdef}); $U_{\Delta,n}$,
$U_\Delta$ and $c_U$ (Definition~\ref{aux:notation-UDelta}, Lemma~\ref{aux:notation-cU}).  For
$h\in\CE$ and a box $B$ we write
\[
 \mathcal D(h,B):=\{\Delta\in\mathfrak S(h):\ \Delta\cap B=\varnothing\}
\]
for the squares outside $B$ charged by $h$; it is finite, and each of its squares meets
$\supp h$ (Lemma~\ref{aux:notation-cutoffs}(iii)).

\begin{auxlemma}[the widest strip exists]\label{aux:d2-widest}
Let $X,Y\subseteq\R^2$ be bounded and non-empty, and suppose that some strip separates $X$
from $Y$.  Then the supremum of the widths of the strips separating $X$ from $Y$ is attained.
In particular, for two localized monomials whose localization squares are separated by some
strip, ``the widest strip bounded by two parallel lines, and separating'' them
\cite[p.~594]{GJS} exists, its width is the supremum $d$ of
Definition~\ref{aux:notation-stripdef}, and every separating strip has width at most $d$.
\end{auxlemma}

Definition~\ref{aux:notation-stripdef} defines the width $d$ as a supremum; the lemma shows that
it is attained, as the wording of p.~594 presupposes.

\begin{proof}
For a unit vector $\nu$, $c\in\R$ and $w\ge0$, the strip bounded by $\{\nu\cdot z=c\}$ and
$\{\nu\cdot z=c+w\}$ has the closed half-planes $H_1=\{\nu\cdot z\le c\}$ and
$H_2=\{\nu\cdot z\ge c+w\}$ (Definition~\ref{aux:notation-stripdef}).  Put
$a(\nu):=\sup_{x\in X}\nu\cdot x$ and $b(\nu):=\inf_{y\in Y}\nu\cdot y$.  Then $X\subseteq H_1$
and $Y\subseteq H_2$ if and only if $a(\nu)\le c$ and $c+w\le b(\nu)$; the other assignment of
the half-planes corresponds to $-\nu$.  So a strip with normal $\nu$, oriented from $X$ to
$Y$, separates them exactly when $b(\nu)-a(\nu)\ge0$, its width is then at most
$b(\nu)-a(\nu)$, and the choice $c=a(\nu)$, $w=b(\nu)-a(\nu)$ attains this.  Hence the
supremum of the separating widths is $\sup_{|\nu|=1}(b(\nu)-a(\nu))$, which is $\ge0$ by
hypothesis.  With $\varrho:=\sup\{|z|:z\in X\cup Y\}<\infty$,
$|a(\nu)-a(\nu')|\le\varrho|\nu-\nu'|$ and $|b(\nu)-b(\nu')|\le\varrho|\nu-\nu'|$, so
$\nu\mapsto b(\nu)-a(\nu)$ is continuous on the unit circle and attains its supremum at some
$\nu_*$; the strip with normal $\nu_*$, $c=a(\nu_*)$ and $w=b(\nu_*)-a(\nu_*)$ separates and has
the maximal width.
\end{proof}

\begin{auxlemma}[square integrability of the pieces]\label{aux:d2-int}
Let $h\in\CE$.  Every $X\in L^4(\dmu)$ lies in $L^2(dq_h)$.  On $L^2(dq_h)$ the covariance
$\Cov_h$ is defined, bilinear and symmetric, and $\Cov_h(X,c)=0$ for constants $c$.  Every finite
sum of localized monomials of type (1.1.5) or (1.1.6) with kernels in $L^2$ and individual degrees
$n_i<\bar n$ lies in $L^4(\dmu)$, hence in $L^2(dq_h)$; so do, in particular, every piece $F_a$
of Lemma~\ref{lem:D2} below and every $U_{\Delta,n}$, $1\le n\le2p'$.
\end{auxlemma}

This also supplies the hypothesis ``$F$ and the finitely many $U_{\Delta,n}$ involved lie in
$L^2(dq_{h'})$'' of Lemma~\ref{aux:notation-cU}(iv), which \S\ref{sec:notation} refers to the
chain of Lemma~\ref{lem:D1}; that chain gives the $dq_h$-integrability of $F$ and of
$\Uc_{h_2-h_1}$, not of a single $U_{\Delta,n}$.

\begin{proof}
By Assumption~A3 of the record \cite[p.~8]{R}, $Z(h)>0$ and $e^{-2\Uc_h}\in L^1(\dmu)$, i.e.\
$e^{-\Uc_h}\in L^2(\dmu)$.  H\"older's inequality with exponents $(2,2)$ gives
\[
 \int X^2e^{-\Uc_h}\dmu\le\|X^2\|_{L^2(\dmu)}\|e^{-\Uc_h}\|_{L^2(\dmu)}
 =\|X\|_{L^4(\dmu)}^2\|e^{-\Uc_h}\|_{L^2(\dmu)}<\infty ;
\]
divide by $Z(h)$.  For $X,Y\in L^2(dq_h)$, $XY\in L^1(dq_h)$
by the Cauchy--Schwarz inequality, and $L^2(dq_h)\subseteq L^1(dq_h)$ because $dq_h$ is a
probability measure; so $\Cov_h(X,Y)$ is defined, and it is bilinear and symmetric by the
linearity of the integral; for a constant $c$, $\ev{Xc}{h}-\ev{X}{h}c=0$.  The membership in
$L^4(\dmu)$ is Lemma~\ref{aux:d1-moments}(b) of \S\ref{sec:d1} (``$Q\in L^q(\dmu)$ for every
$q<\infty$, and so is every $F\in\mathfrak F$''); $U_{\Delta,n}$ is a localized monomial of type
(1.1.6) with kernel $\lambda a_nh\mathbf 1_\Delta\in L^2(\R^2)$ and degree $n<\bar n$
(Lemma~\ref{aux:notation-cU}(i)).
\end{proof}

\subsection{The lemma}\label{sub:d2-lemma}

\begin{dlemma}[localized comparison]\label{lem:D2}
Assume the standing assumption of \S\ref{sub:d2-conv}.  Let $F=\sum_{a\in\mathfrak A}F_a$ be a
finite sum of localized monomials $F_a$ of type (1.1.5) or (1.1.6) \cite[p.~593]{GJS}, with
kernels in $L^2$ and every individual degree $n_i<\bar n$, and put
$\GJS{F}:=\sum_{a\in\mathfrak A}\GJS{F_a}$, the norm (1.1.9) of this decomposition.  Let $R$
be a non-empty finite set of squares containing every localization square of every $F_a$, let
$B\supseteq R$ be a box, let $h\in\CE$, and let $\mathcal D:=\mathcal D(h,B)$.
\begin{enumerate}
\item[\rm(a)] Suppose that to every $\Delta\in\mathcal D$ a strip $S_\Delta$ of positive width
$w_\Delta$ separating $\Delta$ from $\bigcup R$ is assigned.  Then
\[
 \bigl|\ev{F}{h}-\ev{F}{h^B}\bigr|\le C_7\,c_U\,\GJS{F}\sum_{\Delta\in\mathcal D}e^{-m_1w_\Delta} .
\]
Every estimate behind this bound is Source Fact~2.2 of the record (GJS (1.1.10)) for a pair
$(F_a,U_{\Delta,n})$, across the strip $S_\Delta$, at a cutoff $h_\alpha\in\CE$; the
orientation of $S_\Delta$ is the orientation of that application.  If $S_\Delta$ is transverse
to the time-zero line and $F_a$ is of type (1.1.5), that application is the transverse case of
GJS Theorem~1.1.7, used at the printed statement's scope and disclosed by F-T
(\S\ref{sec:purpose}).
\item[\rm(b)] (A form of the statement of \texttt{ugjs/PLAN.md}, Part~1, with positive widths;
the plan's form does not require positive width.)  If every square $\Delta\not\subseteq B$
meeting $\supp h$ is separated from $\bigcup R$ by a strip of positive width, then
\[
 \bigl|\ev{F}{h}-\ev{F}{h^B}\bigr|\le C_7\,c_U\,\GJS{F}
 \sum_{\Delta\not\subseteq B,\ \Delta\cap\supp h\ne\varnothing}e^{-m_1d(\Delta,R)} .
\]
The widest strips used here are not controlled in orientation
(Remark~\ref{aux:d2-uses}(2)); where one is transverse to the time-zero line and some $F_a$
is of type (1.1.5), form (b) uses the transverse case of GJS Theorem~1.1.7 at the printed
statement's scope (F-T, \S\ref{sec:purpose}).
\end{enumerate}
The constants $C_7$ and $m_1$ are those of Theorem~1.1.7 in working coordinates (Remark~3.3 of
the record), and $c_U$ is that of Lemma~\ref{aux:notation-cU}; none depends on $h$, $B$, $R$ or
$F$.
\end{dlemma}

\emph{Status of Lemma~\ref{lem:D2}.}  Proved, conditional on Lemma~C.4 of the record (repairs
R1--R10 of the V3 audit, which include R8).  Lemma~C.4 supplies, in Step~3, the identification of
$\Uc_{h-h^B}=\Uc_h-\Uc_{h^B}$ with the $L^2(\dmu)$-limit of its regularizations, and it supplies
the condition of Lemma~\ref{lem:D1} on the route written in Step~1 of the proof of that lemma (the
status paragraph after Lemma~\ref{lem:D1}).  Repair R8 alone makes $h\mapsto\Uc_h$ linear but
does not give that identification as this note is written.  Lemma~\ref{lem:D2} uses neither
Lemma~C.2 nor Theorem~C.5 of the record.  It applies GJS Theorem~1.1.7 to the localized monomials
of Definition~\ref{aux:d1-obs}; that these are the monomials of p.~593, and that the record's
$\Uc_h$ is GJS's $V(h)$, is the reading of the source of Remark~\ref{aux:notation-gjslattice}(iv)
(Remark~\ref{aux:d1-def}; Step~4).  The orientation of each application is that of the assigned
strip, and a transverse strip with a sharp-time $F_a$ falls under F-T.

\begin{proof}
\emph{Step 1: the path.}  By Lemma~\ref{aux:notation-cutoffs}(i)--(ii), $h^B\in\CE$ and
$h_\alpha=h^B+\alpha(h-h^B)\in\CE$ for every $\alpha\in[0,1]$, with $h_0=h^B$ and $h_1=h$.

\emph{Step 2: the identity.}  $F$ is a finite sum of localized monomials of type (1.1.5) or
(1.1.6) with kernels in $L^2$ and $n_i<\bar n$, which is the class of Lemma~\ref{lem:D1}.
Lemma~\ref{lem:D1} with $h_1=h^B$, $h_2=h$ gives that $\alpha\mapsto\ev{F}{h_\alpha}$ is $C^1$
on $[0,1]$ and
\[
 \ev{F}{h}-\ev{F}{h^B}=-\int_0^1\Cov_{h_\alpha}\bigl(F,\Uc_{h-h^B}\bigr)\,d\alpha ;
\]
the covariance is taken under $dq_{h_\alpha}$, $h_\alpha=h^B+\alpha(h-h^B)$, and
$\Uc_{h-h^B}=\Uc_h-\Uc_{h^B}$ is the observable of Lemma~\ref{lem:D1}.  Here
$h-h^B=h\mathbf 1_{\R^2\setminus B}\ge0$ lies in $\CE$ (Lemma~\ref{aux:notation-cutoffs}(i)); it
enters only as the observable $\Uc_{h-h^B}$, not as a cutoff.

\emph{Step 3: decomposition into squares.}  Put $w:=h-h^B$, a bounded measurable function with
compact support, and, for $\kappa\ge2$, $\Uc^{(\kappa)}_w:=\lambda\int w\,{:}\mathscr
P(\Phi_\kappa){:}_{c_\kappa}dz$, pointwise in $q$, with the record's regularization of
\S\ref{sec:d1}.  For each $q$, $\Uc^{(\kappa)}_w(q)$ is the integral of $w$ against a function
continuous in $z$ (Definition~\ref{aux:d1-obs}: $z\mapsto\Phi_\kappa(z)(q)$ is smooth), so it is
linear in $w$ and unchanged when $w$ is changed on a null set.  In particular
$\Uc^{(\kappa)}_w=\Uc_{h,\kappa}-\Uc_{h^B,\kappa}$, the difference of the record's regularized
interactions (Appendix~C, notation paragraph \cite[p.~39]{R}).  Lemma~C.4 of the record
\cite[p.~40]{R} reads: ``There is $c_2=c_2(\lambda,\mathscr P,\Lambda)$ such that for all
$h\in\CE$ with $\supp h\subseteq\Lambda$ and all $\kappa\ge2$,
$\|\mathcal U_h-\mathcal U_{h,\kappa}\|_{L^2(d\mu_0)}^2\le c_2\,\kappa^{-1}(\log\kappa)^{2p'}$.''
Its hypotheses hold for $h$ and for $h^B$ with the compact set $\Lambda:=\supp h$ (both lie in
$\CE$ and have support in $\Lambda$, Lemma~\ref{aux:notation-cutoffs}(i)) and every $\kappa\ge2$;
it is used conditional on the repairs R1--R10 of the V3 audit.  Hence $\Uc_{h,\kappa}\to\Uc_h$ and
$\Uc_{h^B,\kappa}\to\Uc_{h^B}$ in $L^2(\dmu)$, and
$\Uc^{(\kappa)}_w\to\Uc_h-\Uc_{h^B}=\Uc_{h-h^B}$.  (Step~1 of the proof of Lemma~\ref{lem:D1}
shows from Lemma~\ref{aux:d1-moments}(a) alone that $\Uc^{(\kappa)}_w$ converges; Lemma~C.4
identifies the limit with the record's $\Uc_h-\Uc_{h^B}$.  This identification is the condition
of Lemma~\ref{lem:D2}.)  By Lemma~\ref{aux:notation-cutoffs}(iii),
$w=\sum_{\Delta\in\mathcal D}h\mathbf 1_\Delta$ almost everywhere; with
${:}\mathscr P(y){:}_c=\sum_na_n{:}y^n{:}_c$ and ${:}y^0{:}_c=1$ the linearity in $w$ gives,
pointwise in $q$,
\[
 \Uc^{(\kappa)}_w=\sum_{\Delta\in\mathcal D}\Bigl(\lambda a_0\int_\Delta h\,dz
 +\sum_{n=1}^{2p'}U^{(\kappa)}_{\Delta,n}\Bigr),\qquad
 U^{(\kappa)}_{\Delta,n}:=\lambda a_n\int h\mathbf 1_\Delta\,{:}\Phi_\kappa^n{:}_{c_\kappa}\,dz .
\]
$U^{(\kappa)}_{\Delta,n}$ is the regularized monomial of Definition~\ref{aux:d1-obs}(a) with $r=1$,
degree $n$ and kernel $\lambda a_nh\mathbf 1_\Delta$, so by Lemma~\ref{aux:d1-moments}(a) it
converges in $L^2(\dmu)$ to the localized monomial $U_{\Delta,n}$ of
Definition~\ref{aux:notation-UDelta}.  A finite sum of convergent sequences converges to the sum
of the limits; hence
\[
 \Uc_{h-h^B}=\sum_{\Delta\in\mathcal D}\Bigl(U_{\Delta,0}+\sum_{n=1}^{2p'}U_{\Delta,n}\Bigr)
 \quad\text{in }L^2(\dmu),\qquad U_{\Delta,0}=\lambda a_0\int_\Delta h\,dz .
\]
This is the first identity of Lemma~\ref{aux:notation-cU}(iv), which \S\ref{sec:notation} proves
through Lemma~\ref{aux:notation-cutoffs}(iv), also from Lemma~C.4 of the record.  By
Lemma~\ref{aux:d2-int}, $F$, every
$F_a$ and every $U_{\Delta,n}$ lie in $L^2(dq_{h_\alpha})$, where $\Cov_{h_\alpha}$ is bilinear
and vanishes against the constants $U_{\Delta,0}$; therefore
\[
 \Cov_{h_\alpha}\bigl(F,\Uc_{h-h^B}\bigr)
 =\sum_{\Delta\in\mathcal D}\ \sum_{n=1}^{2p'}\ \sum_{a\in\mathfrak A}
  \Cov_{h_\alpha}\bigl(F_a,U_{\Delta,n}\bigr),
\]
a finite sum.

\emph{Step 4: one term, by Source Fact~2.2.}  Fix $\alpha\in[0,1]$, $\Delta\in\mathcal D$,
$1\le n\le2p'$ and $a\in\mathfrak A$.  Source Fact~2.2 of the record \cite[p.~10]{R} reads:
``For localized monomials $Q_1,Q_2$ with disjoint localization squares separated by a strip of
width $d$ bounded by two parallel lines (any orientation),
$\bigl|\int Q_1Q_2\,dq_h-\int Q_1dq_h\int Q_2dq_h\bigr|\le
C_7e^{-m_1d}\|Q_1\|_{\rm GJS}\|Q_2\|_{\rm GJS}$, with $m_1>0$ and $C_7$ independent of $h$,
$Q_1$, $Q_2$ (so stated in the source).''  The source states: ``\textsc{Theorem} 1.1.7.
\emph{Let $m_0$, $K_2$ and $K_1$ be sufficiently large and let $\lambda$ be sufficiently small.
There are constants $0<m_1$ and $O(1)$ independent of $h$, $Q_1$ and $Q_2$ such that}
(1.1.10) $\bigl|\int Q_1Q_2dq_h-\int Q_1dq_h\int Q_2dq_h\bigr|\leqq O(1)e^{-m_1d}\|Q_1\|\,
\|Q_2\|$'' \cite[p.~594]{GJS}, with $d$ the width of the widest separating strip, which exists by
Lemma~\ref{aux:d2-widest}; every separating strip has width at most that $d$, and $m_1>0$, so the
printed bound implies the bound with the width of any separating strip, which is the record's form
(and the form in which \S\ref{sec:notation} applies it, after
Definition~\ref{aux:notation-stripdef}).  We apply it with:
\begin{itemize}
\item $Q_1:=F_a$: a localized monomial of type (1.1.5) or (1.1.6), as given; kernel in $L^2$;
  every $n_i<\bar n$; localization squares in $R$.
\item $Q_2:=U_{\Delta,n}$: a localized monomial of type (1.1.6) with $r=1$, localized in the
  single square $\Delta$, degree $n\le2p'<\bar n$, kernel $\lambda a_nh\mathbf 1_\Delta\in
  L^2(\R^2)$ (Lemma~\ref{aux:notation-cU}(i)).  Page~593 defines ``localized'' for both forms (``For a
  localized monomial $Q$ or $Q(x_0)$''), and Theorem~1.1.7 is stated for ``two localized
  monomials $Q_1$ and $Q_2$'' \cite[p.~594]{GJS}; the mixed pair is covered by the statement.
\item The objects: $Q_1$ and $Q_2$ are the limits of Definition~\ref{aux:d1-obs}, and the measure
  is the record's $dq_{h_\alpha}$.  That these limits are the monomials of p.~593, which GJS print
  as formal integrals, and that the record's $\Uc_h$ is GJS's $V(h)$, is a reading of the
  source, of the kind the record names in Assumption~A8, Scope (d)
  (Remark~\ref{aux:notation-gjslattice}(iv), Remark~\ref{aux:d1-def}); it is not proved here.
\item The strip: $S_\Delta$ separates $\Delta$ from $\bigcup R$, hence from the localization
  squares of $F_a$, which lie in $\bigcup R$; its width is $w_\Delta>0$ and its orientation is the
  one assigned.  Its two closed half-planes are disjoint because $w_\Delta>0$
  (Definition~\ref{aux:notation-stripdef}), so the localization squares of $F_a$ and $\Delta$ are
  disjoint, even as closed squares.  If $S_\Delta$ is transverse and $F_a$ is of type (1.1.5),
  this is the transverse case of F-T (\S\ref{sec:purpose}).
\item The cutoff: $h_\alpha\in\CE$ (Step~1), the class of \cite[p.~589]{GJS}, over which the
  record reads ``independent of $h$'' (Assumption~A8, Scope (d), main.tex ll.~887--889).
\item The printed hypotheses on $m_0$, $K_2$, $K_1$, $\lambda$: the standing assumption.
\end{itemize}
Source Fact~2.2 gives
\[
 \bigl|\Cov_{h_\alpha}(F_a,U_{\Delta,n})\bigr|\le C_7\,e^{-m_1w_\Delta}\,
 \GJS{F_a}\,\GJS{U_{\Delta,n}} .
\]

\emph{Step 5: summation.}  We sum over $a$ and $n$.  By Lemma~\ref{aux:notation-cU}(ii),
$\sum_n\GJS{U_{\Delta,n}}\le c_U$, and $\sum_a\GJS{F_a}=\GJS{F}$; hence
\begin{align*}
 \bigl|\Cov_{h_\alpha}\bigl(F,\Uc_{h-h^B}\bigr)\bigr|
 &\le\sum_{\Delta\in\mathcal D}C_7e^{-m_1w_\Delta}
   \Bigl(\sum_a\GJS{F_a}\Bigr)\Bigl(\sum_n\GJS{U_{\Delta,n}}\Bigr)\\
 &\le C_7\,c_U\,\GJS{F}\sum_{\Delta\in\mathcal D}e^{-m_1w_\Delta} .
\end{align*}
The right side does not depend on $\alpha$; integrating over $\alpha\in[0,1]$ in Step~2 proves
(a).

\emph{Step 6: the form (b).}  Every $\Delta\in\mathcal D$ is a square outside $B$ meeting $\supp h$
(Lemma~\ref{aux:notation-cutoffs}(iii)), so by hypothesis a strip of positive width separates it
from $\bigcup R$, and by Lemma~\ref{aux:d2-widest} a widest one exists; its width is
$d(\Delta,R)>0$.  Assign it as $S_\Delta$ and apply (a); since every term
$e^{-m_1d(\Delta,R)}$ is positive and
$\mathcal D\subseteq\{\Delta\not\subseteq B:\Delta\cap\supp h\ne\varnothing\}$, the sum over
$\mathcal D$ is at most the sum in (b).
\end{proof}

\begin{auxrem}[what the lemma uses, and the orientation]\label{aux:d2-uses}
\emph{(1)} The lemma uses of the cutoff $h$ exactly: $h\ge0$ (so that $h-h^B\ge0$ and the
charged squares are defined by $\int_\Delta h>0$), $h\le1$ (in $c_U$ and in $h_\alpha\in\CE$),
and compact support (finiteness of $\mathcal D$).  It is a statement for $h\in\CE$ only
(trap~20 of the plan); Theorem~1.1.7 and Lemma~\ref{lem:D1} are used at cutoffs of $\CE$ only.
\emph{(2)} In (b) the widest strips are not controlled in orientation: for a square $\Delta$
far to the side of $R$ in a row near that of $R$, the widest separating strip is transverse to
the time-zero line even when a narrower parallel one exists (Remark~\ref{aux:notation-tilted}:
it may also be tilted).  Where such a strip is used with a sharp-time $F_a$, form (b) rests on
the transverse case of Theorem~1.1.7, used at the printed statement's scope and disclosed by F-T
(\S\ref{sec:purpose}).  Theorems~\ref{thm:D3} and~\ref{thm:D4} therefore use (a), with the
strips assigned explicitly and their orientation stated, so that every application of
Theorem~1.1.7 has a named orientation (gate GA1).  Form (b) is recorded because it is a form of
the statement of the plan; it is not used below.
\emph{(3)} The sum in (a) is over the charged squares $\mathcal D(h,B)$, the squares outside $B$
on which $h$ is not almost everywhere zero; squares that meet $\supp h$ only in a null set carry
$U_\Delta=0$ and need no strip.  This matters in \S\ref{sec:d3}.  An example of the drafter of
\S\ref{sec:d1}: for $g=\mathbf 1_{[0,1/10]}+\mathbf 1_{\{3\}}\in\CH$, $\supp g=[0,\frac1{10}]\cup
\{3\}$ meets column~$3$ in one point only, so $N_g=1$; for $I\subseteq(0,\frac1{10})$ the box
$B_s(g,I)$ of \S\ref{sec:d3} has the single column~$0$, and for $T\ge|l_s|+2$ the squares
$\Delta_{l,3}$ with $|l-l_s|\le1$ meet $\supp h_{T,g}\supseteq[-T,T]\times\{3\}$, lie outside $B$,
and are separated from row $l_s$ only by transverse strips (Lemma~\ref{aux:notation-strip}(iv)).
With the plan's index set $\{\Delta\not\subseteq B:\Delta\cap\supp h\ne\varnothing\}$ they would
enter the sum, and the parallel-only claim and the count $2N_g$ of Theorem~\ref{thm:D3} would
fail; they are not charged ($h_{T,g}\mathbf 1_{\Delta_{l,3}}=0$ almost everywhere), so (a) omits
them.
\end{auxrem}


\section{Phase A: fixed cutoff, parallel strips}\label{sec:d3}
% d3.tex -- drafted by aqft (council member of ugjs/PLAN.md), 2 November 2026 (begun 31 October 2026);
% revised 3 November 2026 after the audits V1 (cqft, mp) and V2: status line; the sharp-time field as
% the covariance-norm extension; Lemma aux:d3-cyl(vi) (the pieces as objects of Definition
% aux:d1-obs); the severities of the cell-count finding; Remark aux:d3-orient (slicing); the boundary
% cases N_g = 0 and 1/2 < T < 1; the inputs.
% Section 5: Theorem D.3 (Phase A; fixed g, strips parallel to the time-zero line).  Merged with
% \S2 (notation.tex, cqft2) on 2 November 2026: the stability constant S_2, the column count and
% the bounds (N.a), and the strip lemma are those of \S2.  Auxiliary statements of this part:
% aux:d3-V3, aux:d3-sharp, aux:d3-cyl, aux:d3-cells, aux:d3-geom and the remarks of \S5.4.

Theorem~\ref{thm:D3} bounds the expectation of a sharp-time cylinder observable under the slab
measures $dq_{h_{T,g}}$ uniformly in $T>0$, for one fixed spatial cutoff $g$, with constants that
depend on $g$; that is all that \S7 of the record consumes (\S\ref{sec:cons}).  Every use of GJS
Theorem~1.1.7 in this section is across a strip parallel to the time-zero line.

\subsection{Setting}\label{sub:d3-setting}

The standing assumption of \S\ref{sub:d2-conv} holds.  Let $g\in\CH$ (Definition~2.1 of the record
\cite[p.~7]{R}: $g:\R\to[0,1]$ measurable with compact support) and, for $T>0$,
$h_{T,g}:=\mathbf 1_{[-T,T]}\otimes g$, which lies in $\CE$ by Lemma~3.1(ii) of the record
\cite[p.~15]{R}.  Let $I\subset\R$ be a bounded interval of length $L>0$, and $s\in\R$, with
$l_s$ the row containing time $s$ (Definition~\ref{aux:notation-squares}).  Put
\begin{gather*}
 K_g:=\bigl\{k\in\Z:\ |\supp g\cap[k-\tfrac12,k+\tfrac12)|>0\bigr\},\qquad N_g:=\#K_g ,\\
 K_I:=\bigl\{k\in\Z:\ |I\cap[k-\tfrac12,k+\tfrac12)|>0\bigr\},\qquad N^{\square}_I=\#K_I ,
\end{gather*}
($|\cdot|$ the Lebesgue measure on $\R$; $N^\square_I$ is the column count of
Remark~\ref{aux:notation-counts}), so that $N_g$ is the number of columns meeting $\supp g$ in
positive measure; both sets are finite, and $K_I\ne\varnothing$ because $L>0$.  Let
$k_-:=\min(K_g\cup K_I)$, $k_+:=\max(K_g\cup K_I)$, $N_B:=k_+-k_-+1$, and
\[
 B=B_s(g,I):=\bigcup\{\Delta_{l,k}:\ |l-l_s|\le1,\ k_-\le k\le k_+\},\qquad
 R_s(I):=\{\Delta_{l_s,k}:\ k\in K_I\} .
\]
So $B$ is rows $l_s-1$ to $l_s+1$ times the smallest interval of $N_B$ columns containing every
column that meets $\supp g\cup I$ in positive measure, $|B|=3N_B$, and the squares of $R_s(I)$
lie in $B$.

A \emph{cylinder monomial} (\S6 of the record \cite[p.~19]{R}) is $Q=\prod_{i=1}^r\varphi_0(f_i)$
with $r\ge1$ and real $f_i\in C_c^\infty(\R)$; it is \emph{supported in} $I$ if every $f_i$ is.
At sharp Euclidean time $s$ it is $Q(s):=\prod_{i=1}^r\Phi(s,f_i)$, with
``$\Phi(\tau,f)\equiv\Phi(\delta_\tau\otimes f)=\int\Phi(\tau,\vec x)f(\vec x)d\vec x$''
\cite[p.~589]{GJS}; GJS print no construction, and $\Phi(s,f)$ is read, as in the proof of
Lemma~\ref{aux:notation-hyper}(i), as the extension of $f\mapsto\Phi(f)$ to $\delta_s\otimes f$ by
continuity in the covariance norm, which is finite there (Lemma~\ref{aux:d3-sharp}).  For a finite
real linear combination $F=\sum_{\beta\in\mathfrak B}c_\beta Q_\beta$ of cylinder monomials,
$F(s):=\sum_\beta c_\beta Q_\beta(s)$.

\subsection{Auxiliary facts}\label{sub:d3-aux}

\begin{auxrem}[Lemma~C.2 and Theorem~C.5 after the V3 audit]\label{aux:d3-V3}
This section is the first of the note to consume Lemma~C.2 and Theorem~C.5 of the record
\cite[pp.~40--41]{R}; Theorem~\ref{thm:D4} and the replaced proofs of \S\ref{sec:cons} consume
them too, and from here on they are load-bearing (trap~12 of the plan).  Their status after the
audit V3 (\texttt{ugjs/council/}\allowbreak\texttt{v3-appendix-c-report.md}, 2 November 2026) is
recorded in the introduction of \S\ref{sec:notation}: the audit found them true and every
displayed computation re-derived, with three MAJOR gaps as written (the uncited Bessel facts of
Lemma~C.4, the undefined constant $C_*$, and the Wick-chaos structure of $\Uc_h$ used without an
import), and wrote out repairs R1--R10 that need no new mathematics.  Accordingly, in this note
Lemma~C.2 is \emph{proved, conditional on repair R8 of the V3 audit}; Theorem~C.5, and with it
every statement about $S_2$, is \emph{proved, conditional on the repairs R1--R10 of the V3
audit}; so is Lemma~C.4, which Lemma~\ref{lem:D2} uses in Step~3 of its proof and
\S\ref{sec:notation} uses for the translation covariance of $\Uc_h$
(Lemma~\ref{aux:notation-transl}(iii)) and hence for the shape-only property of $S_2$
(Proposition~\ref{aux:notation-S2}(iii)); and the structural facts of Gaussian Wick powers (the
chaos decomposition, the Wick isometry, $\Gamma(e^{-t})$ on chaoses) are used conditional on
repair R8.  Every later use says ``as in Remark~\ref{aux:d3-V3}''.
\end{auxrem}

\begin{auxlemma}[the sharp-time field under $\dmu$]\label{aux:d3-sharp}
For every $s\in\R$ and every real $f\in C_c^\infty(\R)$, $\delta_s\otimes f$ has finite
covariance norm, so $\Phi(s,f)$ of \S\ref{sub:d3-setting} is defined.  The family
$(\Phi(s,f))_f$ is under $\dmu$ a centred Gaussian family with covariance $\ip{f}{Cf'}$,
$C=(2\mu)^{-1}$, $\mu=(-\partial_x^2+m_0^2)^{1/2}$ (\S2.1 of the record); it therefore has the law
of $(\varphi_0(f))_f$ under $\dphi$.  In particular $\|F(s)\|_{L^p(\dmu)}=\|F\|_{L^p(\dphi)}$ for
every $p\in[1,\infty)$ and every $F$ as in \S\ref{sub:d3-setting}.
\end{auxlemma}

At $s=0$ this is the identification of the sharp-time field used in
Lemma~\ref{aux:notation-hyper}(i); the lemma records it at every time $s$, which
Theorem~\ref{thm:D3} needs.

\begin{proof}
$\dmu$ is the centred Gaussian measure with covariance $(-\Delta_{\R^2}+m_0^2)^{-1}$
(Assumption~A2 of the record \cite[p.~8]{R}).  The Fourier transform of $\delta_s\otimes f$ is
$e^{-ik_0s}\hat f(k)$, and $|e^{-ik_0s}|=1$, so for real $f,f'\in C_c^\infty(\R)$ the covariance
form at $\delta_s\otimes f$ and $\delta_s\otimes f'$ is
\begin{align*}
 \ip{\delta_s\otimes f}{(-\Delta_{\R^2}+m_0^2)^{-1}\,\delta_s\otimes f'}
 &=\int\frac{dk}{2\pi}\,\overline{\hat f(k)}\hat f'(k)\int\frac{dk_0}{2\pi}\,
 \frac1{k_0^2+\mu(k)^2}\\
 &=\int\frac{dk}{2\pi}\,\overline{\hat f(k)}\hat f'(k)\,\frac1{2\mu(k)}=\ip{f}{Cf'} ,
\end{align*}
by the identity $\int\frac{dk_0}{2\pi}(k_0^2+\mu(k)^2)^{-1}=(2\mu(k))^{-1}$ of Remark~2.5 of the
record \cite[p.~14]{R}, whose computation does not involve the value of $s$; the integrals converge
absolutely because $\hat f,\hat f'$ are Schwartz functions and $(2\mu)^{-1}\le(2m_0)^{-1}$.  With
$f=f'$ this shows that the covariance norm of $\delta_s\otimes f$ is finite.  The extension of the
linear isometry $f\mapsto\Phi(f)$ (from test functions with the covariance norm into $L^2(\dmu)$)
by continuity preserves inner products, so $\E_{\mu_0}[\Phi(s,f)\Phi(s,f')]=\ip f{Cf'}$; and
$L^2$-limits of centred jointly Gaussian variables are centred and jointly Gaussian.  A centred
Gaussian family is determined in law by its covariance, and $\dphi$ is the centred Gaussian measure
with covariance $C$ (\S2.1 of the record).  $F(s)$ is the same polynomial in the family at time $s$
as $F$ is in $(\varphi_0(f))_f$.
\end{proof}

\begin{auxlemma}[cylinder observables as GJS monomials]\label{aux:d3-cyl}
Let $Q=\prod_{i=1}^r\varphi_0(f_i)$ be a cylinder monomial supported in $I$, and $s\in\R$.
\emph{(i)} $Q(s)=\int d\vec x\,w(x)\prod_{i=1}^r{:}\Phi(x_i)^{n_i}{:}$ with $n_i=1$, all times
$x_{i,0}=s$, and kernel $w=f_1\otimes\dots\otimes f_r\in L^2(\R^r)$: a monomial of type (1.1.5)
\cite[p.~593]{GJS}, since ${:}\Phi^1{:}=\Phi$; no Wick reordering is made (Assumption~A8(ii) of the
record: ``products of individually Wick-ordered powers are admissible without reordering'').
Its individual degrees are $n_i=1<\bar n$.
\emph{(ii)} For $\kappa=(\kappa_1,\dots,\kappa_r)\in K_I^r$ let $Q_\kappa(s)$ be the monomial of
type (1.1.5) at time $s$ with kernel
$w_\kappa:=\bigotimes_i f_i\mathbf 1_{[\kappa_i-\frac12,\kappa_i+\frac12)}$.  It is localized in
the squares $\Delta_{l_s,\kappa_1},\dots,\Delta_{l_s,\kappa_r}$, which belong to $R_s(I)$, and
$Q(s)=\sum_{\kappa\in K_I^r}Q_\kappa(s)$.
\emph{(iii)} For the norm $\GJS{Q(s)}:=\sum_\kappa\GJS{Q_\kappa(s)}$ of this decomposition, and for
$Q(s)^2$ (the cylinder monomial with the $2r$ factors $f_1,\dots,f_r,f_1,\dots,f_r$, decomposed
in the same way), the two bounds (N.a) of Remark~\ref{aux:notation-counts}(iii) hold, at every
time $s$.
\emph{(iv)} $N^\square_I\le\lfloor L\rfloor+2$ (Remark~\ref{aux:notation-counts}(i)).
\emph{(v)} For $F=\sum_\beta c_\beta Q_\beta$ as in \S\ref{sub:d3-setting},
$F(s)=\sum_\beta\sum_\kappa c_\beta Q_{\beta,\kappa}(s)$ is a finite sum of localized monomials
of type (1.1.5) with kernels in $L^2$, individual degrees $1<\bar n$ and localization squares in
$R_s(I)$, and the norm of this decomposition is
\[
 \GJS{F(s)}=\sum_\beta|c_\beta|\,\GJS{Q_\beta(s)}
 \le\sum_\beta|c_\beta|\,K_1K_2^{2r_\beta}\,r_\beta!\,(N^\square_I)^{r_\beta/2}
 \prod_{i=1}^{r_\beta}\|f_{\beta,i}\|_2 .
\]
\emph{(vi)} As an element of $L^2(\dmu)$, the localized monomial $Q_\kappa(s)$ of
Definition~\ref{aux:d1-obs}(b) (the object to which Lemma~\ref{lem:D2} is applied) equals
$\prod_i\Phi(s,f_i\mathbf 1_{[\kappa_i-\frac12,\kappa_i+\frac12)})$ almost everywhere, and
$\sum_{\kappa\in K_I^r}Q_\kappa(s)=Q(s)$ almost everywhere, with $Q(s)$ as in
\S\ref{sub:d3-setting}.
\end{auxlemma}

Items (iii)--(iv) are those of Remark~\ref{aux:notation-counts}, applied to the explicit
decomposition (ii); items (i), (ii), (v) and (vi), which name the localized pieces and their
squares for gate GA1 and identify them with the observable, are not in \S\ref{sec:notation}.

\begin{proof}
(i) is the definition (1.1.5) with $n_i=1$.  (ii) A point $(s,\vec x_i)$ with
$\vec x_i\in[k-\frac12,k+\frac12)$ lies in $\Delta_{l_s,k}$, so $w_\kappa$ vanishes outside the
product of the squares $\Delta_{l_s,\kappa_i}$, and the monomial is localized in them
(Remark~\ref{aux:notation-gjslattice}(i)--(ii)).  Since $f_i$ vanishes outside $I$ and $I$ meets
the columns $k\notin K_I$ in null sets, $f_i=\sum_{k\in K_I}f_i\mathbf
1_{[k-\frac12,k+\frac12)}$ almost everywhere; expanding the tensor product gives
$w=\sum_\kappa w_\kappa$ a.e., hence the decomposition (made precise in (vi)).  (iii) The pieces of
(ii) are the decomposition $f_i=\sum_kf_i\mathbf 1_{[k-\frac12,k+\frac12)}$ with which
Remark~\ref{aux:notation-counts}(iii) runs the proof of Proposition~6.2 of the record on the
columns; its bounds use the local degrees and kernel norms of the pieces only, which do not
depend on the row $l_s$.  (iv) is Remark~\ref{aux:notation-counts}(i).  (v) follows from
(i)--(iii) applied to each $Q_\beta$, with $\GJS{cQ_\kappa(s)}=|c|\,\GJS{Q_\kappa(s)}$ by
(1.1.8) (the kernel is multiplied by $c$).  (vi) The kernel $w_\kappa$ is a product of the
functions $f_i\mathbf 1_{[\kappa_i-\frac12,\kappa_i+\frac12)}\in L^2(\R)$, each with bounded support
and vanishing outside the spatial projection of $\Delta_{l_s,\kappa_i}$.  By
Lemma~\ref{aux:d1-moments}(c), $Q_\kappa(s)$ is almost everywhere the product of the one-factor
monomials of degree $1$ at time $s$ with these kernels, and by Lemma~\ref{aux:d1-moments}(f)(1)
each of them is $\Phi(s,f_i\mathbf 1_{[\kappa_i-\frac12,\kappa_i+\frac12)})$.  The map
$f\mapsto\Phi(s,f)$ is linear and, being isometric for the covariance norm, takes almost
everywhere equal functions to the same element; so, with $f_i=\sum_{k\in K_I}f_i\mathbf
1_{[k-\frac12,k+\frac12)}$ a.e.,
$\sum_\kappa\prod_i\Phi(s,f_i\mathbf 1_{[\kappa_i-\frac12,\kappa_i+\frac12)})=
\prod_i\Phi(s,f_i)=Q(s)$.  (Remark~\ref{aux:notation-sharpidrem}(c) records the degree-$1$
identification for $f\in C_c^\infty(\R)$; the pieces have the kernels
$f_i\mathbf 1_{[\kappa_i-\frac12,\kappa_i+\frac12)}$, which are not smooth, and are covered by
Lemma~\ref{aux:d1-moments}(f)(1), and also by Lemma~\ref{aux:notation-sharpid} with $n=1$.)
\end{proof}

\begin{auxrem}[the record's cell count fails in (7) and (8); a finding]\label{aux:d3-cells}
Remark~\ref{aux:notation-counts}(ii)--(iii) shows that the record's cell count $N_I$ and the
column count $N^\square_I$ differ, and that (7) of the record fails with $N_I$ on GJS's squares
($I=(0.4,0.6)$, $Q=\varphi_0(f)$: $\sqrt2K_1K_2^2\|f\|_2$ against $K_1K_2^2\|f\|_2$).  On GJS's
squares the cells $[j,j+1)$ are not localization regions, so the decomposition of the record's
proof of Proposition~6.2 is not available there; the bound fails for the decomposition by columns,
which is what that proof becomes on GJS's squares, and, for $r=1$, for every decomposition of
$\Phi(s,f)$ into one-factor pieces of degree $1$ at time $s$: such pieces are $\Phi(s,f_a)$ with
$f_a$ vanishing outside one column, $\sum_af_a=f$ almost everywhere because $f\mapsto\Phi(s,f)$ is
injective ($C$ is positive definite), so the kernels of the pieces in column $c$ sum to
$f\mathbf 1_c$, and the triangle inequality gives a norm sum
$\ge K_1K_2^2(\|f\mathbf 1_{[-\frac12,\frac12)}\|_2+\|f\mathbf 1_{[\frac12,\frac32)}\|_2)
=\sqrt2K_1K_2^2\|f\|_2$ for $f$ even about $\frac12$.  (Decompositions with pieces of other
degrees are not examined here.)  The same example fails (8), which is the bound Lemma~7.1 of the
record consumes: for $f\in C_c^\infty(I)$ real, even about $\frac12$, with $\|f\|_2=1$, the column
decomposition (ii) of $Q(s)^2=\Phi(s,f)^2$ has the four pieces $\kappa\in\{0,1\}^2$, with norms
$K_1K_2^4\cdot2!\cdot\frac12$ for $\kappa=(0,0)$ and $(1,1)$ and $K_1K_2^4\cdot\frac12$ for
$\kappa=(0,1)$ and $(1,0)$, in total $3K_1K_2^4$, while the right side of (8) with $N_I=1$ is
$K_1K_2^4\cdot2!\cdot1\cdot1=2K_1K_2^4$.  \emph{Severity, as a finding about the record:}
CRITICAL for the display of Proposition~6.2 as stated, on the printed norm of Assumption~A8(iii)
(GJS's squares); the displayed bounds of Lemmas~7.1 and~7.6, whose proofs consume (8), fail as
stated; MINOR downstream, for Theorems~7.7, 7.8 and~12.2, which use finiteness only.
\S\ref{sub:cons-NI} records the edit of the record ($N_I$ redefined by GJS's columns), which
changes the statements of Lemma~6.1 and Proposition~6.2, the notation table and Assumption~A8(ii),
and the displayed constants $\mathfrak n(Q)$ and $K'$.  (The counts and the value $3K_1K_2^4$ for a
bump function are reproduced by \texttt{parts/aqft\_checks.py}, check~6, \emph{numerical}.)
\end{auxrem}

\begin{auxlemma}[geometry of Phase A: parallel strips]\label{aux:d3-geom}
Let $T>0$, $h:=h_{T,g}$, $B:=B_s(g,I)$ and $R:=R_s(I)$.
\emph{(i)} If $\Delta_{l,k}\in\mathcal D(h,B)$, then $k\in K_g$, $|l-l_s|\ge2$, and row $l$ meets
$[-T,T]$ in positive measure.
\emph{(ii)} For every $k\in\Z$ and every $l$ with $|l-l_s|\ge2$, the strip
\[
 P_l:=\bigl\{z:\ \min(l,l_s)+\tfrac12<x^0<\max(l,l_s)-\tfrac12\bigr\},
\]
bounded by two lines of constant time, hence \emph{parallel to the time-zero line}, of width
$|l-l_s|-1\ge1$, separates $\Delta_{l,k}$ from $\bigcup R$ (indeed from the whole row $l_s$).
\emph{(iii)} With $w_{\Delta_{l,k}}:=|l-l_s|-1$,
\[
 \sum_{\Delta\in\mathcal D(h,B)}e^{-m_1w_\Delta}\le2N_g\sum_{j\ge1}e^{-m_1j}
 =\frac{2N_g\,e^{-m_1}}{1-e^{-m_1}} .
\]
\end{auxlemma}

\begin{proof}
(i) By Tonelli's theorem,
$\int_{\Delta_{l,k}}h\,dz=\bigl|[-T,T]\cap[l-\frac12,l+\frac12)\bigr|\cdot
\int_{k-1/2}^{k+1/2}g(x^1)\,dx^1$.  If this is positive, both factors are; the second gives
$|\{g>0\}\cap[k-\frac12,k+\frac12)|>0$, and $\{g>0\}\subseteq\supp g$, so $k\in K_g$, hence
$k_-\le k\le k_+$.  A square of a column in $[k_-,k_+]$ that is not in $B$ has
$l\notin\{l_s-1,l_s,l_s+1\}$, i.e.\ $|l-l_s|\ge2$.
(ii) Lemma~\ref{aux:notation-strip}(i), applied to the rectangle
$S:=\{\Delta_{l_s,k'}:\min K_I\le k'\le\max K_I\}\supseteq R$ (one row, $l_1=l_2=l_s$) and to
$\Delta=\Delta_{l,k}\notin S$, whose row distance from $S$ is $\rho_t=|l-l_s|\ge2$: the strip
bounded by $x^0=l_s+\frac12$ and $x^0=l-\frac12$ (if $l>l_s$), resp.\ $x^0=l+\frac12$ and
$x^0=l_s-\frac12$ (if $l<l_s$), which is $P_l$, is parallel, has width $\rho_t-1$, and separates
$\Delta$ from $\bigcup S\supseteq\bigcup R$.  The same two half-planes contain the whole row $l_s$.
(iii) By (i) every $\Delta_{l,k}\in\mathcal D(h,B)$ has $k\in K_g$ and $|l-l_s|\ge2$; the terms are
positive, so the sum is at most $\sum_{k\in K_g}\sum_{|l-l_s|\ge2}e^{-m_1(|l-l_s|-1)}$, in which
$j:=|l-l_s|-1$ takes every value $j\ge1$ for exactly two $l$.  The geometric series with ratio
$e^{-m_1}\in(0,1)$ gives the closed form.  (Test: $N_g=1$, $e^{-m_1}=\frac12$: both sides equal
$2$.)
\end{proof}

\subsection{The theorem}\label{sub:d3-thm}

\begin{dtheorem}[Phase A: fixed cutoff, parallel strips]\label{thm:D3}
Assume the standing assumption of \S\ref{sub:d2-conv}.  Let $g\in\CH$, let $I\subset\R$ be a
bounded interval of positive length, let $s\in\R$, and let $l_s$, $N_g$, $N_B$, $B=B_s(g,I)$ be
as in \S\ref{sub:d3-setting}.  Let $F=\sum_\beta c_\beta Q_\beta$ be a finite real linear
combination of cylinder monomials supported in $I$, and $F(s)$ its value at sharp time $s$.
Then for every $T>0$,
\[
 \Bigl|\int F(s)\,dq_{h_{T,g}}\Bigr|\le A\,\|F\|_{L^2(\dphi)}+B'\,\GJS{F(s)} ,
\]
\[
 A:=e^{3\lambda|a_0|N_B}\,S_2(B),\qquad
 B':=2\,C_7\,c_U\,N_g\,\frac{e^{-m_1}}{1-e^{-m_1}} ,
\]
where $S_2(B)$ is the constant of Definition~\ref{aux:notation-S2def}.  The norm $\GJS{F(s)}$
is that of the decomposition in Lemma~\ref{aux:d3-cyl}(v), which bounds it in terms of the
functions $f_{\beta,i}$.  In working coordinates (Remark~3.3 of the record), the constant $A$
depends only on $(\lambda,\mathscr P,m_0)$ and on $N_B$, through the unevaluated constant $S_2$;
the constant $B'$ depends only on $C_7$, $m_1$, $K_1$, $K_2$, $\lambda$, $\mathscr P$ and $N_g$.
Both are independent of $T$ and $s$, and $B'$ is independent of $I$.
\end{dtheorem}

\emph{Status of Theorem~\ref{thm:D3}.}  Proved, conditional on: Lemma~\ref{lem:D1}, through
Lemma~\ref{lem:D2} (Lemma~C.4, R1--R10, as stated after Lemma~\ref{lem:D2});
Lemma~C.2 of the record (repair R8); Theorem~C.5 of the record with
Proposition~\ref{aux:notation-S2} (repairs R1--R10), whose part (iii), resting on Lemma~C.4,
gives the independence of $A$ from $s$; and the reading of the source of
Remark~\ref{aux:notation-gjslattice}(iv), which it inherits through Lemma~\ref{lem:D2} (the
pieces of Lemma~\ref{aux:d3-cyl}(vi) and the $U_{\Delta,n}$ are taken as GJS's monomials, and the
record's $\Uc_h$ as GJS's $V(h)$).  It uses GJS Theorem~1.1.7 in the parallel case only, with
a time-integrated partner (Remark~\ref{aux:d3-orient}); it does not use the transverse case, the
GAP recorded after Lemma~\ref{aux:notation-wick}, or GJS Theorem~1.1.11.

\begin{proof}
\emph{Step 1: the observable.}  By Lemma~\ref{aux:d3-cyl}(v), $F(s)=\sum_{\beta,\kappa}c_\beta
Q_{\beta,\kappa}(s)$ is a finite sum of localized monomials of type (1.1.5), with kernels in $L^2$,
individual degrees $n_i=1<\bar n$, and localization squares in $R:=R_s(I)$, and by
Lemma~\ref{aux:d3-cyl}(vi) the sum of these pieces, as objects of Definition~\ref{aux:d1-obs}, is
the observable $F(s)$ almost everywhere; $B$ is a box containing the squares of $R$;
$\GJS{F(s)}=\sum_{\beta,\kappa}|c_\beta|\,\GJS{Q_{\beta,\kappa}(s)}$.  Put $h:=h_{T,g}\in\CE$.

\emph{Step 2: the base term.}  By Lemma~\ref{aux:notation-cutoffs}(i), $h^B\in\CE$ and
$\supp h^B\subseteq\overline B$, a compact set of Lebesgue measure $|B|=3N_B$.  Lemma~C.2 of the
record \cite[p.~40]{R} (proved conditional on repair R8 of the V3 audit,
Remark~\ref{aux:d3-V3}) states: for every $h\in\CE$ with $\supp h\subseteq\Lambda$,
$Z(h)\ge e^{-\lambda|a_0||\Lambda|}>0$; its hypotheses ($h\in\CE$, $\supp h\subseteq\Lambda$,
$\Lambda$ compact as in its notation, main.tex l.~2902) hold for $h^B$ and $\Lambda=\overline B$,
so $Z(h^B)\ge e^{-3\lambda|a_0|N_B}$.  Proposition~\ref{aux:notation-S2}(iv) (proved conditional
on the repairs R1--R10, as in Remark~\ref{aux:d3-V3}) gives $\|e^{-\Uc_{h^B}}\|_{L^2(\dmu)}\le
S_2(B)$.  By Cauchy--Schwarz in $L^2(\dmu)$ and Lemma~\ref{aux:d3-sharp},
\[
 \bigl|\ev{F(s)}{h^B}\bigr|=\frac{\bigl|\int F(s)\,e^{-\Uc_{h^B}}\dmu\bigr|}{Z(h^B)}
 \le e^{3\lambda|a_0|N_B}\,\|F(s)\|_{L^2(\dmu)}\,\|e^{-\Uc_{h^B}}\|_{L^2(\dmu)}
 \le A\,\|F\|_{L^2(\dphi)} .
\]

\emph{Step 3: the comparison term.}  Apply Lemma~\ref{lem:D2}(a) with $F(s)$ and its pieces
(Step~1), $R=R_s(I)$, the box $B=B_s(g,I)$, the cutoff $h=h_{T,g}\in\CE$, and, for
$\Delta=\Delta_{l,k}\in\mathcal D(h,B)$, the strip $S_\Delta:=P_l$ of
Lemma~\ref{aux:d3-geom}(ii), of width $w_\Delta=|l-l_s|-1$ (by Lemma~\ref{aux:d3-geom}(i),
$|l-l_s|\ge2$, so $P_l$ is defined and separates $\Delta$ from $\bigcup R$).  Each estimate inside
that lemma is Source Fact~2.2 of the record with (gate GA1):
\begin{itemize}
\item $Q_1=c_\beta Q_{\beta,\kappa}(s)$: type (1.1.5), $r_\beta$ factors of degree $1<\bar n$ at
  time $s$, localized in $\Delta_{l_s,\kappa_1},\dots,\Delta_{l_s,\kappa_{r_\beta}}\in R_s(I)$,
  kernel $c_\beta\bigotimes_if_{\beta,i}\mathbf 1_{[\kappa_i-\frac12,\kappa_i+\frac12)}\in L^2$;
\item $Q_2=U_{\Delta,n}$, $1\le n\le2p'$: type (1.1.6), localized in the one square
  $\Delta=\Delta_{l,k}$, $k\in K_g$, $|l-l_s|\ge2$, degree $n<\bar n$, kernel
  $\lambda a_nh\mathbf 1_\Delta\in L^2(\R^2)$, norm $\le\lambda|a_n|K_1K_2^{2n}n!$
  (Lemma~\ref{aux:notation-cU}(i));
\item the cutoff $h_\alpha=h^B+\alpha(h-h^B)\in\CE$, $0\le\alpha\le1$
  (Lemma~\ref{aux:notation-cutoffs}(ii));
\item the strip $P_l$: \emph{parallel to the time-zero line}, width $|l-l_s|-1\ge1$;
\item the resulting term, summed over $n$: $C_7\,c_U\,|c_\beta|\,\GJS{Q_{\beta,\kappa}(s)}\,
  e^{-m_1(|l-l_s|-1)}$.
\end{itemize}
Lemma~\ref{lem:D2}(a) and Lemma~\ref{aux:d3-geom}(iii) give
\begin{align*}
 \bigl|\ev{F(s)}{h}-\ev{F(s)}{h^B}\bigr|
 &\le C_7\,c_U\,\GJS{F(s)}\sum_{\Delta\in\mathcal D(h,B)}e^{-m_1w_\Delta}\\
 &\le C_7\,c_U\,\GJS{F(s)}\,\frac{2N_ge^{-m_1}}{1-e^{-m_1}}=B'\,\GJS{F(s)} .
\end{align*}

\emph{Step 4.}  The triangle inequality and Steps~2--3 give the bound.  The constants: $S_2(B)$
depends on $(\lambda,\mathscr P,m_0)$ through Theorem~C.5 and on $B$ only through its shape
$3\times N_B$ (Proposition~\ref{aux:notation-S2}(iii)), so $A$ does not depend on $s$ (which moves
$B$ by whole rows) or on $T$; $B'$ contains $N_g$ and the printed constants and nothing that
depends on $T$, $s$ or $I$.
\end{proof}

\subsection{Remarks on Theorem~\ref{thm:D3}}\label{sub:d3-rem}

\begin{auxrem}[orientation, and what the printed proof reaches]\label{aux:d3-orient}
Every application of Source Fact~2.2 in the proof of Theorem~\ref{thm:D3} pairs a sharp-time
piece at time $s$ with a time-integrated $U_{\Delta,n}$ in a row $l$, $|l-l_s|\ge2$, across the
strip $P_l$, whose bounding lines are lines of constant time.  No strip transverse to the
time-zero line is used in this section.  This is the parallel case, with a time-integrated
partner, which the printed proof reaches only after slicing.  The sentence F-T
(\S\ref{sec:purpose}) and Lemma~7.4 of the record \cite[p.~21]{R} concern pairs of sharp-time
monomials: the printed proof translates the strip and takes ``$A=Q_1$, $B_T=Q_2$ in Theorem
1.1.11'' \cite[p.~597]{GJS}, where ``We assume that $A$ is localized at times $x_{i0}\leqq0$, and
that $B$ is localized at times $x_{i0}\geqq0$'' \cite[p.~596]{GJS}, and Theorem~1.1.11 is printed
for ``monomials of type (1.1.5)'' \cite[p.~596]{GJS}.  The partner $U_{\Delta,n}$ is of type
(1.1.6); the printed proof reaches the pair only after $U_{\Delta,n}$ is sliced in time and the
bound is integrated over the slices, the multi-time reading of (1.1.5) that the plan records
(Part~2) and GJS do not print.  The note uses the printed statement of Theorem~1.1.7, which covers
the pair (Step~4 of the proof of Lemma~\ref{lem:D2}), and does not need to decide whether the
printed proof reaches it.  The strips are parallel because $B$ contains all columns of $\supp g$
in the three rows $l_s-1,l_s,l_s+1$: every charged square outside $B$ lies above or below them
(Lemma~\ref{aux:d3-geom}(i)).  The price is that $A$ depends on $N_B\ge N_g$; this is why
Theorem~\ref{thm:D3} is a fixed-$g$ statement and Theorem~\ref{thm:D4} needs transverse strips.
Remark~\ref{aux:d2-uses}(3) is used here: squares of the rows of $B$ outside its columns that meet
$\supp h_{T,g}$ only in a null set are not charged.
\end{auxrem}

\begin{auxrem}[order of quantifiers]\label{aux:d3-quant}
For every $g\in\CH$ and every bounded interval $I$ of positive length the constants $A$ (through
$N_B$) and $B'$ (through $N_g$) are fixed; then the bound holds for every $s\in\R$, every finite
real linear combination $F$ of cylinder monomials supported in $I$, and every $T>0$.  Neither
constant depends on $s$, $F$ or $T$; $A$ depends on $I$ through $N_B$ only, and $B'$ not at all.
\end{auxrem}

\begin{auxrem}[boundary cases]\label{aux:d3-boundary}
\emph{$T<1$.}  Row $l$ meets $[-T,T]$ in positive measure if and only if $l-\frac12<T$ and
$l+\frac12>-T$.  \emph{(a) $T\le\frac12$:} only row $0$ does.  If moreover $|l_s|\le1$, then
$\mathcal D(h_{T,g},B)=\varnothing$ by Lemma~\ref{aux:d3-geom}(i), so $h_{T,g}-h_{T,g}^B=0$ almost
everywhere (Lemma~\ref{aux:notation-cutoffs}(iii)) and the comparison term vanishes; if
$|l_s|\ge2$, then $h_{T,g}^B=0$ almost everywhere (everywhere unless $T=\frac12$ and $l_s=2$, when the null
set $\{x^0=\frac12\}$ lies in row $1\subseteq B$), the base term is $\int F(s)\dmu$, and
$\mathcal D$ consists of row-$0$ squares in columns of $K_g$, at row distance $|l_s|\ge2$.
\emph{(b) $\frac12<T<1$:} exactly the rows $-1,0,1$ meet $[-T,T]$ in positive measure.  If
$l_s=0$, every charged square lies in rows $-1,0,1$ and in a column of $K_g$, hence in $B$, and
$\mathcal D(h_{T,g},B)=\varnothing$; otherwise the charged squares outside $B$ lie in rows
$l\in\{-1,0,1\}$ with $|l-l_s|\ge2$, and the general case of the proof applies to them.
\emph{$T$ not an integer} enters only through the rows met by the slab.
\emph{$s$ on a lattice line}, $s\in\frac12+\Z$: then $l_s=s+\frac12$, the sharp-time pieces lie on
the lower edge of row $l_s$, and for $l\le l_s-2$ the upper bounding line of $P_l$ is $x^0=s$; the
observable lies on the boundary line, outside the open strip, which is what separation requires.
Assigning such an $s$ to the lower row instead (closed squares) moves $B$ down one row; the proof
is unchanged.
\emph{$I$ meeting one column only}: $K_I=\{k\}$, $R=\{\Delta_{l_s,k}\}$, $N^\square_I=1$.
\emph{$N_g=0$} (which implies $g=0$ a.e.): $h_{T,g}=0$ a.e., $B'=0$, and the bound is the base
term.  Conversely, if $g=0$ a.e., no square is charged and the comparison term vanishes, whatever
$N_g$ is; $B'$ then overestimates it (example: $g=\mathbf 1_{\mathbb Q\cap[0,1]}\in\CH$ vanishes
a.e., while $\supp g=[0,1]$ gives $N_g=2$).
\emph{$N_g=1$}: $B'=2C_7c_Ue^{-m_1}(1-e^{-m_1})^{-1}$.
\emph{The $n=0$ term of $\mathscr P$} enters only through $Z(h^B)$ (the factor
$e^{3\lambda|a_0|N_B}$) and through $\Uc_{h^B}$ inside $S_2$, and is absent from $c_U$
(Lemma~\ref{aux:notation-cU}(iii)).
\end{auxrem}

\begin{auxrem}[what Theorem~\ref{thm:D3} is not, and its inputs]\label{aux:d3-not}
It is not CE1 at slab cutoffs: its bound is $A\|F\|_{L^2}+B'\GJS{F(s)}$, with $A$ growing with
the observable's box (through $N_B$ and $S_2(B)$), not $C_8\GJS{F(s)}$.  The statement ``CE1
follows from CE2'' is not made; what is proved is a bound of this form for observables of bounded
support (trap~8 of the plan).  It is uniform in $T$ and $s$ at fixed $g$ and $I$, and it is not
uniform in $g$: $A$ and $B'$ grow with $N_B$ and $N_g$.  Of GJS only Theorem~1.1.7 is used, in its
parallel case; Theorem~1.1.11 and Assumption~A8(vi) are not used, and none of GJS's proofs is
verified here.  The inputs from the record are: Assumptions~A2, A3 and A8(i)--(iv), Remark~2.5
(Lemma~\ref{aux:d3-sharp}), Lemma~3.1(ii) ($h_{T,g}\in\CE$), Lemma~C.2, Theorem~C.5, and
Lemma~C.4, through Lemma~\ref{lem:D2} (Step~3 of its proof) and through
Proposition~\ref{aux:notation-S2}(iii) (the independence of $A$ from $s$); with Lemmas~\ref{lem:D1}
and~\ref{lem:D2} under their conditions, all as in Remark~\ref{aux:d3-V3} and the status line
after Theorem~\ref{thm:D3}.
\end{auxrem}


\section{Phase B: uniform in the cutoff, transverse strips}\label{sec:d4}
% d4.tex -- drafted by aqft (council member of ugjs/PLAN.md), 2 November 2026 (begun 31 October 2026);
% revised 3 November 2026 after the audits V1 (cqft, mp) and V2: the identification of the two
% constructions of Q_V (V1 cqft finding 1 = V2-1) is cqft2's Lemma aux:notation-sharpid with
% Remark aux:notation-sharpidrem(b), cited here; the status line after the theorem; V_* for the
% constructed V(h'), so that the GAP enters only the transfer to the record (\S8.6).
% Section 6: Theorem D.4 (Phase B; uniform in the cutoff; transverse strips, disclosed by F-T).
% Merged with \S2 (notation.tex, cqft2) on 2 November 2026: the box B_{3x5}, its rings and their
% strips, the Wick isometry with V_*, and the comparison ||Y||_4 <= 3^{p'}||Y||_2 are those of \S2.
% Auxiliary statements of this part: aux:d4-QV and the remarks of \S6.5.

Theorem~\ref{thm:D4} proves the moment bound that Theorem~5.1 of the record consumes, for every
cutoff $h\in\CE$ at once, hence uniformly in $T>0$ and $g\in\CH$ when $h=h_{T,g}$.  Its proof
uses GJS Theorem~1.1.7 also across strips transverse to the time-zero line.  That case is
covered by the printed statement and not reached by its printed proof; it is disclosed by the
sentence F-T of \S\ref{sec:purpose}, and every place below where it is used says so.

\subsection{Setting}\label{sub:d4-setting}

The standing assumption of \S\ref{sub:d2-conv} holds.  Let $h\in\CE$, $a\in\R$, and let $h'$ be a
real measurable function on $\R$ with $\|h'\|_\infty\le1$ and $\supp h'\subseteq[a,a+2]$.

\emph{The observable.}  $V(h'):=\lambda\sum_{n=0}^{2p'}a_nW_n(h')\in L^2(\dphi)$ is the
time-zero observable constructed in Lemma~\ref{aux:notation-wick}(iv), $W_n(h')$ the limit of
Lemma~\ref{aux:notation-wick}(iii).  Its sharp-time version at time $0$ is the element
\[
 V(h')(0):=Y=J_0V(h')\in L^2(\dmu)
\]
of Lemma~\ref{aux:notation-hyper}, with $J_0$ the isometry of the proof of its part (i), and
$Q_V:=V(h')(0)^2=Y^2$.  Formally $V(h')(0)=\lambda\sum_na_n\int h'(x)\,{:}\Phi(0,x)^n{:}\,dx$, the
record's eq.~(4) with $h'$ for its $h$.  By Lemma~\ref{aux:notation-sharpid} and
Remark~\ref{aux:notation-sharpidrem}(b), $Q_V$ is also, almost everywhere, the sum of the
localized monomials of Definition~\ref{aux:d1-obs}(b) to which Lemma~\ref{lem:D2} is applied
(Lemma~\ref{aux:d4-QV}(i) below); so the base term and the comparison term of the proof concern
one and the same element of $L^2(\dmu)$.  Theorem~\ref{thm:D4} is a statement about this
constructed $V(h')$.  The record's $V(h')$ is identified with it by the GAP recorded after
Lemma~\ref{aux:notation-wick} (the record states no construction of the Wick powers in $V(h)$);
Theorem~\ref{thm:D4} does not use that identification, which enters only at the record's
Theorem~5.1, Step~1, and at the transfer of Theorem~\ref{thm:D4} to Theorem~5.1
(\S\ref{sub:cons-51}).

\emph{The box and the rings.}  We use Definition~\ref{aux:notation-ringdef}: $k_a$, the row-$0$
squares $S_a=\{\Delta_{0,k}:k_a\le k\le k_a+2\}$, the box $B_{3\times5}=B_{3\times5}(a)$ (rows
$-1,0,1$, columns $k_a-1,\dots,k_a+3$; $|B_{3\times5}|=15$), and the rings $\mathcal L_j$, $j\ge1$.
By Lemma~\ref{aux:notation-rings}(i) the columns that meet $[a,a+2]$ are exactly
$k_a,k_a+1,k_a+2$, so $[a,a+2]\subseteq[k_a-\frac12,k_a+\frac52)$, and the time $0$ lies in
row $0$.  Put $R:=S_a$.

\emph{The ring constant.}  With $x:=e^{-m_1}\in(0,1)$ let
\[
 \overline\Sigma_{3\times5}:=\sum_{j\ge1}(8j+12)\,e^{-m_1j}=\frac{4x(5-3x)}{(1-x)^2},
\]
the right side of Lemma~\ref{aux:notation-rings}(vi).  That lemma writes $\Sigma_{3\times5}$ for
the sum $\sum_{\Delta\not\subseteq B_{3\times5}}e^{-m_1d(\Delta,R)}$, which depends on $R$, and
bounds it by $\overline\Sigma_{3\times5}$.  Theorem~\ref{thm:D4} is stated and proved with
$\overline\Sigma_{3\times5}$, because its proof applies Theorem~1.1.7 only along the axis-parallel
strips of Lemma~\ref{aux:notation-rings}(iv), whose orientation is named; using the widest strips
instead would give the smaller $\Sigma_{3\times5}$, at the price of tilted (hence transverse)
strips also for squares with $|l|=1+j$ (Remark~\ref{aux:notation-tilted}).

\emph{The Gaussian constant.}  $V_*:=\sup_{h'}\|V(h')\|_{L^2(\dphi)}^2$, the supremum over all $h'$
as above and all $a$, is finite by Lemma~\ref{aux:notation-wick}(iv), with
$V_*\le\overline V_*=\lambda^2\sum_{n=0}^{2p'}a_n^2\,n!\,I_n$ there.  It concerns the
constructed $V(h')$ and does not use the GAP.

\subsection{The observable}\label{sub:d4-aux}

\begin{auxlemma}[the observable $Q_V$]\label{aux:d4-QV}
For $0\le n,n'\le2p'$ and $k,k'\in\{k_a,k_a+1,k_a+2\}$ let
$h'_k:=h'\mathbf 1_{[k-\frac12,k+\frac12)}$, and let $Q_{n,n',k,k'}$ be the localized monomial of
Definition~\ref{aux:d1-obs}(b) at time $0$ with two factors of degrees $n,n'$ and kernel
$\lambda^2a_na_{n'}h'_k\otimes h'_{k'}$, formally
\[
 Q_{n,n',k,k'}=\lambda^2a_na_{n'}\iint h'_k(x)\,h'_{k'}(y)\,{:}\Phi(0,x)^n{:}\,{:}\Phi(0,y)^{n'}{:}
 \,dx\,dy .
\]
\emph{(i)} $Q_V=\sum_{n,n'}\sum_{k,k'}Q_{n,n',k,k'}$ $\dmu$-almost everywhere.
\emph{(ii)} Each $Q_{n,n',k,k'}$ is a monomial of type (1.1.5), in the sense of
Definition~\ref{aux:d1-obs}(b), with $r=2$ factors at time $0$, of individual degrees
$n,n'\le2p'<\bar n$, with kernel $\lambda^2a_na_{n'}h'_k\otimes h'_{k'}\in L^2(\R^2)$, localized
in $\Delta_{0,k}$ and $\Delta_{0,k'}$, which belong to $S_a$.  Its total degree $n+n'\le4p'$ may
exceed $\bar n-1=2p'$; this is admissible because the source bounds only the individual degrees,
``We assume a bound on the $n_i$ in (1.1.5), $n_i<\bar n$'', and adds ``Since there is no
restriction on $r$, $Q$ still may have an arbitrarily large degree'' \cite[p.~593]{GJS}.  This is
where that sentence is used.  No Wick re-expansion of the product is made: as in Step~1 of the
proof of Theorem~5.1 of the record (main.tex ll.~1278--1283), the product of two individually
Wick-ordered powers is itself of type (1.1.5).
\emph{(iii)} For this decomposition, (1.1.9) and (1.1.8) give
\[
 \GJS{Q_V}:=\sum_{n,n',k,k'}\GJS{Q_{n,n',k,k'}}
 \le9\,K_1K_2^{8p'}(4p')!\,\lambda^2\Bigl(\sum_{n=0}^{2p'}|a_n|\Bigr)^2 .
\]
\end{auxlemma}

Part (i) is Remark~\ref{aux:notation-sharpidrem}(b) of \S\ref{sec:notation}, restated with the
present notation; (ii)--(iii) are Steps~1--4 of the proof of Theorem~5.1 of the record,
transcribed to the columns of $S_a$.  That the limits of Definition~\ref{aux:d1-obs} are the
monomials that GJS print as formal integrals is a reading of the source
(Remark~\ref{aux:notation-gjslattice}(iv), Remark~\ref{aux:d1-def}).

\begin{proof}
(i) Each $h'_k$ is real, measurable, bounded by $1$, vanishes outside $[a,a+2]$ and outside the
column $[k-\frac12,k+\frac12)$, so it satisfies the hypotheses of
Lemma~\ref{aux:notation-sharpid} (bounded, compact support) and of
Remark~\ref{aux:notation-sharpidrem}(a) (vanishing outside one column).  Since $h'$ vanishes
outside $[a,a+2]\subseteq[k_a-\frac12,k_a+\frac52)$, $h'=\sum_{k=k_a}^{k_a+2}h'_k$ pointwise.
Remark~\ref{aux:notation-sharpidrem}(b) gives, almost everywhere,
$Q_{n,n',k,k'}=\lambda^2a_na_{n'}\,J_0W_n(h'_k)\,J_0W_{n'}(h'_{k'})$ and $W_n(h')=\sum_kW_n(h'_k)$;
hence, $J_0$ being linear,
\[
 \sum_{n,n',k,k'}Q_{n,n',k,k'}=\Bigl(J_0\,\lambda\sum_na_n\sum_kW_n(h'_k)\Bigr)^2
 =\bigl(J_0V(h')\bigr)^2=Y^2=Q_V .
\]
(ii) The kernel is a product of bounded functions with compact support, hence in $L^2(\R^2)$, and
it vanishes unless $x\in[k-\frac12,k+\frac12)$ and $y\in[k'-\frac12,k'+\frac12)$, i.e.\ unless
$(0,x)\in\Delta_{0,k}$ and $(0,y)\in\Delta_{0,k'}$.  A factor of degree $0$ is ${:}\Phi^0{:}=1$
and contributes $0$ to the local degree.
(iii) This is Steps~3--4 of the proof of Theorem~5.1 of the record (main.tex ll.~1293--1315) with
the cells $[j,j+1)$ replaced by GJS's columns.  For one ordered pair $(k,k')$ and degrees
$(n,n')$: if $k=k'$ the local degree of $\Delta_{0,k}$ is $n+n'$ and the weight of (1.1.8) is
$K_1K_2^{2(n+n')}(n+n')!$; if $k\ne k'$ it is $K_1K_2^{2n}n!\,K_2^{2n'}n'!$.  In both cases the
weight is at most $K_1K_2^{8p'}(4p')!$, because $K_2\ge1$, $n+n'\le4p'$ and
$n!\,n'!\le(n+n')!\le(4p')!$.  The kernel norm is
$\lambda^2|a_n||a_{n'}|\,\|h'_k\|_2\|h'_{k'}\|_2\le\lambda^2|a_n||a_{n'}|$, since
$\|h'_k\|_2^2=\int_{k-1/2}^{k+1/2}h'^2\le1$.  There are at most $3^2=9$ ordered column pairs, and
$\sum_{n,n'}|a_n||a_{n'}|=(\sum_n|a_n|)^2$.  (This is eq.~(5) of the record without its final
equality ``$=M^2/C_8$'', which item n14 of the plan deletes.  Test, $\mathscr P(\xi)=\xi^4$: the
bound is $9\cdot8!\,K_1K_2^{16}\lambda^2=362\,880\,K_1K_2^{16}\lambda^2$.)
\end{proof}

\subsection{The theorem}\label{sub:d4-thm}

\begin{dtheorem}[Phase B: uniform in the cutoff]\label{thm:D4}
Assume the standing assumption of \S\ref{sub:d2-conv}.  For every $h\in\CE$, every $a\in\R$ and
every real measurable $h'$ with $\|h'\|_\infty\le1$ and $\supp h'\subseteq[a,a+2]$,
\[
 \int V(h')(0)^2\,dq_h\le M_B^2 ,
\]
\[
 M_B^2:=e^{15\lambda|a_0|}\,S_2(B_{3\times5})\cdot9^{p'}\,V_*
 +C_7\,c_U\,\overline\Sigma_{3\times5}\cdot9\,K_1K_2^{8p'}(4p')!\,\lambda^2
 \Bigl(\sum_{n=0}^{2p'}|a_n|\Bigr)^2 ,
\]
with $V(h')(0)$, $V_*$ and $\overline\Sigma_{3\times5}$ as in \S\ref{sub:d4-setting}, $c_U$ of
Lemma~\ref{aux:notation-cU}, and $S_2(B_{3\times5})$ the common value of
Definition~\ref{aux:notation-S2def} at all boxes of $3$ rows and $5$ columns
(Proposition~\ref{aux:notation-S2}(iii)).  In working coordinates (Remark~3.3 of the record), $M_B$
depends on $(\lambda,\mathscr P,m_0)$ through $C_7$, $m_1$, $K_1$, $K_2$, the unevaluated finite
constant $S_2(B_{3\times5})$ of Theorem~C.5 of the record, and $V_*$; it does not depend on $h$, $a$
or $h'$.  The proof uses GJS Theorem~1.1.7 across strips transverse to the time-zero line, a case
its printed proof does not reach (F-T, \S\ref{sec:purpose}; \S\ref{sub:d4-transverse}).
\end{dtheorem}

\emph{Status of Theorem~\ref{thm:D4}.}  Proved, conditional on: Lemma~\ref{lem:D1}, through
Lemma~\ref{lem:D2} (Lemma~C.4, R1--R10, as stated after Lemma~\ref{lem:D2});
Lemma~C.2 of the record (repair R8); Theorem~C.5 of the record with
Proposition~\ref{aux:notation-S2}, including the common value of $S_2$ at all $3\times5$ boxes
(repairs R1--R10); repair R8 for the comparison of Lemma~\ref{aux:notation-hyper}(ii); the
transverse case of GJS Theorem~1.1.7 at the printed statement's scope (F-T); and the reading of the
source of Remark~\ref{aux:notation-gjslattice}(iv), which it inherits through Lemma~\ref{lem:D2}
(the pieces $Q_{n,n',k,k'}$ and the $U_{\Delta,n''}$ are taken as GJS's monomials, and the record's
$\Uc_h$ as GJS's $V(h)$).  It does not use the
GAP recorded after Lemma~\ref{aux:notation-wick}, and Lemma~\ref{aux:notation-sharpid} uses no
repair (Remark~\ref{aux:notation-sharpidrem}(d)).

\begin{proof}
\emph{Step 1: the observable and the box.}  By Lemma~\ref{aux:d4-QV}, $Q_V=V(h')(0)^2$ is, almost
everywhere, the finite sum of the localized monomials $Q_{n,n',k,k'}$ of type (1.1.5), with kernels
in $L^2$, individual degrees $\le2p'<\bar n$ and localization squares in $R=S_a$, and
$\GJS{Q_V}\le9K_1K_2^{8p'}(4p')!\lambda^2(\sum_n|a_n|)^2$ for this decomposition.  So Steps~2 and~4
below bound expectations of one and the same element of $L^1(dq_h)$ (Lemma~\ref{aux:d2-int});
almost everywhere equal functions have equal integrals under $dq_h$, which has a density with
respect to $\dmu$.  $B:=B_{3\times5}$ is a box containing the squares of $R$.  Since $Q_V\ge0$,
$\int Q_V\,dq_h=\ev{Q_V}{h^B}+\bigl(\ev{Q_V}{h}-\ev{Q_V}{h^B}\bigr)$ with $\ev{Q_V}{h^B}\ge0$.

\emph{Step 2: the base term.}  $h^B\in\CE$ with $\supp h^B\subseteq\overline B$, a closed $3\times5$
rectangle of Lebesgue measure $15$ (Lemma~\ref{aux:notation-cutoffs}(i)).  Lemma~C.2 of the record
with $\Lambda=\overline B$ (hypotheses: $h^B\in\CE$, $\supp h^B\subseteq\Lambda$, $\Lambda$
compact) gives $Z(h^B)\ge e^{-15\lambda|a_0|}$.  Proposition~\ref{aux:notation-S2}(iv) gives
$\|e^{-\Uc_{h^B}}\|_{L^2(\dmu)}\le S_2(B)$, and by its part (iii) $S_2(B)=S_2(B_{3\times5})$ is one
number for all $a$.  Lemma~\ref{aux:notation-hyper}(ii), applied to $Y=V(h')(0)$, gives
$\|Q_V\|_{L^2(\dmu)}=\|Y^2\|_{L^2(\dmu)}\le9^{p'}V_*$; by Lemma~\ref{aux:notation-sharpid} and
Remark~\ref{aux:notation-sharpidrem}(b) this $Y^2$ is the $Q_V$ of Step~1.  This step concerns
the constructed $V(h')$; the GAP recorded after Lemma~\ref{aux:notation-wick} (the
identification of the record's $V(h)$ with it) is not used here.  (Lemma~C.2, Theorem~C.5
through $S_2$, and the comparison are used as in Remark~\ref{aux:d3-V3}.)  By Cauchy--Schwarz in
$L^2(\dmu)$,
\[
 \ev{Q_V}{h^B}\le e^{15\lambda|a_0|}\,\|Q_V\|_{L^2(\dmu)}\,\|e^{-\Uc_{h^B}}\|_{L^2(\dmu)}
 \le e^{15\lambda|a_0|}\,9^{p'}V_*\,S_2(B_{3\times5}) .
\]

\emph{Step 3: the strips.}  By Lemma~\ref{aux:notation-rings}(ii) every square
$\Delta=\Delta_{l,k}\not\subseteq B$ lies in exactly one ring $\mathcal L_j$, $j\ge1$.  Assign:
\begin{itemize}
\item[(P)] if $|l|=1+j$: the strip $S_\Delta$ bounded by $x^0=\frac12$ and $x^0=l-\frac12$ (if
$l>0$), resp.\ by $x^0=l+\frac12$ and $x^0=-\frac12$ (if $l<0$), which is \emph{parallel to the
time-zero line}, has width $j$, and separates $\Delta$ from $\bigcup S_a$
(Lemma~\ref{aux:notation-rings}(iv), with Lemma~\ref{aux:notation-strip}(i) for $S=S_a$);
\item[(T)] if $|l|\le j$: then $k=k_a-1-j$ or $k=k_a+3+j$, and $S_\Delta$ is the strip bounded by
$x^1=k_a+\frac52$ and $x^1=k-\frac12$, resp.\ by $x^1=k+\frac12$ and $x^1=k_a-\frac12$, which is
\emph{transverse to the time-zero line}, has width $j$, and separates $\Delta$ from $\bigcup S_a$
(Lemma~\ref{aux:notation-rings}(iv), with Lemma~\ref{aux:notation-strip}(ii)).
\end{itemize}
The four corner squares of each ring ($|l|=1+j$ in the two side columns) admit both strips and are
assigned to (P).  So $w_\Delta=j$ for $\Delta\in\mathcal L_j$, and since $|\mathcal L_j|=8j+12$
(Lemma~\ref{aux:notation-rings}(ii)) and the terms are positive,
\[
 \sum_{\Delta\in\mathcal D(h,B)}e^{-m_1w_\Delta}\le\sum_{j\ge1}(8j+12)\,e^{-m_1j}
 =\overline\Sigma_{3\times5} .
\]

\emph{Step 4: the comparison term.}  Apply Lemma~\ref{lem:D2}(a) with $F=\sum Q_{n,n',k,k'}$, which
is $Q_V$ almost everywhere by Lemma~\ref{aux:d4-QV}(i) (Lemma~\ref{aux:notation-sharpid} and
Remark~\ref{aux:notation-sharpidrem}(b)),
and its pieces $Q_{n,n',k,k'}$, $R=S_a$, the box $B=B_{3\times5}$, the cutoff $h\in\CE$, and the
strips of Step~3.  The estimates inside that lemma are Source Fact~2.2 of the record with (gate
GA1):
\begin{itemize}
\item $Q_1=Q_{n,n',k,k'}$: type (1.1.5), two factors at time $0$ of degrees $n,n'\le2p'<\bar n$,
  localized in $\Delta_{0,k},\Delta_{0,k'}\in S_a$, kernel
  $\lambda^2a_na_{n'}h'_k\otimes h'_{k'}\in L^2(\R^2)$;
\item $Q_2=U_{\Delta,n''}$, $1\le n''\le2p'$, $\Delta=\Delta_{l,k''}\in\mathcal D(h,B)$, $\Delta\in
  \mathcal L_j$: type (1.1.6), localized in $\Delta$ alone, degree $n''<\bar n$, kernel
  $\lambda a_{n''}h\mathbf 1_\Delta\in L^2(\R^2)$ (Lemma~\ref{aux:notation-cU}(i));
\item the cutoff $h_\alpha=h^B+\alpha(h-h^B)\in\CE$, $0\le\alpha\le1$
  (Lemma~\ref{aux:notation-cutoffs}(ii));
\item the strip $S_\Delta$ of Step~3, of width $j$: if $\Delta$ is of class (P), \emph{parallel to
  the time-zero line}; if $\Delta$ is of class (T), \emph{transverse to the time-zero line}.
  \emph{Every application to a square of class (T) is the transverse case of GJS
  Theorem~1.1.7, used at the printed statement's scope and disclosed by F-T
  (\S\ref{sec:purpose}); see \S\ref{sub:d4-transverse}.}
\item the resulting term, summed over $n''$: $C_7\,c_U\,\GJS{Q_{n,n',k,k'}}\,e^{-m_1j}$.
\end{itemize}
Lemma~\ref{lem:D2}(a), Lemma~\ref{aux:d4-QV}(iii) and Step~3 give
\[
 \bigl|\ev{Q_V}{h}-\ev{Q_V}{h^B}\bigr|\le C_7\,c_U\,\GJS{Q_V}\,\overline\Sigma_{3\times5}
 \le C_7\,c_U\,\overline\Sigma_{3\times5}\cdot9\,K_1K_2^{8p'}(4p')!\,\lambda^2
 \Bigl(\sum_{n=0}^{2p'}|a_n|\Bigr)^2 ,
\]
where the transverse case enters through the squares of class (T).

\emph{Step 5.}  Steps 1, 2 and 4 give $\int Q_V\,dq_h\le M_B^2$.  None of the constants depends on
$h$, $a$ or $h'$: $S_2(B_{3\times5})$ by Proposition~\ref{aux:notation-S2}(iii), $V_*$ by its
definition as a supremum, and $c_U$, $\overline\Sigma_{3\times5}$, $C_7$, $m_1$, $K_1$, $K_2$ by
construction.
\end{proof}

\subsection{Where the transverse case is used}\label{sub:d4-transverse}

The squares outside $B_{3\times5}$ fall into the two classes of Step~3:
\begin{itemize}
\item \emph{Class (P)}, $|l|=1+j$ for $\Delta\in\mathcal L_j$ (so $|l|\ge2$): separated from $S_a$
by a strip parallel to the time-zero line, of width $j$, which is the width the bound uses.  For
the time-integrated partner $U_{\Delta,n''}$ this is the parallel case that the printed proof
reaches after slicing $U_{\Delta,n''}$ in time (Remark~\ref{aux:d3-orient}), as in \S\ref{sec:d3}.
\item \emph{Class (T) in rows $-1,0,1$}: separated from $S_a$ only by strips that are not
parallel to the time-zero line (Lemma~\ref{aux:notation-rings}(v)); the bound uses the transverse
strip of width $j$.
\item \emph{Class (T) in rows $2\le|l|\le j$}: separated also by parallel strips, but only of width
$\le|l|-1<j$ (Lemma~\ref{aux:notation-rings}(iv)); the bound uses the transverse strip of width
$j$.  Using the parallel strips instead would replace $\overline\Sigma_{3\times5}$ by a sum over
all columns $k$ with $\int_{\Delta_{l,k}}h>0$ of $e^{-m_1(|l|-1)}$, which is not bounded uniformly
in $h\in\CE$ (Remark~\ref{aux:notation-ringrem}(b)).
\end{itemize}
So the transverse case is used for every square of class (T), in rows $-1,0,1$ and beyond, as
Remark~\ref{aux:notation-ringrem}(b) also records; the plan's description (Part~1: ``Squares in rows
$-1$, $0$, $1$ outside $B_{3\times5}$ are separated from the time-zero observable only by strips
transverse to the time-zero line'') names the squares for which no parallel strip exists.  Each of
these applications is Source Fact~2.2 for a sharp-time $Q_1$ (type (1.1.5), at time $0$) and a
time-integrated $Q_2=U_{\Delta,n''}$ (type (1.1.6)) across a strip bounded by two lines of constant
$x^1$.  The printed statement covers it: no orientation is prescribed (``let $d$ be the width of
the widest strip bounded by two parallel lines'', \cite[p.~594]{GJS}), and the record's Source
Fact~2.2 says ``(any orientation)''.  The printed proof does not reach it: it proceeds ``After
rotating and translating the strip separating $Q_1$ and $Q_2$ in Theorem 1.1.7, and after choosing
$A=Q_1$, $B_T=Q_2$ in Theorem 1.1.11'' \cite[p.~597]{GJS}, and Theorem~1.1.11 takes ``monomials of
type (1.1.5)'' localized at times $\leqq0$ and $\geqq0$ \cite[p.~596]{GJS}; a rotation by a right
angle turns the sharp-time piece of $Q_V$ into an object with fixed spatial position and integrated
times, which is not of type (1.1.5) in the rotated frame.  This is the content of the sentence F-T
(\S\ref{sec:purpose}); Theorem~\ref{thm:D4} uses the transverse case at the printed statement's
scope, the standard under which the record imports A1--A8.  Nothing in this note proves that case.

\subsection{Remarks on Theorem~\ref{thm:D4}}\label{sub:d4-rem}

\begin{auxrem}[order of quantifiers]\label{aux:d4-quant}
$M_B$ is fixed first, by $(\lambda,\mathscr P,m_0)$ through the printed constants $C_7$, $m_1$,
$K_1$, $K_2$ and through $S_2(B_{3\times5})$ and $V_*$, in working coordinates (Remark~3.3 of the
record); then the bound holds for every $h\in\CE$, every $a\in\R$ and every admissible $h'$.  In
particular, for $h=h_{T,g}$ it holds for every $g\in\CH$ and every $T>0$ with the same $M_B$.
\end{auxrem}

\begin{auxrem}[boundary cases]\label{aux:d4-boundary}
\emph{$p'=1$}: $\bar n=3$, all individual degrees are $\le2$, and $9^{p'}=9$.
\emph{The position of $a$}: if $a\in\frac12+\Z$, $[a,a+2]$ meets the columns $k_a$ and $k_a+1$ in
positive measure and $k_a+2$ in one point; $B_{3\times5}$ is defined through $k_a$ in every case,
so it has $15$ squares and contains the columns of $h'$ with a margin of at least one column on
each side.
\emph{Ring $1$} (Lemma~\ref{aux:notation-rings}(iii)): $20$ squares, of which the $14$ in rows
$\pm2$ are of class (P) and the $6$ in rows $-1,0,1$ (columns $k_a-2$ and $k_a+4$) of class (T).
\emph{$\supp h\subseteq\overline B_{3\times5}$}: then $\mathcal D(h,B_{3\times5})=\varnothing$, the
comparison term vanishes, and the bound is the base term.
\emph{The coefficient $a_0$}: the measures $dq_h$ do not depend on it
(Lemma~\ref{aux:notation-cU}(iii)), but the bound does, in three places: through $Z(h^B)$ and
$\Uc_{h^B}$ in the base term ($e^{15\lambda|a_0|}$ and $S_2$); through the observable, in $V_*$ (the
term $n=0$ of $V(h')$) and in $\GJS{Q_V}$ (pieces with a factor of degree $0$, covered by
Lemma~\ref{aux:d4-QV}); in $\Uc_{h-h^B}$ it is the constant $\lambda a_0\int(h-h^B)$, which has no
covariance.
\emph{$h'=0$}: both sides are $0$ and $\le M_B^2$.
\end{auxrem}

\begin{auxrem}[what Theorem~\ref{thm:D4} is not, and what it rests on]\label{aux:d4-not}
It does not prove $\UGJS$, and it is not a statement about the status of Theorem~12.2 of the
record: it derives the moment bound that Theorem~5.1 consumes from GJS Theorem~1.1.7 as printed (including
its transverse case) and from the record's Lemma~C.2 and Theorem~C.5; GJS Theorem~1.1.11 is not
used, and none of GJS's proofs is verified (traps~1, 4 and~7 of the plan).  Its conditions are the
status line after the statement.  The transfer to the record's Theorem~5.1 adds one more: the GAP
recorded after Lemma~\ref{aux:notation-wick}, which the record's Theorem~5.1, Step~1, also uses
(\S\ref{sub:cons-51}).  The printed proof of Theorem~1.1.7 takes its constant from Theorem~1.1.8
at the same $h$ \cite[p.~597]{GJS}; that is a fact about the source's proof, which the record does
not re-verify (Scope (c) of Assumption~A8).  $M_B$ is not ``independent of everything but
$(\lambda,\mathscr P,m_0)$'' in an explicit way: it depends on them through $C_7$, $m_1$, $K_1$,
$K_2$ and the unevaluated constant of Theorem~C.5, in working coordinates (Remark~3.3 of the
record) (trap~10).  The Euclidean bound implies the Hamiltonian bound of Theorem~5.1 through
Steps~5--6 of its proof (\S\ref{sec:cons}); the converse is not claimed (trap~9).  The theorem
needs the weak-coupling hypothesis through Theorem~1.1.7 (trap~17).  Lemma~C.2 and Theorem~C.5 are
load-bearing here (trap~12); their independent audit is gate GA4.
\end{auxrem}


\section{Quantifiers and the dependence of the constants}\label{sec:quant}
% quantifiers.tex -- drafted by mp (council member of ugjs/PLAN.md), 2 November 2026.
% Section 7: order of quantifiers, dependence and uniformity of every constant, boundary cases.
% Written from the statements of ugjs/PLAN.md Part 1 and from Sections 2-4 of this note as drafted
% on 2 November 2026 (notation.tex by cqft2, d2.tex by aqft); checked afterwards against d3.tex,
% d4.tex and consequences.tex (aqft) as they stood at 16:04 that day, whose statements agree.

This section states, for every result of the note and for every proof of the record that the
note replaces, the order of the quantifiers and, for every constant, what it depends on and
what it is uniform in. ``Uniform in $X$'' means that one value serves for all $X$ in the stated
range, the remaining data being fixed. The statements are those of
\S\S\ref{sec:d1}--\ref{sec:d4}, which agree with Part~1 of \texttt{ugjs/PLAN.md}.

\subsection{The fixed data}\label{sec:quant-fixed}

The interaction polynomial $\mathscr P(\xi)=\sum_{n=0}^{2p'}a_n\xi^n$, $a_{2p'}>0$,
$p'\ge1$, is fixed. In \S\S\ref{sec:d2}--\ref{sec:d4} and \S\ref{sec:cons} the standing
assumption $0\le\lambda/m_0^2<\varepsilon_{\mathscr P}$ of the record (\S2.6,
\cite[p.~14]{R}) is in force, and the note works in the working coordinates, in which the
pair $(m_0,\lambda)$ satisfies the hypotheses printed with Theorem~1.1.7, ``Let $m_0$, $K_2$
and $K_1$ be sufficiently large and let $\lambda$ be sufficiently small'' \cite[p.~594]{GJS};
Remark~3.3 of the record \cite[p.~16]{R} transports every conclusion back to the original
coordinates. \S\ref{sec:d1} uses neither: Lemma~\ref{lem:D1} holds for every $\lambda\ge0$ and
$m_0>0$. The constants fixed by these data are:
\begin{itemize}
\item[(F1)] $m_1>0$ and $C_7$, the constants of Theorem~1.1.7: ``There are constants $0<m_1$
and $O(1)$ independent of $h$, $Q_1$ and $Q_2$'' \cite[p.~594]{GJS}; the record reads
``independent of $h$'' over the class $\CE$ (Assumption~A8, Scope (d), \cite[p.~13]{R}).
They are fixed once $(m_0,\lambda,\mathscr P)$, $K_1$, $K_2$ and $\bar n$ are fixed by the
choices behind the printed hypotheses; the printed statement makes them independent of $h$,
$Q_1$ and $Q_2$, and says nothing more about their dependence.
\item[(F2)] $K_1$ and $K_2\ge1$, the constants of the printed norms (eq.~(1) of the record,
\cite[p.~10]{R}), and $\bar n=2p'+1$ (Assumption~A8(ii)).
\item[(F3)] The derived constants:
\begin{itemize}
\item $c_U=\lambda\sum_{n=1}^{2p'}|a_n|K_1K_2^{2n}n!$ (Lemma~\ref{aux:notation-cU}),
depending on $\lambda,a_1,\dots,a_{2p'},K_1,K_2$;
\item the stability constant $S_2(B)$ (Definition~\ref{aux:notation-S2def}), finite by
Theorem~C.5 of the record \cite[p.~41]{R}, which is proved conditional on the repairs R1--R10
of the V3 audit (Remark~\ref{aux:d3-V3}), and not evaluated. It depends on $(\lambda,\mathscr P,m_0)$ and on the box $B$
through its shape only (Proposition~\ref{aux:notation-S2}(iii), whose proof uses Lemma~C.4).
Independently of that proposition, the bound that the \emph{proof} of Theorem~C.5 produces
depends on $\Lambda$ only through $|\Lambda|$ and $\operatorname{diam}\Lambda$, monotonically
(V3 audit, items C5.1 and \S6; the statement of Theorem~C.5 does not record this); for a box
this is a bound depending on the shape only;
\item the ring constant
$\overline\Sigma_{3\times5}:=\sum_{j\ge1}(8j+12)e^{-m_1j}=4x(5-3x)(1-x)^{-2}$, $x=e^{-m_1}$
(\S\ref{sec:d4}), depending on $m_1$ only. It bounds the sum
$\Sigma_{3\times5}=\sum_{\Delta\not\subseteq B_{3\times5}}e^{-m_1d(\Delta,R)}$ of
Lemma~\ref{aux:notation-rings}(vi), which depends on $R$, and the sum of $e^{-m_1w_\Delta}$
over the charged squares outside $B_{3\times5}$ with the widths $w_\Delta$ of the axis-parallel
strips assigned in the proof of Theorem~\ref{thm:D4};
\item $V_*$, depending on $(\lambda,\mathscr P,m_0)$, finite (Lemma~\ref{aux:notation-wick}(iv));
the integrability of $C(x-y)^n$ on $[-1,1]^2$ is Lemma~\ref{aux:d1-cov}(ii),(iv).  $V_*$ is
taken for the constructed $V(h')$ of Lemma~\ref{aux:notation-wick}(iv) and needs no GAP.  It
bounds the record's $V(h')$ only through the identification of the record's $V(h)$ with that
object, which is the GAP recorded after Lemma~\ref{aux:notation-wick} (the record constructs no
Wick powers in $V(h)$) and enters at the record's Theorem~5.1, Step~1, and at Corollary~5.2
(\S\ref{sub:cons-51});
\item the factor $9^{p'}=(3^{p'})^2$ of the comparison $\|Y\|_4\le3^{p'}\|Y\|_2$ for $Y$ in the
chaoses of degree $\le2p'$ (proof of Theorem~C.5 of the record; Lemma~\ref{aux:notation-hyper};
the V3 audit reports it as Theorem~I.22 of Simon's book, which this drafter did not read),
depending on $p'$ only; it rests on the chaos structure that the V3 audit found unimported
(its M-3, repair R8), and it is applied to the constructed $Y=J_0V(h')$ of
Lemma~\ref{aux:notation-hyper}, which needs no GAP.
\end{itemize}
\end{itemize}
None of these depends on a cutoff $h$, on $T$, on $g$, on the position of a box, or on an
observable. They are constants of the working coordinates, and Lemma~\ref{lem:D2} and
Theorems~\ref{thm:D3}--\ref{thm:D4} say so (the rider of trap~10).

\subsection{The order of the quantifiers, statement by statement}\label{sec:quant-order}

\emph{Lemma~\ref{lem:D1}.} For all $h_1,h_2\in\CE$ and all $F\in\mathfrak F$ (indeed all
$F\in\bigcup_{p_0>1}L^{p_0}(\dmu)$, Remark~\ref{aux:d1-uses}(a)), the map
$\alpha\mapsto\ev{F}{h_\alpha}$ is $C^1$ on $[0,1]$ and \eqref{eq:d1-identity} holds. There is
no constant; no hypothesis on $\lambda$ or $m_0$ is used.  Status: proved, conditional on repair
R8 of the V3 audit (the linearity of $h\mapsto\Uc_h$ in Step~1), or, on the route written in
Step~1, on Lemma~C.4 of the record (repairs R1--R10); see the status paragraph after
Lemma~\ref{lem:D1}.

\emph{Lemma~\ref{lem:D2}}, in the form (a) of \S\ref{sec:d2}. For every finite family
$(F_a)_{a\in\mathfrak A}$ of localized monomials of type (1.1.5) or (1.1.6) with kernels in
$L^2$ and $n_i<\bar n$, every set $R$ of squares containing their localization squares, every
box $B\supseteq R$, every $h\in\CE$, and every assignment $\Delta\mapsto S_\Delta$,
$\Delta\in\mathcal D(h,B)$, of strips of widths $w_\Delta>0$ separating $\Delta$ from $R$,
\[
 \bigl|\ev{F}{h}-\ev{F}{h^B}\bigr|\le C_7\,c_U\Bigl(\sum_a\GJS{F_a}\Bigr)
 \sum_{\Delta\in\mathcal D(h,B)}e^{-m_1w_\Delta},\qquad F:=\sum_aF_a .
\]
$C_7$, $m_1$ and $c_U$ are those of (F1)--(F3); none depends on $h$, $B$, $R$, the family or the
strips. The right side depends on $h$ only through the finite set $\mathcal D(h,B)$ of the
squares outside $B$ charged by $h$, and on $F$ through the norm of the chosen decomposition
(1.1.9); the bound therefore holds with the infimum over decompositions. Of the cutoff the
lemma uses exactly $h\ge0$, $h\le1$ and compact support (Remark~\ref{aux:d2-uses}(1)): it is a
statement for $h\in\CE$ only, and nothing in the note provides it for cutoffs of variable sign,
for which $h_\alpha\notin\CE$ in general and neither Assumption~A3 nor Theorem~1.1.7 as
imported applies (trap~20 of the plan). The orientation of each application of
Theorem~1.1.7 is that of the strip $S_\Delta$, which is chosen before the bound is applied.
Lemma~\ref{lem:D2} uses Lemma~\ref{lem:D1}, and it is conditional on Lemma~C.4 of the record
(repairs R1--R10, which include R8 and so cover the condition of Lemma~\ref{lem:D1}).  The
condition enters in Step~3 of its proof, which decomposes $\Uc_{h-h^B}$ into the
$U_{\Delta,n}$.  There, Step~1 of the proof of Lemma~\ref{lem:D1} and
Lemma~\ref{aux:d1-moments}(a) give the convergence of the regularizations $\Uc^{(\kappa)}_w$,
$w=h-h^B$, and of their pieces $U^{(\kappa)}_{\Delta,n}$; Lemma~\ref{aux:notation-cutoffs}(iii)
splits $w$ into the squares; and Lemma~C.4 identifies the limit of $\Uc^{(\kappa)}_w$ with
$\Uc_h-\Uc_{h^B}$.  Repair R8 alone makes $h\mapsto\Uc_h$ linear but does not give that
identification as the note is written.

\emph{Theorem~\ref{thm:D3}.} Let $g\in\CH$, let $I\subset\R$ be a bounded interval of positive
length and let $s\in\R$; these determine $l_s$, $N_g$, $N_B$ and $B=B_s(g,I)$, with $|B|=3N_B$, and then
\[
 A:=e^{3\lambda|a_0|N_B}S_2(B),\qquad
 B':=2\,C_7\,c_U\,N_g\,\frac{e^{-m_1}}{1-e^{-m_1}} .
\]
For every finite linear combination $F$ of cylinder monomials (\S6 of the record) supported in
$I$, at the sharp time $s$, and every $T>0$,
$|\int F\,dq_{h_{T,g}}|\le A\|F\|_{L^2(d\varphi_0)}+B'\GJS F$. The order is: $g$ and $I$; then
$A$ and $B'$, which do not depend on $s$ (for $A$ through Proposition~\ref{aux:notation-S2}(iii),
conditional on R1--R10, as in Remark~\ref{aux:d3-quant}); then $s$, $F$ and $T$.
\begin{itemize}
\item $A$ depends on $\lambda$ and $a_0$ through Lemma~C.2 of the record \cite[p.~40]{R},
$Z(h^B)\ge e^{-\lambda|a_0||B|}$ (proved conditional on repair R8 of the V3 audit), and on
$(\lambda,\mathscr P,m_0)$ and $N_B$ through $S_2(B)$ (conditional on R1--R10), which depends
on the shape $3\times N_B$ only. It is uniform in $T>0$, in $s$, in
$F$ and in the position of $B$. It is \emph{not} uniform in $N_B$, which grows with the
diameter of the union of $\supp g$ and $I$, in particular with the distance from $I$ to
$\supp g$ (item m1 of the R3 verification).
\item $B'$ depends on $C_7$, $m_1$, $c_U$ and $N_g$. It is uniform in $T$, in $s$, in $F$
and in the position and the length of $I$; it depends on $g$ only through $N_g$.  The term
$B'\GJS F$ comes from Lemma~\ref{lem:D2}(a), so Theorem~\ref{thm:D3} is conditional on
Lemma~C.4 of the record (repairs R1--R10) through Lemma~\ref{lem:D2}.
\item Both depend on $g$; \S7 of the record fixes $g\in\CH$, so this is all it needs.
\item $\GJS F$ is the norm (1.1.9) of the decomposition of $F$ into its pieces localized in
the squares of row $l_s$ on the columns of \S\ref{sec:notation}
(Lemma~\ref{aux:d3-cyl}(v)), bounded with the column count $N^\square_I$ of
Remark~\ref{aux:notation-counts}; it does not depend on $s$ (Remark~6.3 of the record
\cite[p.~20]{R}).
\item Theorem~\ref{thm:D3} bounds $\int F\,dq_h$ for observables of bounded support, with
the constant $A\|F\|_{L^2}+B'\GJS F$; it is not a form of CE1, since $A$ depends on the
support of $F$ through $N_B$ and both constants depend on $g$ (trap~8).
\end{itemize}

\emph{The consequences of Theorem~\ref{thm:D3} in \S\S4 and 7 of the record.}
\begin{itemize}
\item Lemma~4.3(ii) \cite[p.~17]{R}: at each fixed $T>0$ and $g\in\CH$ the square moments
are finite by Assumption~A3 and H\"older's inequality with the exponents $(2,2)$; no
uniformity in $T$ or $g$ is claimed or used (finding F7 of the plan).
\item Lemma~7.1 \cite[p.~20]{R}: for every $g\in\CH$, every bounded interval $I$ of positive
length and every cylinder monomial $Q$ supported in $I$,
$\sup_{T>0}\|Q\hat\psi_T\|^2\le A\|Q^2\|_{L^2(d\varphi_0)}+B'\mathfrak n(Q)^2$, with $A,B'$ of
Theorem~\ref{thm:D3} at $(g,I,0)$; the same bound holds for $\|Q\Om_g\|^2$. The bound is
uniform in $T$ and depends on $g$ and $I$. Its right side replaces the printed
$C_8\,\mathfrak n(Q)^2$, so the \emph{statement} of Lemma~7.1 changes.
\item Lemmas~7.2--7.3 \cite[pp.~20--21]{R}: the square moments at the time $s$ and at the
times $\pm t/2$ obey the same bound, with $A$ and $B'$ independent of the time and of $T$.
\item Lemma~7.6 \cite[p.~22]{R}: for every $g\in\CH$, every real $f\in C_c^\infty(\R)$,
$f\ne0$, supported in an interval $I$, and every $n\ge1$,
$\E_{\mu_g}[X_f^{2n}]\le A_fK_f^{2n}(2n)!$ with $A_f:=A+B'K_1$ ($A,B'$ at $(g,I,0)$),
$K_f:=\max\{K',\sqrt2\,\sigma_f\}$, $K':=K_2^2(N^\square_I)^{1/2}\|f\|_2$ and
$\sigma_f^2:=\ip{f}{Cf}$. $A_f$ and $K_f$ do not depend on $n$, which is what Theorem~7.7 uses;
$A_f$ depends on $g$ and $I$; $K_f$ depends on $f$ (through $\|f\|_2$, $\sigma_f$ and the column
count of $I$), on $K_2$ and, through $C$, on $m_0$, and not on $g$. The radius $1/K_f$ of the
exponential moment comes from the route through the GJS norm; nothing in the note shows that it
is intrinsic to $\mu_g$ (trap~25). The statement of Lemma~7.6 changes ($C_8K_1$ and $K'$
become $A_f$ and $K_f$).
\item Theorem~7.7 \cite[p.~23]{R}: for given $e_1,\dots,e_k$ the polystrip of its Step~1 can
be taken of half-width $s_0:=\min_{j\le k,\ e_j\ne0}(2kK_{e_j})^{-1}$ ($K_f$ is defined for
$f\ne0$ only); if all $e_j=0$, then $F(z)=\E_{\mu_g}[\overline\zeta]$ is constant and every
$s_0>0$ serves, so $s_0:=+\infty$.  Indeed, by the Cauchy--Schwarz inequality and H\"older's
inequality with $k$ exponents $k$,
\[
 \E_{\mu_g}\Bigl[|\zeta|\,e^{\sum_j|\mathrm{Re}\,z_j||X_{e_j}|}\Bigr]
 \le\|\zeta\|_{L^2(\mu_g)}\prod_{j=1}^k
 \Bigl(\E_{\mu_g}e^{2k|\mathrm{Re}\,z_j||X_{e_j}|}\Bigr)^{1/(2k)},
\]
which is finite when $2k|\mathrm{Re}\,z_j|<1/K_{e_j}$ for every $j$ with $e_j\ne0$ (a factor
with $e_j=0$ is $1$). This $s_0$ depends on $e_1,\dots,e_k$ only, not on $g$. (Part~1 of the plan writes $s_0=1/K_f$, the radius of Lemma~7.6 for one $f$.)
\item Proposition~7.5 and Theorem~7.8 \cite[pp.~22--23]{R}: the gap $m_1$ is the constant of
Theorem~1.1.7 (through Lemma~7.4) and is uniform in $g$; $\widehat C_Q$ depends on $g$, and
only its finiteness is used. Their statements do not change.
\end{itemize}

\emph{Theorem~\ref{thm:D4}.} There is a constant $M_B$, determined by the fixed data
(F1)--(F3),
\[
 M_B^2:=e^{15\lambda|a_0|}S_2(B_{3\times5})\,9^{p'}V_*
 +C_7\,c_U\,\overline\Sigma_{3\times5}\cdot9K_1K_2^{8p'}(4p')!\,\lambda^2
 \Bigl(\sum_{n=0}^{2p'}|a_n|\Bigr)^2,
\]
such that for every $h\in\CE$, every $a\in\R$ and every real measurable $h'$ with
$\|h'\|_\infty\le1$ and $\supp h'\subseteq[a,a+2]$: $\int V(h')(0)^2\,dq_h\le M_B^2$. The order
is: $M_B$ first; then $h$, $a$ and $h'$.
\begin{itemize}
\item $M_B$ depends on $(\lambda,\mathscr P,m_0)$ through $C_7$, $m_1$ (also through
$\overline\Sigma_{3\times5}$), $K_1$, $K_2$, the unevaluated constant $S_2(B_{3\times5})$ of
Theorem~C.5 (conditional on R1--R10) and the finite $V_*$, in the working coordinates
(Remark~3.3 of the record); this is the rider of trap~10, and $M_B$ replaces the record's $M$,
which is not unchanged.  Theorem~\ref{thm:D4} bounds the constructed observable,
$V(h')(0)=J_0V(h')$ with $V(h')$ the limit of Lemma~\ref{aux:notation-wick}(iv); its base term
and its comparison term concern one element of $L^2(\dmu)$, because the two regularizations of a
sharp-time Wick power agree (Lemma~\ref{aux:notation-sharpid} and
Remark~\ref{aux:notation-sharpidrem}(b)).  So $M_B$ and the Euclidean bound need no GAP; through
$9^{p'}$ (Lemma~\ref{aux:notation-hyper}(ii)) and $e^{15\lambda|a_0|}$ they are conditional on
repair R8, and through the comparison term on Lemma~C.4 of the record (repairs R1--R10), through
Lemma~\ref{lem:D2}.  The GAP recorded after Lemma~\ref{aux:notation-wick} enters only where $M_B$ is used
for the record's $V(h)$: at Theorem~5.1, Step~1, and at Corollary~5.2 (\S\ref{sub:cons-51}).
\item It is uniform in $h\in\CE$, hence in $T>0$ and $g\in\CH$ when $h=h_{T,g}$; in $a$,
because the boxes $B_{3\times5}(a)$ (Definition~\ref{aux:notation-ringdef}) all have the shape
$3\times5$ and $S_2$ depends on the shape only (Proposition~\ref{aux:notation-S2}(iii), through
Lemma~C.4, conditional on R1--R10); and in $h'$.
\item The factors:
\begin{itemize}
\item $e^{15\lambda|a_0|}$ is Lemma~C.2 with $|B_{3\times5}|=15$ (conditional on R8);
\item $S_2(B_{3\times5})$ bounds $\|e^{-\Uc_{h^B}}\|_{L^2(\dmu)}$
(Proposition~\ref{aux:notation-S2}(iv));
\item $9^{p'}V_*$ bounds $\|V(h')\|_{L^4(d\varphi_0)}^2$, which equals
$\|V(h')^2\|_{L^2(d\varphi_0)}$, for the constructed $V(h')$ (Lemma~\ref{aux:notation-hyper};
conditional on R8; no GAP);
\item $C_7c_U\overline\Sigma_{3\times5}$ is Lemma~\ref{lem:D2}(a) summed over the rings
(conditional on Lemma~C.4, R1--R10, through Lemma~\ref{lem:D2}),
$\overline\Sigma_{3\times5}$ bounding the sum of $e^{-m_1w_\Delta}$ over the charged squares
outside $B_{3\times5}$;
\item the last factor bounds $\GJS{Q_V}$ as in eq.~(5) of the record \cite[p.~19]{R}, whose
``$=M^2/C_8$'' no longer applies (plan, item n14).
\end{itemize}
\item \emph{Where the transverse case enters.} Every square of rows $-1,0,1$ outside
$B_{3\times5}$ is separated from $R$ only by transverse strips of positive width
(Lemma~\ref{aux:notation-rings}(v)); every side square of the ring $\mathcal L_j$ in a row
$2\le|l|\le j$ admits parallel strips only of width $\le|l|-1<j$, and the bound
$d(\Delta,R)\ge j$ uses a transverse strip for it too
(Remark~\ref{aux:notation-ringrem}(b)). The squares of the rows $|l|=j+1$ of $\mathcal L_j$
are reached by parallel strips of width $j$. Each application with a transverse strip falls
under sentence F-T (\S\ref{sec:purpose}).
\item Theorem~\ref{thm:D4} rests on the printed hypotheses of Theorem~1.1.7; Phase~B does not
remove the weak-coupling restriction (trap~17).
\end{itemize}

\emph{Theorem~5.1 with $M:=M_B$, Corollary~5.2, and \S\S8--14 of the record.}
\begin{itemize}
\item Theorem~5.1 \cite[p.~18]{R}: for every $g\in\CH$ and every real measurable $h'$ with
$\|h'\|_\infty\le1$ and $\supp h'$ in an interval of length $2$ (which lies in some
$[a,a+2]$), $\Om_g\in D(V(h'))$ and $\|V(h')\Om_g\|\le M_B$. Step~5 applies
Theorem~\ref{thm:D4} at $h=h_{T,g}\in\CE$ for every $T>0$; Step~6 lets $T_k\to\infty$ at fixed
$g$. The order ``there is $M_B$ such that for all $g$ and $h'$'' is the content of (M$_2$); the
slab bound implies (M$_2$), and the converse is not claimed (trap~9).  Through
Theorem~\ref{thm:D4}, Theorem~5.1, and with it Corollary~5.2, Lemma~10.2, Theorem~11.1 and
items (1) and (5) of Theorem~12.2, rest on the transverse case of Theorem~1.1.7 at the printed
statement's scope (sentence F-T, \S\ref{sec:purpose}).  They rest on the GAP recorded after
Lemma~\ref{aux:notation-wick} as well, which enters at Step~1 (``$Q_V$ is exactly the
multiplication function $V(h)^2$'', main.tex l.~1278): there the record's $V(h')$ is identified
with the constructed $V(h')$ of Lemma~\ref{aux:notation-wick}(iv), to which
Theorem~\ref{thm:D4} applies (\S\ref{sub:cons-51}).  Lemma~\ref{aux:notation-sharpid} and
Remark~\ref{aux:notation-sharpidrem} identify the two constructions of the note with each other;
they do not reach the record's $V(h)$.
\item Corollary~5.2 \cite[p.~19]{R}: its number $362\,880\,C_8K_1K_2^{16}\lambda^2$ retires
with $C_8$. For $\mathscr P(\xi)=\xi^4$ ($p'=2$, $a_4=1$, the other $a_n=0$) the formula above
gives
\[
 M_B^2=81\,S_2(B_{3\times5})\,V_*+362\,880\,C_7\,c_U\,\overline\Sigma_{3\times5}\,
 K_1K_2^{16}\lambda^2,
\]
with $c_U=24\lambda K_1K_2^8$ and $V_*\le24\lambda^2\iint_{[-1,1]^2}C(x-y)^4\,dx\,dy$.  As a
bound for the record's $V(h)$, Corollary~5.2 inherits from Theorem~5.1 the transverse case (F-T)
and the GAP at Step~1; $M_B$ itself needs no GAP.
\item \S\S8--14 consume the local moment only through Lemma~10.2 \cite[p.~28]{R}, which applies
Theorem~5.1 at the cutoffs $g_s\in\CH$ of an affine path and at $h_j$ supported in
$[j-1,j+1]$; with $M:=M_B$, which is uniform in $g$, Lemma~10.2 holds verbatim (cf.\ trap~11).
The constants $A$, $B'$, $A_f$, $K_f$ and $s_0$ do not leave \S7 (gate GA6,
\S\ref{sec:cons}).
\item In the original coordinates the conclusions are transported by Remark~3.3 of the
record; for the local moment this means supports of length $2/s$, or a partition of unity into
at most $\lceil s\rceil+1$ pieces and the constant $(\lceil s\rceil+1)M_B$, where
$s=m_0/M_*$.
\end{itemize}

\subsection{Boundary cases}\label{sec:quant-boundary}

Squares, rows, columns, $l_s$ and the charged squares $\mathcal D(h,B)$ are those of
\S\S\ref{sec:notation}--\ref{sec:d2}: $\Delta_{l,k}=[l-\frac12,l+\frac12)\times
[k-\frac12,k+\frac12)$, $l_s=\lfloor s+\frac12\rfloor$, and $\mathcal D(h,B)$ the squares
$\Delta\not\subset B$ with $\int_\Delta h>0$.

\emph{(1) $T<1$: a slab thinner than one row.} $h_{T,g}=\mathbf 1_{[-T,T]}\otimes g$ charges at
most the rows $l$ with $|[l-\frac12,l+\frac12)\cap[-T,T]|>0$, i.e.\ $|l|<T+\frac12$: row $0$
alone for $T\le\frac12$, rows $-1,0,1$ for $\frac12<T<1$. In Theorem~\ref{thm:D3} with
$l_s=0$, the box $B$ (rows $-1,0,1$, all columns meeting $\supp g\cup I$ in positive measure)
then contains every charged square, $\mathcal D(h_{T,g},B)=\varnothing$, $h^B=h_{T,g}$, and the
bound is the base term $A\|F\|_{L^2}$ alone. For other $l_s$ the charged squares outside $B$ lie
at row distance $\ge2$ from $l_s$ and the parallel strips of the general case apply.
Theorem~\ref{thm:D3} needs no lower bound on $T$. In Theorem~\ref{thm:D4} with $T\le\frac12$
the only charged row is $0$, and every charged square outside $B_{3\times5}$ lies in row $0$
and is reached only by a transverse strip: the transverse case is used even for the thinnest
slabs. (The fallback Lemma~L of the plan needs $T\ge t_*=1/m_0$; Theorems~\ref{thm:D3}
and~\ref{thm:D4} do not.)

\emph{(2) $T$ not an integer.} The rows at the edge of the slab are charged only in part, so
$h_{T,g}\mathbf 1_\Delta\ne g$ on such squares; $c_U$ uses only
$\|h\mathbf 1_\Delta\|_2\le1$ (Lemma~\ref{aux:notation-cU}(i)), and no constant depends on $T$.
If $T\in\frac12+\Z$, the row $l=T+\frac12$ meets $[-T,T]$ only in the null set $\{x^0=T\}$, is
not charged, and carries $U_\Delta=0$; it needs no strip (Remark~\ref{aux:d2-uses}(3)).

\emph{(3) $s$ on a lattice line.} If $s=l-\frac12$ with $l\in\Z$, then $l_s=l$
(Definition~\ref{aux:notation-squares}): the observable is localized in row $l$, although the
line $x^0=s$ is also the upper edge of row $l-1$. The box $B$ of Theorem~\ref{thm:D3} contains
rows $l-1$, $l$, $l+1$, hence both rows touching the line. A charged square in row $l'$ with
$|l'-l|\ge2$ is separated from the row-$l$ squares by the parallel strip between
$x^0=\min\{l,l'\}+\frac12$ and $x^0=\max\{l,l'\}-\frac12$, of width $|l'-l|-1$; it separates the
closed squares, in particular the line $x^0=s$ itself. For $l'\le l-2$ the line lies on the
boundary of that strip; for $l'\ge l+2$ the true distance is one unit larger than the width
used. The constants are unchanged.

\emph{(4) $N_g=1$, and $N_g=0$.} If $\supp g$ meets one column in positive measure, the charged
squares outside $B$ lie in that column, in rows $l$ with $|l-l_s|\ge2$, at most two at each
strip width $j=|l-l_s|-1\ge1$, and
$B'=2C_7c_Ue^{-m_1}(1-e^{-m_1})^{-1}$. If $g=0$ almost everywhere, no square is charged,
$dq_{h_{T,g}}=\dmu$ (since $\Uc_{h_{T,g}}=0$ and $Z=1$), the comparison term vanishes, and the
theorem is the bound $|\int F\,d\mu_0|\le\|F\|_{L^2}\le A\|F\|_{L^2}$ ($S_2\ge1$). The count
must be taken with the charged squares of Lemma~\ref{lem:D2}(a). With the index set
$\{\Delta\not\subset B:\Delta\cap\supp h\ne\varnothing\}$ of its form (b), a column that meets
$\supp g$ only in a null set would contribute squares in the rows $l_s-1$, $l_s$, $l_s+1$
outside $B$, which only transverse strips separate from $R$; for instance
$g:=\mathbf 1_{[0,1/10]}+\mathbf 1_{\{3\}}\in\CH$ has $N_g=1$ while column $3$ meets
$\supp g=[0,\frac1{10}]\cup\{3\}$. Such columns carry $U_\Delta=0$, and form (a) omits them
(Remark~\ref{aux:d2-uses}(3), which records this example).

\emph{(5) $I$ meeting one column only.} Then $N^\square_I=1$, every cylinder monomial of $F$ has
a single localized piece, all its factors in one square $\Delta_{l_s,k}$, and
$\GJS F$ is bounded by (N.a) of \S\ref{sec:notation} with count $1$. If $\supp g$ meets only the
same column, $N_B=1$, $|B|=3$ and $A=e^{3\lambda|a_0|}S_2(B)$. If $I$ touches a second column in
a null set (e.g.\ $I=[k-\frac12,k+\frac12]$), the pieces there vanish in $L^2$ and give no
localization square. On the record's cells $[j,j+1)$ the count can differ
(Remark~\ref{aux:notation-counts}(ii)); the note counts on the columns.

\emph{(6) $p'=1$.} $\bar n=3$; the pieces $U_{\Delta,n}$ have $n\in\{1,2\}$, the factors of $Q_V$
have degrees $\le2<3$, and its local degrees are $\le4$, which (1.1.7) permits (the bound is on
the individual $n_i$). $c_U=\lambda K_1(|a_1|K_2^2+2|a_2|K_2^4)$, $9^{p'}=9$, and
$9K_1K_2^{8p'}(4p')!=216\,K_1K_2^8$. Nothing degenerates.

\emph{(7) Ring $k=1$.} $\mathcal L_1$ has $8\cdot1+12=20$ squares
(Lemma~\ref{aux:notation-rings}(iii)): the fourteen squares of rows $\pm2$ in the columns
$k_a-2,\dots,k_a+4$, reached by parallel strips of width $1$; and the six squares of rows
$-1,0,1$ in the columns $k_a-2$ and $k_a+4$, reached by transverse strips of width $1$ and by no
parallel strip of positive width. The four corners (rows $\pm2$ in the two side columns) admit
both. The ring contributes the term $20e^{-m_1}$ of $\overline\Sigma_{3\times5}$, of which
$6e^{-m_1}$ rest on the transverse case alone.

\emph{(8) $a$ such that $[a,a+2]$ meets three columns.} $[a,a+2]$ meets exactly the three
columns $k_a,k_a+1,k_a+2$ (Lemma~\ref{aux:notation-rings}(i)); it meets all three in positive
measure unless $a\in\frac12+\Z$, when it meets $k_a+2$ only in the point $a+2$. In that case
the pieces of $Q_V$ in column $k_a+2$ vanish in $L^2$ and $R$ lies in two columns, while
$B_{3\times5}(a)$ is built on all three, which only enlarges the box: the base term and the ring
count are unchanged. Under the ``positive measure'' reading the box would be $3\times4$ for
$a\in\frac12+\Z$; its rings have $(3+2j)(4+2j)-(1+2j)(2+2j)=8j+10\le8j+12$ squares, each at
strip width $\ge j$ from $R$ by Lemma~\ref{aux:notation-strip}, so that
$\overline\Sigma_{3\times5}$ still bounds their sum; its base
term has $e^{12\lambda|a_0|}\le e^{15\lambda|a_0|}$, and its $S_2$ is at most $S_2(B_{3\times5})$
(Proposition~\ref{aux:notation-S2}(ii), the $3\times4$ box lying in a $3\times5$ box). So $M_B$
is valid under both readings, and the record's count in Theorem~5.1, Step~2 (``at most three
half-open unit spatial cells \dots\ in positive measure'') holds on the columns as well.  (The
ring counts $8j+10$ with strip widths $\ge j$ for $j\le12$, and the rows charged by a slab in case
(1), are also checked in exact integer and rational arithmetic by \texttt{parts/check\_d1.py},
parts (4)--(5); \emph{numerical}, run record \texttt{parts/}\allowbreak\texttt{check\_d1.output.txt}, 3 November 2026.)

\emph{(9) The term $n=0$ of $\mathscr P$.} $U_{\Delta,0}=\lambda a_0\int_\Delta h$ is a
constant, with zero covariance under every $dq_{h'}$ (Lemma~\ref{aux:notation-cU}(iii)), and
$c_U$ starts at $n=1$. The coefficient $a_0$ does not change $dq_h$. As a coefficient of the
interaction it enters $A$ and $M_B$ only through $e^{3\lambda|a_0|N_B}$ and $e^{15\lambda|a_0|}$
(Lemma~C.2, conditional on R8) and through $S_2$, an artefact of bounding
$\|e^{-\Uc_{h^B}}\|_2$ and $Z(h^B)$ separately; with $\mathscr P-a_0$ in place of $\mathscr P$
every $dq_h$ is the same and those factors disappear. As a coefficient of the observable it
enters $M_B$ through $V_*$ and $\GJS{Q_V}$, where the term $n=0$ stays: $V(h')$ contains
$\lambda a_0\int h'$, and the sums $\sum_{n=0}^{2p'}$ include it.

\emph{(10) $\lambda=0$.} Then $\Uc_h=0$, $dq_h=\dmu$ for every $h$, both sides of
Lemma~\ref{lem:D2} vanish, $c_U=V_*=0$, $S_2=1$, and $M_B=0=\|V(h')\Om_g\|$.

\emph{(11) Lemma~\ref{lem:D1} at its edges.} For $h_1=h_2$ both sides of
\eqref{eq:d1-identity} vanish ($\Uc_{h_2-h_1}=0$); for constant $F$ the covariance vanishes and
$\ev{F}{h}$ is constant; at $\alpha\in\{0,1\}$ the derivatives are one-sided, and Step~3 of the
proof covers them, the dominator not depending on $\alpha$.


\section{Consequences for the record: the replaced proofs}\label{sec:cons}
% consequences.tex -- drafted by aqft (council member of ugjs/PLAN.md), 2 November 2026 (begun 31 October 2026);
% revised 3 November 2026 after the audits V1 (cqft, mp) and V2.
% Section 8: the replaced proofs of the record.  Each replacement text is enclosed between the
% comment lines "% BEGIN-REPLACEMENT <name>" and "% END-REPLACEMENT", so that gate GA5 can be
% checked by extracting exactly those blocks (\S\ref{sub:cons-gates}).  Auxiliary statements
% aux:consequences-*.  Line numbers: part_i/main.tex as of 2 November 2026 (3181 lines).

This section writes, for each result of the record that consumes Assumption~A8(vi) (equation
$(\UGJS)$ of the record) or its constant $C_8$, the record's current text with its line numbers in
\texttt{part\_i/main.tex}, the replacement text, and the checks of gates GA5 and GA6 of the plan.
Phase~A (Theorem~\ref{thm:D3}) replaces the proofs of Lemma~4.3(ii) and Lemmas~7.1, 7.2, 7.3
and~7.6 and fixes the radius in the proof of Theorem~7.7; Phase~B (Theorem~\ref{thm:D4}) replaces
Step~5 of the proof of Theorem~5.1, with $M:=M_B$, and Corollary~5.2.  The replaced proofs need
the count $N_I$ on GJS's columns (\S\ref{sub:cons-NI}).  Record numbers are those of the label
table of 31 October 2026 (unchanged in the build of 2 November 2026).

\emph{Which statements change.}  Through the edit of \S\ref{sub:cons-NI}: the statement of
Lemma~6.1 (its counterexample sentence, ll.~1364--1366) and the content of Proposition~6.2,
eqs.~(7)--(8), through the meaning of $N_I$; the text of the notation table (l.~448) and of
Assumption~A8(ii) (ll.~659--663) changes too.  Through Theorems~\ref{thm:D3} and~\ref{thm:D4}:
the statement of Lemma~7.1 changes its form, gaining the term $A\|Q^2\|_{L^2(d\phi_0)}$, with
$B'$ in place of $C_8$; the statement of Lemma~7.6 has $A_f$ and $K_f$ in place of $C_8K_1$ and
$K'$; the statements of Theorem~5.1 and Corollary~5.2 keep their form, with the value
$M:=M_B$ (main.tex ll.~1262, 1346, 1430--1431, 1609--1617 contain $C_8$).  The statement of
Theorem~7.7 is unchanged; the radius of the polystrip in its proof becomes explicit
(\S\ref{sub:cons-76}).  The imports paragraph of Theorem~12.2 changes at integration
(\S\ref{sub:cons-rest}).  Every other numbered statement of the record is unchanged; the
constants of \S\S5 and~7 that reach statements of \S\S8--14 are $m_1$ (through Theorem~7.8,
unchanged) and $M$, whose value becomes $M_B$ (\S\ref{sub:cons-gates}, gate GA6).

\emph{Dependence on Appendix~C, on the transverse case and on the GAP.}  The replaced proofs of
Lemmas~7.1--7.3, Lemma~7.6 (and through it the radius in the proof of Theorem~7.7) and
Theorem~5.1 consume Lemma~C.2 and Theorem~C.5 of the record, and Lemmas~\ref{lem:D1}
and~\ref{lem:D2}, through Theorems~\ref{thm:D3} and~\ref{thm:D4}, as in Remark~\ref{aux:d3-V3}:
Lemma~C.2 is proved conditional on repair R8 of the V3 audit, Lemma~\ref{lem:D2} (with
Lemma~\ref{lem:D1}) conditional on Lemma~C.4, and Theorem~C.5 and Lemma~C.4 conditional on its
repairs R1--R10; with the reading of Remark~\ref{aux:notation-gjslattice}(iv), the conditions are
listed in the status lines after Lemma~\ref{lem:D2} and Theorems~\ref{thm:D3} and~\ref{thm:D4}.  The replaced proof of Theorem~5.1 also uses, through
Theorem~\ref{thm:D4}, the transverse case of GJS Theorem~1.1.7 (F-T, \S\ref{sec:purpose}), and the
GAP recorded after Lemma~\ref{aux:notation-wick}: Theorem~\ref{thm:D4} bounds the constructed
$V(h)$ of Lemma~\ref{aux:notation-wick}(iv), and the record's $V(h)$ is identified with it by that
GAP, on which the record's Step~1 (``$Q_V$ is exactly the multiplication function $V(h)^2$'',
ll.~1278--1283), which is kept, rests as well.  The replaced proof of Lemma~4.3(ii) uses
Assumption~A3, H\"older's inequality, the sharp-time law of cylinder monomials
(Lemma~\ref{aux:d3-sharp}, from Assumption~A2 and Remark~2.5 of the record) and Assumption~A5(i);
it applies its moment bound to the cylinder monomials of \S7 only and keeps A5(i) for the
interaction-type multiplications, so it does not use the GAP.

\subsection{An edit needed by the replaced proofs: the count $N_I$}\label{sub:cons-NI}

The record defines $N_I:=\#\{j\in\Z:|I\cap[j,j+1)|>0\}$ (main.tex l.~1359, and in the notation
table, l.~448), counting the cells $[j,j+1)$, and proves Proposition~6.2 by decomposing over those
cells (ll.~1393--1409); Theorem~5.1, Steps~2--4 (ll.~1285--1315), count the same cells.  On GJS's
squares the cells are not localization regions; the same computation gives (7)--(8) with the
number of columns $[k-\frac12,k+\frac12)$ meeting $I$ in positive measure, and the two counts
differ for some $I$ (Remark~\ref{aux:d3-cells}, with $I=(0.4,0.6)$).  The replaced proofs below
need (8) for GJS's norm, so at integration:
% BEGIN-REPLACEMENT sec6-NI
\begin{quote}\small
\emph{Notation table, l.~448, replaced:} $N_I$ \& number of columns $[k-\frac12,k+\frac12)$ of the
squares of p.~593 meeting the interval $I$ in positive measure.
\emph{Assumption~A8(ii), ll.~659--663, the sentence beginning ``The cells'' replaced:} The spatial
cells used in \S\S5--6 are the columns $[k-\frac12,k+\frac12)$ of these squares; the count bounds
used below (at most three columns for a support of length $2$, Theorem~5.1, Step~2;
$N_I\le\lfloor L\rfloor+2$, Lemma~6.1) are proved on them.
\emph{\S6, l.~1359, replaced:} $N_I:=\#\{k\in\Z:\ |I\cap[k-\tfrac12,k+\tfrac12)|>0\}$, the number
of columns of the squares of \cite[p.~593]{GJS} that meet $I$ in positive measure.
\emph{Lemma~6.1, ll.~1364--1366, the counterexample replaced:} the interval $I=(0.4,0.6)$ has
length $0.2$ and meets both $[-\frac12,\frac12)$ and $[\frac12,\frac32)$ in positive measure.
\emph{Proof of Lemma~6.1 and of Proposition~6.2:} the cells $[j,j+1)$ replaced by
$[k-\frac12,k+\frac12)$ throughout (the proofs are otherwise verbatim;
Remark~\ref{aux:notation-counts}(i),(iii)).
\emph{Theorem~5.1, Steps~2--4:} the cells $[j,j+1)$ replaced by the columns
$[k-\frac12,k+\frac12)$ throughout; the arguments are otherwise verbatim (in Step~2, four columns
met in positive measure would give infimum $<k_{\min}+\frac12$ and supremum
$>k_{\max}-\frac12\ge k_{\min}+\frac52$, hence length $>2$).
\end{quote}
% END-REPLACEMENT
With this edit $N_I=N^\square_I$ of Remark~\ref{aux:notation-counts}, (8) is the second bound of
(N.a) there (applied in Lemma~\ref{aux:d3-cyl}(iii)), the bound $N_I\le\lfloor L\rfloor+2$ of
Lemma~6.1 and the Appendix~A row for $N_I$ are unchanged, and Steps~2--4 of Theorem~5.1 give the
bound of Lemma~\ref{aux:d4-QV}(iii).  As the sentence of Assumption~A8(ii) stands, its first
clause (the cells ``used in \S\S5--6'') would be false after the edit of \S6, and its conclusion
``so the same computation on the squares of p.~593 gives the same constants'' is false for
$\mathfrak n(Q)$ and $K'$ (finding, \S\ref{sub:cons-findings}); it is therefore replaced, not kept.

\subsection{Lemma 4.3(ii) (Phase A; finding F7 of the plan)}\label{sub:cons-43}

\emph{Statement of (ii)} (main.tex ll.~1188--1202; Lemma~4.3 of the record \cite[p.~17]{R}):
unchanged.

\emph{Current proof of (ii), the lines that cite $(\UGJS)$} (ll.~1207--1218):
\begin{quote}\small
``(ii) The Hilbert-space limit is dominated convergence as displayed. Membership
$\psi_a\in D(F)$ is supplied at each application: finiteness of the corresponding $t=0$ identity
of part (i) over the $a$-slab, together with the slab bound $(\UGJS)$, forces $F\psi_a\in\HH$
(this is how it is verified in Theorem~5.1, Step~5, and in Lemma~7.1); for the interaction-type
multiplications it also follows independently from Assumption~A5(i). On path space,
$F_i^{[\ell]}(\Phi(s_i,\cdot))\to F_i(\Phi(s_i,\cdot))$ pointwise with the displayed majorant,
whose integrability is exactly the finiteness of the two square moments; the cutoff-uniform
estimate $(\UGJS)$ supplies those square moments at every later application (\S5, \S7). Dominated
convergence then passes Lemma~4.2 to the limit on both sides.''
\end{quote}

\emph{Replacement.}
% BEGIN-REPLACEMENT lemma-4.3-ii
\begin{quote}\small
(ii) The Hilbert-space limit is dominated convergence as displayed.  Membership
$\psi_a\in D(F)$ is supplied at each application.  For the interaction-type multiplications it
follows from Assumption~A5(i).  For the cylinder monomials of \S7, part (i) at $t=0$ over the
$a$-slab gives $\|F\psi_a\|^2/\|\psi_a\|^2=\int F(\Phi(0,\cdot))^2\,dq_{h_{a,g}}$ in $[0,\infty]$;
at every sharp time $s$ a cylinder monomial has under $d\mu_0$ the law of the same polynomial in
the Gaussian variables $\varphi_0(f_i)$ under $d\phi_0$ (Appendix~D, \S5), so
$F(\Phi(s,\cdot))\in L^4(d\mu_0)$, and H\"older's inequality with exponents $(2,2)$ and
Assumption~A3 give, for every $h\in\CE$,
\[
 \int F(\Phi(s,\cdot))^2\,dq_h\le Z(h)^{-1}\,\|F(\Phi(s,\cdot))\|_{L^4(d\mu_0)}^2\,
 \|e^{-\mathcal U_h}\|_{L^2(d\mu_0)}<\infty ,
\]
since $Z(h)>0$ and $e^{-2\mathcal U_h}\in L^1(d\mu_0)$ by Assumption~A3.  With $h=h_{a,g}$ and
$s=0$ this gives $F\psi_a\in\HH$.  On path space,
$F_i^{[\ell]}(\Phi(s_i,\cdot))\to F_i(\Phi(s_i,\cdot))$ pointwise with the displayed majorant,
whose integrability is exactly the finiteness of the two square moments; for the cylinder
monomials of \S7, at each fixed $T$ the same bound with $h=h_{T,g}$ supplies them, and no
uniformity in $T$ is used here.  Dominated convergence then passes Lemma~4.2 to the limit on both
sides.
\end{quote}
% END-REPLACEMENT
\emph{Check.}  The replaced proof uses Assumption~A3 ($Z(h)>0$; $e^{-\varrho\,\mathcal U_h}\in L^1$
with $\varrho=2$), H\"older with exponents $(2,2)$ on $\int F^2e^{-\mathcal U_h}d\mu_0\le
\|F^2\|_2\|e^{-\mathcal U_h}\|_2$, the sharp-time law of Lemma~\ref{aux:d3-sharp} (Assumption~A2
and Remark~2.5 of the record), and Assumption~A5(i) as before; it cites neither $(\UGJS)$ nor
$C_8$.  The record applies part (ii) only to cylinder monomials (Lemmas~7.2 and~7.3, main.tex
ll.~1462 and~1501); for $|V(h)|$ in \S5 it uses part (i) only (Theorem~5.1, Step~5, l.~1320), so
the restriction of the moment bound to cylinder monomials loses nothing, and the GAP recorded after
Lemma~\ref{aux:notation-wick} is not used.  Lemma~4.3(ii) needs finiteness at each $T$ only, as
finding F7 of the plan says; the $T$-uniform bounds of \S\S5 and~7 are supplied where they are
used, by Theorems~\ref{thm:D3} and~\ref{thm:D4}.

\subsection{Lemma 7.1 (Phase A)}\label{sub:cons-71}

\emph{Current statement} (ll.~1428--1432; Lemma~7.1 \cite[p.~20]{R}):
``$\sup_{T>0}\|Q\hat\psi_T\|^2\le C_8\,\mathfrak n(Q)^2$; moreover $\Om_g\in D(Q)$ and
$\|Q\Om_g\|^2\le C_8\,\mathfrak n(Q)^2$.''
\emph{Current proof} (ll.~1434--1450): ``Apply (2) at $t=0$ with $F_1=F_2=|Q|$ (Lemma~4.3(i)); the
Euclidean integrand is $Q^2$ at time $0$, so $(\UGJS)$, Remark~6.3 and (8) give
$\|Q\hat\psi_T\|^2=\int Q(0)^2\,dq_{h_{T,g}}\le C_8\|Q^2\|_{\rm GJS}\le C_8\,\mathfrak n(Q)^2$ for
every $T$.'' [weak-graph passage] ``\dots weak lower semicontinuity gives
$\|Q\Om_g\|\le\liminf_\alpha\|Q\widehat\psi_{T_\alpha}\|$le$C_8^{1/2}\mathfrak n(Q)$.''  (The
``le'' is the source line's \verb+\|le+, finding F2 of the plan, l.~1448.)

\emph{Replacement.}
% BEGIN-REPLACEMENT lemma-7.1
\begin{quote}\small
\textsc{Lemma 7.1} (moment transfer).  \emph{If the interval $I$ of \S6 has length $L>0$, let $A$
and $B'$ be the constants of Theorem~D.3 for the cutoff $g$, the interval $I$ and the time $s=0$;
if $L=0$, put $A:=B':=0$ (then every $f_i$ vanishes and $Q=0$).  Then
$\sup_{T>0}\|Q\hat\psi_T\|^2\le A\,\|Q^2\|_{L^2(d\phi_0)}+B'\,\mathfrak n(Q)^2$; moreover
$\Om_g\in D(Q)$ and $\|Q\Om_g\|^2\le A\,\|Q^2\|_{L^2(d\phi_0)}+B'\,\mathfrak n(Q)^2$.}

\emph{Proof.}  If $L=0$, every $f_i\in C_c^\infty(\R)$ vanishes outside a set of at most one
point, hence everywhere, so $Q=0$ and there is nothing to prove.  Let $L>0$.  Apply (2) at $t=0$
with $F_1=F_2=|Q|$ (Lemma~4.3(i)); the Euclidean integrand is $Q^2$ at time $0$.  $Q^2$ is a
cylinder monomial supported in $I$ (its factors are $f_1,\dots,f_r,f_1,\dots,f_r$), so
Theorem~D.3 with $F=Q^2$ and $s=0$, and (8), give for every $T>0$
\begin{align*}
 \|Q\hat\psi_T\|^2=\int Q(0)^2\,dq_{h_{T,g}}
 &\le A\,\|Q^2\|_{L^2(d\phi_0)}+B'\,\|Q(0)^2\|_{\rm GJS}\\
 &\le A\,\|Q^2\|_{L^2(d\phi_0)}+B'\,\mathfrak n(Q)^2 ,
\end{align*}
the right side being finite by the Gaussian moments of the free field.  Choose from the bounded
family $Q\hat\psi_T$ a weakly convergent subnet $Q\hat\psi_{T_\alpha}\rightharpoonup y$.  Since
$\hat\psi_{T_\alpha}\to\Om_g$ strongly (Lemma~4.4) and the graph of the closed multiplication
operator $Q$ is a closed linear subspace of $\HH\oplus\HH$, hence weakly closed, $(\Om_g,y)$ lies
in that graph.  Thus $\Om_g\in D(Q)$, $y=Q\Om_g$, and weak lower semicontinuity gives
$\|Q\Om_g\|\le\liminf_\alpha\|Q\hat\psi_{T_\alpha}\|\le\bigl(A\|Q^2\|_{L^2(d\phi_0)}+
B'\mathfrak n(Q)^2\bigr)^{1/2}$.
\end{quote}
% END-REPLACEMENT
\emph{Check.}  For $L>0$, Theorem~\ref{thm:D3} applies: $g\in\CH$, $I$ bounded of positive
length, $F=Q^2$ a cylinder monomial supported in $I$, $s=0$, every $T>0$; its constants are
$A=e^{3\lambda|a_0|N_B}S_2(B_0(g,I))$ and $B'=2C_7c_UN_ge^{-m_1}(1-e^{-m_1})^{-1}$, independent of
$T$.  $\GJS{Q(0)^2}\le\mathfrak n(Q)^2$ is (8) with $N_I$ as in \S\ref{sub:cons-NI}
(Lemma~\ref{aux:d3-cyl}(iii)); Remark~6.3 is no longer needed, since the bound of
Lemma~\ref{aux:d3-cyl}(iii) holds at every time.  The case $L=0$ is outside the hypotheses of
Theorem~\ref{thm:D3} and is trivial.  The replaced text cites neither $(\UGJS)$ nor $C_8$, and
the misprint F2 is gone.

\subsection{Lemmas 7.2 and 7.3 (Phase A)}\label{sub:cons-7273}

\emph{Lemma 7.2, current lines} (ll.~1462--1464; \cite[p.~20]{R}): ``For $T>|s|$, (3) (extended to
$Q$ by Lemma~4.3; the square moments $\int Q(s)^2\,dq_{h_{T,g}}\le C_8\|Q^2\|_{\rm GJS}$ at the
needed times follow from $(\UGJS)$ with Remark~6.3, as in Lemma~7.1) gives''.
% BEGIN-REPLACEMENT lemma-7.2
\begin{quote}\small
For $T>|s|$, (3) (extended to $Q$ by Lemma~4.3; the square moments $\int Q(s)^2\,dq_{h_{T,g}}$ at
the needed times are finite at each $T$ by the H\"older bound in the proof of Lemma~4.3(ii), and
Theorem~D.3 with $F=Q^2$ at sharp time $s$ bounds them by
$A\|Q^2\|_{L^2(d\phi_0)}+B'\mathfrak n(Q)^2$ uniformly in $T$, with $A$ and $B'$ as in Lemma~7.1)
gives
\end{quote}
% END-REPLACEMENT
The rest of the proof of Lemma~7.2 is unchanged; its sentence ``Thus (10) is bounded in either case,
by Lemma~7.1'' now cites the replaced Lemma~7.1.  The constant $A$ of Theorem~\ref{thm:D3} at time
$s$ equals that at time $0$: the box $B_s(g,I)$ is $B_0(g,I)$ moved by $l_s$ whole rows, and
$S_2$ is translation invariant (Proposition~\ref{aux:notation-S2}(iii)).

\emph{Lemma 7.3, current lines} (ll.~1500--1503; \cite[p.~21]{R}): ``With
$u_T:=Q\psi_{T-t/2}/\|\psi_T\|$, the extension of (2) to $Q$ (Lemma~4.3; square moments at the two
times by $(\UGJS)$ and Remark~6.3, as in Lemma~7.1) is the exact equality''.
% BEGIN-REPLACEMENT lemma-7.3
\begin{quote}\small
With $u_T:=Q\psi_{T-t/2}/\|\psi_T\|$, the extension of (2) to $Q$ (Lemma~4.3; square moments at
the two times finite at each $T$ by the H\"older bound in the proof of Lemma~4.3(ii), and bounded
uniformly in $T$ by Theorem~D.3 at $s=\mp t/2$, as in Lemma~7.2) is the exact equality
\end{quote}
% END-REPLACEMENT
The rest is unchanged; ``The ratio bound of Lemma~4.4 and Lemma~7.1 make $(u_T)_T$ bounded''
(l.~1507) cites the replaced Lemma~7.1.  \emph{Check.}  Neither replaced text cites $(\UGJS)$ or
$C_8$.  Only finiteness at each $T$ is needed for the extension by Lemma~4.3; the $T$-uniform bound
of Theorem~\ref{thm:D3} at times $s\ne0$ is recorded but not consumed (the record's \S7 consumes
$T$-uniformity only through Lemma~7.1 at time $0$).

\subsection{Lemma 7.6 and the radius in the proof of Theorem 7.7 (Phase A)}\label{sub:cons-76}

\begin{auxlemma}[a Gaussian moment inequality]\label{aux:consequences-moment}
Let $X$ be a centred Gaussian random variable with variance $\sigma^2$.  Then
$\E[X^{2m}]=\sigma^{2m}(2m-1)!!$ for $m\ge0$, and for $n\ge1$
\[
 \|X^{2n}\|_{L^2}=\sigma^{2n}\bigl((4n-1)!!\bigr)^{1/2}\le(\sqrt2\,\sigma)^{2n}\,(2n)! .
\]
\end{auxlemma}

\begin{proof}
$\E[e^{tX}]=e^{\sigma^2t^2/2}$ for real $t$, and $\E[e^{|tX|}]<\infty$, so by dominated
convergence $\E[e^{tX}]=\sum_kt^k\E[X^k]/k!$; comparing the coefficients of $t^{2m}$ with
$e^{\sigma^2t^2/2}=\sum_m\sigma^{2m}t^{2m}/(2^mm!)$ gives
$\E[X^{2m}]=\sigma^{2m}(2m)!/(2^mm!)=\sigma^{2m}(2m-1)!!$.  With $m=2n$,
$\|X^{2n}\|_{L^2}=(\E[X^{4n}])^{1/2}=\sigma^{2n}((4n-1)!!)^{1/2}$.  The inequality is equivalent to
$(4n-1)!!\le4^n((2n)!)^2$, i.e.\ (using $(4n-1)!!=(4n)!/(2^{2n}(2n)!)$) to
$\binom{4n}{2n}\le16^n\,(2n)!$, which holds because $\binom{4n}{2n}\le\sum_j\binom{4n}{j}=16^n$
and $(2n)!\ge1$.  (Tests: $n=1$, $\sqrt3\le4$; $n=2$, $\sqrt{105}\le96$; the integer
inequality $(4n-1)!!\le4^n((2n)!)^2$ is checked for $n\le40$ by \texttt{parts/aqft\_checks.py},
check~4.)
\end{proof}

\emph{Current statement} (ll.~1603--1618; Lemma~7.6 \cite[p.~22]{R}): ``Let
$d\mu_g:=|\Om_g|^2\,d\phi_0$ and $X_f:=\varphi_0(f)$ for real $f\in C_c^\infty(\R)$, $f\ne0$,
supported in an interval of length $L$, and put $K':=K_2^{2}N_I^{1/2}\|f\|_2>0$. Then
$\E_{\mu_g}[X_f^{2n}]\le C_8K_1\,(K')^{2n}(2n)!$ $(n\ge1)$, and consequently
$\E_{\mu_g}e^{s|X_f|}\le2\max\{1,C_8K_1\}\sum_{n=0}^\infty(sK')^{2n}<\infty$ $(0\le s<1/K')$. The
factor $C_8K_1$ affects only the prefactor, not the radius $1/K'$.''  \emph{Current proof}
(ll.~1620--1635) begins ``$\E_{\mu_g}[X_f^{2n}]=\|X_f^n\Om_g\|^2\le C_8\,\mathfrak n(X_f^n)^2$ by
Lemma~7.1 applied to $Q=X_f^n$''.

\emph{Replacement.}
% BEGIN-REPLACEMENT lemma-7.6
\begin{quote}\small
\textsc{Lemma 7.6} (exponential moments).  \emph{Let $d\mu_g:=|\Om_g|^2\,d\phi_0$ and
$X_f:=\varphi_0(f)$ for real $f\in C_c^\infty(\R)$, $f\ne0$, supported in an interval $I$ of
length $L$, and put $K':=K_2^{2}N_I^{1/2}\|f\|_2>0$, $\sigma_f^2:=\ip{f}{Cf}$,
$K_f:=\max\{K',\sqrt2\,\sigma_f\}$ and $A_f:=A+B'K_1$, with $A$ and $B'$ as in Lemma~7.1 for this
$I$ (which has $L>0$, since $f\ne0$).  Then}
\[
 \E_{\mu_g}\bigl[X_f^{2n}\bigr]\le A_f\,K_f^{2n}\,(2n)!\qquad(n\ge1),
\]
\emph{and consequently}
\[
 \E_{\mu_g}e^{s|X_f|}\le2\max\{1,A_f\}\sum_{n=0}^\infty(sK_f)^{2n}<\infty\qquad(0\le s<1/K_f).
\]

\emph{Proof.}  $\E_{\mu_g}[X_f^{2n}]=\|X_f^n\Om_g\|^2\le A\,\|X_f^{2n}\|_{L^2(d\phi_0)}+B'\,
\mathfrak n(X_f^n)^2$ by Lemma~7.1 applied to $Q=X_f^n$ ($r=n$, all $f_i=f$).  By (8),
$\mathfrak n(X_f^n)^2=K_1K_2^{4n}(2n)!\,N_I^n\|f\|_2^{2n}=K_1(K')^{2n}(2n)!$.  Under $d\phi_0$,
$X_f$ is centred Gaussian with variance $\sigma_f^2$, so
$\|X_f^{2n}\|_{L^2(d\phi_0)}=\sigma_f^{2n}((4n-1)!!)^{1/2}\le(\sqrt2\,\sigma_f)^{2n}(2n)!$, because
$(4n-1)!!=(4n)!/(2^{2n}(2n)!)$ and $\binom{4n}{2n}\le16^n$.  Hence
$\E_{\mu_g}[X_f^{2n}]\le(A+B'K_1)K_f^{2n}(2n)!$.  For the exponential bound: $e^{s|x|}\le
e^{sx}+e^{-sx}=2\cosh(sx)$, and by monotone convergence
\[
 \E_{\mu_g}\cosh(sX_f)=\sum_{n\ge0}\frac{s^{2n}}{(2n)!}\E_{\mu_g}[X_f^{2n}]
 \le\max\{1,A_f\}\sum_{n\ge0}(sK_f)^{2n},
\]
where the maximum covers the $n=0$ term $\E_{\mu_g}[1]=1$.  The series converges for $sK_f<1$.
\end{quote}
% END-REPLACEMENT
\emph{Remark on the radius} (outside the statement; trap~25 of the plan).  The radius $1/K_f$
comes from the route through Theorem~\ref{thm:D3}: the component $K'$ from the GJS norm in its
comparison term, the component $\sqrt2\,\sigma_f$ from its base term $A\|X_f^{2n}\|_{L^2(d\phi_0)}$.
Whether $\E_{\mu_g}e^{s|X_f|}<\infty$ for larger $s$ is not decided here.

\emph{Check.}  The step uses the replaced Lemma~7.1, (8) with $N_I$ as in \S\ref{sub:cons-NI}, and
Lemma~\ref{aux:consequences-moment}; it cites neither $(\UGJS)$ nor $C_8$.  $A_f$ and $K_f$ are
independent of $n$, which is what the proof of Theorem~7.7 needs (the plan, Part~1).  $A_f$
depends on $g$ (through $N_g$ and $N_B$) and on $I$ (through $N_B$); $K_f$ depends on $f$
(through $\|f\|_2$, $\sigma_f$ and $N_I$), on $K_2$ and on $m_0$ (through $C$), and not on $g$.
Special case $n=1$: the bound reads $\E_{\mu_g}[X_f^2]\le2A_fK_f^2$.

\emph{Theorem 7.7: the radius.}  The statement of Theorem~7.7 (ll.~1637--1642) is unchanged and
so is the text of its proof; the proof's polystrip (ll.~1652--1654: ``By Cauchy--Schwarz,
H\"older, and Lemma~7.6, $F$ is defined and analytic in a nonempty polystrip
$\{|\mathrm{Re}\,z_j|<s_0\}$ about $i\R^k$'') now has the explicit radius
\[
 s_0:=\min_{1\le j\le k,\ e_j\ne0}\ \frac1{2k\,K_{e_j}},\qquad K_{e_j}:=K_f\ \text{at}\ f=e_j ,
\]
and any $s_0>0$ if all $e_j=0$ (then $F$ is the constant $\E_{\mu_g}[\overline\zeta]$), as in
Remark~\ref{aux:constants-s0}.  Indeed $|\exp(\sum_jz_jX_{e_j})|\le\prod_je^{|\mathrm{Re}\,z_j|
|X_{e_j}|}$, and by Cauchy--Schwarz in $L^2(d\mu_g)$ and then H\"older with $k$ exponents equal
to $k$,
\begin{align*}
 \E_{\mu_g}\Bigl[|\zeta|\prod_je^{|\mathrm{Re}\,z_j||X_{e_j}|}\Bigr]
 &\le\|\zeta\|_{L^2(d\mu_g)}\Bigl(\E_{\mu_g}\prod_je^{2|\mathrm{Re}\,z_j||X_{e_j}|}\Bigr)^{1/2}\\
 &\le\|\zeta\|_{L^2(d\mu_g)}\prod_j\Bigl(\E_{\mu_g}e^{2k|\mathrm{Re}\,z_j||X_{e_j}|}\Bigr)^{1/(2k)},
\end{align*}
which is finite by Lemma~7.6 when $2k|\mathrm{Re}\,z_j|<1/K_{e_j}$ for every $j$ with $e_j\ne0$
(factors with $e_j=0$ equal $1$).  The plan's ``$s_0=1/K_f$'' (Part~1) is the radius of Lemma~7.6
for one $f$; the polystrip for $k$ functions has the radius above.  Before the change the same
computation gave $\min_j(2kK'_{e_j})^{-1}$.  Differentiation under the integral is justified, as in
the record (ll.~1653--1654), by the exponential dominator of Lemma~7.6, quantified as follows: for
$0<s_1<s_2<s_0$, real $x$ and integers $\ell\ge0$, $|x|^\ell e^{s_1|x|}\le\ell!\,(s_2-s_1)^{-\ell}
e^{s_2|x|}$, so on the polystrip $\{|\mathrm{Re}\,z_j|\le s_1\}$ the integrand of every
derivative of order $(\ell_1,\dots,\ell_k)$ is dominated by $\prod_j\ell_j!\,(s_2-s_1)^{-\ell_j}$
times $|\zeta|\prod_je^{s_2|X_{e_j}|}$, which is integrable by the display with $s_2$ in place of
$|\mathrm{Re}\,z_j|$, since $2ks_2<1/K_{e_j}$.

\subsection{Theorem 5.1, Step 5, with $M:=M_B$, and Corollary 5.2 (Phase B)}\label{sub:cons-51}

\emph{Current statement} (ll.~1255--1267; Theorem~5.1 \cite[p.~18]{R}): for every $g\in\CH$ and
every real measurable $h$ with $\|h\|_\infty\le1$ and $\supp h$ contained in an interval of
length $2$, ``$\|V(h)\Om_g\|\le M$, $M^2:=9\,C_8\,\mathfrak K(p')\,\lambda^2(\sum_{n=0}^{2p'}|a_n|)^2$,
$\mathfrak K(p'):=K_1K_2^{8p'}(4p')!$''.  \emph{Current lines of the proof citing $(\UGJS)$ or
$C_8$}: eq.~(5) (l.~1313), ``$\|Q_V\|_{\rm GJS}\le9K_1K_2^{8p'}(4p')!\lambda^2(\sum_{n=0}^{2p'}
|a_n|)^2=M^2/C_8$''; Step~5 (ll.~1319--1327), ``For every $T>0$, apply (2) at $t=0$ with the two
individually nonnegative functions $F_1=F_2=|V(h)|$ (Lemma~4.3(i)); their product is
$V(h)^2=Q_V$. Then $(\UGJS)$ and (5) give $\|V(h)\hat\psi_T\|^2=\int V(h)^2\,dq_{h_{T,g}}\le
C_8\|Q_V\|_{\rm GJS}\le M^2$. (6) In particular the left side is finite, so
$\psi_T\in D(V(h))$.''

\emph{Replacement.}
% BEGIN-REPLACEMENT theorem-5.1
\begin{quote}\small
\emph{In the statement:} $\|V(h)\Om_g\|\le M$, with
\begin{gather*}
 M^2:=M_B^2=e^{15\lambda|a_0|}\,S_2(B_{3\times5})\,9^{p'}\,V_*
 +9\,C_7\,c_U\,\overline\Sigma_{3\times5}\,\mathfrak K(p')\,\lambda^2\Bigl(\sum_{n=0}^{2p'}|a_n|\Bigr)^2,\\
 \mathfrak K(p'):=K_1K_2^{8p'}(4p')! ,
\end{gather*}
the constants $S_2(B_{3\times5})$, $V_*$, $c_U$ and $\overline\Sigma_{3\times5}$ being those of
Theorem~D.4.

\emph{Step 4, eq.~(5):} $\|Q_V\|_{\rm GJS}\le9\,K_1K_2^{8p'}(4p')!\,\lambda^2
(\sum_{n=0}^{2p'}|a_n|)^2$.

\emph{Step 5: transfer at time zero.}  For every $T>0$, apply (2) at $t=0$ with the two
individually nonnegative functions $F_1=F_2=|V(h)|$ (Lemma~4.3(i)); their product is
$V(h)^2=Q_V$.  Choose $a\in\R$ with $\supp h\subseteq[a,a+2]$.  Theorem~D.4, applied with the
Euclidean cutoff $h_{T,g}\in\CE$ (Lemma~3.1(ii)) and with $h':=h$, gives
\[
 \|V(h)\hat\psi_T\|^2=\int V(h)^2\,dq_{h_{T,g}}\le M_B^2=M^2 .
 \tag{6}
\]
In particular the left side is finite, so $\psi_T\in D(V(h))$.  (Theorem~D.4 is a bound for the
$V(h)$ constructed in Appendix~D, \S2, which is identified with the $V(h)$ of \S2.1 only through
the gap recorded there, as in Step~1; and it uses GJS Theorem~1.1.7 also across strips
transverse to the time-zero line, a case its printed proof does not reach, F-T.)
\end{quote}
% END-REPLACEMENT
Steps~1--4 and~6 are kept (Steps~2--4 with the columns of \S\ref{sub:cons-NI}); in Step~4 the
final equality of eq.~(5), which reads ``$=M^2/C_8$'', is deleted (item n14 of the plan).
\emph{Check.}  Theorem~\ref{thm:D4} applies: $h_{T,g}\in\CE$ by Lemma~3.1(ii) of the record; $h$
is real measurable with $\|h\|_\infty\le1$ and support in an interval of length~$2$, hence in a
closed interval $[a,a+2]$; the bound is uniform in $T>0$ and $g\in\CH$ because $M_B$ does not
depend on the cutoff.  Theorem~\ref{thm:D4} bounds $\int Y^2\,dq_{h_{T,g}}$ for the constructed
$Y=V(h)(0)=J_0V(h)$ of Lemma~\ref{aux:notation-hyper}, and by Lemma~\ref{aux:notation-sharpid}
with Remark~\ref{aux:notation-sharpidrem}(b) this $Y^2$ is also the sum of the localized pieces to
which its comparison term applies; the identity $\|V(h)\hat\psi_T\|^2=\int V(h)^2\,dq_{h_{T,g}}$
of (2) concerns the record's $V(h)$, and the two are identified by the GAP recorded after
Lemma~\ref{aux:notation-wick}, which Step~1 of the record's proof uses in the same way (its
sentence ``$Q_V$ is exactly the multiplication function $V(h)^2$'').  Step~6 (the weak-graph
passage, ll.~1329--1340) uses only the $T$-uniform bound (6) and Lemma~4.4, so it gives
$\|V(h)\Om_g\|\le M_B$.  The replaced statement and proof cite neither $(\UGJS)$ nor $C_8$.
Steps~2--4, which bound $\GJS{Q_V}$, are no longer consumed by Step~5 (Theorem~\ref{thm:D4}
bounds $\GJS{Q_V}$ on GJS's columns itself, Lemma~\ref{aux:d4-QV}(iii)); they are kept, on the
columns, with the same constant.  Through Theorem~\ref{thm:D4}, Step~5 uses GJS Theorem~1.1.7 in
its transverse case (F-T, \S\ref{sec:purpose}; \S\ref{sub:d4-transverse}).

\emph{Corollary 5.2, current statement} (ll.~1342--1347; \cite[p.~19]{R}): ``For
$\mathscr P(\xi)=\xi^4$ ($p'=2$, $\sum_n|a_n|=1$; degrees exactly $4+4=8$),
$M^2=9\cdot8!\;C_8K_1K_2^{16}\lambda^2=362\,880\;C_8K_1K_2^{16}\lambda^2$.''
% BEGIN-REPLACEMENT corollary-5.2
\begin{quote}\small
\textsc{Corollary 5.2} ($\varphi^4$ instance).  \emph{For $\mathscr P(\xi)=\xi^4$ ($p'=2$, $a_4=1$,
$a_0=\dots=a_3=0$, $\sum_n|a_n|=1$; degrees exactly $4+4=8$),}
\begin{gather*}
 M^2=81\,S_2(B_{3\times5})\,V_*+8\,709\,120\;C_7K_1^2K_2^{24}\,\overline\Sigma_{3\times5}\,\lambda^3,\\
 V_*\le24\,\lambda^2\int_{-1}^1\!\int_{-1}^1C(x-y)^4\,dx\,dy .
\end{gather*}
(The value is that of Theorem~5.1, through Theorem~D.4: it rests on the transverse case of GJS
Theorem~1.1.7, F-T, and on the gap note of Appendix~D, \S2.)
\end{quote}
% END-REPLACEMENT
\emph{Derivation.}  Here $e^{15\lambda|a_0|}=1$ and $9^{p'}=81$.  By
Lemma~\ref{aux:notation-cU},
\begin{gather*}
 c_U=\lambda|a_4|K_1K_2^{8}\,4!=24\lambda K_1K_2^8 ,\\
 9\,\mathfrak K(2)\lambda^2\Bigl(\sum_n|a_n|\Bigr)^2=9\cdot8!\,K_1K_2^{16}\lambda^2
 =362\,880\,K_1K_2^{16}\lambda^2 ,
\end{gather*}
and the product
\[
 C_7\cdot24\lambda K_1K_2^8\cdot\overline\Sigma_{3\times5}\cdot362\,880\,K_1K_2^{16}\lambda^2
 =8\,709\,120\,C_7K_1^2K_2^{24}\,\overline\Sigma_{3\times5}\,\lambda^3
\]
($24\cdot362\,880=8\,709\,120$); and $V_*\le\overline V_*=\lambda^2a_4^2\,4!\,I_4=24\lambda^2I_4$ by Lemma~\ref{aux:notation-wick}(iv),
with $I_4=\iint_{[-1,1]^2}C(x-y)^4\,dx\,dy$ of Lemma~\ref{aux:notation-kernel}(iv).  The number
$362\,880\,C_8K_1K_2^{16}\lambda^2$ retires (trap~10 of the plan).  Alternatively the corollary is
retired with its number kept (the plan, Part~3, Phase~B).

\subsection{Theorem 12.2, the other sections, and Appendix A}\label{sub:cons-rest}

\emph{Theorem 12.2.}  Its proof (ll.~2427--2446) is unchanged: ``(5) is Theorem~7.8 and
Theorem~5.1. These supply (G$_\gamma$) with $\gamma=m_1$ and (M$_2$), so Theorem~11.1 and
Proposition~11.2 give (1)--(2)'' holds verbatim with the replaced Theorem~5.1.  Its statement
keeps its text except the imports paragraph (ll.~2415--2424), whose rewrite at integration is
listed in the plan (Part~6) and is not drafted here; the value of $M$ in items (1) (through
$\Psi_{m_1}$, Lemma~10.5) and (5) is $M_B$ after Phase~B (trap~10).  The plan's Part~6 predates
the V3 audit and the GAP; the rewritten paragraph must name, besides GJS Theorem~1.1.7 at its
printed scope including the transverse case (F-T), the record's Lemma~C.2, Lemma~C.4 and
Theorem~C.5 under the V3 repairs (R8 importing the construction of $\Uc_h$), and the construction
of the time-zero Wick powers that closes the GAP recorded after Lemma~\ref{aux:notation-wick}.
\emph{Disclosure.}  After Phase~B, Theorem~5.1, and through $M=M_B$ the statements of Lemma~10.2,
Lemma~10.3, Proposition~10.4, Lemma~10.5, Theorem~11.1 and Theorem~12.2 (items (1) and (5)),
use the transverse case of GJS Theorem~1.1.7 (F-T, \S\ref{sec:purpose}) and the GAP; after
Phase~A alone, Lemmas~7.1--7.3 and~7.6 use the parallel case only.  Proposition~7.5 and
Theorem~7.8 keep statements and proofs ($\widehat C_Q$ contains $\GJS{Q}$ and $\|\psi_Q\|$ only,
the latter finite by the replaced Lemma~7.1).  Lemma~10.2 (proof, l.~2110) and \S11 (l.~2254)
cite Theorem~5.1's $M$ by name: text unchanged, value $M_B$.

\emph{Appendix A of the record} (ll.~2655--2685).  Rows that change: $K_1,K_2,C_7,C_8,m_1$ (l.~2668:
$C_8$ dropped once A8(vi) is retired after Phase~B); $N_I$ (l.~2670: defined by GJS's columns,
\S\ref{sub:cons-NI}, bound unchanged); $M^2$ (l.~2674: $M_B^2$); $M^2$ ($\varphi^4$) (l.~2675:
Corollary~5.2 as replaced, or retired); $K'$ (l.~2676: same formula, $N_I$ by columns); the
$\Psi_\gamma$ rows (ll.~2681--2682: same formulas; the $M$ in them is $M_B$).  Unchanged:
$\varepsilon_{\mathscr P}$, $\mathfrak K(p')$, $\widehat C_Q$ (its formula; $\GJS Q$ is computed on
GJS's squares), $\mathfrak m$, $d_0$, the shell count.  New rows: $c_U$, $S_2(B)$, $A$, $B'$,
$A_f$, $K_f$ (with $\sigma_f$), $s_0$ of the proof of Theorem~7.7, $V_*$,
$\overline\Sigma_{3\times5}$ and $M_B$; $S_2(B)$, and through it $A$, $A_f$ and $M_B$, contain the
finite constant of Theorem~C.5, which is not evaluated, and Appendix~A must say so (gate GB3).
This list agrees with the closing paragraph of Appendix~A.4 of this note.

\subsection{Gates GA5 and GA6}\label{sub:cons-gates}

\emph{GA5 (no proof in \S\S4, 6, 7 cites $(\UGJS)$).}  Command, in the repository (run on
2 and 3 November 2026):
\begin{quote}\footnotesize\begin{verbatim}
grep -n -e 'eq:UGJS' -e 'C_8' part_i/main.tex | awk -F: '$1>=1133 && $1<1722'
\end{verbatim}\end{quote}
Output (line numbers only): 1210, 1216 (Lemma~4.3(ii)); 1262, 1313, 1321, 1324 (Theorem~5.1);
1346 (Corollary~5.2); 1430, 1431, 1436, 1439, 1440, 1449 (Lemma~7.1); 1463, 1464 (Lemma~7.2); 1502
(Lemma~7.3); 1609, 1614, 1617, 1621, 1631 (Lemma~7.6).  Every hit lies in a text replaced above.
The replacement texts themselves are the eight blocks between the markers of this file
(\texttt{sec6-NI}, \texttt{lemma-4.3-ii}, \texttt{lemma-7.1}, \texttt{lemma-7.2},
\texttt{lemma-7.3}, \texttt{lemma-7.6}, \texttt{theorem-5.1}, \texttt{corollary-5.2}); the command
\begin{quote}\footnotesize\begin{verbatim}
awk '/% BEGIN-REPLACEMENT/,/% END-REPLACEMENT/' parts/consequences.tex \
  | grep -n -e 'UGJS' -e 'C_8' -e 'eq:UGJS'
\end{verbatim}\end{quote}
prints nothing (run on this file, 3 November 2026).  So after the replacement no proof of
\S\S4--7 cites $(\UGJS)$ or $C_8$.  Outside \S\S4--7 the record still cites $(\UGJS)$; the
command (3 November 2026)
\begin{quote}\footnotesize\begin{verbatim}
grep -n -e 'eq:UGJS' -e 'U_{\rm GJS}' -e 'UGJS' part_i/main.tex
\end{verbatim}\end{quote}
finds it, outside \S\S4--7, in the status notes after the title (ll.~68, 72, 76, 81), in \S1.2 (ll.~184, 210), in \S1.4 (ll.~290, 302), in
Assumption~A8 (ll.~710--803, 868), in Theorem~12.2's imports paragraph (l.~2417), in Appendices~B
and~C (ll.~2737, 2776) and in Proposition~C.7 (l.~3075); this list is complete for that grep.
Their edits are listed in the plan (Part~6) and are not part of this section.

\emph{GA6 (the new constants occur in no statement of \S\S8--14).}  Command:
\begin{quote}\footnotesize\begin{verbatim}
awk 'NR>=1722 && NR<=2654 {print NR": "$0}' part_i/main.tex \
 | grep -E -e 'C_8|UGJS|U_\{\\rm GJS\}|mathfrak n' \
     -e "K'|widehat C_Q|s_0" \
     -e 'lem:transfer|lem:onepoint|lem:twopoint|lem:expmom' \
     -e 'prop:semidecay|thm:dense|eq:M2|thm:M2|thm:G[^a-z]|eq:G[^a-z]'
\end{verbatim}\end{quote}
(run in the repository on 2 and 3 November 2026).  Output, line numbers and what they are: 1731
(\S8 text, ``by Theorem~7.8 this holds with $\gamma=m_1$''); 1859--1860 ($s_0$ here is the path
parameter in the proof of Proposition~8.2(ii), not the radius of Theorem~7.7); 2110 (proof of
Lemma~10.2, ``Theorem~5.1 applies to the cutoff $g_s$ and gives $\|V_j\Om_s\|\le M$''); 2252, 2254
(\S11 text, $\gamma=m_1$ by Theorem~7.8 and $\Psi_{m_1}$ with $M$ the constant of Theorem~5.1); 2394
(a source comment); 2417 (Theorem~12.2's imports paragraph, $(\UGJS)$); 2429--2430 (proof of
Theorem~12.2); 2486, 2530 (\S13 text and proof, Theorem~7.8).  Hence $C_8$, $\mathfrak n$, $K'$,
$\widehat C_Q$ and the labels of Lemmas~7.1--7.3, 7.6, Proposition~7.5 and Theorem~7.7 do not occur
in \S\S8--14 at all, and the new constants $A$, $B'$, $A_f$, $K_f$ and the radius $s_0$ of
Theorem~7.7 cannot occur there.  The only constants of \S\S5 and~7 that reach \S\S8--14 are $m_1$
(through Theorem~7.8, unchanged) and $M$ (through Theorem~5.1).  The statements that contain $M$
are those of Lemma~10.2 (eq.~(25), l.~2099), Lemma~10.3 (eq.~(26), l.~2126), Proposition~10.4
(eq.~(27), l.~2155), Lemma~10.5 (eq.~(29), l.~2214, and its constant $C_{N,\gamma,M}$), and,
through $\Psi_{m_1}$, Theorem~11.1 and Theorem~12.2, items (1) and (5); the grep matches labels and
constants, not the letter $M$, so this list was made by reading ll.~2088--2232 of the record.  The
value of $M$ becomes $M_B$ after Phase~B; that change is the plan's trap~10, not a GA6 failure.
After Phase~A alone, $M$ is unchanged.

\subsection{Findings about the record and the plan made while drafting}\label{sub:cons-findings}

\begin{itemize}
\item \emph{The count $N_I$} (main.tex l.~1359, the notation table l.~448, and A8(ii),
ll.~659--663): see \S\ref{sub:cons-NI} and Remark~\ref{aux:d3-cells}.  On GJS's squares (7)--(8)
hold with the column count; with the record's cell count, (7) fails for $\Phi(s,f)$ and (8) for
the column decomposition of $\Phi(s,f)^2$ when $I=(0.4,0.6)$ ($\sqrt2K_1K_2^2\|f\|_2$ against
$K_1K_2^2\|f\|_2$; $3K_1K_2^4$ against $2K_1K_2^4$).  Severity: CRITICAL for the display of
Proposition~6.2 as stated (on the printed norm of A8(iii)); the displayed bounds of Lemmas~7.1
and~7.6, whose proofs consume (8), fail as stated; MINOR downstream (Theorems~7.7, 7.8 and~12.2
use finiteness only).  The edit of \S\ref{sub:cons-NI} repairs all of these.
\item \emph{Statements that change}: the statements of Lemma~6.1 and Proposition~6.2 (through
$N_I$), of Lemma~7.1 (its form) and Lemma~7.6, Theorem~5.1 and Corollary~5.2 (their constants);
Theorem~7.7's statement has no radius, which lives in its proof.  Severity: MINOR (wording of the
note, corrected in the introduction of this section).
\item \emph{The GAP at Theorem~5.1, Step~1} (ll.~1278--1283): ``$Q_V$ is exactly the multiplication
function $V(h)^2$'' identifies the record's $V(h)$, of which the record states no construction,
with an object expanded into Wick monomials; this is the GAP recorded after
Lemma~\ref{aux:notation-wick}, used by the kept Step~1 and by the replaced Step~5.
\item \emph{Wick re-expansion} (drafting brief, d4 item): no Wick re-expansion of $V(h')(0)^2$ is
needed or made; the record's own Step~1 (ll.~1278--1283) treats the product as a (1.1.5) monomial,
and Lemma~\ref{aux:d4-QV} follows it.
\item \emph{Transverse uses} (plan Part~1, ``Where the transverse case enters''): besides the
squares of rows $-1,0,1$, the squares of class (T) in rows $|l|\ge2$ use the transverse strip in
Theorem~\ref{thm:D4}'s bound (\S\ref{sub:d4-transverse}).
\item \emph{The radius $s_0$} (plan Part~1, ``Theorem 7.7 then has $s_0=1/K_f$''): the polystrip
of Theorem~7.7 has radius $\min_{j:\,e_j\ne0}(2kK_{e_j})^{-1}$ (\S\ref{sub:cons-76}).
\item \emph{The form of Lemma~D.2} (plan Part~1): a form of the plan's statement, with the widest
strips $d(\Delta,R)$, positive widths and the sum over squares meeting $\supp h$, is
Lemma~\ref{lem:D2}(b); the theorems use form (a), with explicitly assigned strips (so that each
application has a named orientation) and the sum over the squares on which $h$ is not almost
everywhere zero (Remark~\ref{aux:d2-uses}).
\item \emph{$\Sigma_{3\times5}$} (plan Part~1, and Lemma~\ref{aux:notation-rings}(vi) of
\S\ref{sec:notation}, define it as $\sum_{\Delta\not\subseteq B_{3\times5}}e^{-m_1d(\Delta,R)}$,
which depends on $R$): Theorem~\ref{thm:D4} is stated with the constant
$\overline\Sigma_{3\times5}:=\sum_{k\ge1}(8k+12)e^{-m_1k}=4x(5-3x)(1-x)^{-2}$, the right side of
that lemma, because its proof uses the axis-parallel strips of named orientation
(\S\ref{sub:d4-setting}); $M_B$ therefore does not depend on $h'$, and the plan's $M_B$ with the
widest strips would need tilted strips also for squares of class (P).
\end{itemize}


\appendix
\section{Constants recomputed by hand (gate GA3)}\label{app:constants}
% constants.tex -- drafted by cqft2 (council of ugjs/PLAN.md), 2 November 2026.
% Recomputed by hand from ugjs/PLAN.md Part 1, the record and Section 2 of this note,
% independently of the drafts of Sections 4-6 and 8 (not on disk when this was written;
% the comparison with them, made afterwards, is subsection A.6).

This appendix recomputes, line by line, every constant of Theorems~\ref{thm:D3} and
\ref{thm:D4} and of the consequences of \S\ref{sec:cons}: $c_U$, $S_2(B)$, $A$, $B'$, $A_f$,
$K_f$, $s_0$, $\overline\Sigma_{3\times5}$, $V_*$ and $M_B$, together with the auxiliary quantities they
are built from.  The statements of \texttt{ugjs/PLAN.md}, Part~1, are the targets.  The inputs
are the printed constants of \S\ref{sec:notation-printed}; Lemma~\ref{lem:D1}; Source Fact~2.2 of
the record (Assumption~A8(iv) \cite[p.~10]{R}); Lemma~C.2 (the base terms), Lemma~C.4 (through
Lemmas~\ref{aux:notation-cutoffs}(iv), \ref{aux:notation-cU}(iv) and~\ref{aux:notation-transl},
i.e.\ the decomposition of $\Uc_{h-h^B}$ and the shape invariance of $S_2$), Theorem~C.5 (finiteness
of $S_2$) and the comparison step in the proof of Theorem~C.5 (the factor $9^{p'}$), with the V3
dependences of \S\ref{sec:notation}; Remark~2.5, Assumptions~A2, A3 and A4 and Lemma~3.1(ii) of the
record; and the lemmas of \S\ref{sec:notation}.  The recomputation was made independently of the text of
\S\S\ref{sec:d2}--\ref{sec:d4} and \S\ref{sec:cons}; \S A.5 lists every discrepancy found
with the plan, and \S A.6 those found afterwards with those sections.  Every statement here is \emph{conjectured} until
audited.

\subsection*{A.1\quad The constant of the comparison}

\begin{auxlemma}[the comparison constant is $C_7c_U$]\label{aux:constants-comparison}
Let $h\in\CE$, let $B$ be a box, and let $F=\sum_{a}F_a$ be a finite sum of localized
monomials of type (1.1.5) or (1.1.6) with $L^2$ kernels and factor degrees $<\bar n$, all of
whose localization squares lie in a set $R\subseteq B$, with $\GJS F=\sum_a\GJS{F_a}$ by
(1.1.9).  Assume that $F$ satisfies the hypotheses of Lemma~\ref{lem:D1} with $h_1=h^B$,
$h_2=h$, and that to every $\Delta\in\mathfrak S(h)$ with $\Delta\cap B=\varnothing$ a strip
$S_\Delta$ of width $w_\Delta>0$ separating $\Delta$ from $\bigcup R$ is assigned.  Then
\[
 \bigl|\ev Fh-\ev F{h^B}\bigr|\le C_7\,c_U\,\GJS F\sum_{\Delta\in\mathfrak S(h),\ \Delta\cap
 B=\varnothing}e^{-m_1w_\Delta} .
\]
Every estimate behind this bound is Source Fact~2.2 for a pair $(F_a,U_{\Delta,n})$ across
$S_\Delta$, so the orientation of $S_\Delta$ is the orientation of that application.  Since
$w_\Delta$ may be taken arbitrarily close to $d(\Delta,R)$, the bound also holds with
$d(\Delta,R)$ in place of $w_\Delta$; the orientations are then not controlled
(Remark~\ref{aux:notation-tilted}): where a strip so chosen is transverse to the time-zero line
and some $F_a$ is of type (1.1.5), this form uses the transverse case of GJS Theorem~1.1.7 at the
printed statement's scope (F-T, \S\ref{sec:purpose}).
\end{auxlemma}

\begin{proof}
By Lemma~\ref{lem:D1} and Lemma~\ref{aux:notation-cutoffs}(ii),
$\ev Fh-\ev F{h^B}=-\int_0^1\Cov_{h_\alpha}(F,\Uc_{h-h^B})\,d\alpha$ with $h_\alpha\in\CE$; by
Lemma~\ref{aux:notation-cU}(iv) the covariance is
$\sum_{\Delta}\sum_{n=1}^{2p'}\sum_a\Cov_{h_\alpha}(F_a,U_{\Delta,n})$, the sum over
$\Delta\in\mathfrak S(h)$ with $\Delta\cap B=\varnothing$.  Each pair $(F_a,U_{\Delta,n})$
consists of two localized monomials (types (1.1.5)/(1.1.6) and (1.1.6),
Lemma~\ref{aux:notation-cU}(i)) with disjoint localization squares (those of $F_a$ lie in
$R\subseteq B$, and $\Delta\cap B=\varnothing$), factor degrees $<\bar n$, and the strip
$S_\Delta$ separates $\Delta$ from $\bigcup R$, hence from the squares of $F_a$.  Source
Fact~2.2, at the cutoff $h_\alpha\in\CE$, in the weakened form of \S\ref{sec:notation-strips}
with the strip $S_\Delta$, gives
$|\Cov_{h_\alpha}(F_a,U_{\Delta,n})|\le C_7e^{-m_1w_\Delta}\GJS{F_a}\GJS{U_{\Delta,n}}$.  Sum:
$\sum_a\GJS{F_a}=\GJS F$, $\sum_n\GJS{U_{\Delta,n}}\le c_U$ (Lemma~\ref{aux:notation-cU}(ii)),
and $\int_0^1d\alpha=1$.  For the last sentence: given $\varepsilon>0$ choose strips with
$w_\Delta>d(\Delta,R)-\varepsilon$; the sum is finite, so letting $\varepsilon\to0$ gives the
bound with $d(\Delta,R)$.
\end{proof}

The orientation of each application is that of the assigned strip;
Propositions~\ref{aux:constants-D3} and~\ref{aux:constants-D4} assign the strips explicitly.  The
constant of the comparison is $C_7c_U$ and nothing else; the geometry enters only through the
sum.

\subsection*{A.2\quad The constants of Theorems~\ref{thm:D3} and~\ref{thm:D4}}

\begin{auxprop}[the constants of Theorem~\ref{thm:D3}]\label{aux:constants-D3}
Let $g\in\CH$, $I\subset\R$ a bounded interval, $s\in\R$, $N_g$ the number of columns meeting
$\supp g$ in positive measure, and $B=B_s(g,I)$ the rows $l_s-1,l_s,l_s+1$ times the smallest
interval of $N_B$ columns containing every column that meets $\supp g$ or $I$ (in positive
measure, or set-theoretically: both readings work below).  Let $F$ be a finite linear
combination of cylinder monomials at sharp time $s$, supported in $I$, satisfying the
hypotheses of Lemma~\ref{lem:D1}.  Then for every $T>0$
\[
\begin{gathered}
 \Bigl|\int F\,dq_{h_{T,g}}\Bigr|\le A\,\|F\|_{L^2(\dphi)}+B'\,\GJS F,\\
 A:=e^{3\lambda|a_0|N_B}S_2(B),\qquad B':=\frac{2\,C_7\,c_U\,N_g\,e^{-m_1}}{1-e^{-m_1}} .
\end{gathered}
\]
$A$ depends on $(\lambda,\mathscr P,m_0)$ and $N_B$ and is uniform in $T>0$ and $s\in\R$;
$B'$ depends on $C_7,m_1,K_1,K_2,\lambda,a_1,\dots,a_{2p'}$ and $N_g$, and is uniform in $T$,
$s$ and the position of $I$.  All strips used are parallel.
\end{auxprop}

\begin{proof}
\emph{Base term.}  $h:=h_{T,g}\in\CE$ (Lemma~3.1(ii) \cite[p.~15]{R}); $h^B\in\CE$ and $\supp
h^B\subseteq\overline B$ (Lemma~\ref{aux:notation-cutoffs}(i)).  Cauchy--Schwarz in
$L^2(\dmu)$:
$|\ev F{h^B}|\le Z(h^B)^{-1}\|F\|_{L^2(\dmu)}\|e^{-\Uc_{h^B}}\|_{L^2(\dmu)}$.  Lemma~C.2
\cite[p.~40]{R}, ``For every $h\in\CE$ with $\supp h\subseteq\Lambda$:
$Z(h)\ge e^{-\lambda|a_0|\,|\Lambda|}>0$'', applied with $\Lambda=\overline B$ (compact,
$|\overline B|=|B|=3N_B$; hypotheses checked above) gives $Z(h^B)^{-1}\le e^{3\lambda|a_0|N_B}$
(Lemma~C.2 is used conditional on the repair R8 of the V3 audit, which the V3 report, \S2, finds
sufficient for it; \S\ref{sec:notation});
Proposition~\ref{aux:notation-S2}(iv) gives $\|e^{-\Uc_{h^B}}\|_2\le S_2(B)$; and
$\|F\|_{L^2(\dmu)}=\|F\|_{L^2(\dphi)}$ because the sharp-time field has the law of
$\varphi_0$ (Remark~2.5 \cite[p.~14]{R}, as in the proof of
Lemma~\ref{aux:notation-hyper}).  So $|\ev F{h^B}|\le A\|F\|_{L^2(\dphi)}$.  Uniformity: $A$
contains neither $T$ nor $s$; $B_s(g,I)$ has shape $(3,N_B)$ for every $s$, so $S_2(B)$ is
independent of $s$ (Proposition~\ref{aux:notation-S2}(iii)).

\emph{Comparison term.}  The localization squares of $F$ (Remark~\ref{aux:notation-gjslattice}(ii),
pieces $f_i\mathbf 1_{[k-1/2,k+1/2)}$) lie in $R\subseteq S:=\{\Delta_{l_s,k}:k\text{ a column
of }B\}$, a rectangle in $B$.  If $\Delta_{l,k}\in\mathfrak S(h_{T,g})$, then
$\int_{\Delta_{l,k}}h_{T,g}=|[-T,T]\cap[l-\frac12,l+\frac12)|\int_{k-1/2}^{k+1/2}g>0$, so
column $k$ meets $\{g>0\}\subseteq\supp g$ in positive measure: $k$ is one of the $N_g$
columns, all of which are columns of $B$.  If moreover $\Delta_{l,k}\cap B=\varnothing$, then
$|l-l_s|\ge2$, so $\rho_t=|l-l_s|\ge2$ and $\rho_x=0$ with respect to $S$, and
Lemma~\ref{aux:notation-strip}(i) gives the \emph{parallel} strip $S_\Delta$ bounded by the
lines $x^0=l_s+\frac12$ and $x^0=l-\frac12$ (if $l>l_s$), resp.\ $x^0=l+\frac12$ and
$x^0=l_s-\frac12$ (if $l<l_s$), of width $w_\Delta=|l-l_s|-1\ge1$, separating $\Delta$ from
$\bigcup S\supseteq\bigcup R$.  Hence, with $x=e^{-m_1}$ and $\sum_{j\ge1}x^j=x/(1-x)$ (proof of
Lemma~\ref{aux:notation-rings}(vi)),
\[
 \sum_{\Delta\in\mathfrak S(h_{T,g}),\ \Delta\cap B=\varnothing}e^{-m_1w_\Delta}
 \le N_g\sum_{|l-l_s|\ge2}e^{-m_1(|l-l_s|-1)}=2N_g\sum_{j\ge1}x^j=\frac{2N_ge^{-m_1}}{1-e^{-m_1}} .
\]
Lemma~\ref{aux:constants-comparison} with these strips gives
$|\ev F{h_{T,g}}-\ev F{h^B}|\le B'\GJS F$.  The
bound contains neither $T$, nor $s$, nor $I$.  For $T<\frac12$ only row $0$ is charged; the
bound is the same.
\end{proof}

\begin{auxprop}[the constants of Theorem~\ref{thm:D4}]\label{aux:constants-D4}
Let $h\in\CE$, $a\in\R$, and $h'$ real measurable with $\|h'\|_\infty\le1$ and $\supp
h'\subseteq[a,a+2]$, and assume that $Q_V:=V(h')(0)^2$ satisfies the hypotheses of
Lemma~\ref{lem:D1}.  Then $\int V(h')(0)^2\,dq_h\le M_B^2$, where
\[
 M_B^2:=e^{15\lambda|a_0|}\,S_2(B_{3\times5})\,9^{p'}V_*
 +C_7\,c_U\,\overline\Sigma_{3\times5}\,\mathfrak Q,
 \qquad
 \mathfrak Q:=9K_1K_2^{8p'}(4p')!\,\lambda^2\Bigl(\sum_{n=0}^{2p'}|a_n|\Bigr)^2 ,
\]
where $S_2(B_{3\times5})$ is the common value for all $a$ (Proposition~\ref{aux:notation-S2}(iii))
and $\overline\Sigma_{3\times5}=\sum_{j\ge1}(8j+12)e^{-m_1j}=4e^{-m_1}(5-3e^{-m_1})
(1-e^{-m_1})^{-2}$ is the ring constant of Lemma~\ref{aux:notation-rings}(vi).  Fully explicit up
to $S_2$:
\[
 M_B^2\le\overline M_B^2:=e^{15\lambda|a_0|}\,S_2(B_{3\times5})\,9^{p'}\,\overline V_*
 +C_7\,c_U\,\overline\Sigma_{3\times5}\,\mathfrak Q ,
\]
$\overline V_*=\lambda^2\sum_{n=0}^{2p'}a_n^2n!\,I_n$ (Lemma~\ref{aux:notation-wick}).  $M_B$
depends on $(\lambda,\mathscr P,m_0)$ through $C_7,m_1,K_1,K_2$, $V_*$ and the unevaluated
$S_2(B_{3\times5})$, and is uniform in $h\in\CE$ (hence in $T>0$ and $g\in\CH$), in $a$ and in
$h'$.  The comparison term uses the strips of Lemma~\ref{aux:notation-rings}(iv): parallel ones
for the squares with $|l|=1+j$ (the corner squares included) and transverse ones for the squares
with $|l|\le j$ (Remark~\ref{aux:notation-ringrem}(b)).
\end{auxprop}

\begin{proof}
\emph{The norm $\mathfrak Q$.}  By eq.~(4) of the record (Theorem~5.1, Step~1
\cite[p.~18]{R}), $Q_V=\lambda^2\sum_{n,n'=0}^{2p'}a_na_{n'}\iint h'(x)h'(y)
{:}\Phi(0,x)^n{:}\,{:}\Phi(0,y)^{n'}{:}\,dx\,dy$.  Split $h'=\sum_kh'\mathbf
1_{[k-1/2,k+1/2)}$ over the columns meeting $[a,a+2]$ in positive measure, at most three, all
among $k_a,k_a+1,k_a+2$ (Lemma~\ref{aux:notation-rings}(i)).  The piece $(n,n',k,k')$ is a
localized monomial of type (1.1.5) at time $0\in[-\frac12,\frac12)$, i.e.\ in row $0$, with
kernel $w=\lambda^2a_na_{n'}h'\mathbf 1_{k}\otimes h'\mathbf 1_{k'}$,
$\|w\|_2\le\lambda^2|a_n||a_{n'}|$ (unit columns, $|h'|\le1$), and factor degrees
$n,n'\le2p'<\bar n$.  Its weight in (1.1.8) is $K_1K_2^{2(n+n')}(n+n')!$ if $k=k'$ and
$K_1K_2^{2n}n!\,K_2^{2n'}n'!$ if $k\ne k'$; since $n+n'\le4p'$, $K_2\ge1$ and
$n!\,n'!\le(n+n')!$, both are $\le K_1K_2^{8p'}(4p')!$.  By (1.1.9), summing over at most $9$
pairs $(k,k')$ and over $(n,n')$: $\GJS{Q_V}\le\mathfrak Q$.  (This is the record's eq.~(5)
\cite[p.~19]{R}, recomputed on the squares of Definition~\ref{aux:notation-squares}; the
record's count of three cells holds for the columns by Lemma~\ref{aux:notation-rings}(i).)

\emph{Base term.}  $B:=B_{3\times5}(a)$ is a box of $15$ squares containing $R$, the
localization squares of $Q_V$, which lie in $S_a$.  As in Proposition~\ref{aux:constants-D3}:
$|\ev{Q_V}{h^B}|\le e^{15\lambda|a_0|}S_2(B)\|Q_V\|_{L^2(\dmu)}$, and
$\|Q_V\|_{L^2(\dmu)}=\|Y^2\|_{L^2(\dmu)}\le9^{p'}V_*$ by Lemma~\ref{aux:notation-hyper}.  The
$Q_V$ of the base term and the $Q_V$ whose pieces enter the comparison term below are the same
element of $L^2(\dmu)$ by Lemma~\ref{aux:notation-sharpid} and
Remark~\ref{aux:notation-sharpidrem}(b).

\emph{Comparison term.}  Lemma~\ref{aux:constants-comparison} with $F=Q_V$, $R=S_a$, and, for
$\Delta=\Delta_{l,k}\in\mathcal L_j$, the strip $S_\Delta$ of Lemma~\ref{aux:notation-rings}(iv) of
width $w_\Delta=j$: parallel if $|l|=1+j$ (the corner squares included), transverse if
$|l|\le j$.  Every application to a square with $|l|\le j$ is the transverse case of GJS
Theorem~1.1.7, used at the printed statement's scope (sentence F-T, \S\ref{sec:purpose}).  The
terms are positive and the charged squares outside $B$ are among all squares $\Delta\not\subseteq
B$, so by Lemma~\ref{aux:notation-rings}(ii),(vi) the sum is at most
$\sum_{\Delta\not\subseteq B}e^{-m_1w_\Delta}=\overline\Sigma_{3\times5}$, and
$|\ev{Q_V}h-\ev{Q_V}{h^B}|\le C_7c_U\,\mathfrak Q\,\overline\Sigma_{3\times5}$.  Add the two
terms; $\ev{Q_V}h\ge0$.  Uniformity: no
term contains $h$, $T$, $g$ or $h'$, and $a$ enters only through $S_2(B_{3\times5}(a))$, which
does not depend on $a$.
\end{proof}

\begin{auxrem}[the $\varphi^4$ instance; Corollary~5.2 recomputed]\label{aux:constants-phi4}
For $\mathscr P(\xi)=\xi^4$: $p'=2$, $a_4=1$, all other $a_n=0$, so $a_0=0$ and
$\sum_n|a_n|=1$.  Then $c_U=\lambda K_1K_2^8\,4!=24\,\lambda K_1K_2^8$;
$\mathfrak Q=9\cdot8!\,K_1K_2^{16}\lambda^2=362\,880\,K_1K_2^{16}\lambda^2$;
$\overline V_*=\lambda^2\,4!\,I_4=24\,\lambda^2I_4$; $9^{p'}=81$; and
\begin{align*}
 M_B^2&\le81\cdot24\,\lambda^2I_4\,S_2(B_{3\times5})
       +24\cdot362\,880\,C_7K_1^2K_2^{24}\lambda^3\,\overline\Sigma_{3\times5}\\
      &=1944\,\lambda^2I_4\,S_2(B_{3\times5})+8\,709\,120\,C_7K_1^2K_2^{24}\lambda^3\,
       \overline\Sigma_{3\times5} .
\end{align*}
(Arithmetic: $9\cdot40\,320=362\,880$; $81\cdot24=1944$; $24\cdot362\,880=8\,709\,120$.)  The
record's number $362\,880\,C_8K_1K_2^{16}\lambda^2$ (Corollary~5.2 \cite[p.~19]{R}) is
replaced, as the plan says; the new value is not a number, since $S_2$ is not evaluated, and
$I_4$ depends on $m_0=M_*$.
\end{auxrem}

\subsection*{A.3\quad The constants of the consequences}

\begin{auxlemma}[Gaussian moments of $X_f$]\label{aux:constants-gauss}
Let $f\in C_c^\infty(\R)$ be real, $X_f:=\varphi_0(f)$, $\sigma_f^2:=\ip f{Cf}=\iint
f(x)C(x-y)f(y)\,dx\,dy$ (Lemma~\ref{aux:notation-kernel}(ii)).  Then for $n\ge0$
\[
 \|X_f^{2n}\|_{L^2(\dphi)}=\sigma_f^{2n}\bigl((4n-1)!!\bigr)^{1/2}\le(\sqrt2\,\sigma_f)^{2n}(2n)! ,
\]
with $(-1)!!:=1$.
\end{auxlemma}

\begin{proof}
$X_f$ is centred Gaussian with variance $\sigma_f^2$, so $\E e^{sX_f}=e^{s^2\sigma_f^2/2}$;
comparing coefficients of $s^{2m}$, $\E[X_f^{2m}]/(2m)!=\sigma_f^{2m}/(2^mm!)$, i.e.\
$\E[X_f^{2m}]=(2m-1)!!\,\sigma_f^{2m}$.  With $m=2n$,
$\|X_f^{2n}\|_2=(\E X_f^{4n})^{1/2}=\sigma_f^{2n}((4n-1)!!)^{1/2}$.  Next,
$(4n-1)!!=(4n)!/(2^{2n}(2n)!)$ and $(4n)!=\binom{4n}{2n}((2n)!)^2\le2^{4n}((2n)!)^2$, so
$(4n-1)!!\le2^{2n}(2n)!$ and $((4n-1)!!)^{1/2}\le2^n((2n)!)^{1/2}\le2^n(2n)!$.  (Test $n=1$:
$\|X_f^2\|_2=\sqrt3\,\sigma_f^2\le4\sigma_f^2$.  Numerical check of $(4n-1)!!\le4^n(2n)!$ for
$n\le40$ in exact integers: \texttt{parts/cqft2\_ringcheck.py}, run record in
Remark~\ref{aux:notation-ringrem}(c); nothing is deduced from it.)
\end{proof}

\begin{auxcor}[Lemma~7.6 with the constants of Theorem~\ref{thm:D3}]\label{aux:constants-76}
Let $g\in\CH$, let $f\in C_c^\infty(\R)$ be real, $f\ne0$, supported in a bounded interval $I$
of length $L$, and let $A$, $B'$ be the constants of Proposition~\ref{aux:constants-D3} for
$g$, $I$ and $s=0$.  Put
\[
 K'_\square:=K_2^2\,(N^{\square}_I)^{1/2}\|f\|_2,\qquad K_f:=\max\{K'_\square,\sqrt2\,\sigma_f\},
 \qquad A_f:=A+B'K_1 .
\]
Assume the transfer
\[
 \E_{\mu_g}[X_f^{2n}]\le\sup_{T>0}\int X_f(0)^{2n}\,dq_{h_{T,g}}
\]
(the weak-graph passage of Lemma~7.1 of the record, as replaced in \S\ref{sec:cons}).  Then for
every $n\ge1$ and every $0\le s<1/K_f$,
\[
 \E_{\mu_g}[X_f^{2n}]\le A_f\,K_f^{2n}(2n)!,\qquad
 \E_{\mu_g}e^{s|X_f|}\le2\max\{1,A_f\}\sum_{n\ge0}(sK_f)^{2n}<\infty .
\]
The constants $A_f$ and $K_f$ do not depend on $n$.
\end{auxcor}

\begin{proof}
Proposition~\ref{aux:constants-D3} with $F=X_f^{2n}$ (a cylinder monomial at time $0$ with $2n$
factors, supported in $I$) gives
$\int X_f(0)^{2n}dq_{h_{T,g}}\le A\|X_f^{2n}\|_{L^2(\dphi)}+B'\GJS{X_f^{2n}}$ for every $T>0$.
By Lemma~\ref{aux:constants-gauss}, $\|X_f^{2n}\|_2\le(\sqrt2\sigma_f)^{2n}(2n)!$; by (N.a) of
Remark~\ref{aux:notation-counts} with $r=2n$ and all $f_i=f$,
$\GJS{X_f^{2n}}\le K_1K_2^{4n}(2n)!\,(N^{\square}_I)^n\|f\|_2^{2n}=K_1(K'_\square)^{2n}(2n)!$.
Both are $\le K_f^{2n}(2n)!$ times $1$, resp.\ $K_1$; add.  For the exponential bound:
$e^{s|x|}\le2\cosh(sx)$, and by monotone convergence
$\E_{\mu_g}\cosh(sX_f)=\sum_{n\ge0}s^{2n}\E_{\mu_g}[X_f^{2n}]/(2n)!\le1+A_f\sum_{n\ge1}(sK_f)^{2n}
\le\max\{1,A_f\}\sum_{n\ge0}(sK_f)^{2n}$, a geometric series with ratio $(sK_f)^2<1$.
\end{proof}

\begin{auxrem}[the radius $s_0$ of Theorem~7.7]\label{aux:constants-s0}
In the proof of Theorem~7.7 of the record \cite[p.~23]{R}, $F(z)=\E_{\mu_g}[\overline\zeta
\exp(\sum_{j=1}^kz_jX_{e_j})]$ must be defined and analytic in a polystrip
$\{|\mathrm{Re}\,z_j|<s_0\}$.  By Cauchy--Schwarz and H\"older with $k$ equal exponents,
\[
 \E_{\mu_g}\Bigl[|\zeta|\,e^{\sum_j|\mathrm{Re}\,z_j||X_{e_j}|}\Bigr]\le\|\zeta\|_{L^2(\mu_g)}
 \prod_{j=1}^k\Bigl(\E_{\mu_g}e^{2k|\mathrm{Re}\,z_j||X_{e_j}|}\Bigr)^{1/(2k)} ,
\]
finite, locally uniformly, when $2k|\mathrm{Re}\,z_j|<1/K_{e_j}$ for every $j$ with
$e_j\ne0$ (Corollary~\ref{aux:constants-76}, which assumes $f\ne0$; a factor with $e_j=0$ is
$1$).  So
\[
 s_0:=\min_{1\le j\le k,\ e_j\ne0}\frac1{2k\,K_{e_j}}
\]
works (any $s_0>0$ if all $e_j=0$); $1/K_f$ is the one-variable radius of Corollary~\ref{aux:constants-76}, and any
$s_0>0$ suffices for Theorem~7.7.  $s_0$ depends on $e_1,\dots,e_k$ and $K_2$, $m_0$, not on
$g$ (the constants $A_{e_j}$ do, but they do not enter $s_0$).
\end{auxrem}

\begin{auxrem}[Lemmas~4.3, 7.1--7.3, Proposition~7.5]\label{aux:constants-other}
(a) Lemma~4.3(ii) \cite[p.~17]{R} needs, at each fixed $T$, only the finiteness of
\[
 \int F(\Phi(s,\cdot))^2\,dq_{h_{T,g}}\le Z(h_{T,g})^{-1}\,\|F^2\|_{L^2(\dmu)}\,
 \|e^{-\Uc_{h_{T,g}}}\|_{L^2(\dmu)},
\]
which holds by Assumption~A3 ($e^{-2\Uc}\in L^1$, $0<Z$) and Gaussian moments of $F$; no
constant is needed.  (b) Lemma~7.1 becomes
\[
 \sup_{T>0}\|Q\hat\psi_T\|^2\le A\|Q^2\|_{L^2(\dphi)}+B'\,\mathfrak n_\square(Q)^2,\qquad
 \mathfrak n_\square(Q)^2:=K_1K_2^{4r}(2r)!\,(N^{\square}_I)^r\prod_i\|f_i\|_2^2 ,
\]
the second bound of (N.a), and Lemmas~7.2--7.3 use the same constants at the times $s$ and $\mp t/2$ (the
box $B_s(g,I)$ has the shape $(3,N_B)$ for every $s$).  Here
\[
 \|Q^2\|_{L^2(\dphi)}\le\bigl((4r-1)!!\bigr)^{1/2}\prod_i\sigma_{f_i}^2 ,
\]
by H\"older with $r$ equal exponents and the moment formula in the proof of
Lemma~\ref{aux:constants-gauss}:
\[
 \E\prod_iX_{f_i}^4\le\prod_i\bigl(\E X_{f_i}^{4r}\bigr)^{1/r}=(4r-1)!!\prod_i\sigma_{f_i}^4 .
\]
(c) Proposition~7.5 keeps
$\widehat C_Q=\max\{C_7e^{2m_1}\GJS Q^2,e^{3m_1}\|\psi_Q\|^2\}$, with $\GJS Q$ computed on the
squares of Definition~\ref{aux:notation-squares}; the gap $m_1$ of Theorem~7.8 is unchanged.
\end{auxrem}

\subsection*{A.4\quad The table}

All constants are those of the working coordinates (\S2.6 and Remark~3.3 of the record
\cite[pp.~14, 16]{R}).  ``Cond.'' marks a constant whose finiteness or bound uses Lemma~C.2,
Lemma~C.4, Theorem~C.5 or the comparison step in the proof of Theorem~C.5, hence is proved
conditional on repairs of the V3 audit (\S\ref{sec:notation}): on R1--R10 where $S_2$ enters or
the comparison bound of Lemma~\ref{aux:constants-comparison} is used (through Lemma~C.4, in
Lemmas~\ref{aux:notation-cU}(iv) and~\ref{aux:notation-cutoffs}(iv), and through the condition of
Lemma~\ref{lem:D1}), on R8 alone where only $Z$ or the factor $9^{p'}$ enters.  ``F-T'' marks a constant whose
derivation uses GJS Theorem~1.1.7 across strips transverse to the time-zero line (sentence F-T
of \S\ref{sec:purpose}).  ``GAP'' marks a constant that bounds the record's operator $V(h)$ and
so rests on the GAP after Lemma~\ref{aux:notation-wick}: no row as a constant of this note;
$M_B$ when it is used for the record's Theorem~5.1 (\S\ref{sec:cons}).  The rows $V_*$,
$9^{p'}$ and $M_B$ concern the constructed objects, which Lemma~\ref{aux:notation-sharpid}
identifies with the record's regularized ones, and need no GAP as constants.

\begin{center}
\footnotesize
\begin{tabular}{@{}p{0.10\textwidth}p{0.34\textwidth}p{0.18\textwidth}p{0.14\textwidth}p{0.12\textwidth}@{}}
\toprule
constant & formula & depends on & uniform in & derived in\\
\midrule
$\bar n$ & $2p'+1$ & $p'$ & --- & \S\ref{sec:notation-printed}; A8(ii)\\
$K_1$, $K_2$ & printed, (1.1.7)--(1.1.8); $K_2\ge1$ & fixed with $(m_0,\lambda)$ of \S2.6 &
$h$, $Q$ & \S\ref{sec:notation-printed}\\
$C_7$, $m_1$ & printed, (1.1.10) & fixed with $(m_0,\lambda,K_1,K_2)$; not printed &
$h$, $Q_1$, $Q_2$ & \S\ref{sec:notation-printed}\\
$c_U$ & $\lambda\sum_{n=1}^{2p'}|a_n|K_1K_2^{2n}n!$ & $\lambda$, $a_1,\dots,a_{2p'}$, $K_1$,
$K_2$ & $\Delta$, $h\in\CE$ & L.~\ref{aux:notation-cU}\\
$C_7c_U$ & constant of the comparison; cond. & $C_7$, $c_U$ & $h$, $B$, $F$, $\alpha$ &
L.~\ref{aux:constants-comparison}\\
$S_2(B)$ & $\sup\{\|e^{-\Uc_{h'}}\|_{L^2(\dmu)}:h'\in\CE,\ \supp h'\subseteq\overline B\}$;
$1\le S_2<\infty$; not evaluated; cond. & $\lambda$, $\mathscr P$, $m_0$; shape of $B$ &
position of $B$ & P.~\ref{aux:notation-S2}\\
$N_g$, $N_B$ & columns meeting $\supp g$ in positive measure; columns of $B_s(g,I)$ & $g$;
$g$ and $I$ & $s$, $T$ & P.~\ref{aux:constants-D3}\\
$A$ & $e^{3\lambda|a_0|N_B}S_2(B_s(g,I))$; cond. & $\lambda$, $\mathscr P$, $m_0$, $N_B$ &
$T>0$, $s\in\R$ & P.~\ref{aux:constants-D3}\\
$B'$ & $2C_7c_UN_ge^{-m_1}(1-e^{-m_1})^{-1}$; cond. & $C_7$, $m_1$, $c_U$, $N_g$ & $T$, $s$, position
of $I$ & P.~\ref{aux:constants-D3}\\
\bottomrule
\end{tabular}
\end{center}

\begin{center}
\footnotesize
\begin{tabular}{@{}p{0.10\textwidth}p{0.34\textwidth}p{0.18\textwidth}p{0.14\textwidth}p{0.12\textwidth}@{}}
\toprule
constant & formula & depends on & uniform in & derived in\\
\midrule
$N^{\square}_I$ & $\#\{k:|I\cap[k-\frac12,k+\frac12)|>0\}\le\lfloor L\rfloor+2$ & $I$ & --- &
R.~\ref{aux:notation-counts}\\
$\mathfrak n_\square(Q)^2$ & $K_1K_2^{4r}(2r)!(N^{\square}_I)^r\prod_i\|f_i\|_2^2$ &
$K_1$, $K_2$, $r$, $I$, $f_i$ & $s$ & R.~\ref{aux:notation-counts}\\
$\sigma_f$ & $\bigl(\iint f(x)C(x-y)f(y)\,dx\,dy\bigr)^{1/2}$ & $f$, $m_0$ & --- &
L.~\ref{aux:notation-kernel}\\
$K'_\square$ & $K_2^2(N^{\square}_I)^{1/2}\|f\|_2$ & $K_2$, $I$, $f$ & $n$ & C.~\ref{aux:constants-76}\\
$K_f$ & $\max\{K'_\square,\sqrt2\,\sigma_f\}$ & $K_2$, $m_0$, $I$, $f$ & $n$, $g$ &
C.~\ref{aux:constants-76}\\
$A_f$ & $A+B'K_1$ ($A$, $B'$ for $g$, $I\supseteq\supp f$, $s=0$); cond. & as $A$, $B'$;
$K_1$ & $n$, $T$ & C.~\ref{aux:constants-76}\\
$s_0$ & $\min_{j\le k,\,e_j\ne0}(2kK_{e_j})^{-1}$; one-variable radius $1/K_f$ & $e_1,\dots,e_k$,
$K_2$, $m_0$ & $g$ & R.~\ref{aux:constants-s0}\\
$\mathfrak Q$ & $9K_1K_2^{8p'}(4p')!\lambda^2(\sum_{n=0}^{2p'}|a_n|)^2\ge\GJS{Q_V}$ &
$\lambda$, $\mathscr P$, $K_1$, $K_2$ & $h'$, $a$ & P.~\ref{aux:constants-D4}\\
$\overline\Sigma_{3\times5}$ & $\sum_{j\ge1}(8j+12)e^{-m_1j}=4e^{-m_1}(5-3e^{-m_1})
(1-e^{-m_1})^{-2}\ge\Sigma_{3\times5}$ & $m_1$ & $a$, $h$, $R\subseteq S_a$ &
L.~\ref{aux:notation-rings}\\
$b_0$, $I_n$ & $\log4+e^{-1}+m_0^{-2}$; $I_0=4$,
$I_n\le(4\pi)^{-n}2^{n+1}(2b_0^n+2^n(n!+(\log2)^n))$ & $m_0$, $n$ & --- &
L.~\ref{aux:notation-kernel}\\
$V_*$ & $\sup_{h'}\|V(h')\|^2_{L^2(\dphi)}\le\overline V_*:=\lambda^2\sum_{n=0}^{2p'}a_n^2n!\,I_n$,
with $V(h')=\lambda\sum_na_nW_n$ of L.~\ref{aux:notation-wick}(iv)
& $\lambda$, $\mathscr P$, $m_0$ & $h'$, $a$ & L.~\ref{aux:notation-wick}\\
$9^{p'}$ & $\|Y\|_{L^4(\dmu)}^2\le9^{p'}\|Y\|_{L^2(\dmu)}^2$; cond. & $p'$ & $h'$, $a$ &
L.~\ref{aux:notation-hyper}\\
$M_B^2$ & $e^{15\lambda|a_0|}S_2(B_{3\times5})\,9^{p'}V_*+C_7c_U\overline\Sigma_{3\times5}\mathfrak Q$;
cond.; F-T; GAP for Thm~5.1 & $(\lambda,\mathscr P,m_0)$ via $C_7$, $m_1$, $K_1$, $K_2$, $V_*$, $S_2(B_{3\times5})$ &
$h\in\CE$ (so $T$, $g$), $a$, $h'$ & P.~\ref{aux:constants-D4}\\
$M_B^2$ ($\varphi^4$) & $\le1944\lambda^2I_4S_2(B_{3\times5})+8\,709\,120\,C_7K_1^2K_2^{24}
\lambda^3\overline\Sigma_{3\times5}$; cond.; F-T; GAP for Cor.~5.2 & $\lambda$, $m_0$, $C_7$, $m_1$, $K_1$, $K_2$ & as $M_B$ &
R.~\ref{aux:constants-phi4}\\
\bottomrule
\end{tabular}
\end{center}

In the record's Appendix~A \cite[p.~35]{R} the rows that change are, as \S\ref{sec:cons} lists
them: $K_1,K_2,C_7,C_8,m_1$ ($C_8$ dropped after Phase~B); $N_I$ (counted on GJS's columns, bound
$\lfloor L\rfloor+2$ unchanged, Remark~\ref{aux:notation-counts}(i)); $M^2$ (now $M_B^2$); $M^2$
($\varphi^4$); $K'$ (same formula, $N_I$ by columns); and the $\Psi_\gamma$ rows (the $M$ in them
is $M_B$).  Unchanged: $\varepsilon_{\mathscr P}$, $\mathfrak K(p')$, $\widehat C_Q$ (its formula;
$\GJS Q$ is computed on GJS's squares), $\mathfrak m$, $d_0$ and the shell count.  New rows: $c_U$,
$S_2(B)$, $A$, $B'$, $A_f$, $K_f$ (with $\sigma_f$), $s_0$, $V_*$, $\overline\Sigma_{3\times5}$ and
$M_B$; $S_2(B)$, and through it $A$, $A_f$ and $M_B$, contain the unevaluated finite constant of
Theorem~C.5.

\subsection*{A.5\quad Discrepancies found}

Compared with \texttt{ugjs/PLAN.md}, Part~1 (the targets).  Each item is a finding for the
audit; none changes a theorem of the record.
\begin{enumerate}
\item[(D1)] \emph{Counts on GJS squares} (MINOR in the plan; MAJOR in \S\ref{sec:cons} if it
cites the record's (7)--(8) or $K'$ with $N_I$ for norms on the squares of
Definition~\ref{aux:notation-squares}).  The plan's $K_f:=\max\{K',\sqrt2\sigma_f\}$ and its
Lemma~7.1 bound ``$B'\mathfrak n(Q)^2$'' use the record's count $N_I$ of cells $[j,j+1)$.  The
norms of Theorem~\ref{thm:D3} are computed on GJS's squares, where the bound can fail with
$N_I$: Remark~\ref{aux:notation-counts}(iii) gives $\GJS{\varphi_0(f)}=\sqrt2K_1K_2^2\|f\|_2$ for
$f$ even about $\frac12$ and supported in $(0.4,0.6)$, against the record's bound
$K_1K_2^2\|f\|_2$.  Repair: $N^{\square}_I$ in place of $N_I$ (this appendix), or the common bound
$\lfloor L\rfloor+2$.  For the record itself the same example makes Proposition~6.2, eq.~(7),
false as stated when $\GJS\cdot$ is computed on GJS's squares from the column decomposition; the
V1 attack classifies this as CRITICAL for Proposition~6.2, MAJOR for the displayed bounds of
Lemmas~7.1 and~7.6, and MINOR for Theorems~7.7, 7.8 and~12.2, which need finiteness only.  This
appendix adopts that split; \S\ref{sec:cons} repairs the record.
\item[(D2)] \emph{The radius $s_0$} (MINOR).  The plan's ``Theorem~7.7 then has $s_0=1/K_f$'' is
the one-variable radius of Lemma~7.6.  The polystrip of Theorem~7.7 has $k$ variables, and
Cauchy--Schwarz with H\"older give the half-width $\min_{j,\,e_j\ne0}(2kK_{e_j})^{-1}$
(Remark~\ref{aux:constants-s0}); any positive value suffices.
\item[(D3)] \emph{Where the transverse case enters} (MINOR, disclosure).  The plan (Part~1,
Theorem~D.4) names ``Squares in rows $-1,0,1$ outside $B_{3\times5}$''.  Those are the squares with
no parallel strip of positive width, but the uniform ring bound uses the transverse strip for
every side square of $\mathcal L_j$ with $|l|\le j$ (Remark~\ref{aux:notation-ringrem}(b));
sentence F-T must accompany all of them.
\item[(D4)] \emph{The columns of $B_{3\times5}$} (no change of constants).  ``The at most three
columns meeting $[a,a+2]$'' are exactly three, $k_a,k_a+1,k_a+2$, when ``meeting'' is
set-theoretic (Lemma~\ref{aux:notation-rings}(i)); in positive measure only two meet when
$a\in\frac12+\Z$, and the box would be $3\times4$, with ring count $8j+10$.  Every constant
stays valid: $S_2$ is monotone (Proposition~\ref{aux:notation-S2}(ii)), $e^{12\lambda|a_0|}\le
e^{15\lambda|a_0|}$, and $8j+10\le8j+12$.  This appendix uses the three columns always.
\item[(D5)] \emph{Which squares are summed} (consistency).  The plan's D.2 sums over squares
meeting $\supp h$; Theorem~D.3 counts $N_g$ columns meeting $\supp g$ in positive measure.  These
match if the sum is over the charged squares $\mathfrak S(h)$, as here; if it is over squares
meeting $\supp h$ set-theoretically, $N_g$ must be read set-theoretically too, or the extra
squares (with $U_\Delta=0$) dropped.
\item[(D6)] \emph{The coefficient $a_0$} (observation).  $dq_h$ does not depend on $a_0$
(Lemma~\ref{aux:notation-cU}(iii)), so the factors $e^{3\lambda|a_0|N_B}$ in $A$ and
$e^{15\lambda|a_0|}$ in $M_B$ could be dropped by computing $S_2$ for $\mathscr P-a_0$; the plan's
constants are valid as they stand.
\item[(D7)] \emph{Agreement.}  $c_U$, $A$, $B'$, the bound $\Sigma_{3\times5}\le\overline\Sigma_{3\times5}=\sum_j(8j+12)e^{-m_1j}$,
the bound on $V_*$, the factor $9^{p'}$, $M_B$, $A_f$ and the inequality
$\|X_f^{2n}\|_{L^2(\dphi)}\le(\sqrt2\sigma_f)^{2n}(2n)!$ agree with the plan.  New here: the
closed form of $\overline\Sigma_{3\times5}$, the explicit bound on $I_n$, and the $\varphi^4$ value of
Remark~\ref{aux:constants-phi4}.
\item[(D8)] \emph{Status of Appendix~C} (not a constant).  The plan counts Lemmas~C.2--C.5 among
``the record's own proofs'' (Part~7).  After the V3 audit, every constant marked ``cond.'' rests
on them conditional on the repairs R1--R10 (on R8 alone where only Lemma~C.2 or the comparison
step enters).
\item[(D9)] \emph{The ring constant in $M_B$} (deviation, stated).  The plan writes $M_B$ with
$\Sigma_{3\times5}=\sum_{\Delta\not\subseteq B_{3\times5}}e^{-m_1d(\Delta,R)}$, over the widest
strips.  This appendix and \S\ref{sec:d4} use the closed form
$\overline\Sigma_{3\times5}\ge\Sigma_{3\times5}$ (Lemma~\ref{aux:notation-rings}(vi)), because the
comparison is applied along assigned axis-parallel strips whose orientation is named
(Lemma~\ref{aux:constants-comparison}); the widest strips may be tilted
(Remark~\ref{aux:notation-tilted}).  $M_B$ is therefore an upper bound for the plan's value.
\end{enumerate}

Compared with \S\S\ref{sec:d2}--\ref{sec:d4} and \S\ref{sec:cons}: \S A.6.

\subsection*{A.6\quad Comparison with \S\S\ref{sec:d1}--\ref{sec:cons} as drafted}

\S\S A.1--A.5 were written before the drafts of \S\S\ref{sec:d1}--\ref{sec:cons} existed.  Those
drafts were read afterwards (2 November 2026, versions of 16:00--16:06) and again after their
duplicates were merged into \S\ref{sec:notation} (versions of 16:24--16:31), and compared item by
item.
\begin{enumerate}
\item[(C1)] \emph{Agreement of every constant.}  $c_U$ (Lemma~\ref{aux:notation-cU}); $A$ and $B'$
(Theorem~\ref{thm:D3}); the ring constant $\overline\Sigma_{3\times5}=4x(5-3x)(1-x)^{-2}$, with
which \S\ref{sec:d4} states Theorem~\ref{thm:D4} (the bound of
Lemma~\ref{aux:notation-rings}(vi)); the bound on
$V_*$ (Lemma~\ref{aux:notation-wick}); the factor $9^{p'}$ (Lemma~\ref{aux:notation-hyper}); $M_B$
(Theorem~\ref{thm:D4}); $A_f$, $K_f$ and $s_0$ (\S\ref{sec:cons}); and the $\varphi^4$ value of
Corollary~5.2 as replaced in \S\ref{sec:cons},
$81\,S_2(B_{3\times5})V_*+8\,709\,120\,C_7K_1^2K_2^{24}\overline\Sigma_{3\times5}\lambda^3$, which with
$V_*\le24\lambda^2I_4$ is Remark~\ref{aux:constants-phi4}.  No numerical discrepancy was found.
\item[(C2)] \emph{The count.}  \S\ref{sec:d3} found the discrepancy (D1) independently
(Remark~\ref{aux:d3-cells}, with $\Phi(s,f)^2$ and the values $3K_1K_2^4$ against $2K_1K_2^4$),
and \S\ref{sub:cons-NI} repairs it by redefining the record's $N_I$ by GJS's columns; the
redefined $N_I$ is the $N^{\square}_I$ of Remark~\ref{aux:notation-counts}.
\item[(C3)] \emph{The transverse case.}  \S\ref{sec:d4} (its subsection on where the transverse
case is used) states (D3) independently: the squares of class (T) in rows $|l|\ge2$ also use the
transverse strip.  It assigns the strips of Lemma~\ref{aux:notation-rings}(iv), the corner
squares of each ring ($|l|=1+j$ in the side columns) to the parallel class;
Proposition~\ref{aux:constants-D4} does the same, so every square of $\mathcal L_j$ carries width
$j$ and the bound is $\overline\Sigma_{3\times5}$.
\item[(C4)] \emph{The radius $s_0$.}  \S\ref{sec:cons} restricts the minimum to $e_j\ne0$; this
appendix now does too (Remark~\ref{aux:constants-s0}).  \S\ref{sec:quant} writes
$\min_{j\le k}(2kK_{e_j})^{-1}$ without the restriction (MINOR; same repair: for $e_j=0$,
$K_{e_j}=0$).
\item[(C5)] \emph{Duplicates.}  After the merge, \S\S\ref{sec:d2}--\ref{sec:d4} cite
\S\ref{sec:notation} for the cutoffs, the strips, $c_U$, $S_2$, the rings, the comparison
$9^{p'}$ and $V_*$.  The remaining overlap is the kernel bound of \S\ref{sec:d1}
(Lemma~\ref{aux:d1-cov}, which bounds the Euclidean kernel $G$ and the time-zero kernel
$C=G(0,\cdot)$ with other, also valid, constants than Lemma~\ref{aux:notation-kernel}).  The two
differences recorded before the merge (a bound on $V_*$ from the record's uncited ``classical
Bessel facts'', and a comparison run on $L^2(\dphi)$) are gone: \S\ref{sec:d4} now uses
Lemmas~\ref{aux:notation-wick} and~\ref{aux:notation-hyper}.
\item[(C6)] \emph{The form of Lemma~D.2.}  Lemma~\ref{lem:D2}(a) and
Lemma~\ref{aux:constants-comparison} have the same form: explicitly assigned strips of positive
width $w_\Delta$, one for each charged square outside the box, and the bound
$C_7c_U\GJS F\sum e^{-m_1w_\Delta}$.  In Theorems~\ref{thm:D3} and~\ref{thm:D4} every assigned
strip has width $\ge1$.
\item[(C7)] \emph{The two constructions of $Q_V$} (audit V1, finding~1; audit V2, item~V2-1).
Theorem~\ref{thm:D4} bounds its base term through Lemma~\ref{aux:notation-hyper} (the $J_0$
construction) and its comparison term through Lemma~\ref{lem:D2} (the pieces of
Definition~\ref{aux:d1-obs}(b), built with the record's mollifier).
Lemma~\ref{aux:notation-sharpid} and Remark~\ref{aux:notation-sharpidrem}(b) identify the two, so
Proposition~\ref{aux:constants-D4} and Theorem~\ref{thm:D4} bound one observable.  With $V_*$
taken for the constructed $V(h')$ of Lemma~\ref{aux:notation-wick}(iv), the Euclidean bound of
Theorem~\ref{thm:D4} needs no GAP; the GAP enters at the transfer to the record's Theorem~5.1
(\S\ref{sec:cons}).
\end{enumerate}


\section{Gate record}\label{app:gates}
% gates.tex -- gate record of the note (orchestrator, 2 November 2026; updated after each round).
\subsection*{Status of the note}
The note as a whole has status \emph{conjectured} until the audits of \texttt{ugjs/PLAN.md}
Part~5 and gates GA1--GA8 and GB1--GB3 (Part~3) have passed; this appendix records the rounds.
The status lines inside the sections are the drafters' claims after the disposition of rounds V1
and V2 (3 November 2026), under re-check: Lemma~\ref{lem:D1} proved, conditional on repair R8 of
the V3 audit (which supplies the linearity of $h\mapsto\Uc_h$ and, with Lemma~C.4 on the
Step-1 route (R1--R10), the identification of the record's $\Uc_h$ with the $L^2$-limit of its
regularizations $\Uc_{h,\kappa}$) (mp); Lemma~\ref{lem:D2} proved, conditional on Lemma~\ref{lem:D1} and on Lemma~C.4 (R1--R10), which
identifies $\Uc_{h-h^B}$ with the limit of its regularizations in its Step~3, using neither
Lemma~C.2 nor Theorem~C.5 (aqft); Theorem~\ref{thm:D3} proved, conditional on
Lemma~\ref{lem:D1} and Lemma~C.4 (both through \ref{lem:D2}; R1--R10), Lemma~C.2 (R8) and
Theorem~C.5 with the shape property of $S_2$ (R1--R10), parallel strips only (aqft); Theorem~\ref{thm:D4} proved,
conditional on the same, on R8 for the comparison step of Lemma~N.21(ii), and on the transverse
case of Theorem~1.1.7 at the printed statement's scope (F-T), and not using the GAP of
\S\ref{sec:notation} (aqft); the auxiliary statements of \S\ref{sec:notation} and
Appendix~\ref{app:constants} proved or conditional as marked there (cqft2), among them the
identification Lemma~N.22 that closed the one MAJOR finding of V1 and V2. The one GAP mark of the
note stands in \S\ref{sec:notation}, after the Wick-isometry lemma: the record never constructs
the Wick powers in its $V(h)$ (\S2.1 of the record), so the identification of the record's
$V(h)$ with the constructed object is assumed; it enters only at Step~1 of the record's
Theorem~5.1 and at Corollary~5.2 (\S\ref{sec:cons}), not in the Euclidean Theorem~\ref{thm:D4}.
It is the time-zero face of the V3 audit's finding M-3, whose repair R8 would supply the
construction. The status lines of Lemma~\ref{lem:D2} and Theorems~\ref{thm:D3}--\ref{thm:D4} also rest on
the reading, named in Remark~N.2(iv) of \S\ref{sec:notation}, that identifies the note's
localized pieces, one-square interaction terms and $V(h')(0)^2$ with GJS's monomials of
p.~593; the printed text that reading uses is the object of C5 item~10, an open human reading. Nothing in this note is a theorem of the record
until integration (Checkpoint~2).

\subsection*{Deviations from PLAN Part 1 (findings for the audits)}
Lemma~\ref{lem:D2} is stated in a square-by-square form (a), summing over the squares where $h$ is
not almost everywhere zero, from which a form (b) of the plan's statement, with strips of
positive width, follows; the plan's index set
$\{\Delta\not\subset B:\Delta\cap\supp h\neq\emptyset\}$ would break Theorem~\ref{thm:D3}'s
parallel-only claim and its $2N_g$ count (mp's example $g=\mathbf 1_{[0,1/10]}+\mathbf 1_{\{3\}}$).
$\Sigma_{3\times5}$ is the constant $4x(5-3x)/(1-x)^2$, $x=e^{-m_1}$, bounding the plan's
$R$-dependent sum. Transverse strips are used for every square whose column distance exceeds $|l|$,
not only in rows $-1,0,1$. No Wick re-expansion of $V(h')(0)^2$ is needed (the record's Step~1).
Theorem~7.7's radius is $s_0=\min_{j:\,e_j\neq0}(2kK_{e_j})^{-1}$, not the plan's $1/K_f$. The
$3^{p'}$ comparison is derived under $\dmu$ (and on $L^2(\dphi)$) from A4. $B_{3\times5}$ always has
15 squares, built from the column containing $a$. The count $N_I$ of the record's \S6 must be
redefined on GJS columns (\S\ref{sec:cons}; PLAN finding F3 was inexact on this point).

\subsection*{Source facts recorded by the drafters (misprints kept verbatim)}
GJS p.~594 prints the strip-definition sentence before Theorem~1.1.7 without a closing period
(cqft2, aqft, 300\,dpi; confirmed by the V1 attacker at 600\,dpi). GJS p.~595 prints
``Feyman-Kac'' [sic] (mp). GJS (1.1.11) on p.~594 is printed ``0(1)$\,e^{-m_1d/2}$'' with an
upright digit zero, the glyph of the label's ``0'', while (1.1.10) on the same page prints the
italic ``$O(1)$'' (mp at 300\,dpi; the V1 attacker at 600\,dpi with glyph metrics; the
orchestrator on those crops); the record's A8(vii$'$)(a), l.~814, transcribes (1.1.11) with
``$O(1)$'', a silent correction to be marked [sic] at integration (KB correction node
\texttt{a8-vii-prime-silent-correction-of-gjs-0-1}).

\subsection*{Audits}
\emph{V1 (2--3 November 2026; the frozen version of 2 November, 16:38, 59 pp, sha256 of the PDF
beginning 7e393b9a440fe6a2).} mp, quantifiers of \S\S4--6 and \S8: no CRITICAL or MAJOR
finding, 14 MINOR (conditions not stated in the statements of D.2--D.4; the GAP unnamed at D.4's
Step~2; D.2's dependence misdescribed; the replaced Lemma~4.3(ii) uses more than A3; $L=0$;
the ``artefact'' sentence inside Lemma~7.6's statement; the polystrip threshold and $s_0$ when
all $e_j=0$; $K_f$ does not depend on $g$; $N_g=0$ one direction only; ``$T<1$'' covers $T\le\frac12$
only; the integration sentence of Theorem~12.2; the master's ``radius of Theorem~7.7''; \S7 and
Appendix~A omit the GAP from $M_B$'s dependence). cqft, attacker: D.1--D.3 hold and can be
recorded as proved conditional on the V3 repairs (D.1: the identification of $\Uc_h$ with the
$L^2$-limit, R8 with Lemma~C.4; D.2: Lemma~C.4; D.3: Lemma~C.2 at R8 and Theorem~C.5 at R1--R10,
parallel strips only); every application of Theorem~1.1.7 in \S\S4--6 checks out and every
transverse one is flagged; D.4 has one MAJOR gap: it treats two constructions of $V(h')(0)^2$
(time-zero Wick powers with a one-dimensional mollifier, carried by $J_0$, in Step~2; the limits
with the record's two-dimensional mollifier $\delta_\kappa$ in Step~4) as one object, which
\S8's transfer to Theorem~5.1 needs; the repair is the cross-covariance computation giving
$\|X_n-J_0W_n\|^2=0$; 14 MINOR (status labels; the $N_I$ repair incomplete at l.~448, A8(ii),
Theorem~5.1 Steps~2--4, Lemma~6.1 and Proposition~6.2; the severity of the cell-count error;
``the case the printed proof reaches'' overstated for a sharp-time/time-integrated pair; F-T
missing at Theorem~12.2, Lemma~10.2 and Corollary~5.2; the ``$\|$le'' line; D.2(b)'s ``positive
width''; the input lists; Lemma~N.18(iii) and $a_0$; F-T adapted rather than verbatim; \S7's
$N^{\mathrm{col}}_I$ and $s_0$; the ``0(1)'' transcription). The reports are
\texttt{ugjs/council/v1-cqft-report.md} and \texttt{ugjs/council/v1-mp-report.md}.
Disposition of every finding: done by the drafters on 3 November 2026 (\texttt{ugjs/council/disposition-aqft.md}, \texttt{-mp.md}, \texttt{-cqft2.md}); the MAJOR is closed by Lemma~N.22 of \S\ref{sec:notation}, a new claim under re-check.
\emph{V2 (3 November 2026; same frozen version; a fresh auditor that read neither the memos nor
the V1 reports).} No CRITICAL finding and no false mathematical statement in D.1--D.4; the same
MAJOR as V1 (V2-1, the two constructions of $Q_V$), 21 MINOR (status and conditionality of
D.1; R8 against R1--R10 at one place; ``conjectured'' and ``proved conditional'' applied to the
same statements; the $N_I$ edit incomplete at l.~448 and A8(ii); the replaced Lemma~4.3(ii)'s
inputs; eight replacement blocks, not ten, and the GA5 list; GA6's list omits Lemma~10.3 and
Proposition~10.4; the GA3 sentence; F-T not verbatim and missing at D.2(b), Lemma~N.46 and the
Theorem~5.1 block; the ``artefact'' sentence; $\mathcal H_r$ against R8's $\Gamma_r$; \S8.9's
stale quote; $N^{\mathrm{col}}_I$ undefined; $s_0$; Appendix~A.4 against \S8.7; the input list;
Lemma~N.18(iii); script run records; the V3 report under two paths; locator slips; the
``0(1)'' transcription). Every constant, strip width and orientation recomputed by hand
agrees with the note; all GJS quotations match the images. Report:
\texttt{ugjs/council/v2-report.md}.

\subsection*{Gates GA and GB (PLAN Part 3), as of the re-check rounds V1b and V2b}
\begin{itemize}\setlength{\itemsep}{1pt}
\item GA1 (every use of (1.1.10) names $Q_1$, $Q_2$ with their types, their squares, the cutoff with
  its membership in $\CE$, the strip with orientation and width; a second agent re-derives each):
  \emph{passed}. The V1 attacker checked every application in \S\S\ref{sec:d2}--\ref{sec:d4}
  (types (1.1.5)/(1.1.6), $n_i<\bar n$, $L^2$ kernels, disjoint closed squares, orientation and
  width, $h_\alpha\in\CE$, covariance under $dq_{h_\alpha}$, the constant) and the V2 auditor
  re-derived every strip width and orientation; every transverse use is flagged.
\item GA2 (the integrability chain of D.1 cites only A3, Lemma~C.2, Theorem~C.5, the Gaussian
  moments, H\"older and dominated convergence, every exponent written): \emph{passed with one
  addition}: the chain also uses the linearity of $h\mapsto\Uc_h$ (repair R8), which the
  plan's list omitted (mp); exponents written; checked by V1 (cqft) and V2.
\item GA3 ($A$, $B'$, $A_f$, $K_f$ recomputed by hand and entered in Appendix~A): \emph{passed};
  Appendix~\ref{app:constants} (cqft2, independent of aqft's text), re-recomputed by V2 and
  V2b: every constant agrees with \S\S\ref{sec:d3}--\ref{sec:d4} and \S\ref{sec:cons}; the
  deviations from PLAN Part~1 are listed in Appendix~A.5.
\item GA4 (an independent audit of Lemmas~C.2--C.4 and Theorem~C.5, including the
  translation-invariance sentence): \emph{run, with findings}: round V3, whose report is
  \texttt{ugjs/council/}\allowbreak\texttt{v3-appendix-c-report.md}; the four statements are true and every
  display re-derives, but three load-bearing inputs are neither proved nor imported (MAJOR);
  repairs R1--R10 written out there, not applied. Until they are, every use of Lemma~C.2,
  Lemma~C.4 and Theorem~C.5 in this note is conditional, as the status lines say; the
  translation invariance of the bound is shown in Theorem~C.5's proof and that of $S_2(B)$ in
  \S\ref{sec:notation}.
\item GA5 (no proof in \S\S4, 6, 7 of the record cites \texttt{eq:UGJS}): \emph{passed for the
  replacement texts}: the eight replacement blocks of \S\ref{sec:cons} contain no
  \texttt{UGJS} and no $C_8$ (grep recorded in \S8.8, re-run by V1 and V2); the record itself
  still carries 21 hits in \S\S4--7, all inside the texts to be replaced at integration.
\item GA6 ($A$, $B'$, $\widehat C_Q$, $s_0$ appear in no statement of \S\S8--14 of the record):
  \emph{passed}: grep recorded in \S8.8 (re-run by V1 and V2); $M$ appears in Lemmas~10.3 and
  10.5, Proposition~10.4 and the later statements and becomes $M_B$.
\item GA7 (a G6-style adversarial pass with zero unresolved findings): \emph{see the audit
  record above}: rounds V1 (two auditors) and V2 (fresh), re-checks V1b and V2b; no CRITICAL
  finding, the one MAJOR closed by Lemma~N.22 and confirmed by both re-checks; the residue of
  the re-checks (status wording of D.2--D.4, Remark~N.23(c), Remark~N.2(iv), Appendix~A.4's
  marks, one F-T, \S7's description of D.2's route, the run-record path) is listed in the
  disposition files and its closure is recorded below.
\item GA8 (the C5 checklist contains items~1--4 and 10 with pages and text): \emph{passed}
  (\texttt{part\_i/C5\_reading\_record.md}, items 1--4 and 10--13, committed 4166883).
\item GB1 (the surviving \texttt{U\_GJS}/\texttt{UGJS}/\texttt{C\_8} hits in the record are
  exactly the listed lines): \emph{not applicable before integration}; the record is unchanged
  by this note. The lines that would survive are listed in \S8.8.
\item GB2 (C5 items 1--4 and 10--13 exist with pages, images and text; the status sentences
  name them as open): \emph{passed} (items 10--13 added 31 October 2026; named as open in the
  record's Appendix~B, the README and \texttt{docs/}).
\item GB3 (Appendix~A's $M$ rows rewritten; $M$ contains $S_2$, unevaluated, and Appendix~A
  says so): \emph{prepared, not applied}: \S8.7 lists the rows; applied at integration.
\end{itemize}
\emph{Closure of the re-check residue:} confirmed by the attacker's closing pass V1c (4 November
2026, on the frozen version of 18:30, 72 pp): every item of V1b and V2b closed with no new
defect; final word: D.1--D.4 can be recorded exactly as their status lines and the status
paragraph above say; two wording leftovers in the master's \S1 (the labelling of the inputs
on working coordinates, V1c-1; the qualifier on ``no new import'') were applied by the
orchestrator after V1c, wording only. For the repository the two numerics output files
with the suffix \texttt{.out} were renamed to the suffix \texttt{.output.txt}, since that pattern is
git-ignored as a LaTeX artefact; the citations follow.
Of V1c's four MINOR leftovers, two were applied at once (V1c-1 and the qualifier), the
quadrature tolerance of the ring check is now stated in Remark~N.12(c), and the cosmetic
one-line summaries of the reading are left as they are. GA7 is therefore recorded as passed
for this version, with every status-affecting finding resolved and every status conditional
as stated; the note's own status stays \emph{conjectured} until integration.


\begin{thebibliography}{9}
\bibitem{R} A.~Stottmeister, \emph{Removal of the spatial cutoff for weakly coupled
$P(\varphi)_2$: norm convergence of the ground-state net on the quasi-local algebra. Part I ---
unified manuscript}, document of record \texttt{part\_i/main.pdf}, 42 pp., version of
31 October 2026, \url{https://github.com/alexander-stottmeister/pphi2-cutoff-removal}.
\bibitem{GJS} J.~Glimm, A.~Jaffe, T.~Spencer,
\emph{The Wightman axioms and particle structure in the $\mathscr P(\phi)_2$ quantum field
model},
Ann.\ of Math.\ \textbf{100} (1974) 585--632.
\end{thebibliography}
\end{document}
